% Options for packages loaded elsewhere
\PassOptionsToPackage{unicode}{hyperref}
\PassOptionsToPackage{hyphens}{url}
\PassOptionsToPackage{dvipsnames,svgnames,x11names}{xcolor}
\documentclass[
  chinese,
  9pt,
  a4paper,
]{ctexart}
\usepackage{xcolor}
\usepackage{amsmath,amssymb}
\setcounter{secnumdepth}{-\maxdimen} % remove section numbering
\usepackage{iftex}
\ifPDFTeX
  \usepackage[T1]{fontenc}
  \usepackage[utf8]{inputenc}
  \usepackage{textcomp} % provide euro and other symbols
\else % if luatex or xetex
  \usepackage{unicode-math} % this also loads fontspec
  \defaultfontfeatures{Scale=MatchLowercase}
  \defaultfontfeatures[\rmfamily]{Ligatures=TeX,Scale=1}
\fi
\usepackage{lmodern}
\ifPDFTeX\else
  % xetex/luatex font selection
\fi
% Use upquote if available, for straight quotes in verbatim environments
\IfFileExists{upquote.sty}{\usepackage{upquote}}{}
\IfFileExists{microtype.sty}{% use microtype if available
  \usepackage[]{microtype}
  \UseMicrotypeSet[protrusion]{basicmath} % disable protrusion for tt fonts
}{}
\makeatletter
\@ifundefined{KOMAClassName}{% if non-KOMA class
  \IfFileExists{parskip.sty}{%
    \usepackage{parskip}
  }{% else
    \setlength{\parindent}{0pt}
    \setlength{\parskip}{6pt plus 2pt minus 1pt}}
}{% if KOMA class
  \KOMAoptions{parskip=half}}
\makeatother
\ifLuaTeX
\usepackage[bidi=basic,provide=*]{babel}
\else
\usepackage[bidi=default,provide=*]{babel}
\fi
% get rid of language-specific shorthands (see #6817):
\let\LanguageShortHands\languageshorthands
\def\languageshorthands#1{}
\setlength{\emergencystretch}{3em} % prevent overfull lines
\providecommand{\tightlist}{%
  \setlength{\itemsep}{0pt}\setlength{\parskip}{0pt}}
% 这个文件是 Pandoc 生成 LaTeX 时注入的排版头文件。
% 目标不是做“论文式漂亮排版”，而是做“纸质开卷考试速查排版”：
% 1. 字体要清楚；
% 2. 页数要尽量少；
% 3. 标题、列表、段落的间距要紧凑；
% 4. 每页要有页眉和页码；
% 5. 目录必须显示真实页码，方便打印后查找。

\usepackage[a4paper,left=12mm,right=12mm,top=13mm,bottom=14mm,headheight=10pt,headsep=3mm,footskip=8mm,includeheadfoot]{geometry}
\usepackage{setspace}
\usepackage{enumitem}
\usepackage{titlesec}
\usepackage{fancyhdr}
\usepackage{tocloft}
\usepackage{needspace}
\usepackage{xcolor}
\usepackage{etoolbox}

% 中文字体：ctexart 在 Windows/XeLaTeX 下会自动使用系统中文字体。
% 这里不强行写死字体文件，避免不同机器字体名略有差异导致编译失败。
\setCJKmainfont{SimSun}[AutoFakeBold=2.5]
\setCJKsansfont{Microsoft YaHei}[AutoFakeBold=2.0]
\setCJKmonofont{Microsoft YaHei}

% 正文字号和行距：9pt + 1.03 倍行距，适合近距离打印阅读。
\linespread{1.03}
\setlength{\parindent}{1.4em}
\setlength{\parskip}{0.12em}

% 列表是 Markdown 笔记最占页数的部分，所以这里显著压缩列表上下间距。
\setlist{nosep,leftmargin=1.6em,itemsep=0.04em,topsep=0.08em,parsep=0pt,partopsep=0pt}
\setlist[enumerate]{labelsep=0.45em}
\setlist[itemize]{labelsep=0.45em}

% 标题间距：保证能一眼看出层级，但不浪费大片空白。
\titleformat{\section}{\large\bfseries\sffamily}{\thesection}{0.6em}{}
\titleformat{\subsection}{\normalsize\bfseries\sffamily}{\thesubsection}{0.6em}{}
\titleformat{\subsubsection}{\small\bfseries\sffamily}{\thesubsubsection}{0.6em}{}
\titleformat{\paragraph}{\small\bfseries}{\theparagraph}{0.6em}{}
\titlespacing*{\section}{0pt}{0.7em}{0.32em}
\titlespacing*{\subsection}{0pt}{0.45em}{0.18em}
\titlespacing*{\subsubsection}{0pt}{0.32em}{0.12em}
\titlespacing*{\paragraph}{0pt}{0.22em}{0.6em}

% 避免标题孤零零落在页底，减少打印时的断裂感。
\pretocmd{\section}{\Needspace{5\baselineskip}}{}{}
\pretocmd{\subsection}{\Needspace{3\baselineskip}}{}{}
\pretocmd{\subsubsection}{\Needspace{2\baselineskip}}{}{}

% 页眉页脚：页眉左侧显示当前章节，右侧显示课程名；页脚居中显示页码。
\pagestyle{fancy}
\fancyhf{}
\fancyhead[L]{\small\nouppercase{\leftmark}}
\fancyhead[R]{\small 心理危机干预与预防}
\fancyfoot[C]{\small 第 \thepage 页}
\renewcommand{\headrulewidth}{0.25pt}
\renewcommand{\footrulewidth}{0pt}

% 目录压缩：保留真实页码，同时减少目录本身占用页数。
\setcounter{tocdepth}{4}
\renewcommand{\contentsname}{打印速查目录}
\setlength{\cftbeforesecskip}{0.08em}
\setlength{\cftbeforesubsecskip}{0.02em}
\setlength{\cftbeforesubsubsecskip}{0pt}
\renewcommand{\cftsecfont}{\small\bfseries}
\renewcommand{\cftsubsecfont}{\footnotesize}
\renewcommand{\cftsubsubsecfont}{\scriptsize}
\renewcommand{\cftparafont}{\scriptsize}
\renewcommand{\cftsecpagefont}{\small\bfseries}
\renewcommand{\cftsubsecpagefont}{\footnotesize}
\renewcommand{\cftsubsubsecpagefont}{\scriptsize}
\renewcommand{\cftparapagefont}{\scriptsize}

% 让表格和代码块尽量不要把正文撑得过宽。
\setlength{\emergencystretch}{3em}
\sloppy
\usepackage{bookmark}
\IfFileExists{xurl.sty}{\usepackage{xurl}}{} % add URL line breaks if available
\urlstyle{same}
\hypersetup{
  pdftitle={心理危机干预与预防开卷考试速查版},
  pdflang={zh-CN},
  colorlinks=true,
  linkcolor={Maroon},
  filecolor={Maroon},
  citecolor={Blue},
  urlcolor={Blue},
  pdfcreator={LaTeX via pandoc}}

\title{心理危机干预与预防开卷考试速查版}
\author{}
\date{2026-04-21}

\begin{document}
\maketitle

{
\hypersetup{linkcolor=}
\setcounter{tocdepth}{4}
\tableofcontents
}
\begin{quote}
说明：本 PDF 由 \texttt{心理危机干预与预防/Notes/} 下 7 份 Markdown
笔记合并生成。目录页码由 XeLaTeX 编译得到，适合打印后按页码查找。
\end{quote}

\section{第1讲：心理危机导论}\label{ux7b2c1ux8bb2ux5fc3ux7406ux5371ux673aux5bfcux8bba}

\begin{quote}
说明：以下笔记严格按老师讲课推进顺序整理，只整理到你这次发来的转写片段末尾。个别专有名词在人名转写中前后写法不一致时，按同一案例处理；涉及新闻、统计或个案处，以下均按\textbf{老师课堂表述}记录。
\end{quote}

\begin{center}\rule{0.5\linewidth}{0.5pt}\end{center}

\subsection{一、课程导入：这门课为什么重要}\label{ux4e00ux8bfeux7a0bux5bfcux5165ux8fd9ux95e8ux8bfeux4e3aux4ec0ux4e48ux91cdux8981}

老师开场说明：今天开始讲一门新课------\textbf{心理危机干预与预防}（\textbf{psychological
crisis intervention and prevention}）。

老师先用很口语化、很直接的方式指出：
大学生的心理危机在极端情况下，常表现为两类后果：

\begin{enumerate}
\def\labelenumi{\arabic{enumi}.}
\tightlist
\item
  \textbf{自杀（suicide）}：伤害自己、结束生命。
\item
  \textbf{伤人/杀人（homicide / violence toward
  others）}：把危机转向他人。
\end{enumerate}

虽然这类事件总体上并不常见，但一旦发生，社会影响非常恶劣。
因此，这门课不只是完成一门课程，更重要的是：

\begin{itemize}
\tightlist
\item
  学会\textbf{识别（identification）}
\item
  学会\textbf{调整（adjustment）}
\item
  学会\textbf{助人自助（help others and help oneself）}
\end{itemize}

核心目的还是``健康平安''。

\begin{center}\rule{0.5\linewidth}{0.5pt}\end{center}

\subsection{二、任课教师背景与校内外求助资源}\label{ux4e8cux4efbux8bfeux6559ux5e08ux80ccux666fux4e0eux6821ux5185ux5916ux6c42ux52a9ux8d44ux6e90}

老师介绍：任课教师来自学校\textbf{精神卫生与心理学系}，其中一部分老师同时也是执业医师，在大学医院从事\textbf{心理诊疗（psychological
counseling / clinical psychological service）}工作。

\subsubsection{【求助信息】}\label{ux6c42ux52a9ux4fe1ux606f}

老师特别提到毛老师的门诊时间，提醒大家记住，以备自己或亲友将来需要时使用：

\begin{itemize}
\tightlist
\item
  \textbf{每周四上午}：医科大学总医院 \textbf{临床心理科（clinical
  psychology）}
\item
  \textbf{每周六下午}：一大二院
\end{itemize}

老师提示：如果以后自己或亲友出现如下情况，可以尽早就诊、加号咨询：

\begin{itemize}
\tightlist
\item
  睡不着觉
\item
  焦虑、失眠（anxiety, insomnia）
\item
  抑郁（depression）
\end{itemize}

\begin{center}\rule{0.5\linewidth}{0.5pt}\end{center}

\subsection{三、老师展示个人科普文章：大学生抑郁症要重在早识别、早预防}\label{ux4e09ux8001ux5e08ux5c55ux793aux4e2aux4ebaux79d1ux666eux6587ux7ae0ux5927ux5b66ux751fux6291ux90c1ux75c7ux8981ux91cdux5728ux65e9ux8bc6ux522bux65e9ux9884ux9632}

老师提到左侧一本杂志，是中国心理卫生协会主办的科普杂志《心理与健康》。
其中某年（老师说为\textbf{2023年5月}）邀请其撰写封面文章，主题是：

\textbf{大学生抑郁症早期识别与预防}（early identification and prevention
of depression in college students）

\subsubsection{老师强调的背景判断}\label{ux8001ux5e08ux5f3aux8c03ux7684ux80ccux666fux5224ux65ad}

疫情以后，大学生\textbf{抑郁症（depression）}已经成为大学生\textbf{休学、退学}的重要原因，甚至是主要原因之一。

老师回顾说：

\begin{itemize}
\tightlist
\item
  过去学生休学退学原因中，较常见的是传染病、外伤等；
\item
  现在心理问题尤其突出；
\item
  抑郁症一旦发展起来，\textbf{不容易恢复}，因此更应强调\textbf{预防（prevention）}。
\end{itemize}

\subsubsection{【学习建议】}\label{ux5b66ux4e60ux5efaux8bae}

凡是``不好治''的病，都要强调预防。
对抑郁症而言，关键不是等严重了再处理，而是：

\begin{itemize}
\tightlist
\item
  稍有苗头就要注意；
\item
  稍有症状就应及时发现；
\item
  及时调整，避免进一步加重。
\end{itemize}

\begin{center}\rule{0.5\linewidth}{0.5pt}\end{center}

\subsection{四、政策背景：健康第一，尤其要重视``心理脆弱''}\label{ux56dbux653fux7b56ux80ccux666fux5065ux5eb7ux7b2cux4e00ux5c24ux5176ux8981ux91cdux89c6ux5fc3ux7406ux8106ux5f31}

老师提到，国家层面已明确强调要重视\textbf{心理健康（mental
health）}与\textbf{精神卫生（mental hygiene / mental health service）}。

并提到教育系统当前强调``\textbf{健康第一}''，不是``学习第一''。

\subsubsection{老师提到的``三小''问题}\label{ux8001ux5e08ux63d0ux5230ux7684ux4e09ux5c0fux95eeux9898}

当前学生群体，尤其是中小学生中，突出有``三小''：

\begin{enumerate}
\def\labelenumi{\arabic{enumi}.}
\tightlist
\item
  \textbf{小眼镜}：近视问题突出；
\item
  \textbf{小胖墩儿}：肥胖问题增加；
\item
  \textbf{小脆心儿}：心理脆弱、抗挫能力差。
\end{enumerate}

老师的意思是：
过去很多家长、学生把``学习''放在第一位，甚至牺牲健康去换成绩；现在越来越清楚，\textbf{健康没了，学习最后也守不住}。所以必须把``健康第一''摆到前面。

\begin{center}\rule{0.5\linewidth}{0.5pt}\end{center}

\subsection{五、从``危机（crisis）''讲起：危机的广义概念}\label{ux4e94ux4eceux5371ux673acrisisux8bb2ux8d77ux5371ux673aux7684ux5e7fux4e49ux6982ux5ff5}

老师先从一般性的``危机''讲起。
危机可理解为：某种事件使社会偏离正常运行轨道，对公共安全、稳定造成重大影响的一种\textbf{非均衡状态（state
of disequilibrium）}。

\subsubsection{危机按来源大致可分三类}\label{ux5371ux673aux6309ux6765ux6e90ux5927ux81f4ux53efux5206ux4e09ux7c7b}

\begin{enumerate}
\def\labelenumi{\arabic{enumi}.}
\tightlist
\item
  \textbf{天灾}：如地震、洪水、山火等；
\item
  \textbf{人祸}：如战争；
\item
  \textbf{自然因素与人为因素叠加}：人为因素加重自然危机。
\end{enumerate}

老师这里举了若干国际和自然灾害新闻作为引入，目的不是让大家背新闻本身，而是让大家理解：
\textbf{危机状态意味着安全感被打破，日常秩序被中断，人会进入高度不确定与高压力状态。}

\subsubsection{老师的引申}\label{ux8001ux5e08ux7684ux5f15ux7533}

大家能坐在教室里正常上课、没有寒冷和生命危险，其实已经是一种相对稳定、安全的状态。
而在地球另一些地方，人可能随时面对死亡威胁，这就是危机情境。

\begin{center}\rule{0.5\linewidth}{0.5pt}\end{center}

\subsection{六、【考务】本课考试形式与平时成绩}\label{ux516dux8003ux52a1ux672cux8bfeux8003ux8bd5ux5f62ux5f0fux4e0eux5e73ux65f6ux6210ux7ee9}

老师在这里插入了课程考核说明。

\subsubsection{【考务】}\label{ux8003ux52a1}

\begin{itemize}
\tightlist
\item
  \textbf{期末为开卷考试（open-book examination）}
\item
  但\textbf{不能互相抄袭}
\item
  考试时可以看自己准备的材料、自己的笔记
\item
  不能看别人的试卷
\end{itemize}

\subsubsection{【考务】平时成绩}\label{ux8003ux52a1ux5e73ux65f6ux6210ux7ee9}

\begin{itemize}
\item
  大约有\textbf{30～40分平时成绩}
\item
  具体形式视任课老师习惯而定，可能包括：

  \begin{itemize}
  \tightlist
  \item
    随堂测验
  \item
    签到
  \item
    课堂表现等
  \end{itemize}
\end{itemize}

\subsubsection{【备考建议】}\label{ux5907ux8003ux5efaux8bae}

总原则是：

\begin{itemize}
\tightlist
\item
  平时成绩为辅；
\item
  期末开卷考试为主；
\item
  重点是\textbf{自己做好笔记}；
\item
  考前要做基本复习，不要完全空着去考。
\end{itemize}

\subsubsection{老师说通常什么情况会不及格}\label{ux8001ux5e08ux8bf4ux901aux5e38ux4ec0ux4e48ux60c5ux51b5ux4f1aux4e0dux53caux683c}

\begin{enumerate}
\def\labelenumi{\arabic{enumi}.}
\tightlist
\item
  经常点名不在；
\item
  平时成绩丢得太多；
\item
  期末卷面基本空白。
\end{enumerate}

\subsubsection{【考点提示】}\label{ux8003ux70b9ux63d0ux793a}

老师特别提醒：
凡是\textbf{考试可能涉及的内容}，毛老师通常会在课件或板书中画\textbf{红三角}。
同学们把这些重点拍下来、记下来即可；实在懒的话，课件下课可由学习委员拷走。

\begin{center}\rule{0.5\linewidth}{0.5pt}\end{center}

\subsection{七、【考点】心理危机（psychological
crisis）的定义}\label{ux4e03ux8003ux70b9ux5fc3ux7406ux5371ux673apsychological-crisisux7684ux5b9aux4e49}

老师明确说：\textbf{今天第一个重点，就是心理危机的概念。}

\subsubsection{【考点】心理危机定义}\label{ux8003ux70b9ux5fc3ux7406ux5371ux673aux5b9aux4e49}

\textbf{心理危机（psychological crisis）}：
指个体面临突然的或重大的生活困境，\textbf{无法用既有方式解决}，从而出现\textbf{心理失衡（psychological
disequilibrium）}的状态。

\subsubsection{【考点】心理危机成立的三条标准}\label{ux8003ux70b9ux5fc3ux7406ux5371ux673aux6210ux7acbux7684ux4e09ux6761ux6807ux51c6}

老师说，判断心理危机，要同时抓住三条：

\paragraph{1. 有一个``事儿''------即存在生活事件（life
event）}\label{ux6709ux4e00ux4e2aux4e8bux513fux5373ux5b58ux5728ux751fux6d3bux4e8bux4ef6life-event}

必须先有一个外部或内部触发的重大事件、困境或打击。

\paragraph{2. 出现严重反应}\label{ux51faux73b0ux4e25ux91cdux53cdux5e94}

包括多方面的严重心身反应：

\begin{itemize}
\tightlist
\item
  \textbf{情绪反应（emotional response）}
\item
  \textbf{认知反应（cognitive response）}
\item
  \textbf{身体反应（physical / somatic response）}
\item
  \textbf{行为反应（behavioral response）}
\end{itemize}

\paragraph{3.
不能有效自我缓解}\label{ux4e0dux80fdux6709ux6548ux81eaux6211ux7f13ux89e3}

也就是：

\begin{itemize}
\tightlist
\item
  不能靠现有经验、办法、知识把问题解决掉；
\item
  ``老办法解决不了新问题''；
\item
  已超出个体固有能力边界。
\end{itemize}

\subsubsection{老师总结式表述}\label{ux8001ux5e08ux603bux7ed3ux5f0fux8868ux8ff0}

心理危机本质上就是： \textbf{有重大事件 + 有严重反应 + 自己缓不过来。}

\begin{center}\rule{0.5\linewidth}{0.5pt}\end{center}

\subsection{八、大学生常见心理危机的极端形式：以自杀为主}\label{ux516bux5927ux5b66ux751fux5e38ux89c1ux5fc3ux7406ux5371ux673aux7684ux6781ux7aefux5f62ux5f0fux4ee5ux81eaux6740ux4e3aux4e3b}

老师再次回到大学生群体，指出大学生常见心理危机的极端形式之一就是\textbf{自杀（suicide）}。

\subsubsection{老师对本校与兄弟院校的比较性评论}\label{ux8001ux5e08ux5bf9ux672cux6821ux4e0eux5144ux5f1fux9662ux6821ux7684ux6bd4ux8f83ux6027ux8bc4ux8bba}

老师说，近几年本校没有自杀成功案例，是比较幸运的。
但其他高校并不都如此，尤其一些学校学生多、竞争强、卷得厉害，相关问题更突出。

\subsubsection{老师的解释}\label{ux8001ux5e08ux7684ux89e3ux91ca}

并不是医学生``天生更健康''，而是：

\begin{itemize}
\tightlist
\item
  学校健康教育抓得相对较好；
\item
  医学生多学了一些健康知识；
\item
  更有利于自我识别和自我保护。
\end{itemize}

\begin{center}\rule{0.5\linewidth}{0.5pt}\end{center}

\subsection{九、【举例 /
临床关联】天津师范大学女大学生因乙肝``大三阳''被孤立后自杀案例}\label{ux4e5dux4e3eux4f8b-ux4e34ux5e8aux5173ux8054ux5929ux6d25ux5e08ux8303ux5927ux5b66ux5973ux5927ux5b66ux751fux56e0ux4e59ux809dux5927ux4e09ux9633ux88abux5b64ux7acbux540eux81eaux6740ux6848ux4f8b}

老师说，单讲抽象概念大家容易困，所以开始举案例。
他说明：许多真实案例涉及工作机密，不能讲，因此选了一个网上公开案例。

\subsubsection{案例基本情况}\label{ux6848ux4f8bux57faux672cux60c5ux51b5}

\begin{itemize}
\tightlist
\item
  学校：天津师范大学
\item
  对象：一名大一女生
\item
  转写中名字前后有异写，这里按同一案例整理
\end{itemize}

\subsubsection{事件经过}\label{ux4e8bux4ef6ux7ecfux8fc7}

该女生入学后，在一次偶然检查中发现自己是\textbf{乙肝大三阳}。
老师指出：她本人可能当时并无明显症状，但被检查出有病毒、且具有传染性后，同学开始疏远、孤立她。

\subsubsection{【医学相关】}\label{ux533bux5b66ux76f8ux5173}

老师这里提到：

\begin{itemize}
\tightlist
\item
  乙肝大三阳与\textbf{乙型肝炎（hepatitis B）}相关；
\item
  传播主要与\textbf{体液传播（body fluid transmission）}有关；
\item
  并不是像呼吸道传染病那样通过日常呼吸轻易传播。
\end{itemize}

老师的意思是：
真正传播风险未必像大家想象得那样大，但因为同学不懂医学知识，所以恐惧、疏远仍然发生了。

\subsubsection{被孤立后的变化}\label{ux88abux5b64ux7acbux540eux7684ux53d8ux5316}

\begin{itemize}
\tightlist
\item
  吃饭、上课、生活中被刻意拉开距离；
\item
  逐渐被同学群体排斥；
\item
  心理压力越来越大。
\end{itemize}

\subsubsection{学校与辅导员的处理}\label{ux5b66ux6821ux4e0eux8f85ux5bfcux5458ux7684ux5904ux7406}

辅导员曾劝其先休学治病。 但该女生认为：

\begin{enumerate}
\def\labelenumi{\arabic{enumi}.}
\tightlist
\item
  自己没症状；
\item
  大三阳不是短时间就能治好的；
\item
  家里贫困，休学意味着拖延毕业、增加经济负担；
\item
  患病不能成为剥夺其受教育权的理由。
\end{enumerate}

\subsubsection{学校后续安排}\label{ux5b66ux6821ux540eux7eedux5b89ux6392}

在宿舍紧张情况下，学校为``照顾''她，给了一个\textbf{单间宿舍}。

老师指出：
表面上看，像是``优待''；但实际上，这种安排也可能进一步加重其\textbf{隔离感、孤独感}。

\subsubsection{结局}\label{ux7ed3ux5c40}

她于3月入住单间，一个月后在宿舍内\textbf{烧炭自杀}。

\subsubsection{遗书内容与心理状态}\label{ux9057ux4e66ux5185ux5bb9ux4e0eux5fc3ux7406ux72b6ux6001}

老师转述其遗书意思：

\begin{itemize}
\tightlist
\item
  对未来彻底绝望；
\item
  感到人生漫长而看不到尽头；
\item
  想从当前``腐坏''的生活中尽快逃离；
\item
  也知道``不应效仿''，但自己已救不了自己。
\end{itemize}

\subsubsection{老师对其心理逻辑的概括}\label{ux8001ux5e08ux5bf9ux5176ux5fc3ux7406ux903bux8f91ux7684ux6982ux62ec}

她感到：

\begin{itemize}
\tightlist
\item
  因疾病被孤立；
\item
  眼前看不到希望；
\item
  未来读书、就业、社会交往似乎都会继续遭受歧视；
\item
  于是选择``尽快逃离''。
\end{itemize}

\subsubsection{【易错 / 识别点】}\label{ux6613ux9519-ux8bc6ux522bux70b9}

老师补充说：

\begin{itemize}
\tightlist
\item
  她不是没有努力过，也尝试向朋友解释疾病并不可怕；
\item
  但新生群体对传染病不了解，恐惧和流言仍在扩散；
\item
  自杀当天其母亲已有不祥预感，曾联系辅导员；
\item
  但因误判她人在图书馆，错过了抢救关键时机。
\end{itemize}

\subsubsection{引发的争议}\label{ux5f15ux53d1ux7684ux4e89ux8bae}

事件后来引起很大争议，焦点包括：

\begin{itemize}
\tightlist
\item
  患有传染病的学生应如何安排；
\item
  学校在保障其权利与保护其他学生之间如何平衡；
\item
  辅导员、学校是否已尽责；
\item
  同学的排斥与恐惧在其中起了多大作用。
\end{itemize}

\subsubsection{老师的评价}\label{ux8001ux5e08ux7684ux8bc4ux4ef7}

老师认为：

\begin{itemize}
\tightlist
\item
  这是一个由疾病引发、再被孤立放大的心理危机；
\item
  其中有老师和学校的责任；
\item
  但客观说，学校也并非完全不作为，而是提供了某种支持；
\item
  \textbf{真正困难的是：心理危机的预防并不容易。}
\end{itemize}

\subsubsection{【助人者视角】}\label{ux52a9ux4ebaux8005ux89c6ux89d2}

老师引用该校心理中心主任在事件后的朋友圈感慨，大意是：
自杀像地震，很多征兆往往只有事后追溯时才显得具有预测意义。
每个生命的逝去，都会让助人者感到悲伤、内疚、悔恨，以及无奈、无助。

\begin{center}\rule{0.5\linewidth}{0.5pt}\end{center}

\subsection{十、【考点】危机干预（crisis
intervention）的概念、目的与意义}\label{ux5341ux8003ux70b9ux5371ux673aux5e72ux9884crisis-interventionux7684ux6982ux5ff5ux76eeux7684ux4e0eux610fux4e49}

讲完案例后，老师转入理论。

\subsubsection{【考点】危机干预的提出}\label{ux8003ux70b9ux5371ux673aux5e72ux9884ux7684ux63d0ux51fa}

老师说，\textbf{危机干预（crisis
intervention）}概念最早由美国心理学家\textbf{凯普兰（Caplan）}于\textbf{1954年}提出，并进行系统研究。

\subsubsection{【考点】危机干预的定义}\label{ux8003ux70b9ux5371ux673aux5e72ux9884ux7684ux5b9aux4e49}

危机干预属于广义的\textbf{心理治疗（psychotherapy）}范畴。
其核心是运用心理治疗的手段，帮助处于危机状态的当事人：

\begin{itemize}
\tightlist
\item
  处理迫在眉睫的问题；
\item
  恢复心理平衡；
\item
  安全度过危机；
\item
  重新适应生活。
\end{itemize}

\subsubsection{【考点】危机干预的两个主要目的}\label{ux8003ux70b9ux5371ux673aux5e72ux9884ux7684ux4e24ux4e2aux4e3bux8981ux76eeux7684}

\begin{enumerate}
\def\labelenumi{\arabic{enumi}.}
\item
  \textbf{避免自伤或伤人}

  \begin{itemize}
  \tightlist
  \item
    防止自杀、自残
  \item
    防止伤及他人
  \end{itemize}
\item
  \textbf{恢复心理平衡与心理动力}

  \begin{itemize}
  \tightlist
  \item
    使人重新``运转起来''
  \item
    不再陷于崩溃状态
  \end{itemize}
\end{enumerate}

\subsubsection{【考点】危机干预成功的意义}\label{ux8003ux70b9ux5371ux673aux5e72ux9884ux6210ux529fux7684ux610fux4e49}

如果危机干预成功，个体可以：

\begin{enumerate}
\def\labelenumi{\arabic{enumi}.}
\tightlist
\item
  更好地把握自己的现状；
\item
  重新认识危机事件；
\item
  学会未来再遇到类似危机时更好的应对策略。
\end{enumerate}

\begin{center}\rule{0.5\linewidth}{0.5pt}\end{center}

\subsection{十一、世界范围内的自杀情况：老师用于引入的流行病学比较}\label{ux5341ux4e00ux4e16ux754cux8303ux56f4ux5185ux7684ux81eaux6740ux60c5ux51b5ux8001ux5e08ux7528ux4e8eux5f15ux5165ux7684ux6d41ux884cux75c5ux5b66ux6bd4ux8f83}

老师接着从全球范围谈自杀问题，作为宏观背景。

\subsubsection{老师课堂中提到的比较}\label{ux8001ux5e08ux8bfeux5802ux4e2dux63d0ux5230ux7684ux6bd4ux8f83}

\begin{itemize}
\tightlist
\item
  自杀率较高的国家：韩国、匈牙利、日本等
\item
  自杀率较低的国家：冰岛、西班牙、希腊等
\item
  自杀是一个\textbf{概率事件}，没有哪个国家是绝对零发生
\end{itemize}

\subsubsection{老师进一步讲的趋势}\label{ux8001ux5e08ux8fdbux4e00ux6b65ux8bb2ux7684ux8d8bux52bf}

从纵向看，西方国家自杀率在历史上有：

\begin{itemize}
\tightlist
\item
  两次低谷：两次世界大战时期
\item
  两次高峰：经济大萧条相关时期
\end{itemize}

老师借此强调一个朴素判断：
\textbf{经济压力、贫困、失业与心理危机密切相关。}

\subsubsection{老师的口语化总结}\label{ux8001ux5e08ux7684ux53e3ux8bedux5316ux603bux7ed3}

``钱不是万能的，但没有钱是万万不能的。'' 他用这个意思强调：
收入和现实生活条件，确实能缓解大量现实压力和心理问题。

\begin{center}\rule{0.5\linewidth}{0.5pt}\end{center}

\subsection{十二、中国大学生自杀率：老师强调``总体不高，但趋势需警惕''}\label{ux5341ux4e8cux4e2dux56fdux5927ux5b66ux751fux81eaux6740ux7387ux8001ux5e08ux5f3aux8c03ux603bux4f53ux4e0dux9ad8ux4f46ux8d8bux52bfux9700ux8b66ux60d5}

老师提到：

\begin{itemize}
\tightlist
\item
  美国总体自杀率居中；
\item
  美国大学生自杀率大致低于其全人群平均；
\item
  中国总体自杀率也有自己的水平；
\item
  中国大学生自杀率相对更低。
\end{itemize}

\subsubsection{老师的解释}\label{ux8001ux5e08ux7684ux89e3ux91ca-1}

中国高校自杀率相对较低，老师将其归因为：

\begin{itemize}
\tightlist
\item
  学校对心理健康工作的高度重视；
\item
  有专门心理老师；
\item
  有辅导员体系；
\item
  对学生有更强的组织性和日常管理。
\end{itemize}

\subsubsection{【复习 / 观点】}\label{ux590dux4e60-ux89c2ux70b9}

老师顺带比较了中国大学与国外大学的学生管理差异：

\begin{itemize}
\tightlist
\item
  国外大学多是个人化、分散式管理；
\item
  没有中国式的班级、班委、辅导员等强集体支持系统；
\item
  中国高校的班集体、宿舍集体，在心理危机预防上反而是一种保障。
\end{itemize}

老师要表达的是：
\textbf{一个人出问题时，若身处集体，更容易被发现、被支持。}

\begin{center}\rule{0.5\linewidth}{0.5pt}\end{center}

\subsection{十三、天津市学生自杀的趋势：疫情后出现低龄化}\label{ux5341ux4e09ux5929ux6d25ux5e02ux5b66ux751fux81eaux6740ux7684ux8d8bux52bfux75abux60c5ux540eux51faux73b0ux4f4eux9f84ux5316}

老师提到，天津市每年具体自杀人数属于保密信息，但他给出了几个趋势判断：

\subsubsection{趋势}\label{ux8d8bux52bf}

\begin{enumerate}
\def\labelenumi{\arabic{enumi}.}
\item
  \textbf{总体逐年增长}
\item
  \textbf{疫情后结构变化明显}

  \begin{itemize}
  \tightlist
  \item
    疫情前：非正常死亡学生中，研究生和本科生占比更高
  \item
    疫情后：中小学生占比明显上升
  \end{itemize}
\item
  \textbf{低龄化}

  \begin{itemize}
  \tightlist
  \item
    老师提到最小案例已到小学四年级
  \end{itemize}
\end{enumerate}

\subsubsection{老师结论}\label{ux8001ux5e08ux7ed3ux8bba}

当前形势不容乐观。

\begin{center}\rule{0.5\linewidth}{0.5pt}\end{center}

\subsection{十四、北京某调查总结的大学生自杀规律}\label{ux5341ux56dbux5317ux4eacux67d0ux8c03ux67e5ux603bux7ed3ux7684ux5927ux5b66ux751fux81eaux6740ux89c4ux5f8b}

老师说，最新数据未必掌握，但举了一个较早的北京研究作参考。
该研究收集了\textbf{200例大学生自杀相关个案}，从中总结了一些规律。

\subsubsection{1）自杀地点分布}\label{ux81eaux6740ux5730ux70b9ux5206ux5e03}

老师提到，地点以以下场所为主：

\begin{itemize}
\tightlist
\item
  \textbf{宿舍}最多
\item
  其次是\textbf{教学楼}
\item
  还有\textbf{科研楼、家属楼}等
\end{itemize}

\subsubsection{解释}\label{ux89e3ux91ca}

说明自杀行为常发生在学生\textbf{最熟悉、最常活动}的场所。

\subsubsection{2）跳楼楼层分布}\label{ux8df3ux697cux697cux5c42ux5206ux5e03}

老师说，自杀跳楼楼层较集中在：

\begin{itemize}
\tightlist
\item
  5楼
\item
  6楼
\item
  7楼
\item
  8楼
\end{itemize}

其中\textbf{7楼最多}。

老师用建筑标准与致死概率作口语化解释：

\begin{itemize}
\tightlist
\item
  太低可能摔不死；
\item
  七层楼又常是学校建筑的典型高度。
\end{itemize}

\subsubsection{3）最主要原因：不适应（maladjustment）}\label{ux6700ux4e3bux8981ux539fux56e0ux4e0dux9002ux5e94maladjustment}

老师强调，``不适应''是非常核心的原因，而且是\textbf{各种不适应叠加}：

\begin{itemize}
\tightlist
\item
  学业不适应
\item
  人际不适应
\item
  气候不适应
\item
  饮食不适应
\item
  环境不适应
\item
  专业不适应
\end{itemize}

\paragraph{【举例】}\label{ux4e3eux4f8b}

\begin{itemize}
\tightlist
\item
  从小学到初中、高中到大学，角色变化导致失落；
\item
  原来在原学校很优秀，换环境后优势消失；
\item
  第一本志愿没录上，被调剂到不喜欢的专业；
\item
  想退学重考，又没勇气，长期纠结。
\end{itemize}

\subsubsection{4）抑郁（depression）}\label{ux6291ux90c1depression}

老师说，很多学生已经处于抑郁状态，但：

\begin{itemize}
\tightlist
\item
  自己不知道；
\item
  或者不承认；
\item
  硬撑、逞强，最后越来越严重。
\end{itemize}

\subsubsection{5）危机发生时间}\label{ux5371ux673aux53d1ux751fux65f6ux95f4}

老师强调一个容易与直觉不符的点：

\begin{itemize}
\tightlist
\item
  大家往往以为夜里更容易出事；
\item
  但调查提示：\textbf{白天比晚上多}
\item
  尤其是\textbf{上午}更容易发生。
\end{itemize}

\subsubsection{6）自杀方式}\label{ux81eaux6740ux65b9ux5f0f}

最显著高于其他方式的是：\textbf{跳楼}

老师给的解释是：

\begin{itemize}
\tightlist
\item
  不需要额外工具；
\item
  致死率高；
\item
  操作相对``简单''。
\end{itemize}

\subsubsection{7）年龄与年级分布}\label{ux5e74ux9f84ux4e0eux5e74ux7ea7ux5206ux5e03}

老师指出：

\begin{itemize}
\tightlist
\item
  中小学阶段：高年级更突出
\item
  大学阶段：\textbf{随着年级升高，危机数量逐年增加}
\item
  \textbf{毕业年级最高发}
\end{itemize}

\subsubsection{老师对毕业年级高发的解释}\label{ux8001ux5e08ux5bf9ux6bd5ux4e1aux5e74ux7ea7ux9ad8ux53d1ux7684ux89e3ux91ca}

毕业班问题交织最复杂：

\begin{itemize}
\tightlist
\item
  能否毕业
\item
  找工作困难
\item
  考研压力
\item
  答辩、盲审压力
\item
  恋爱与分手问题
\item
  多重现实压力集中爆发
\end{itemize}

\subsubsection{8）性别差异}\label{ux6027ux522bux5deeux5f02}

老师说，北京这项研究提示：

\begin{itemize}
\item
  男女比例基本接近 \textbf{1:1}
\item
  自杀方式差异也不大
\item
  只是时间选择上略有不同：

  \begin{itemize}
  \tightlist
  \item
    男生夜间稍多
  \item
    女生白天稍多
  \end{itemize}
\end{itemize}

\begin{center}\rule{0.5\linewidth}{0.5pt}\end{center}

\subsection{十五、费立鹏关于中国人自杀原因的总结}\label{ux5341ux4e94ux8d39ux7acbux9e4fux5173ux4e8eux4e2dux56fdux4ebaux81eaux6740ux539fux56e0ux7684ux603bux7ed3}

老师接着提到加拿大心理学家\textbf{费立鹏（Michael R.
Phillips）}在中国开展研究，并总结出中国人自杀的八大相关因素。

\subsubsection{八大原因}\label{ux516bux5927ux539fux56e0}

\begin{enumerate}
\def\labelenumi{\arabic{enumi}.}
\tightlist
\item
  \textbf{抑郁程度重}
\item
  \textbf{有自杀未遂史}
\item
  \textbf{应激强度大（high stress intensity）}
\item
  \textbf{生命质量低}
\item
  \textbf{慢性心理压力大}
\item
  \textbf{死亡前两天内有人际冲突}
\item
  \textbf{有血缘关系的人存在自杀行为}
\item
  \textbf{熟人有自杀行为}
\end{enumerate}

\subsubsection{【考点理解】}\label{ux8003ux70b9ux7406ux89e3}

老师特别强调第6条很有``中国特色''：
中国人的自杀往往与\textbf{亲密关系中的激烈冲突}有关，很多是``吵起来了''\,``怼起来了''之后的\textbf{激情性自杀}。

\begin{center}\rule{0.5\linewidth}{0.5pt}\end{center}

\subsection{十六、中国人与国外自杀模式的差异}\label{ux5341ux516dux4e2dux56fdux4ebaux4e0eux56fdux5916ux81eaux6740ux6a21ux5f0fux7684ux5deeux5f02}

老师根据费立鹏研究，讲了几组对比。

\subsubsection{1）性别差异不同}\label{ux6027ux522bux5deeux5f02ux4e0dux540c}

老师的课堂表述是：

\begin{itemize}
\tightlist
\item
  \textbf{中国}：女性自杀多于男性
\item
  \textbf{国外}：男性自杀多于女性
\end{itemize}

老师解释认为： 中国总体数据受\textbf{农村女性自杀率较高}影响很大。

\subsubsection{2）老师对农村女性高自杀率的解释}\label{ux8001ux5e08ux5bf9ux519cux6751ux5973ux6027ux9ad8ux81eaux6740ux7387ux7684ux89e3ux91ca}

他认为与以下因素有关：

\begin{itemize}
\tightlist
\item
  女性社会地位相对较低；
\item
  封建观念残余；
\item
  家庭矛盾无法有效解决；
\item
  遇到不公和挫折时缺乏支持与出路。
\end{itemize}

\subsubsection{【举例】}\label{ux4e3eux4f8b-1}

老师以\textbf{彩礼}为例，从其角度批评这是一种对女性的物化和歧视。
并指出在某些家庭情境中，女性在婚后面临：

\begin{itemize}
\tightlist
\item
  婆媳矛盾
\item
  夫妻冲突
\item
  娘家、婆家双重压力
\item
  离也离不了，过也过不下去
\end{itemize}

最后可能走向极端。

\subsubsection{3）年龄差异不同}\label{ux5e74ux9f84ux5deeux5f02ux4e0dux540c}

老师说：

\begin{itemize}
\tightlist
\item
  国外自杀年龄偏大；
\item
  中国青少年、自青年龄段更突出；
\item
  中国15～24岁是高风险年龄段之一；
\item
  平均自杀年龄偏年轻。
\end{itemize}

\subsubsection{老师的解释}\label{ux8001ux5e08ux7684ux89e3ux91ca-2}

年轻人之所以更容易自杀，是因为：

\begin{itemize}
\tightlist
\item
  年龄阶段本身更敏感；
\item
  长期读书，理想主义色彩强；
\item
  接触现实后，发现现实与想象不一致；
\item
  本该``调整主观去适应客观''，但有些年轻人做不到，容易走向极端。
\end{itemize}

\begin{center}\rule{0.5\linewidth}{0.5pt}\end{center}

\subsection{十七、精神障碍与躯体疾病也是自杀的重要原因}\label{ux5341ux4e03ux7cbeux795eux969cux788dux4e0eux8eafux4f53ux75beux75c5ux4e5fux662fux81eaux6740ux7684ux91cdux8981ux539fux56e0}

老师接着补充，自杀与\textbf{精神障碍（mental
disorders）}、\textbf{躯体疾病（physical illness / somatic
disease）}密切相关。

\subsubsection{1）抑郁症（depression）}\label{ux6291ux90c1ux75c7depression}

老师强调，抑郁症不是简单的``想不开''，而是一种深度的\textbf{精神耗竭（mental
exhaustion）}状态。

典型体验包括：

\begin{itemize}
\tightlist
\item
  持续痛苦
\item
  大脑像被掏空
\item
  转不动
\item
  没动力（loss of motivation）
\item
  没兴趣（anhedonia）
\item
  认知越来越消极
\item
  觉得活着只有痛苦，没有出头之日
\end{itemize}

对正常人来说，``死''是可怕的；
但对重度抑郁患者来说，死亡可能被体验为``唯一解脱''。

\subsubsection{2）精神分裂症（schizophrenia）}\label{ux7cbeux795eux5206ux88c2ux75c7schizophrenia}

老师说，\textbf{精神分裂症（schizophrenia）}患者也可能自杀，因为会受：

\begin{itemize}
\tightlist
\item
  \textbf{幻觉（hallucination）}
\item
  \textbf{妄想（delusion）}
\end{itemize}

支配。
老师还引用电影《追捕》里的情节，帮助理解``被诱导走向死亡''的体验。

\subsubsection{3）严重躯体疾病}\label{ux4e25ux91cdux8eafux4f53ux75beux75c5}

老师提到，一些严重躯体疾病也会增加自杀风险，例如：

\begin{itemize}
\tightlist
\item
  癌症（cancer）
\item
  艾滋病（AIDS / HIV infection）等
\end{itemize}

其相关原因包括：

\begin{itemize}
\tightlist
\item
  疾病痛苦大；
\item
  治疗困难；
\item
  经济负担重；
\item
  心态崩溃；
\item
  常伴发抑郁。
\end{itemize}

\begin{center}\rule{0.5\linewidth}{0.5pt}\end{center}

\subsection{十八、与危机事件相关的常见诱因}\label{ux5341ux516bux4e0eux5371ux673aux4e8bux4ef6ux76f8ux5173ux7684ux5e38ux89c1ux8bf1ux56e0}

老师归纳说，与心理危机相关联的事件，常有以下几类特点。

\subsubsection{1）持续压力（chronic
stress）}\label{ux6301ux7eedux538bux529bchronic-stress}

例如：

\begin{itemize}
\tightlist
\item
  家庭经济困难
\item
  学业困难
\item
  单亲家庭压力
\end{itemize}

\subsubsection{【举例】单亲家庭}\label{ux4e3eux4f8bux5355ux4eb2ux5bb6ux5ead}

老师特别讲到单亲家庭，尤其单身母亲抚养孩子时，容易出现：

\begin{itemize}
\tightlist
\item
  家长把全部希望投射到孩子身上；
\item
  对孩子要求过高；
\item
  孩子一开始努力讨好；
\item
  之后逐渐达不到期待；
\item
  最终矛盾升级。
\end{itemize}

甚至现实中会出现极端的\textbf{弑亲}新闻。

\subsubsection{2）财产、健康、名誉、尊严丧失}\label{ux8d22ux4ea7ux5065ux5eb7ux540dux8a89ux5c0aux4e25ux4e27ux5931}

老师举到：

\begin{itemize}
\tightlist
\item
  网络诈骗（telecom / online fraud）
\item
  健康受损
\item
  名誉受损
\item
  尊严被羞辱
\end{itemize}

\subsubsection{【举例】网络诈骗}\label{ux4e3eux4f8bux7f51ux7edcux8bc8ux9a97}

老师提醒：网络诈骗的两个高发受害群体是：

\begin{itemize}
\tightlist
\item
  老年人
\item
  大学生
\end{itemize}

大学生容易成为目标，是因为：

\begin{itemize}
\tightlist
\item
  社会经验不足；
\item
  容易轻信；
\item
  容易被``兼职''\,``退款''``小便宜''诱骗。
\end{itemize}

\subsubsection{【学习建议】}\label{ux5b66ux4e60ux5efaux8bae-1}

一定要提高警惕。
有的人损失一点钱无所谓，但对另一些人来说，经济打击可能直接把人推入生存危机。

\subsubsection{3）感情受挫}\label{ux611fux60c5ux53d7ux632b}

如：

\begin{itemize}
\tightlist
\item
  失恋
\item
  情感挫败
\item
  情感冲突
\end{itemize}

\subsubsection{4）环境改变（environmental
change）}\label{ux73afux5883ux6539ux53d8environmental-change}

如：

\begin{itemize}
\tightlist
\item
  入学
\item
  转学
\item
  搬家
\item
  留学
\end{itemize}

这些都可能成为危机高发节点。

\subsubsection{5）持续的人际关系紧张}\label{ux6301ux7eedux7684ux4ebaux9645ux5173ux7cfbux7d27ux5f20}

包括：

\begin{itemize}
\tightlist
\item
  夫妻关系
\item
  家庭矛盾
\item
  同学关系
\item
  宿舍关系
\end{itemize}

\subsubsection{老师的总结}\label{ux8001ux5e08ux7684ux603bux7ed3}

中国人的自杀，常有\textbf{明显的人际关系因素}参与。
这是与一些西方社会相对不同的地方。

\begin{center}\rule{0.5\linewidth}{0.5pt}\end{center}

\subsection{十九、人格（personality）与心理危机的关系}\label{ux5341ux4e5dux4ebaux683cpersonalityux4e0eux5fc3ux7406ux5371ux673aux7684ux5173ux7cfb}

老师进一步说，心理危机不只和外部事件有关，也和人格基础有关。

\subsubsection{人格的概念}\label{ux4ebaux683cux7684ux6982ux5ff5}

\textbf{人格（personality）}是一个人心理特征的总和，像一个相对稳定的``脸谱''。

人格是双刃剑：

\begin{itemize}
\tightlist
\item
  有优势的一面；
\item
  也有脆弱、容易出问题的一面。
\end{itemize}

\subsubsection{容易出现心理危机者常见人格特点}\label{ux5bb9ux6613ux51faux73b0ux5fc3ux7406ux5371ux673aux8005ux5e38ux89c1ux4ebaux683cux7279ux70b9}

\begin{enumerate}
\def\labelenumi{\arabic{enumi}.}
\tightlist
\item
  注意力明显缺乏
\item
  过分内省（excessive introspection）
\item
  情绪稳定性差
\item
  解决问题能力差
\end{enumerate}

\subsubsection{另一项研究归纳的心理特点}\label{ux53e6ux4e00ux9879ux7814ux7a76ux5f52ux7eb3ux7684ux5fc3ux7406ux7279ux70b9}

容易出现危机的人，往往还有以下特点：

\begin{itemize}
\item
  不能客观评价自己，也不能客观评价外界；
\item
  容易夸大外部危险；
\item
  对自己评价要么过于自大，要么过度自卑；
\item
  存在多种\textbf{非理性认知（irrational cognition）}：

  \begin{itemize}
  \tightlist
  \item
    过度概括（overgeneralization）
  \item
    非此即彼（all-or-none thinking）
  \item
    以偏概全
  \item
    负性推导
  \item
    过度绝对化、肯定化
  \end{itemize}
\item
  人际关系不良；
\item
  人格不成熟；
\item
  适应能力差；
\item
  不会求助；
\item
  不会应对困难局面。
\end{itemize}

\begin{center}\rule{0.5\linewidth}{0.5pt}\end{center}

\subsection{二十、【举例】两个典型案例：一个杀人，一个自杀}\label{ux4e8cux5341ux4e3eux4f8bux4e24ux4e2aux5178ux578bux6848ux4f8bux4e00ux4e2aux6740ux4ebaux4e00ux4e2aux81eaux6740}

老师说，前面理论有点抽象，接下来用两个典型案例串起来理解：

\begin{enumerate}
\def\labelenumi{\arabic{enumi}.}
\tightlist
\item
  \textbf{马加爵}：杀人案例
\item
  \textbf{杨元元}：自杀案例
\end{enumerate}

\begin{center}\rule{0.5\linewidth}{0.5pt}\end{center}

\subsection{二十一、【案例一】马加爵：从优秀贫困生到宿舍杀人案}\label{ux4e8cux5341ux4e00ux6848ux4f8bux4e00ux9a6cux52a0ux7235ux4eceux4f18ux79c0ux8d2bux56f0ux751fux5230ux5bbfux820dux6740ux4ebaux6848}

\subsubsection{基本情况}\label{ux57faux672cux60c5ux51b5}

\begin{itemize}
\tightlist
\item
  籍贯：广西
\item
  学校：云南大学
\item
  事件：大四最后一个学期前后，杀害4名同学，后逃到三亚，被捕并判死刑
\end{itemize}

\subsubsection{老师先强调的一点}\label{ux8001ux5e08ux5148ux5f3aux8c03ux7684ux4e00ux70b9}

马加爵并不是人们想象中的``青面獠牙''式坏人，从外表看反而非常朴实。

\subsubsection{入学前后的优秀经历}\label{ux5165ux5b66ux524dux540eux7684ux4f18ux79c0ux7ecfux5386}

\begin{itemize}
\tightlist
\item
  从小学到高中成绩都非常优秀；
\item
  初中参加全国数学竞赛获二等奖；
\item
  高中考入云南大学，成为村镇骄傲；
\item
  大一入学后积极报名当班委；
\item
  热心班级事务，参加活动，承担任务，是积极分子。
\end{itemize}

\begin{center}\rule{0.5\linewidth}{0.5pt}\end{center}

\subsection{二十二、【案例二】杨元元：从优秀学生干部到研究生宿舍自杀}\label{ux4e8cux5341ux4e8cux6848ux4f8bux4e8cux6768ux5143ux5143ux4eceux4f18ux79c0ux5b66ux751fux5e72ux90e8ux5230ux7814ux7a76ux751fux5bbfux820dux81eaux6740}

\subsubsection{基本情况}\label{ux57faux672cux60c5ux51b5-1}

\begin{itemize}
\tightlist
\item
  2009级，上海海事大学法学系硕士生
\item
  研一入学后不久，于宿舍卫生间自杀
\end{itemize}

\subsubsection{自杀方式}\label{ux81eaux6740ux65b9ux5f0f-1}

老师提到她是用两条毛巾固定后，坐地将自己勒死。
老师用这个细节强调：这说明其求死意志已经非常坚决，因为过程中理论上存在中断机会，但她没有停止。

\subsubsection{入学前的优秀经历}\label{ux5165ux5b66ux524dux7684ux4f18ux79c0ux7ecfux5386}

\begin{itemize}
\tightlist
\item
  从小学习刻苦，成绩优秀；
\item
  从高中到大学一直是学生干部；
\item
  大学期间加入中国共产党；
\item
  还能帮助同学调解矛盾，是标准的``优秀学生干部''形象。
\end{itemize}

\begin{center}\rule{0.5\linewidth}{0.5pt}\end{center}

\subsection{二十三、两案共同基础之一：贫困（poverty）}\label{ux4e8cux5341ux4e09ux4e24ux6848ux5171ux540cux57faux7840ux4e4bux4e00ux8d2bux56f0poverty}

老师说，自己的一家之言是： 这两个人有一个共同基础------\textbf{穷}。

\subsubsection{马加爵的贫困表现}\label{ux9a6cux52a0ux7235ux7684ux8d2bux56f0ux8868ux73b0}

\begin{itemize}
\tightlist
\item
  吃饭时常只买素菜；
\item
  有时只靠馒头和免费汤；
\item
  家里寄钱晚了就借钱吃饭；
\item
  借不到钱就躺着不动，甚至以睡觉来减少消耗；
\item
  又穷又好面子，不愿申请贫困资助。
\end{itemize}

\subsubsection{杨元元的贫困表现}\label{ux6768ux5143ux5143ux7684ux8d2bux56f0ux8868ux73b0}

\begin{itemize}
\tightlist
\item
  六岁丧父；
\item
  母亲无固定工作，还要照顾弟弟；
\item
  她一边上学，一边承担家庭责任；
\item
  本科毕业时因欠学校贷款，毕业证一度被扣；
\item
  因没钱甚至卖掉手机，与同学联系中断。
\end{itemize}

\begin{center}\rule{0.5\linewidth}{0.5pt}\end{center}

\subsection{二十四、两案共同基础之二：学业受挫}\label{ux4e8cux5341ux56dbux4e24ux6848ux5171ux540cux57faux7840ux4e4bux4e8cux5b66ux4e1aux53d7ux632b}

\subsubsection{马加爵的学业受挫}\label{ux9a6cux52a0ux7235ux7684ux5b66ux4e1aux53d7ux632b}

\begin{itemize}
\tightlist
\item
  大一时非常努力，想拿奖学金；
\item
  但英语基础差，成为短板；
\item
  尽管专业课不错，却因英语拖后腿，多次与奖学金失之交臂；
\item
  原本想通过奖学金同时解决经济和尊严问题，结果两头落空；
\item
  之后开始躺平、松散。
\end{itemize}

\subsubsection{杨元元的学业受挫}\label{ux6768ux5143ux5143ux7684ux5b66ux4e1aux53d7ux632b}

\begin{itemize}
\tightlist
\item
  本科业务成绩第一；
\item
  但综合测评中其他加分因素影响，最终没能保研；
\item
  她对这种结果感到委屈、不公平；
\item
  毕业后又花两年准备考研，但一直没考上；
\item
  最终是在反复受挫中继续往前顶。
\end{itemize}

\begin{center}\rule{0.5\linewidth}{0.5pt}\end{center}

\subsection{二十五、两案共同基础之三：感情受挫}\label{ux4e8cux5341ux4e94ux4e24ux6848ux5171ux540cux57faux7840ux4e4bux4e09ux611fux60c5ux53d7ux632b}

\subsubsection{马加爵}\label{ux9a6cux52a0ux7235}

\begin{itemize}
\tightlist
\item
  喜欢班里一名女同学；
\item
  在男同学鼓动下尝试表白；
\item
  打电话被拒；
\item
  写情书又被当面撕毁；
\item
  受到强烈打击和羞辱感。
\end{itemize}

\subsubsection{杨元元}\label{ux6768ux5143ux5143}

\begin{itemize}
\tightlist
\item
  大学时有男朋友；
\item
  因毕业、考研、压力等因素，关系受挫，最终分手。
\end{itemize}

\begin{center}\rule{0.5\linewidth}{0.5pt}\end{center}

\subsection{二十六、两案爆发前的关键人际冲突}\label{ux4e8cux5341ux516dux4e24ux6848ux7206ux53d1ux524dux7684ux5173ux952eux4ebaux9645ux51b2ux7a81}

\subsubsection{马加爵的爆发点}\label{ux9a6cux52a0ux7235ux7684ux7206ux53d1ux70b9}

老师讲到两个关键事件：

\paragraph{1.
被最看重的朋友排除在外}\label{ux88abux6700ux770bux91cdux7684ux670bux53cbux6392ux9664ux5728ux5916}

他认为某位老乡兼贫困生同学是自己最好的朋友。
但这位同学过生日时请了几乎所有人，唯独没请马加爵。 这件事对他伤害极大。

\paragraph{2.
打牌时被指``作弊''，并被拿``人品差''羞辱}\label{ux6253ux724cux65f6ux88abux6307ux4f5cux5f0aux5e76ux88abux62ffux4ebaux54c1ux5deeux7f9eux8fb1}

另一名同学打牌输了后，不仅说他作弊，还拿前面``生日不请他''这件事嘲讽他，说这是因为他``人品不行''。

老师认为，这一下``捅到了肺管子''，彻底触发了他的杀意。

\subsubsection{之后的发展}\label{ux4e4bux540eux7684ux53d1ux5c55}

\begin{itemize}
\tightlist
\item
  他买了锤子准备报复；
\item
  原本计划针对两人；
\item
  后来另外两名同学``误入现场''，被牵连杀害；
\item
  还有一名曾在他困难时请过他吃饭的同学，被他放过。
\end{itemize}

\subsubsection{【举例中的老师提示】}\label{ux4e3eux4f8bux4e2dux7684ux8001ux5e08ux63d0ux793a}

老师借此强调一句很口语的话： ``勿以善小而不为。''
有时候一个很小的善意，真的可能在关键时刻改变走向。

\begin{center}\rule{0.5\linewidth}{0.5pt}\end{center}

\subsection{二十七、杨元元案中的关键人际冲突：贫困、母亲、宿舍与羞辱感}\label{ux4e8cux5341ux4e03ux6768ux5143ux5143ux6848ux4e2dux7684ux5173ux952eux4ebaux9645ux51b2ux7a81ux8d2bux56f0ux6bcdux4eb2ux5bbfux820dux4e0eux7f9eux8fb1ux611f}

\subsubsection{事件背景}\label{ux4e8bux4ef6ux80ccux666f}

杨元元上大学时曾把母亲带到学校附近生活，学校后勤还曾给母亲安排临时保洁工作。
到了上海读研后，她也想继续这样做，让母亲和自己一起生活。

\subsubsection{冲突升级}\label{ux51b2ux7a81ux5347ux7ea7}

但研究生宿舍管理严格，学校不同意。
母亲曾趁宿管不注意进入宿舍，与她同住。
其他同学反映后，学院领导与宿管介入，要求母亲离开。

\subsubsection{她的主观体验}\label{ux5979ux7684ux4e3bux89c2ux4f53ux9a8c}

\begin{itemize}
\tightlist
\item
  她希望学校能同情她、体谅她；
\item
  但学校采取公事公办态度；
\item
  宿管老师多次驱离其母亲；
\item
  她感到母亲脸面受辱，也感到自己无能；
\item
  与宿管、学校之间的矛盾不断加深。
\end{itemize}

\begin{center}\rule{0.5\linewidth}{0.5pt}\end{center}

\subsection{二十八、老师对两案的心理学分析（一）：欲求的绝对性}\label{ux4e8cux5341ux516bux8001ux5e08ux5bf9ux4e24ux6848ux7684ux5fc3ux7406ux5b66ux5206ux6790ux4e00ux6b32ux6c42ux7684ux7eddux5bf9ux6027}

老师认为，两人一步步走上危机道路，首先与\textbf{欲求的绝对性}有关。

\subsubsection{欲求的绝对性是什么意思}\label{ux6b32ux6c42ux7684ux7eddux5bf9ux6027ux662fux4ec0ux4e48ux610fux601d}

就是对自己、对他人、对事情抱有过于理想化、绝对化的要求。 典型想法包括：

\begin{itemize}
\tightlist
\item
  人都应该讲理；
\item
  事情都应该公平；
\item
  别人应该按我的期待回应我。
\end{itemize}

\subsubsection{在马加爵身上的表现}\label{ux5728ux9a6cux52a0ux7235ux8eabux4e0aux7684ux8868ux73b0}

\begin{itemize}
\tightlist
\item
  我打牌厉害，别人就应该承认并尊重我；
\item
  我把你当朋友，你就必须也把我当朋友；
\item
  如果你没有这样做，就是严重背叛与侮辱。
\end{itemize}

\subsubsection{在杨元元身上的表现}\label{ux5728ux6768ux5143ux5143ux8eabux4e0aux7684ux8868ux73b0}

\begin{itemize}
\tightlist
\item
  坚信``知识改变命运''；
\item
  认为只要努力学习，就应该得到改变命运的结果；
\item
  明明资源不够、条件不够，却执着于一定要考北上广的研究生；
\item
  认为学校理应同情、帮助自己和母亲；
\item
  一旦现实没有按理想运行，就难以接受。
\end{itemize}

\subsubsection{老师结论}\label{ux8001ux5e08ux7ed3ux8bba-1}

表面看像是``别人逼她/逼他''，但深层常常也是\textbf{自己用理想化要求把自己逼到了墙角}。

\begin{center}\rule{0.5\linewidth}{0.5pt}\end{center}

\subsection{二十九、老师对两案的心理学分析（二）：感知的歪曲性}\label{ux4e8cux5341ux4e5dux8001ux5e08ux5bf9ux4e24ux6848ux7684ux5fc3ux7406ux5b66ux5206ux6790ux4e8cux611fux77e5ux7684ux6b6aux66f2ux6027}

第二个特点是\textbf{感知的歪曲性}，即因为贫困、自卑、敏感等基础，在解释外界行为时更容易感到：

\begin{itemize}
\tightlist
\item
  屈辱
\item
  被歧视
\item
  被贬低
\item
  被看不起
\end{itemize}

\subsubsection{马加爵身上的例子}\label{ux9a6cux52a0ux7235ux8eabux4e0aux7684ux4f8bux5b50}

老师举了若干例子：

\begin{itemize}
\tightlist
\item
  请老乡吃饭时，饭馆老板更重视老乡，不把他当主客；
\item
  去商场买鞋，被店员默认只能买最便宜的；
\item
  去银行等候，被保安当成蹭空调的人往外赶；
\item
  当班委时因常去找团支书汇报工作，被男同学嘲笑想``癞蛤蟆吃天鹅肉''。
\end{itemize}

老师的意思是：
这些事在一般人看来未必都是``致命打击''，但对一个本就高敏感、高自尊的人来说，会不断叠加成屈辱记忆。

\begin{center}\rule{0.5\linewidth}{0.5pt}\end{center}

\subsection{三十、老师对两案的心理学分析（三）：思维的消极性}\label{ux4e09ux5341ux8001ux5e08ux5bf9ux4e24ux6848ux7684ux5fc3ux7406ux5b66ux5206ux6790ux4e09ux601dux7ef4ux7684ux6d88ux6781ux6027}

老师说，他们在受挫后，思维会迅速走向消极和极端。

\subsubsection{表现}\label{ux8868ux73b0}

\begin{itemize}
\tightlist
\item
  对自己、对他人、对事情都越来越悲观；
\item
  一旦一件事没达到预期，就容易全面否定；
\item
  从``很理想主义''一下变成``彻底否定现实''。
\end{itemize}

\subsubsection{例子}\label{ux4f8bux5b50}

\begin{itemize}
\tightlist
\item
  马加爵奖学金拿不到后，开始躺平、逃课、上网；
\item
  杨元元长期坚持读书考研，发现并未改变命运后，开始彻底怀疑自己全部努力的意义。
\end{itemize}

\begin{center}\rule{0.5\linewidth}{0.5pt}\end{center}

\subsection{三十一、老师对两案的心理学分析（四）：应对的非理性}\label{ux4e09ux5341ux4e00ux8001ux5e08ux5bf9ux4e24ux6848ux7684ux5fc3ux7406ux5b66ux5206ux6790ux56dbux5e94ux5bf9ux7684ux975eux7406ux6027}

第四个问题是\textbf{应对非理性（irrational coping）}。
即在面对压力和挫折时，不会调整路线，不会换方法，不会求助，而是用极端方式解决问题。

\subsubsection{马加爵}\label{ux9a6cux52a0ux7235-1}

被捕后自述大意是：

\begin{itemize}
\tightlist
\item
  贫困和艰辛我能忍；
\item
  别人的歧视、蔑视我也能忍；
\item
  但最好的朋友践踏我的人格尊严，我不能忍；
\item
  既然这样，就``玉石俱焚''，要给那些歧视贫穷、歧视苦人的人一个教训。
\end{itemize}

老师强调： 其核心是\textbf{自尊受伤}后，以极端暴力方式``讨回尊严''。

\subsubsection{杨元元}\label{ux6768ux5143ux5143-1}

老师转述其遗书中的一句核心态度：
\textbf{``人可以被毁灭，但不可以被打败。''}

老师借此指出：
这类人往往把``宁折不弯''当成尊严，但如果不会调整和转弯，就可能把自己逼向绝路。

\begin{center}\rule{0.5\linewidth}{0.5pt}\end{center}

\subsection{三十二、老师总结两案共有的人格与心理特点}\label{ux4e09ux5341ux4e8cux8001ux5e08ux603bux7ed3ux4e24ux6848ux5171ux6709ux7684ux4ebaux683cux4e0eux5fc3ux7406ux7279ux70b9}

老师把两人的共性概括为三个词：

\subsubsection{1. 高敏感}\label{ux9ad8ux654fux611f}

\begin{itemize}
\tightlist
\item
  对人、对事太在乎；
\item
  容易把小事放大；
\item
  对羞辱、轻视体验格外强烈。
\end{itemize}

\subsubsection{2. 高自尊}\label{ux9ad8ux81eaux5c0a}

\begin{itemize}
\tightlist
\item
  内里可能自卑；
\item
  但越自卑，越想通过成绩、表现、荣誉来证明自己；
\item
  追求高自尊的过程，反而不断受伤。
\end{itemize}

\subsubsection{3. 执着、倔强}\label{ux6267ux7740ux5014ux5f3a}

\begin{itemize}
\tightlist
\item
  一条路走不通，也不愿转弯；
\item
  资源不够，却设定很高目标；
\item
  遇到困难不知调整；
\item
  坚持到最后，把自己``绷断''。
\end{itemize}

\subsubsection{老师的概括性结论}\label{ux8001ux5e08ux7684ux6982ux62ecux6027ux7ed3ux8bba}

本质上是：

\begin{itemize}
\tightlist
\item
  \textbf{资源不足}
\item
  却定了过于理想的目标
\item
  遇到挫折不懂调整
\item
  最后把自己逼崩了
\end{itemize}

老师认为，这种教训是很惨痛的。

\begin{center}\rule{0.5\linewidth}{0.5pt}\end{center}

\subsection{三十三、【举例 /
本校案例】过度要强、宿舍关系差、学习方法僵化，最终精神失常边缘}\label{ux4e09ux5341ux4e09ux4e3eux4f8b-ux672cux6821ux6848ux4f8bux8fc7ux5ea6ux8981ux5f3aux5bbfux820dux5173ux7cfbux5deeux5b66ux4e60ux65b9ux6cd5ux50f5ux5316ux6700ux7ec8ux7cbeux795eux5931ux5e38ux8fb9ux7f18}

老师接着举了一个本校女同学的危险案例，用来提醒大家重视现实中的前兆。

\subsubsection{基本情况}\label{ux57faux672cux60c5ux51b5-2}

\begin{itemize}
\tightlist
\item
  女同学
\item
  本硕连读
\item
  高中时是学霸
\item
  家长在原学校还是主任级人物
\item
  从前荣誉很多，长期处于``优秀学生''位置
\end{itemize}

\subsubsection{入学后的目标}\label{ux5165ux5b66ux540eux7684ux76eeux6807}

她进入医大后，给自己偷偷立了一个很高的标杆：

\begin{itemize}
\tightlist
\item
  要拿奖学金
\item
  要当班干部
\item
  要再创辉煌
\item
  要继续追求卓越
\end{itemize}

老师说： 这种``想变好''本身没错，但她有明显短板。

\subsubsection{短板一：看不起别人，宿舍关系恶化}\label{ux77edux677fux4e00ux770bux4e0dux8d77ux522bux4ebaux5bbfux820dux5173ux7cfbux6076ux5316}

\begin{itemize}
\tightlist
\item
  宿舍里谁都看不顺眼；
\item
  对室友多有意见；
\item
  一个宿舍里同时得罪多数人，结果就是自己被孤立。
\end{itemize}

\subsubsection{短板二：学习方法僵化}\label{ux77edux677fux4e8cux5b66ux4e60ux65b9ux6cd5ux50f5ux5316}

\begin{itemize}
\item
  想继续用高中模式对付大学；
\item
  面对大学大量内容和不同老师讲课方式，不会调整方法；
\item
  开始靠``蛮力''学：

  \begin{itemize}
  \tightlist
  \item
    少吃饭
  \item
    少睡觉
  \item
    拼命学
  \end{itemize}
\end{itemize}

\subsubsection{结果}\label{ux7ed3ux679c}

形成恶性循环：

\begin{itemize}
\tightlist
\item
  越睡不好
\item
  越吃不好
\item
  上课效率越差
\item
  情绪越不稳定
\end{itemize}

最终有一天夜里失控：

\begin{itemize}
\tightlist
\item
  别人都睡了，她睡不着；
\item
  半夜喊口号给自己打气；
\item
  严重影响舍友；
\item
  大家意识到``不对劲''，赶紧上报老师。
\end{itemize}

\subsubsection{后续处理}\label{ux540eux7eedux5904ux7406}

学校通知家长，希望家长陪同处理。
但家长一开始并不理解，反而先追究是不是别人欺负了她，试图找``幕后黑手''。

老师的判断是：

\begin{itemize}
\item
  宿舍小摩擦肯定有；
\item
  但更核心的问题还是她自身：

  \begin{itemize}
  \tightlist
  \item
    过度要强
  \item
    性格短板
  \item
    不会处理人际关系
  \item
    学习遇挫时不会调整方法
  \end{itemize}
\end{itemize}

\subsubsection{【临床关联 /
风险提醒】}\label{ux4e34ux5e8aux5173ux8054-ux98ceux9669ux63d0ux9192}

老师明确提醒： 如果这种情况继续发展，确实存在进一步升级的风险。
虽然是女同学，不一定发展成持刀等暴力，但医学院环境里有：

\begin{itemize}
\tightlist
\item
  实验室
\item
  有毒有害物质
\end{itemize}

如果人在极端状态下失控，后果可能非常严重。

\subsubsection{【举例】复旦投毒案}\label{ux4e3eux4f8bux590dux65e6ux6295ux6bd2ux6848}

老师顺带提到复旦大学研究生宿舍矛盾引发投毒致死的案例，用以说明：

\begin{itemize}
\tightlist
\item
  宿舍矛盾若严重激化；
\item
  又碰上具备实验条件和危险物质；
\item
  后果可能是致命的。
\end{itemize}

\begin{center}\rule{0.5\linewidth}{0.5pt}\end{center}

\subsection{三十四、本段结尾处老师的未完提醒}\label{ux4e09ux5341ux56dbux672cux6bb5ux7ed3ux5c3eux5904ux8001ux5e08ux7684ux672aux5b8cux63d0ux9192}

到你发来的这段转写末尾，老师的话停在：

\begin{quote}
``所以大家呢一定呢，哎，这个这个有矛盾啊，咱呢\ldots\ldots{}''
\end{quote}

也就是说，这一段最后显然还会继续讲：
\textbf{遇到矛盾时应如何处理、如何避免激化、如何及时求助}。
但由于当前转写片段在这里中断，笔记也先整理到这里。

\begin{center}\rule{0.5\linewidth}{0.5pt}\end{center}

\subsection{本节英文术语索引}\label{ux672cux8282ux82f1ux6587ux672fux8bedux7d22ux5f15}

\begin{itemize}
\tightlist
\item
  心理危机 ------ \textbf{psychological crisis}
\item
  危机 ------ \textbf{crisis}
\item
  危机干预 ------ \textbf{crisis intervention}
\item
  心理危机干预与预防 ------ \textbf{psychological crisis intervention
  and prevention}
\item
  心理治疗 ------ \textbf{psychotherapy}
\item
  心理平衡 ------ \textbf{psychological balance}
\item
  心理失衡 ------ \textbf{psychological disequilibrium}
\item
  生活事件 ------ \textbf{life event}
\item
  情绪反应 ------ \textbf{emotional response}
\item
  认知反应 ------ \textbf{cognitive response}
\item
  身体反应 / 躯体反应 ------ \textbf{physical response / somatic
  response}
\item
  行为反应 ------ \textbf{behavioral response}
\item
  心理健康 ------ \textbf{mental health}
\item
  精神卫生 ------ \textbf{mental hygiene / mental health service}
\item
  抑郁 / 抑郁症 ------ \textbf{depression}
\item
  焦虑 ------ \textbf{anxiety}
\item
  失眠 ------ \textbf{insomnia}
\item
  自杀 ------ \textbf{suicide}
\item
  伤人 / 杀人 ------ \textbf{homicide / violence toward others}
\item
  临床心理科 ------ \textbf{clinical psychology}
\item
  乙型肝炎 ------ \textbf{hepatitis B}
\item
  乙肝大三阳 ------ \textbf{hepatitis B ``big three positive'' /
  HBeAg-positive status}
\item
  体液传播 ------ \textbf{body fluid transmission}
\item
  精神障碍 ------ \textbf{mental disorder}
\item
  精神分裂症 ------ \textbf{schizophrenia}
\item
  幻觉 ------ \textbf{hallucination}
\item
  妄想 ------ \textbf{delusion}
\item
  人格 ------ \textbf{personality}
\item
  非理性认知 ------ \textbf{irrational cognition}
\item
  过度概括 ------ \textbf{overgeneralization}
\item
  非此即彼思维 ------ \textbf{all-or-none thinking}
\item
  慢性压力 ------ \textbf{chronic stress}
\item
  应激 ------ \textbf{stress / stress response}
\item
  适应不良 / 不适应 ------ \textbf{maladjustment}
\item
  自杀未遂史 ------ \textbf{history of suicide attempt}
\item
  生命质量 ------ \textbf{quality of life}
\item
  躯体疾病 ------ \textbf{physical illness / somatic disease}
\item
  求助 ------ \textbf{help-seeking}
\end{itemize}

你把后半段转写继续发来，我可以直接接在这份笔记下面，继续按同样格式无缝往下整理。
https://resourcemap.lol/sub?token=eE-E2UXL8bvajtGCpyj3t-cyooYr\_js6m9JDqUEluas\&type=clash\&protocol=ss2022\&use=domestic\&rule=default
库洛米 720 prepend: - `DOMAIN-SUFFIX,dmsguild.com,🇺🇸 United States' -
`DOMAIN-SUFFIX,openrouter.ai,EXFLUX' -
`DOMAIN-SUFFIX,withgoogle.com,Google' -
`DOMAIN,speed.cloudflare.com,Test' - `DOMAIN-KEYWORD,pikpak,GlobalMedia'
- `IP-CIDR,202.113.53.231/32,TMU' -
`PROCESS-NAME,AcroCEF.exe,REJECT-DROP' -
`PROCESS-NAME,AcrobatProPortable.exe,REJECT-DROP' -
`PROCESS-NAME,Acrobat.exe,REJECT-DROP' -
`DOMAIN-SUFFIX,apkpure.com,EXFLUX' - `DOMAIN-SUFFIX,openai.com,EXFLUX' -
`PROCESS-NAME,ChatGPT.exe,EXFLUX' - `DOMAIN-SUFFIX,chatgpt.com,EXFLUX' -
`DOMAIN-SUFFIX,tmu.edu.cn,TMU' -
`DOMAIN-SUFFIX,tmuvod.smartclass.cn,TMU' -
`DOMAIN-SUFFIX,tmu.smartclass.cn,TMU' -
`DOMAIN-SUFFIX,googletagmanager.com,EXFLUX' -
`DOMAIN-SUFFIX,auth0.com,EXFLUX' - `DOMAIN-SUFFIX,jsdelivr.net,EXFLUX' -
`DOMAIN-KEYWORD,tmu,TMU' - `DOMAIN-SUFFIX,noodlemagazine.com,Porn'
append: {[}{]} delete: {[}{]} prepend: - name: `TW-HiNet-SS' type: `ss'
server: `tw-wap-ac-hinet.ddns.zjhstudio.com' port: 35102 cipher:
`aes-128-gcm' password: `82K2ORjiDMenG1K/zbEvKPSJ2DxEt4Hzl+Sk+yVPmRc='
dialer-proxy: `Auto' - name: `DiTW-HiNet-SS' type: `ss' server:
`tw-wap-ac-hinet.ddns.zjhstudio.com' port: 35102 cipher: `aes-128-gcm'
password: `82K2ORjiDMenG1K/zbEvKPSJ2DxEt4Hzl+Sk+yVPmRc=' - name:
`twTW-HiNet-SS' type: `ss' server: `tw-wap-ac-hinet.ddns.zjhstudio.com'
port: 35102 cipher: `aes-128-gcm' password:
`82K2ORjiDMenG1K/zbEvKPSJ2DxEt4Hzl+Sk+yVPmRc=' dialer-proxy: `🇨🇳 Taiwan'
append: - name: `台湾 12' type: `ss' server:
`1253be5372bd821df0c14814ec430772.lv12.exzxwmrbnedyoldwod.com' port:
7411 cipher: `chacha20-ietf-poly1305' password: `L0CjU3' udp: true -
name: `tmu' type: `ss' server: `59.110.9.119' port: 26001 cipher:
`aes-128-gcm' password: `OLL7bLkKXaKP8CMaw+gDNMRHx0IhzptxRHb8RT0c1Hg='
delete: {[}{]} prepend: - type: `select' name: `Test' interval: 300
timeout: 5000 max-failed-times: 5 lazy: true proxies: - `Kuromis' -
`EXFLUX' - `tmu' - `TW-HiNet-SS' - `DiTW-HiNet-SS' - `twTW-HiNet-SS' -
type: `select' name: `EXFLUX' interval: 300 timeout: 5000
max-failed-times: 5 lazy: true proxies: - `台湾 12' - `TW-HiNet-SS' -
`Kuromis' - `🇭🇰 Hong Kong 01' - `🇭🇰 Hong Kong 02' - `🇭🇰 Hong Kong 03' -
`🇭🇰 Hong Kong 04' - `🇭🇰 Hong Kong 05' - `🇭🇰 Hong Kong 06' - `🇭🇰 Hong
Kong 07' - `🇭🇰 Hong Kong 08' - `🇭🇰 Hong Kong 09' append: - type:
`select' name: `TMU' interval: 300 timeout: 5000 max-failed-times: 5
lazy: true proxies: - `tmu' - type: `select' name: `Porn' interval: 300
timeout: 5000 max-failed-times: 5 lazy: true proxies: - `🇭🇰 Hong Kong' -
`🇸🇬 Singapore' - `🇯🇵 Japan' - `🇺🇸 United States' - `🇨🇳 Taiwan' - `🇩🇪
Germany' - `🇬🇧 United Kingdom' - `🇰🇷 South Korea' - `🇦🇺 Australia' - `🇫🇷
France' - `🇨🇭 Switzerland' - `🇳🇱 Netherlands' - `🇮🇳 India' delete:
{[}{]} /** * ========================= * 全局扩展脚本（JavaScript） *
作用： * 1. 给配置追加你自己的固定节点 * 2. 追加你自己的代理组 * 3.
在规则最前面插入你自己的前置规则 * 4. 尽量避免重复添加 *
========================= \emph{ } 用法建议： * - 这份脚本在你自己的 Git
仓库里维护 * - Clash Verge Rev 里只粘贴``运行副本'' * - 改动时永远先改
Git 里的源文件 */

/* ------------------------- * 你自己的固定节点 * 说明： * 1.
所有代理节点至少都要有 name / type / server / port * 2.
不同协议还会有各自的专属字段 * 3. 下面先给一个 ss 示例 *
------------------------- */ const customProxies = {[} \{ name:
``TW-HiNet-SS'', type: ``ss'', server:
``tw-wap-ac-hinet.ddns.zjhstudio.com'', port: 35102, cipher:
``aes-128-gcm'', password:
``82K2ORjiDMenG1K/zbEvKPSJ2DxEt4Hzl+Sk+yVPmRc='', `dialer-proxy':
``Auto'', udp: true \}, \{ name: ``tmu'', type: ``ss'', server:
``59.110.9.119'', port: 26001, cipher: ``aes-128-gcm'', password:
``OLL7bLkKXaKP8CMaw+gDNMRHx0IhzptxRHb8RT0c1Hg='', udp: true \}, //
你可以继续往下加更多节点 // \{ // name: ``自建-US-01'', // type:
``trojan'', // server: ``example.com'', // port: 443, // password:
``你的密码'', // sni: ``example.com'', // skip-cert-verify: false //
\}{]};

/* ------------------------- * 你自己的代理组 * 说明： * 1.
这里既可以显式写 proxies，也可以用 include-all * 2. Mihomo 支持 select /
url-test 等组类型 * ------------------------- */ const customGroups =
{[} \{ name: ``ChatGPT'', type: ``select'', proxies: {[}
``TW-HiNet-SS'', ``Auto'' {]}, url: ``https://chatgpt.com'', interval:
300, timeout: 5000, `max-failed-times': 5, tolerance: 50 \}, \{ name:
``TMU'', type: ``select'', proxies: {[} ``tmu'' {]}, url:
``http://tyjw.tmu.edu.cn/'', interval: 300, timeout: 5000,
`max-failed-times': 5, tolerance: 50 \}, \{ name: ``Porn'', type:
``select'', proxies: {[} ``🇭🇰 Hong Kong'', ``🇸🇬 Singapore'', ``🇯🇵
Japan'', ``🇺🇸 United States'', `🇨🇳 Taiwan', `🇩🇪 Germany', `🇬🇧 United
Kingdom', `🇰🇷 South Korea', `🇦🇺 Australia', `🇫🇷 France', `🇨🇭
Switzerland', `🇳🇱 Netherlands', `🇮🇳 India' {]} \}{]};

/* ------------------------- * 你自己的前置规则 * 说明： * 1.
这些规则会被插到原始 rules 前面 * 2.
规则从上到下匹配，所以前置规则优先级更高 * ------------------------- */
const customRules = {[} `DOMAIN-SUFFIX,dmsguild.com,🇺🇸 United States',
`DOMAIN-SUFFIX,openrouter.ai,EXFLUX',
`DOMAIN-SUFFIX,withgoogle.com,Google',
`DOMAIN,speed.cloudflare.com,Test', `DOMAIN-KEYWORD,pikpak,GlobalMedia',
`IP-CIDR,202.113.53.231/32,TMU', `PROCESS-NAME,AcroCEF.exe,REJECT-DROP',
`PROCESS-NAME,AcrobatProPortable.exe,REJECT-DROP',
`PROCESS-NAME,Acrobat.exe,REJECT-DROP',
`DOMAIN-SUFFIX,apkpure.com,EXFLUX', `DOMAIN-SUFFIX,openai.com,EXFLUX',
`PROCESS-NAME,ChatGPT.exe,EXFLUX', `DOMAIN-SUFFIX,chatgpt.com,EXFLUX',
`DOMAIN-SUFFIX,tmu.edu.cn,TMU',
`DOMAIN-SUFFIX,tmuvod.smartclass.cn,TMU',
`DOMAIN-SUFFIX,tmu.smartclass.cn,TMU',
`DOMAIN-SUFFIX,googletagmanager.com,EXFLUX',
`DOMAIN-SUFFIX,auth0.com,EXFLUX', `DOMAIN-SUFFIX,jsdelivr.net,EXFLUX',
`DOMAIN-KEYWORD,tmu,TMU', `DOMAIN-SUFFIX,noodlemagazine.com,Porn'{]};

/* ------------------------- * 工具函数：保证字段是数组 *
------------------------- */ function ensureArray(obj, key) \{ if
(!Array.isArray(obj{[}key{]})) \{ obj{[}key{]} = {[}{]}; \} \}

/* ------------------------- * 工具函数：按 name 去重插入 *
如果已存在同名项，就替换；否则追加 * 适用于 proxies / proxy-groups *
------------------------- */ function upsertByName(arr, item) \{ const
index = arr.findIndex(x =\textgreater{} x \&\& x.name === item.name); if
(index \textgreater= 0) \{ arr{[}index{]} = item; \} else \{
arr.push(item); \} \}

/* ------------------------- * 工具函数：前置插入规则并去重 *
------------------------- */ function prependUniqueRules(oldRules,
newRules) \{ const oldList = Array.isArray(oldRules) ? oldRules :
{[}{]}; const oldSet = new Set(oldList);

// 只保留那些旧规则里还没有的新增规则 const uniqueNewRules =
newRules.filter(rule =\textgreater{} !oldSet.has(rule));

return uniqueNewRules.concat(oldList); \}

/* ------------------------- * 主入口函数 * profileName
是当前处理的配置文件名 * 你可以按订阅名称做条件判断 *
------------------------- */ function main(config, profileName) \{ //
===== 可选：只对某个订阅生效 ===== //
如果你想让这份脚本只作用于某个主订阅， //
就把下面两行取消注释，并把``我的主订阅''改成你的订阅名。 // // if
(profileName !== ``我的主订阅'') \{ // return config; // \}

ensureArray(config, ``proxies''); ensureArray(config, ``proxy-groups'');
ensureArray(config, ``rules'');

// 1) 注入你自己的固定节点 for (const proxy of customProxies) \{
upsertByName(config{[}``proxies''{]}, proxy); \}

// 2) 注入你自己的代理组 for (const group of customGroups) \{
upsertByName(config{[}``proxy-groups''{]}, group); \}

// 3) 把你的规则放到最前面 config{[}``rules''{]} =
prependUniqueRules(config{[}``rules''{]}, customRules);

return config; \}

\clearpage

\section{第2讲：心理危机的常见原因}\label{ux7b2c2ux8bb2ux5fc3ux7406ux5371ux673aux7684ux5e38ux89c1ux539fux56e0}

\subsection{主题：心理危机的常见原因（第 2
次课）}\label{ux4e3bux9898ux5fc3ux7406ux5371ux673aux7684ux5e38ux89c1ux539fux56e0ux7b2c-2-ux6b21ux8bfe}

\begin{center}\rule{0.5\linewidth}{0.5pt}\end{center}

\subsection{一、开场与本次课程主题}\label{ux4e00ux5f00ux573aux4e0eux672cux6b21ux8bfeux7a0bux4e3bux9898}

老师先说明了本节课内容：这是\textbf{《心理危机干预与预防》第二次课}，本次主要讨论\textbf{心理危机（psychological
crisis）的常见原因}。

\begin{center}\rule{0.5\linewidth}{0.5pt}\end{center}

\subsection{二、【复习】危机（crisis）的基本概念}\label{ux4e8cux590dux4e60ux5371ux673acrisisux7684ux57faux672cux6982ux5ff5}

老师先回顾了上节课可能已经讲过的``危机''概念，并强调危机有\textbf{两层含义}：

\subsubsection{1.
危机可以指``突发事件''本身}\label{ux5371ux673aux53efux4ee5ux6307ux7a81ux53d1ux4e8bux4ef6ux672cux8eab}

即一个\textbf{出乎人们意料而发生的事件}，例如：

\begin{itemize}
\tightlist
\item
  自然灾害（natural disasters）
\item
  技术性灾害（technological disasters）
\item
  恐怖袭击（terrorist attacks）
\item
  战争（war）等
\end{itemize}

老师举例说，人们会说``伊朗危机''，这里的``危机''就是指这个事件本身。

\subsubsection{2.
危机也可以指``人所处的紧急失衡状态''}\label{ux5371ux673aux4e5fux53efux4ee5ux6307ux4ebaux6240ux5904ux7684ux7d27ux6025ux5931ux8861ux72b6ux6001}

当个体遭遇重大问题或变化，感到：

\begin{itemize}
\tightlist
\item
  难以解决
\item
  难以把握
\item
  原有平衡被打破
\item
  正常生活受干扰
\item
  内心紧张不断积蓄
\end{itemize}

进一步就可能出现：

\begin{itemize}
\tightlist
\item
  无所适从
\item
  思维和行为紊乱
\item
  进入一种\textbf{失衡状态（state of disequilibrium）}
\end{itemize}

这就是心理学上讲的\textbf{危机状态（crisis state）}。

\subsubsection{3.
危机的核心含义}\label{ux5371ux673aux7684ux6838ux5fc3ux542bux4e49}

老师强调：

\begin{itemize}
\tightlist
\item
  \textbf{危机意味着平衡与稳定被破坏。}
\item
  危机之所以叫``危机''，就是因为它打破了人长期依赖的平衡，引起混乱和不安。
\end{itemize}

\subsubsection{4.
危机并不单由事件本身决定，而取决于个体的应对能力}\label{ux5371ux673aux5e76ux4e0dux5355ux7531ux4e8bux4ef6ux672cux8eabux51b3ux5b9aux800cux53d6ux51b3ux4e8eux4e2aux4f53ux7684ux5e94ux5bf9ux80fdux529b}

老师特别强调一个判断标准：

\begin{quote}
危机的出现，是因为个体意识到某个事件或情境\textbf{超过了自己的应付能力（coping
capacity）}，而不只是因为事件本身发生了。
\end{quote}

也就是说：

\begin{itemize}
\tightlist
\item
  同一件事，对有些人构成危机；
\item
  对另一些人，可能并不构成危机。
\end{itemize}

\begin{center}\rule{0.5\linewidth}{0.5pt}\end{center}

\subsection{三、【举例】以考试说明``危机感''的形成}\label{ux4e09ux4e3eux4f8bux4ee5ux8003ux8bd5ux8bf4ux660eux5371ux673aux611fux7684ux5f62ux6210}

老师用医学生熟悉的考试场景做例子。

\subsubsection{1.
过去的生活中，很多``特别棘手的事''都可能让人感到危机}\label{ux8fc7ux53bbux7684ux751fux6d3bux4e2dux5f88ux591aux7279ux522bux68d8ux624bux7684ux4e8bux90fdux53efux80fdux8ba9ux4ebaux611fux5230ux5371ux673a}

例如：

\begin{itemize}
\tightlist
\item
  英语四级、六级
\item
  高考
\item
  未来的\textbf{执业医师考试（physician licensing examination）}
\end{itemize}

\subsubsection{2.
【课堂提醒】执业医师考试对很多医学生来说会构成危机}\label{ux8bfeux5802ux63d0ux9192ux6267ux4e1aux533bux5e08ux8003ux8bd5ux5bf9ux5f88ux591aux533bux5b66ux751fux6765ux8bf4ux4f1aux6784ux6210ux5371ux673a}

老师提到：

\begin{itemize}
\tightlist
\item
  临床、儿科、麻醉、影像等专业学生以后都要面对执业医师考试。
\item
  这对很多人来说是一道``坎''。
\end{itemize}

老师还提到学校执业医师通过率``不是那么喜人''，并借此说明为什么考试会构成危机：

\begin{itemize}
\tightlist
\item
  今年没过，明年还得继续考；
\item
  明年又会加入新的考生；
\item
  ``盘子更大''，竞争更强；
\item
  在通过率限制下，个体主观上会感到希望更渺茫。
\end{itemize}

老师借此说明：\textbf{危机感本质上来自于``我觉得自己应付不过来''}。

\begin{center}\rule{0.5\linewidth}{0.5pt}\end{center}

\subsection{四、心理危机的表现与潜伏性}\label{ux56dbux5fc3ux7406ux5371ux673aux7684ux8868ux73b0ux4e0eux6f5cux4f0fux6027}

老师接着说明，\textbf{心理危机（psychological
crisis）}可以表现为很多层面：

\begin{itemize}
\tightlist
\item
  心理状态严重失调
\item
  心理矛盾激烈冲突、难以解决
\item
  精神濒临崩溃，甚至精神失常
\item
  出现心理障碍（psychological disorder）
\end{itemize}

\subsubsection{1.
当事人不一定总能立刻察觉自己处于心理危机中}\label{ux5f53ux4e8bux4ebaux4e0dux4e00ux5b9aux603bux80fdux7acbux523bux5bdfux89c9ux81eaux5df1ux5904ux4e8eux5fc3ux7406ux5371ux673aux4e2d}

老师指出：

\begin{itemize}
\tightlist
\item
  有的人会及时察觉；
\item
  也有人是在``不知不觉中''进入危机状态。
\end{itemize}

\subsubsection{2.
表面看似``正常''的行为模式下，也可能潜伏危机}\label{ux8868ux9762ux770bux4f3cux6b63ux5e38ux7684ux884cux4e3aux6a21ux5f0fux4e0bux4e5fux53efux80fdux6f5cux4f0fux5371ux673a}

老师举例：

\begin{itemize}
\tightlist
\item
  社会上很多人每天照常上班、下班；
\item
  表面看起来是正常的上班族；
\item
  但``压垮骆驼的是最后一根稻草''。
\end{itemize}

也就是说，一个人即便长期维持某种习惯化、看似稳定的行为模式，也可能潜藏着严重的心理危机。

\begin{center}\rule{0.5\linewidth}{0.5pt}\end{center}

\subsection{五、【重点】成瘾行为背后的心理危机}\label{ux4e94ux91cdux70b9ux6210ux763eux884cux4e3aux80ccux540eux7684ux5fc3ux7406ux5371ux673a}

老师由``隐癖''引入成瘾问题，强调：

\begin{itemize}
\tightlist
\item
  严重不良隐癖（如物质滥用）背后，往往潜伏心理危机；
\item
  当隐癖被去除时，心理危机就会暴露出来。
\end{itemize}

\subsubsection{1.
【课堂提问】最常见的精神活性物质滥用是什么}\label{ux8bfeux5802ux63d0ux95eeux6700ux5e38ux89c1ux7684ux7cbeux795eux6d3bux6027ux7269ux8d28ux6ee5ux7528ux662fux4ec0ux4e48}

老师提问后给出答案：

\begin{itemize}
\tightlist
\item
  最常见、也最容易引发法律问题的精神活性物质（psychoactive
  substance）滥用是\textbf{酒精（alcohol）}。
\end{itemize}

并提醒医学生：

\begin{itemize}
\tightlist
\item
  这是临床上必须知道的常识；
\item
  酒驾、醉驾就涉及酒精滥用带来的现实法律问题。
\end{itemize}

\subsubsection{2.
匿名戒酒者协会}\label{ux533fux540dux6212ux9152ux8005ux534fux4f1a}

老师提到一个组织：

\begin{itemize}
\tightlist
\item
  \textbf{匿名戒酒者协会（Alcoholics Anonymous, AA）}
\end{itemize}

并说明其特点：

\begin{itemize}
\tightlist
\item
  它不完全等同于正规心理治疗组织；
\item
  更接近一种\textbf{互助组织（mutual-help group）}；
\item
  会员采取匿名制；
\item
  定期组织活动，帮助成员戒酒。
\end{itemize}

\subsubsection{3.
【核心观点】成瘾往往只是表面现象}\label{ux6838ux5fc3ux89c2ux70b9ux6210ux763eux5f80ux5f80ux53eaux662fux8868ux9762ux73b0ux8c61}

老师强调这一点非常重要：

\begin{itemize}
\tightlist
\item
  酒精成瘾也好
\item
  手机成瘾也好
\end{itemize}

这些都往往只是\textbf{表面现象}，其背后真正掩盖的可能是：

\begin{itemize}
\tightlist
\item
  无法解决的内心冲突
\item
  情绪问题
\item
  家庭矛盾
\item
  其他长期心理危机来源
\end{itemize}

\subsubsection{4.
只``戒行为''往往治标不治本}\label{ux53eaux6212ux884cux4e3aux5f80ux5f80ux6cbbux6807ux4e0dux6cbbux672c}

老师举精神科病房中的例子：

\begin{itemize}
\tightlist
\item
  有些反复住院戒酒的患者，在住院期间当然喝不到酒；
\item
  但他们会在医院里商量出院后去哪儿``大喝一顿''。
\end{itemize}

老师借此说明：

\begin{itemize}
\tightlist
\item
  单纯把成瘾行为压下去，并不等于问题解决；
\item
  关键是识别\textbf{成瘾行为背后的心理危机}。
\end{itemize}

\begin{center}\rule{0.5\linewidth}{0.5pt}\end{center}

\subsection{六、【复习】心理危机原因的多元性：生物---心理---社会医学模式}\label{ux516dux590dux4e60ux5fc3ux7406ux5371ux673aux539fux56e0ux7684ux591aux5143ux6027ux751fux7269ux5fc3ux7406ux793eux4f1aux533bux5b66ux6a21ux5f0f}

老师提醒大家，医学第一课通常就会讲到：

\begin{itemize}
\tightlist
\item
  \textbf{生物---心理---社会医学模式（biopsychosocial medical model）}
\end{itemize}

意思是：人的问题通常不是单一来源，而是由多个层面共同构成。

老师把危机原因概括为三个大方面：

\begin{enumerate}
\def\labelenumi{\arabic{enumi}.}
\tightlist
\item
  \textbf{生物/生理层面（biological/physiological factors）}
\item
  \textbf{心理层面（psychological factors）}
\item
  \textbf{社会文化层面（social and cultural factors）}
\end{enumerate}

\begin{center}\rule{0.5\linewidth}{0.5pt}\end{center}

\subsection{七、生物生理层面：季节、光照与情绪变化}\label{ux4e03ux751fux7269ux751fux7406ux5c42ux9762ux5b63ux8282ux5149ux7167ux4e0eux60c5ux7eeaux53d8ux5316}

老师首先从\textbf{生物生理学层面}讲危机原因，并以``冬天和夏天的差异''为例。

\subsubsection{1.
冬天常见的身心变化}\label{ux51acux5929ux5e38ux89c1ux7684ux8eabux5fc3ux53d8ux5316}

老师提问并总结：

冬天相较夏天，人往往更容易出现：

\begin{itemize}
\tightlist
\item
  更想吃碳水化合物（carbohydrates）
\item
  早晨起不来床
\item
  总犯困
\item
  情绪低落
\item
  整体更``打蔫儿''
\end{itemize}

\subsubsection{2.
原因：北半球冬季光照减少}\label{ux539fux56e0ux5317ux534aux7403ux51acux5b63ux5149ux7167ux51cfux5c11}

老师解释：

\begin{itemize}
\tightlist
\item
  天津位于北半球；
\item
  北半球进入冬天后，\textbf{光照逐渐减少}；
\item
  到冬至时，光照最少；
\item
  夏至则昼最长、夜最短，冬至相反。
\end{itemize}

\subsubsection{3.
光照减少对脑内神经递质的影响}\label{ux5149ux7167ux51cfux5c11ux5bf9ux8111ux5185ux795eux7ecfux9012ux8d28ux7684ux5f71ux54cd}

老师从神经生化角度解释：

\begin{itemize}
\item
  在光照减少时，大脑中的\textbf{松果体（pineal gland）}活动相关变化会使

  \begin{itemize}
  \tightlist
  \item
    \textbf{5-羟色胺（5-hydroxytryptamine,
    5-HT；serotonin，血清素）}减少
  \item
    \textbf{褪黑素（melatonin）}增加
  \end{itemize}
\end{itemize}

由此带来：

\begin{itemize}
\tightlist
\item
  褪黑素高 → 更容易困、总想睡觉
\item
  5-HT 低 → 情绪活跃性下降，容易低落
\end{itemize}

老师借此说明：

\begin{itemize}
\tightlist
\item
  冬天人总想睡觉、总显得没夏天那么活跃，是有\textbf{生物化学机制}的。
\end{itemize}

\subsubsection{4. 【考点味很浓】季节性情感障碍
SAD}\label{ux8003ux70b9ux5473ux5f88ux6d53ux5b63ux8282ux6027ux60c5ux611fux969cux788d-sad}

老师引出：

\begin{itemize}
\tightlist
\item
  \textbf{季节性情感障碍（Seasonal Affective Disorder, SAD）}
\end{itemize}

高纬度地区更常见，表现包括：

\begin{itemize}
\tightlist
\item
  冬季情绪低落
\item
  烦躁
\item
  嗜睡
\item
  每天睡十几个小时，甚至更久
\item
  喜欢吃碳水化合物
\item
  体重增加
\end{itemize}

老师说明：

\begin{itemize}
\tightlist
\item
  有些患者如果不做干预，也可能在春天自然缓解；
\item
  这说明它与\textbf{自然节律（natural rhythm）}关系密切。
\end{itemize}

\subsubsection{5.
高纬度地区患病率更高}\label{ux9ad8ux7eacux5ea6ux5730ux533aux60a3ux75c5ux7387ux66f4ux9ad8}

老师举美国不同地区为例说明：

\begin{itemize}
\tightlist
\item
  佛罗里达州：约 1\%
\item
  华盛顿特区：约 4\%
\item
  阿拉斯加：约 10\%
\end{itemize}

老师借此得出结论：

\begin{itemize}
\tightlist
\item
  随着\textbf{纬度升高}，SAD 发生率上升；
\item
  说明人的情绪状态会受到自然环境显著影响。
\end{itemize}

\begin{center}\rule{0.5\linewidth}{0.5pt}\end{center}

\subsection{八、【机制讲解】为什么情绪不好时容易想吃碳水}\label{ux516bux673aux5236ux8bb2ux89e3ux4e3aux4ec0ux4e48ux60c5ux7eeaux4e0dux597dux65f6ux5bb9ux6613ux60f3ux5403ux78b3ux6c34}

老师顺势解释``情绪不好就想吃东西''的生化基础。

\subsubsection{1.
情绪差时，很多人偏爱碳水食物}\label{ux60c5ux7eeaux5deeux65f6ux5f88ux591aux4ebaux504fux7231ux78b3ux6c34ux98dfux7269}

老师提醒大家观察生活现象：

\begin{itemize}
\tightlist
\item
  情绪不好的人，常会``化郁闷为食量''；
\item
  尤其偏爱富含糖和淀粉的食物。
\end{itemize}

\subsubsection{2. 生化机制}\label{ux751fux5316ux673aux5236}

老师按生理生化逻辑解释：

\begin{itemize}
\tightlist
\item
  吃碳水后 → \textbf{胰岛素（insulin）}分泌增高；
\item
  胰岛素会促进除\textbf{色氨酸（tryptophan）}以外的其他竞争性氨基酸从血液进入组织；
\item
  血液中色氨酸相对比例上升；
\item
  氨基酸通过\textbf{血脑屏障（blood-brain barrier, BBB）}需要载体转运；
\item
  色氨酸比例升高后，更容易抢占载体进入脑内；
\item
  色氨酸是\textbf{5-羟色胺（5-HT）}的原料。
\end{itemize}

\subsubsection{3. 结论}\label{ux7ed3ux8bba}

从这个意义上说：

\begin{itemize}
\tightlist
\item
  情绪低落时吃碳水，\textbf{在一定程度上确实可能改善抑郁情绪}；
\item
  这种做法是有生化机制支持的。
\end{itemize}

但老师也强调副作用：

\begin{itemize}
\tightlist
\item
  容易肥胖
\item
  体重增加
\item
  增加代谢负担
\item
  长期大量甜食对胰岛不利
\end{itemize}

\begin{center}\rule{0.5\linewidth}{0.5pt}\end{center}

\subsection{九、社会文化层面：人不仅是生物学的人，也是社会文化的人}\label{ux4e5dux793eux4f1aux6587ux5316ux5c42ux9762ux4ebaux4e0dux4ec5ux662fux751fux7269ux5b66ux7684ux4ebaux4e5fux662fux793eux4f1aux6587ux5316ux7684ux4eba}

老师接着强调：

\begin{itemize}
\tightlist
\item
  人不能只作为一个生物学个体来理解；
\item
  人还处在\textbf{社会文化（social and cultural）}之中；
\item
  文化与历史不断在人身上留下``烙印''。
\end{itemize}

\subsubsection{1.
【举例】男女衬衫开襟方向不同}\label{ux4e3eux4f8bux7537ux5973ux886cux886bux5f00ux895fux65b9ux5411ux4e0dux540c}

老师举了一个很日常但很多人没注意过的例子：

\begin{itemize}
\tightlist
\item
  男装衬衫通常是\textbf{左压右}
\item
  女装衬衫通常是\textbf{右压左}
\end{itemize}

老师解释其历史文化原因：

\paragraph{男装为什么左压右}\label{ux7537ux88c5ux4e3aux4ec0ux4e48ux5de6ux538bux53f3}

\begin{itemize}
\tightlist
\item
  欧洲最早有纽扣衣服时，穿得起带纽扣衣服的多是贵族；
\item
  贵族男性常佩剑；
\item
  剑一般挂在左侧；
\item
  如果衣服不是左压右，而是右压左，拔剑时剑的护手容易勾住衣服；
\item
  为避免决斗时影响动作，男装形成左压右。
\end{itemize}

\paragraph{女装为什么相反}\label{ux5973ux88c5ux4e3aux4ec0ux4e48ux76f8ux53cd}

\begin{itemize}
\tightlist
\item
  女性一般不佩剑；
\item
  贵族女性的穿衣常由仆人协助；
\item
  反向开襟更方便侍者操作。
\end{itemize}

老师借此说明：\textbf{看似理所当然的日常习惯，很多都是文化长期塑造的结果。}

\subsubsection{2.
【举例】东京与大阪自动扶梯左右站位不同}\label{ux4e3eux4f8bux4e1cux4eacux4e0eux5927ux962aux81eaux52a8ux6276ux68afux5de6ux53f3ux7ad9ux4f4dux4e0dux540c}

老师又举日本的例子：

\begin{itemize}
\tightlist
\item
  东京：站左边，把右边让出来
\item
  大阪：站右边，把左边让出来
\end{itemize}

解释如下：

\paragraph{东京（江户）为什么靠左}\label{ux4e1cux4eacux6c5fux6237ux4e3aux4ec0ux4e48ux9760ux5de6}

\begin{itemize}
\tightlist
\item
  东京历史上是政治中心；
\item
  武士佩刀，刀通常挂左侧；
\item
  若两人都靠右走，刀容易相碰，引发冲突；
\item
  因此形成靠左通行的习惯。
\end{itemize}

\paragraph{大阪为什么靠右}\label{ux5927ux962aux4e3aux4ec0ux4e48ux9760ux53f3}

\begin{itemize}
\tightlist
\item
  大阪是商业中心，商人多；
\item
  商人常把贵重物品放在右侧，便于右手看护；
\item
  因此习惯靠右站立，以减少被顺手牵羊的风险。
\end{itemize}

老师借此再次强调：\textbf{文化历史会长期塑造人的行为模式。}

\begin{center}\rule{0.5\linewidth}{0.5pt}\end{center}

\subsection{十、心理层面：人格、选择与``后悔''}\label{ux5341ux5fc3ux7406ux5c42ux9762ux4ebaux683cux9009ux62e9ux4e0eux540eux6094}

接着老师谈到心理层面，尤其是\textbf{人格与应对方式（personality and
coping style）}。

\subsubsection{1.
关于``后悔过去的决定''}\label{ux5173ux4e8eux540eux6094ux8fc7ux53bbux7684ux51b3ux5b9a}

老师问大家：

\begin{itemize}
\tightlist
\item
  在过去 20 多年的生活中，有没有做过让今天的自己后悔的决定？
\end{itemize}

甚至开玩笑举例：``比如学医。''

\subsubsection{2.
老师的观点：很多后悔是``事后诸葛亮视角''}\label{ux8001ux5e08ux7684ux89c2ux70b9ux5f88ux591aux540eux6094ux662fux4e8bux540eux8bf8ux845bux4eaeux89c6ux89d2}

老师提出一个思考框架：

\begin{itemize}
\tightlist
\item
  假如用时光机器把你送回当年；
\item
  但你对未来发生的事\textbf{一无所知}；
\item
  在当时那个情境里，你大概率还是会做出和当年相似的选择。
\end{itemize}

由此老师认为：

\begin{itemize}
\tightlist
\item
  我们今天常常是站在``事后''评价当初的自己；
\item
  但人在当时的信息条件下，通常已经做出了\textbf{当时自己能做出的最好选择}。
\end{itemize}

\subsubsection{3. 结论}\label{ux7ed3ux8bba-1}

老师的结论是：

\begin{itemize}
\tightlist
\item
  没有必要过度沉溺于对过去决定的后悔；
\item
  因为再来一遍，在同样信息条件下，大概率还是会那么选。
\end{itemize}

并由此引出一句话：

\begin{itemize}
\tightlist
\item
  ``人格决定命运。''
\end{itemize}

\begin{center}\rule{0.5\linewidth}{0.5pt}\end{center}

\subsection{十一、心理危机原因的呈现方式：两种典型模型}\label{ux5341ux4e00ux5fc3ux7406ux5371ux673aux539fux56e0ux7684ux5448ux73b0ux65b9ux5f0fux4e24ux79cdux5178ux578bux6a21ux578b}

老师用两个图像比喻危机来源的两种呈现方式。

\subsubsection{1.
左图：山上巨石突然滚落}\label{ux5de6ux56feux5c71ux4e0aux5de8ux77f3ux7a81ux7136ux6edaux843d}

象征：

\begin{itemize}
\tightlist
\item
  突如其来
\item
  毫无预兆
\item
  灾难性强
\item
  破坏性大
\item
  难以提前预测和应对
\end{itemize}

即：\textbf{突发性灾难事件}。

\subsubsection{2.
右图：鞋里不断进沙子}\label{ux53f3ux56feux978bux91ccux4e0dux65adux8fdbux6c99ux5b50}

象征：

\begin{itemize}
\tightlist
\item
  生活中看似平凡琐碎的小事
\item
  不断摩擦、积累
\item
  持续消耗个体心理资源
\end{itemize}

老师强调：

\begin{itemize}
\tightlist
\item
  大多数人一生中未必会遭遇巨石般的灾难；
\item
  但更常见的是``鞋里的沙子''------也就是日常生活中不断积累的小困扰、小冲突、小压力。
\end{itemize}

\subsubsection{3.
【举例】中年人的日常``沙子''}\label{ux4e3eux4f8bux4e2dux5e74ux4ebaux7684ux65e5ux5e38ux6c99ux5b50}

老师举例：

\begin{itemize}
\tightlist
\item
  孩子作业写错得离谱（如``800-700=900''）
\item
  父母被保健品推销套路
\item
  伴侣追问``你是不是不爱我了''
\end{itemize}

老师借此说明：

\begin{itemize}
\tightlist
\item
  这些看似日常的小事，同样可能成为心理危机的重要来源。
\end{itemize}

\begin{center}\rule{0.5\linewidth}{0.5pt}\end{center}

\subsection{十二、心理危机原因总分类：外部原因与内部原因}\label{ux5341ux4e8cux5fc3ux7406ux5371ux673aux539fux56e0ux603bux5206ux7c7bux5916ux90e8ux539fux56e0ux4e0eux5185ux90e8ux539fux56e0}

老师把心理危机原因概括为两大类：

\subsubsection{1. 外部原因}\label{ux5916ux90e8ux539fux56e0}

\begin{itemize}
\tightlist
\item
  生态环境（ecological environment）
\item
  社会环境（social environment）
\end{itemize}

\subsubsection{2. 内部原因}\label{ux5185ux90e8ux539fux56e0}

\begin{itemize}
\tightlist
\item
  健康状况（health status）
\item
  人格及应对方式（personality and coping style）
\end{itemize}

接下来老师按这个框架展开。

\begin{center}\rule{0.5\linewidth}{0.5pt}\end{center}

\subsection{十三、生态环境导致的心理危机}\label{ux5341ux4e09ux751fux6001ux73afux5883ux5bfcux81f4ux7684ux5fc3ux7406ux5371ux673a}

\subsubsection{1. 主要类型}\label{ux4e3bux8981ux7c7bux578b}

老师列举生态环境层面的常见危机来源：

\begin{itemize}
\tightlist
\item
  自然灾害：地震、洪水、海啸、雪灾
\item
  传染性疾病：SARS、COVID-19、甲流等
\item
  技术性灾害：飞机坠毁、核设施泄漏、建筑物倒塌等
\end{itemize}

\subsubsection{2. 共同特点}\label{ux5171ux540cux7279ux70b9}

老师强调这些事件的特点：

\begin{itemize}
\tightlist
\item
  多数\textbf{不可预见}
\item
  趋向\textbf{不可预测}
\item
  趋向\textbf{不可控制}
\end{itemize}

\begin{center}\rule{0.5\linewidth}{0.5pt}\end{center}

\subsection{十四、【临床关联】流行病高峰与儿科医疗系统压力}\label{ux5341ux56dbux4e34ux5e8aux5173ux8054ux6d41ux884cux75c5ux9ad8ux5cf0ux4e0eux513fux79d1ux533bux7597ux7cfbux7edfux538bux529b}

老师结合医疗工作经验举例：

\begin{itemize}
\tightlist
\item
  秋冬季通常是儿科最忙的时候；
\item
  新冠放开后某一年甲流大暴发，医院压力巨大。
\end{itemize}

\subsubsection{1.
【举例】老师路过天津老儿童医院的观察}\label{ux4e3eux4f8bux8001ux5e08ux8defux8fc7ux5929ux6d25ux8001ux513fux7ae5ux533bux9662ux7684ux89c2ux5bdf}

老师说自己曾亲眼看到：

\begin{itemize}
\tightlist
\item
  老儿童医院门口排队车辆一直排到友谊路并拐弯；
\item
  现场有人举牌引导停车；
\item
  刚过去的 2025 年，天津市儿童医院曾创下一天接诊 \textbf{13000
  人次}的纪录。
\end{itemize}

\subsubsection{2.
【临床观察】儿科急诊系统承压}\label{ux4e34ux5e8aux89c2ux5bdfux513fux79d1ux6025ux8bcaux7cfbux7edfux627fux538b}

老师还提到：

\begin{itemize}
\tightlist
\item
  自己所在私立医院儿科急诊，晚上看诊也要等 2～3 小时；
\item
  这说明整个医疗体系正在承受严重挤兑。
\end{itemize}

\subsubsection{3.
【临床关联】天津与深圳儿童医院分诊体系对比}\label{ux4e34ux5e8aux5173ux8054ux5929ux6d25ux4e0eux6df1ux5733ux513fux7ae5ux533bux9662ux5206ux8bcaux4f53ux7cfbux5bf9ux6bd4}

老师对比两种模式：

\paragraph{天津常见情况}\label{ux5929ux6d25ux5e38ux89c1ux60c5ux51b5}

\begin{itemize}
\tightlist
\item
  分诊处多为护士；
\item
  先分科，再等待就诊；
\item
  见到医生后才开检查；
\item
  检查后又得排队，效率低。
\end{itemize}

\paragraph{深圳儿童医院模式}\label{ux6df1ux5733ux513fux7ae5ux533bux9662ux6a21ux5f0f}

\begin{itemize}
\tightlist
\item
  分诊处由儿科医生先做初步判断；
\item
  提前开必要的检查化验；
\item
  家长先带孩子做检查；
\item
  回来时差不多轮到正式就诊；
\item
  医生能更快根据检查结果做出初步诊断并开药。
\end{itemize}

老师认为这种模式效率更高，也更有利于患儿尽早接受处理。

\begin{center}\rule{0.5\linewidth}{0.5pt}\end{center}

\subsection{十五、【举例】重大生态灾难与创伤反应}\label{ux5341ux4e94ux4e3eux4f8bux91cdux5927ux751fux6001ux707eux96beux4e0eux521bux4f24ux53cdux5e94}

\subsubsection{1.
唐山大地震（1976）}\label{ux5510ux5c71ux5927ux5730ux97071976}

老师提到：

\begin{itemize}
\tightlist
\item
  唐山地震将整座城市夷为平地；
\item
  铁轨都被拧成麻花，显示其破坏力极大。
\end{itemize}

老师转述老一辈人的描述：

\begin{itemize}
\tightlist
\item
  地震后几个月，幸存者之间见面会互相问``你们家几个''，意思是``家里死了几个人''。
\end{itemize}

老师指出，这种\textbf{情感麻木、淡漠}是典型的：

\begin{itemize}
\tightlist
\item
  \textbf{创伤后应激障碍（post-traumatic stress disorder,
  PTSD）}表现之一。
\end{itemize}

\subsubsection{2. 切尔诺贝利核泄漏（Chernobyl
disaster）}\label{ux5207ux5c14ux8bfaux8d1dux5229ux6838ux6cc4ux6f0fchernobyl-disaster}

老师讲到：

\begin{itemize}
\tightlist
\item
  1986 年 4 月 26 日，乌克兰切尔诺贝利 4 号反应堆爆炸；
\item
  这是人类历史上最严重的核泄漏事故之一。
\end{itemize}

老师重点强调了灾后处置中的极端危险：

\begin{itemize}
\tightlist
\item
  直升机向反应堆上空投硼砂/硼酸和沙土；
\item
  辐射极强，机器与人都难以承受；
\item
  机器人失效后，只能招募人员身穿厚重铅防护服上屋顶人工铲除放射性碎片；
\item
  每个人暴露时间必须极短。
\end{itemize}

老师借此说明：

\begin{itemize}
\tightlist
\item
  技术性灾害不仅是``灾难''，还可能迅速转化为严重的群体性心理创伤。
\end{itemize}

\subsubsection{3. 【举例】SARS（2003）}\label{ux4e3eux4f8bsars2003}

老师回忆 2003 年非典：

\begin{itemize}
\tightlist
\item
  当时防疫条件远不如今天；
\item
  甚至没有今天熟悉的 \textbf{N95 口罩}概念；
\item
  防疫常靠多层纱布口罩。
\end{itemize}

老师还提到：

\begin{itemize}
\tightlist
\item
  当年许多医学院研究生往返医院和学校做课题；
\item
  疫情影响下，很多研究和毕业安排被打乱；
\item
  学校甚至延后了毕业答辩时间。
\end{itemize}

\subsubsection{4. 汶川地震（2008）}\label{ux6c76ux5dddux5730ux97072008}

老师简单提到：

\begin{itemize}
\tightlist
\item
  汶川地震十分惨烈；
\item
  大家虽可能年纪太小印象不深，但这是重要的国家级灾难事件。
\end{itemize}

\subsubsection{5.
东日本大地震（2011）与福岛核泄漏}\label{ux4e1cux65e5ux672cux5927ux5730ux97072011ux4e0eux798fux5c9bux6838ux6cc4ux6f0f}

老师提到：

\begin{itemize}
\tightlist
\item
  2011 年 3 月 11 日东日本大地震是人类观测史上极严重的地震之一；
\item
  岩手、宫城、福岛受灾尤重；
\item
  此后又引发\textbf{福岛核电站泄漏}。
\end{itemize}

\begin{center}\rule{0.5\linewidth}{0.5pt}\end{center}

\subsection{十六、【举例】福岛核泄漏后的``抢盐''现象：群体性荒谬反应}\label{ux5341ux516dux4e3eux4f8bux798fux5c9bux6838ux6cc4ux6f0fux540eux7684ux62a2ux76d0ux73b0ux8c61ux7fa4ux4f53ux6027ux8352ux8c2cux53cdux5e94}

老师借福岛事件讲了一个经典社会心理例子。

\subsubsection{1.
当年出现``抢盐潮''}\label{ux5f53ux5e74ux51faux73b0ux62a2ux76d0ux6f6e}

老师回忆：

\begin{itemize}
\tightlist
\item
  2011 年春天，很多地方的盐被抢购一空；
\item
  有人一口气囤很多盐，但其实自己也说不清理由。
\end{itemize}

\subsubsection{2.
谣言的演变过程}\label{ux8c23ux8a00ux7684ux6f14ux53d8ux8fc7ux7a0b}

老师描述谣言是如何``进化''的：

\begin{itemize}
\tightlist
\item
  最初：担心海水受污染，所以海盐不安全；
\item
  后来进一步演变为：\textbf{吃加碘盐可以防辐射}。
\end{itemize}

\subsubsection{3.
【课堂提醒】从医学角度看待``吃盐防辐射''}\label{ux8bfeux5802ux63d0ux9192ux4eceux533bux5b66ux89d2ux5ea6ux770bux5f85ux5403ux76d0ux9632ux8f90ux5c04}

老师提醒医学生：

\begin{itemize}
\tightlist
\item
  理论上碘在某些特定情况下与甲状腺放射性核素防护有关；
\item
  但\textbf{``抛开剂量谈疗效''是没有意义的}。
\end{itemize}

老师举出极端对比：

\begin{itemize}
\tightlist
\item
  若真靠食盐摄入碘达到防护剂量，所需摄入量极大；
\item
  而盐本身过量摄入早就先带来严重风险。
\end{itemize}

\subsubsection{4.
【重点】心理危机可用荒谬形式爆发}\label{ux91cdux70b9ux5fc3ux7406ux5371ux673aux53efux7528ux8352ux8c2cux5f62ux5f0fux7206ux53d1}

老师借此说明：

\begin{itemize}
\tightlist
\item
  在灾难和群体焦虑背景下，
\item
  人会以非常荒谬、失去判断力的方式做出反应；
\item
  这也是危机状态的一种社会性表现。
\end{itemize}

\begin{center}\rule{0.5\linewidth}{0.5pt}\end{center}

\subsection{十七、【举例】812
天津港爆炸事故}\label{ux5341ux4e03ux4e3eux4f8b812-ux5929ux6d25ux6e2fux7206ux70b8ux4e8bux6545}

老师接着讲到较近的本地重大事故：

\begin{itemize}
\tightlist
\item
  天津港危险化学品违规储存并起火；
\item
  早期进入现场的消防员不清楚仓库存放物质；
\item
  喷水降温诱发剧烈爆炸；
\item
  造成多名年轻消防员牺牲。
\end{itemize}

老师还提到：

\begin{itemize}
\tightlist
\item
  爆炸点附近小区玻璃几乎全部震碎；
\item
  大量伤员被送往塘沽及市区各医院；
\item
  那晚医疗系统也面临突发性巨大压力。
\end{itemize}

老师借此再次说明：

\begin{itemize}
\tightlist
\item
  技术性灾难往往来得突然、破坏巨大，也会迅速造成严重的群体心理冲击。
\end{itemize}

\begin{center}\rule{0.5\linewidth}{0.5pt}\end{center}

\subsection{十八、【举例】新冠疫情及其后果}\label{ux5341ux516bux4e3eux4f8bux65b0ux51a0ux75abux60c5ux53caux5176ux540eux679c}

老师提到：

\begin{itemize}
\tightlist
\item
  2020 年 1 月新冠疫情爆发；
\item
  在中国持续约三年；
\item
  到 2022 年底才全面放开。
\end{itemize}

\subsubsection{1.
放开后的群体性影响}\label{ux653eux5f00ux540eux7684ux7fa4ux4f53ux6027ux5f71ux54cd}

老师的课堂描述包括：

\begin{itemize}
\tightlist
\item
  ``每个人是自己健康第一责任人''；
\item
  整个过程某种意义上变成了全民免疫过程；
\item
  放开初期带来显著的群体性冲击与社会系统压力。
\end{itemize}

\subsubsection{2.
老师的社会观察}\label{ux8001ux5e08ux7684ux793eux4f1aux89c2ux5bdf}

老师用天津本地的丧葬行业举例，强调：

\begin{itemize}
\tightlist
\item
  放开后的冬天，死亡相关事务陡增；
\item
  说明疫情冲击不只是医学层面的，也深刻影响社会心理与生活秩序。
\end{itemize}

\begin{center}\rule{0.5\linewidth}{0.5pt}\end{center}

\subsection{十九、社会环境导致的心理危机}\label{ux5341ux4e5dux793eux4f1aux73afux5883ux5bfcux81f4ux7684ux5fc3ux7406ux5371ux673a}

老师接下来进入\textbf{社会环境（social
environment）}层面，列出常见来源：

\begin{itemize}
\tightlist
\item
  恐怖袭击
\item
  人质事件
\item
  暴乱
\item
  经济危机
\item
  社会变革
\item
  学业压力
\item
  工作压力
\item
  家庭矛盾
\end{itemize}

\begin{center}\rule{0.5\linewidth}{0.5pt}\end{center}

\subsection{二十、【课堂举例】恐怖袭击、人质事件与战争}\label{ux4e8cux5341ux8bfeux5802ux4e3eux4f8bux6050ux6016ux88adux51fbux4ebaux8d28ux4e8bux4ef6ux4e0eux6218ux4e89}

\subsubsection{1. 9·11 事件}\label{ux4e8bux4ef6}

老师提到 9·11
事件，并说课堂上有人还会从纪录片角度提出若干``疑点''与质疑。

这里的重点不是历史定论，而是老师借此强调：

\begin{itemize}
\tightlist
\item
  \textbf{重大恐怖袭击事件会构成强烈社会性危机来源}；
\item
  其后续政治、战争、社会舆论与群体心理反应，会进一步放大危机。
\end{itemize}

\begin{quote}
注：老师课堂中提到的部分纪录片内容属于其转述的``质疑性观点/材料''，可作为课堂举例理解，不宜直接当作已证实结论。
\end{quote}

\subsubsection{2.
别斯兰人质事件（2004）}\label{ux522bux65afux5170ux4ebaux8d28ux4e8bux4ef62004}

老师提到：

\begin{itemize}
\tightlist
\item
  俄罗斯北奥塞梯别斯兰第一小学开学典礼上发生人质事件；
\item
  武装分子劫持师生；
\item
  最终造成大量人质死亡，绝大多数为儿童。
\end{itemize}

老师借此说明：

\begin{itemize}
\tightlist
\item
  人质事件、恐怖袭击不仅造成直接伤亡，也会引起深重的社会创伤。
\end{itemize}

\subsubsection{3.
【时事举例】伊朗局势与霍尔木兹海峡}\label{ux65f6ux4e8bux4e3eux4f8bux4f0aux6717ux5c40ux52bfux4e0eux970dux5c14ux6728ux5179ux6d77ux5ce1}

老师以近期中东局势为例，强调：

\begin{itemize}
\tightlist
\item
  国际局势并不一定``离我们很远''；
\item
  霍尔木兹海峡是全球重要能源运输通道；
\item
  国际冲突会影响石油、天然气供应；
\item
  最终通过油价、运输成本、物价等方式影响普通人的日常生活。
\end{itemize}

老师借此说明：

\begin{itemize}
\tightlist
\item
  国际政治冲突也可能成为普通人的心理压力源和危机来源。
\end{itemize}

\begin{center}\rule{0.5\linewidth}{0.5pt}\end{center}

\subsection{二十一、【举例】恶性无差别伤害事件与社会安全感受损}\label{ux4e8cux5341ux4e00ux4e3eux4f8bux6076ux6027ux65e0ux5deeux522bux4f24ux5bb3ux4e8bux4ef6ux4e0eux793eux4f1aux5b89ux5168ux611fux53d7ux635f}

老师列举了一系列恶性无差别伤害事件，说明社会环境如何成为心理危机的重要来源。

\subsubsection{1.
福建南平实验小学案（2010）}\label{ux798fux5efaux5357ux5e73ux5b9eux9a8cux5c0fux5b66ux68482010}

老师提到：

\begin{itemize}
\tightlist
\item
  该案针对的是校门口聚集的小学生；
\item
  之后全国出现多起模仿式校园无差别伤害事件。
\end{itemize}

\subsubsection{2.
老师的个人记忆}\label{ux8001ux5e08ux7684ux4e2aux4ebaux8bb0ux5fc6}

老师回忆：

\begin{itemize}
\tightlist
\item
  当年很多小学、幼儿园上下学时段，民警持枪站岗；
\item
  社会整体安全感明显受到冲击。
\end{itemize}

\subsubsection{3.
近年的无差别袭击事件}\label{ux8fd1ux5e74ux7684ux65e0ux5deeux522bux88adux51fbux4e8bux4ef6}

老师又举例：

\begin{itemize}
\tightlist
\item
  沈阳铁西区多起无差别伤害案件
\item
  山东德州驾车冲撞放学小学生
\item
  珠海体育中心驾车冲撞事件
\end{itemize}

老师借此说明：

\begin{itemize}
\tightlist
\item
  这类事件会显著破坏公众的安全感；
\item
  社会环境中的不确定性、暴力性也会构成心理危机来源。
\end{itemize}

\begin{center}\rule{0.5\linewidth}{0.5pt}\end{center}

\subsection{二十二、学业与就业压力：多数人更常见的危机来源}\label{ux4e8cux5341ux4e8cux5b66ux4e1aux4e0eux5c31ux4e1aux538bux529bux591aux6570ux4ebaux66f4ux5e38ux89c1ux7684ux5371ux673aux6765ux6e90}

老师强调：

\begin{itemize}
\tightlist
\item
  大多数人未必遇到战争、恐袭、灾难；
\item
  但更常遇到的是\textbf{就业、工作、经济与家庭压力}。
\end{itemize}

\subsubsection{1.
【课堂提醒】医学并不意味着``天然好就业''}\label{ux8bfeux5802ux63d0ux9192ux533bux5b66ux5e76ux4e0dux610fux5473ux7740ux5929ux7136ux597dux5c31ux4e1a}

老师直言：

\begin{itemize}
\tightlist
\item
  学医不等于一定轻松找到工作；
\item
  医院更希望招来就能独当一面的人；
\item
  不希望从零开始慢慢培养。
\end{itemize}

\subsubsection{2.
【学习建议】医学生在校期间就要完成从``学生''到``医生''的转换}\label{ux5b66ux4e60ux5efaux8baeux533bux5b66ux751fux5728ux6821ux671fux95f4ux5c31ux8981ux5b8cux6210ux4eceux5b66ux751fux5230ux533bux751fux7684ux8f6cux6362}

老师强调：

\begin{itemize}
\tightlist
\item
  大家在校期间就必须尽快完成这一角色转变；
\item
  否则到临床后会很被动。
\end{itemize}

\begin{center}\rule{0.5\linewidth}{0.5pt}\end{center}

\subsection{二十三、【临床关联】医生职业的工作压力}\label{ux4e8cux5341ux4e09ux4e34ux5e8aux5173ux8054ux533bux751fux804cux4e1aux7684ux5de5ux4f5cux538bux529b}

老师用医院工作流程举例，说明工作压力如何成为危机来源。

\subsubsection{1.
医生的工作节奏很辛苦}\label{ux533bux751fux7684ux5de5ux4f5cux8282ux594fux5f88ux8f9bux82e6}

老师举总医院为例：

\begin{itemize}
\tightlist
\item
  早晨 7:30 交班；
\item
  住院医生应更早到达科室；
\item
  到岗后要先看病人、看夜间病情变化、看临时医嘱、看检查结果。
\end{itemize}

\subsubsection{2.
【课堂提醒】汇报病例不是照着病历念}\label{ux8bfeux5802ux63d0ux9192ux6c47ux62a5ux75c5ux4f8bux4e0dux662fux7167ux7740ux75c5ux5386ux5ff5}

老师强调：

\begin{itemize}
\tightlist
\item
  汇报病例要把主诉、现病史等内容真正掌握；
\item
  重要化验值也要记住；
\item
  不能只是拿着病历照本宣科。
\end{itemize}

\subsubsection{3.
医生常常下班不由自己}\label{ux533bux751fux5e38ux5e38ux4e0bux73edux4e0dux7531ux81eaux5df1}

老师举例：

\begin{itemize}
\tightlist
\item
  快下班时急诊再收一个病人；
\item
  就意味着还要完成查体、化验、输液、病历等一整套工作。
\end{itemize}

\subsubsection{4. 老师的概括}\label{ux8001ux5e08ux7684ux6982ux62ec}

老师用一句很直白的话形容高强度职业：

\begin{itemize}
\tightlist
\item
  \textbf{``人肉换猪肉。''}
\end{itemize}

意思是：

\begin{itemize}
\tightlist
\item
  高收入背后往往是用身体、健康、时间、精力去换的。
\end{itemize}

\begin{center}\rule{0.5\linewidth}{0.5pt}\end{center}

\subsection{二十四、【举例】996、007
与高压劳动}\label{ux4e8cux5341ux56dbux4e3eux4f8b996007-ux4e0eux9ad8ux538bux52b3ux52a8}

老师把医生工作进一步放到更广泛的劳动环境里讨论。

\subsubsection{1. 996}\label{section}

\begin{itemize}
\tightlist
\item
  早 9 点到晚 9 点
\item
  一周工作 6 天
\end{itemize}

\subsubsection{2. 007}\label{section-1}

\begin{itemize}
\tightlist
\item
  0 点到 0 点
\item
  一周 7 天无缝衔接
\end{itemize}

\subsubsection{3. 老师观点}\label{ux8001ux5e08ux89c2ux70b9}

老师借此说明：

\begin{itemize}
\tightlist
\item
  无论是互联网大厂还是其他高压行业，
\item
  高收入背后都可能意味着严重的身体透支与心理消耗。
\end{itemize}

\begin{center}\rule{0.5\linewidth}{0.5pt}\end{center}

\subsection{二十五、【举例】富士康事件与
EAP}\label{ux4e8cux5341ux4e94ux4e3eux4f8bux5bccux58ebux5eb7ux4e8bux4ef6ux4e0e-eap}

老师又举富士康跳楼事件，说明制度化高压环境如何导致心理危机。

\subsubsection{1. EAP}\label{eap}

老师提到企业后来引入：

\begin{itemize}
\tightlist
\item
  \textbf{员工帮助计划（Employee Assistance Program, EAP）}
\end{itemize}

\subsubsection{2.
老师指出的问题}\label{ux8001ux5e08ux6307ux51faux7684ux95eeux9898}

老师认为，富士康暴露出的并不只是个体心理问题，而是制度性、环境性问题，例如：

\begin{itemize}
\tightlist
\item
  车间``静音模式''
\item
  吃饭时间极短
\item
  工作与生活高度机械化
\item
  宿舍工友因作息完全错开，彼此甚至连名字都说不全
\end{itemize}

\subsubsection{3. 结论}\label{ux7ed3ux8bba-2}

老师借此强调：

\begin{itemize}
\tightlist
\item
  心理危机并不总是个体``脆弱''；
\item
  很多时候是环境本身在持续制造危机。
\end{itemize}

\begin{center}\rule{0.5\linewidth}{0.5pt}\end{center}

\subsection{二十六、经济压力与现实生活困境}\label{ux4e8cux5341ux516dux7ecfux6d4eux538bux529bux4e0eux73b0ux5b9eux751fux6d3bux56f0ux5883}

老师接下来谈经济压力。

\subsubsection{1.
收入与生活压力}\label{ux6536ux5165ux4e0eux751fux6d3bux538bux529b}

老师提到天津市职工平均工资数据，并借此说明：

\begin{itemize}
\tightlist
\item
  很多岗位想月薪过万并不容易；
\item
  往往要付出极大的体力、时长或劳动强度。
\end{itemize}

老师列举：

\begin{itemize}
\tightlist
\item
  快递员
\item
  外卖员
\item
  装卸工
\item
  网约车司机
\item
  电子厂工人等
\end{itemize}

意在说明：

\begin{itemize}
\tightlist
\item
  经济压力是非常现实、非常普遍的危机来源。
\end{itemize}

\subsubsection{2. 房地产问题}\label{ux623fux5730ux4ea7ux95eeux9898}

老师提到：

\begin{itemize}
\tightlist
\item
  房企爆雷、楼盘烂尾
\item
  房价下跌
\item
  买房后跌破买价
\item
  贷款压力巨大
\end{itemize}

老师强调：

\begin{itemize}
\tightlist
\item
  对很多家庭来说，这会成为持续性的心理危机。
\end{itemize}

\subsubsection{3.
【举例】真实``蜗居''}\label{ux4e3eux4f8bux771fux5b9eux8717ux5c45}

老师还提到：

\begin{itemize}
\tightlist
\item
  日本胶囊酒店
\item
  日本网吧常住人群
\end{itemize}

用来说明：

\begin{itemize}
\tightlist
\item
  经济压力会直接转化为居住困境、尊严受损和长期压抑。
\end{itemize}

\begin{center}\rule{0.5\linewidth}{0.5pt}\end{center}

\subsection{二十七、【举例】股市暴跌与金融压力}\label{ux4e8cux5341ux4e03ux4e3eux4f8bux80a1ux5e02ux66b4ux8dccux4e0eux91d1ux878dux538bux529b}

老师讲到股市暴跌引发的心理崩溃案例。

\subsubsection{1. 杠杆风险}\label{ux6760ux6746ux98ceux9669}

老师提醒：

\begin{itemize}
\tightlist
\item
  不要轻易加杠杆；
\item
  杠杆放大收益，也同样放大亏损；
\item
  一旦触发强制平仓（forced liquidation），可能``什么都没了''。
\end{itemize}

\subsubsection{2. 【举例】2015
年股灾}\label{ux4e3eux4f8b2015-ux5e74ux80a1ux707e}

老师回忆自己亲眼见到与股灾相关的极端后果，借此说明：

\begin{itemize}
\tightlist
\item
  金融风险也会迅速演变成心理危机、甚至极端事件。
\end{itemize}

\begin{center}\rule{0.5\linewidth}{0.5pt}\end{center}

\subsection{二十八、婚恋压力与人际失败}\label{ux4e8cux5341ux516bux5a5aux604bux538bux529bux4e0eux4ebaux9645ux5931ux8d25}

老师提到，当代年轻人与中年人的另一个高频危机来源是：

\begin{itemize}
\tightlist
\item
  找对象难
\item
  婚恋压力大
\end{itemize}

\subsubsection{1.
【举例】天津中心公园``婚恋市场''}\label{ux4e3eux4f8bux5929ux6d25ux4e2dux5fc3ux516cux56edux5a5aux604bux5e02ux573a}

老师用``婚姻人口期货交易市场''作调侃式描述，指出：

\begin{itemize}
\tightlist
\item
  许多相亲场景中，真正到场的反而不是当事人，而是父母；
\item
  从``盘面''上看，供求失衡明显。
\end{itemize}

\subsubsection{2.
老师借出生率变化说明}\label{ux8001ux5e08ux501fux51faux751fux7387ux53d8ux5316ux8bf4ux660e}

老师提到近年出生率下降、节日婚车减少等现象，想说明：

\begin{itemize}
\tightlist
\item
  婚恋困难已成为现实而广泛的社会问题。
\end{itemize}

\subsubsection{3.
【举例】由单身羞辱引发的极端事件}\label{ux4e3eux4f8bux7531ux5355ux8eabux7f9eux8fb1ux5f15ux53d1ux7684ux6781ux7aefux4e8bux4ef6}

老师提到：

\begin{itemize}
\tightlist
\item
  有人因长期单身、被周围人轻视嘲讽而产生强烈报复心理，最终实施投毒。
\end{itemize}

老师借此强调：

\begin{itemize}
\tightlist
\item
  不要轻视婚恋挫败、羞辱感带来的心理打击；
\item
  有时这会成为危机爆发的诱因。
\end{itemize}

\begin{center}\rule{0.5\linewidth}{0.5pt}\end{center}

\subsection{二十九、家庭矛盾：最常见也最容易被忽视的危机来源}\label{ux4e8cux5341ux4e5dux5bb6ux5eadux77dbux76feux6700ux5e38ux89c1ux4e5fux6700ux5bb9ux6613ux88abux5ffdux89c6ux7684ux5371ux673aux6765ux6e90}

老师指出：

\begin{itemize}
\tightlist
\item
  结不了婚的愁，结了婚也可能更愁；
\item
  家庭中的冲突非常常见，也常常最能消耗一个人。
\end{itemize}

\subsubsection{1.
【举例】辅导作业导致家长崩溃}\label{ux4e3eux4f8bux8f85ux5bfcux4f5cux4e1aux5bfcux81f4ux5bb6ux957fux5d29ux6e83}

老师举例：

\begin{itemize}
\tightlist
\item
  小女儿把``800-700''算成``900''；
\item
  父亲情绪崩溃，甚至哭求警察``把自己抓走''。
\end{itemize}

老师借此说明：

\begin{itemize}
\tightlist
\item
  养育、教育、家庭责任会成为很多中年人崩溃的临界点。
\end{itemize}

\subsubsection{2.
网络段子中的家庭情绪调节}\label{ux7f51ux7edcux6bb5ux5b50ux4e2dux7684ux5bb6ux5eadux60c5ux7eeaux8c03ux8282}

老师引用流行说法，说明现实中很多家长会靠自我调侃撑住：

\begin{itemize}
\tightlist
\item
  ``亲生的、亲生的\ldots\ldots{}''
\item
  ``孩子爸是我选的\ldots\ldots{}''
\end{itemize}

虽然带有玩笑成分，但本质上反映的还是：

\begin{itemize}
\tightlist
\item
  家庭教育压力真实存在。
\end{itemize}

\subsubsection{3.
【学习建议/教育理念】不要总替孩子改作业}\label{ux5b66ux4e60ux5efaux8baeux6559ux80b2ux7406ux5ff5ux4e0dux8981ux603bux66ffux5b69ux5b50ux6539ux4f5cux4e1a}

老师转述一位小学数学老师的教育理念：

\begin{itemize}
\tightlist
\item
  家长不要总把孩子写错的作业改正确；
\item
  保留原始错误，老师才能判断孩子的问题出在哪里；
\item
  过度代改作业，反而遮蔽了真正的教育信息。
\end{itemize}

这段虽然是家庭教育话题，但也体现出：

\begin{itemize}
\tightlist
\item
  正确处理家庭中的压力与角色边界，是降低危机的重要方式。
\end{itemize}

\begin{center}\rule{0.5\linewidth}{0.5pt}\end{center}

\subsection{三十、健康状况与心理危机}\label{ux4e09ux5341ux5065ux5eb7ux72b6ux51b5ux4e0eux5fc3ux7406ux5371ux673a}

老师接着讲内部原因中的\textbf{健康状况（health status）}。

\subsubsection{1. 相关类型}\label{ux76f8ux5173ux7c7bux578b}

包括：

\begin{itemize}
\tightlist
\item
  不治之症
\item
  慢性疾病（chronic disease）
\item
  躯体伤残
\item
  功能障碍
\item
  精神疾病（mental illness）
\item
  行为问题（behavioral problems）
\end{itemize}

\subsubsection{2.
【举例】艾滋病防治宣传}\label{ux4e3eux4f8bux827eux6ecbux75c5ux9632ux6cbbux5ba3ux4f20}

老师提到学校厕所中的公益海报``防艾''，借此提醒：

\begin{itemize}
\tightlist
\item
  艾滋病（AIDS / HIV infection）并不遥远；
\item
  健康问题本身就是心理危机的重要来源。
\end{itemize}

\subsubsection{3.
【举例】隐瞒精神病史引发的极端案件}\label{ux4e3eux4f8bux9690ux7792ux7cbeux795eux75c5ux53f2ux5f15ux53d1ux7684ux6781ux7aefux6848ux4ef6}

老师提到 2021 年天津一例轰动案件：

\begin{itemize}
\tightlist
\item
  当事人有精神病史，却在恋爱、结婚过程中隐瞒；
\item
  甚至自行停药；
\item
  后来精神疾病复发，在冲突中对妻子实施严重暴力。
\end{itemize}

老师借此提醒：

\begin{itemize}
\tightlist
\item
  精神疾病与停药、隐瞒病史、婚恋关系破裂等，都是高风险危机场景。
\end{itemize}

\begin{center}\rule{0.5\linewidth}{0.5pt}\end{center}

\subsection{三十一、人格与应对方式：心理危机的内部基础}\label{ux4e09ux5341ux4e00ux4ebaux683cux4e0eux5e94ux5bf9ux65b9ux5f0fux5fc3ux7406ux5371ux673aux7684ux5185ux90e8ux57faux7840}

老师进入最后一个大部分：\textbf{人格（personality）与应对方式（coping
style）}。

\subsubsection{1. 人格的定义}\label{ux4ebaux683cux7684ux5b9aux4e49}

老师说：

\begin{itemize}
\tightlist
\item
  人格过去也叫``个性''；
\item
  指一个人较固定的行为模式，以及日常待人处事的习惯方式；
\item
  是全部心理特征的综合。
\end{itemize}

\subsubsection{2. 人格的特征}\label{ux4ebaux683cux7684ux7279ux5f81}

老师总结人格具有以下特点：

\paragraph{（1）社会性}\label{ux793eux4f1aux6027}

\begin{itemize}
\tightlist
\item
  不同社会中，人格表现受到不同社会规范影响。
\end{itemize}

\paragraph{（2）独特性}\label{ux72ecux7279ux6027}

\begin{itemize}
\tightlist
\item
  每个人都是独特个体。
\end{itemize}

\paragraph{（3）稳定性}\label{ux7a33ux5b9aux6027}

\begin{itemize}
\tightlist
\item
  人格短时间内不会发生巨大变化；
\item
  若一个人``昨天和今天判若两人''，就需要警惕问题。
\end{itemize}

\paragraph{（4）可塑性}\label{ux53efux5851ux6027}

\begin{itemize}
\tightlist
\item
  人格并非完全固定；
\item
  会随着时间和经历逐渐变化。
\end{itemize}

\begin{center}\rule{0.5\linewidth}{0.5pt}\end{center}

\subsection{三十二、人格的构成：气质、性格、能力}\label{ux4e09ux5341ux4e8cux4ebaux683cux7684ux6784ux6210ux6c14ux8d28ux6027ux683cux80fdux529b}

老师把人格构成分为三部分：

\begin{enumerate}
\def\labelenumi{\arabic{enumi}.}
\tightlist
\item
  \textbf{气质（temperament）}：先天部分，与遗传素质有关
\item
  \textbf{性格（character）}：后天部分，与生活环境和教育有关
\item
  \textbf{能力（ability）}：与个人实践活动、智力发展等相关
\end{enumerate}

\subsubsection{【比喻】}\label{ux6bd4ux55bb}

老师把人格比作一幅名人字画：

\begin{itemize}
\tightlist
\item
  气质 = 画纸/画布（先天基础）
\item
  性格 = 画了什么（后天塑造）
\item
  能力 = 画得怎么样、值不值钱（实践水平）
\end{itemize}

\begin{center}\rule{0.5\linewidth}{0.5pt}\end{center}

\subsection{三十三、【复习】四种基本气质类型}\label{ux4e09ux5341ux4e09ux590dux4e60ux56dbux79cdux57faux672cux6c14ux8d28ux7c7bux578b}

老师讲了经典的``四个人去看戏''故事，说明四种基本气质。

\subsubsection{1. 胆汁质（choleric
temperament）}\label{ux80c6ux6c41ux8d28choleric-temperament}

特点：

\begin{itemize}
\tightlist
\item
  外向
\item
  情绪不稳定
\end{itemize}

情境反应：

\begin{itemize}
\tightlist
\item
  因为迟到被拦，直接和检票员吵起来、打起来。
\end{itemize}

\subsubsection{2. 多血质（sanguine
temperament）}\label{ux591aux8840ux8d28sanguine-temperament}

特点：

\begin{itemize}
\tightlist
\item
  外向
\item
  情绪稳定
\end{itemize}

情境反应：

\begin{itemize}
\tightlist
\item
  趁别人混乱时自己溜进去。
\end{itemize}

\subsubsection{3. 黏液质（phlegmatic
temperament）}\label{ux9ecfux6db2ux8d28phlegmatic-temperament}

特点：

\begin{itemize}
\tightlist
\item
  内向
\item
  情绪稳定
\end{itemize}

情境反应：

\begin{itemize}
\tightlist
\item
  既然说中场休息能进，那就去咖啡厅等会儿。
\end{itemize}

\subsubsection{4. 抑郁质（melancholic
temperament）}\label{ux6291ux90c1ux8d28melancholic-temperament}

特点：

\begin{itemize}
\tightlist
\item
  内向
\item
  情绪不稳定
\end{itemize}

情境反应：

\begin{itemize}
\tightlist
\item
  觉得``今天干啥都不顺''，索性不看了，撕票回家。
\end{itemize}

老师借此说明：

\begin{itemize}
\tightlist
\item
  这是人格中\textbf{先天气质}的差异。
\end{itemize}

\begin{center}\rule{0.5\linewidth}{0.5pt}\end{center}

\subsection{三十四、性格与环境：后天塑造的重要性}\label{ux4e09ux5341ux56dbux6027ux683cux4e0eux73afux5883ux540eux5929ux5851ux9020ux7684ux91cdux8981ux6027}

老师强调：

\begin{itemize}
\tightlist
\item
  后天性格受从小到大接触过的人的影响；
\item
  ``近朱者赤，近墨者黑''；
\item
  ``孟母三迁''本质上也是为了环境塑造。
\end{itemize}

老师指出：

\begin{itemize}
\tightlist
\item
  今天家长拼学区房，本质上也是在争取更好的环境塑造条件。
\end{itemize}

\begin{center}\rule{0.5\linewidth}{0.5pt}\end{center}

\subsection{三十五、人格发展的特点}\label{ux4e09ux5341ux4e94ux4ebaux683cux53d1ux5c55ux7684ux7279ux70b9}

老师总结：

\subsubsection{1.
随年龄呈抛物线式变化}\label{ux968fux5e74ux9f84ux5448ux629bux7269ux7ebfux5f0fux53d8ux5316}

\begin{itemize}
\tightlist
\item
  人格逐渐发展成熟；
\item
  到一定阶段后可能走向退化；
\item
  极端情况下会涉及痴呆等问题。
\end{itemize}

\subsubsection{2. 变化速度慢}\label{ux53d8ux5316ux901fux5ea6ux6162}

\begin{itemize}
\tightlist
\item
  短期效应一般不明显；
\item
  头天和今天相比，不会突然发生根本变化。
\end{itemize}

\subsubsection{3.
受环境影响很大}\label{ux53d7ux73afux5883ux5f71ux54cdux5f88ux5927}

老师强调：

\begin{itemize}
\tightlist
\item
  经常接触社会、与不同人打交道，会促进人格成熟；
\item
  若总是不接触社会、一意孤行，则人格更容易出现问题。
\end{itemize}

\begin{center}\rule{0.5\linewidth}{0.5pt}\end{center}

\subsection{三十六、人格素质、人格问题与人格障碍}\label{ux4e09ux5341ux516dux4ebaux683cux7d20ux8d28ux4ebaux683cux95eeux9898ux4e0eux4ebaux683cux969cux788d}

\subsubsection{1.
人格素质本身无绝对好坏}\label{ux4ebaux683cux7d20ux8d28ux672cux8eabux65e0ux7eddux5bf9ux597dux574f}

老师认为：

\begin{itemize}
\tightlist
\item
  各种人格素质都有积极面和消极面；
\item
  关键看比例是否均衡。
\end{itemize}

\subsubsection{2. 【概念纠正】强迫特质 ≠
强迫症}\label{ux6982ux5ff5ux7ea0ux6b63ux5f3aux8febux7279ux8d28-ux5f3aux8febux75c7}

老师用手术举例说明：

\begin{itemize}
\tightlist
\item
  手术结束前必须清点器械、敷料、针线、纱布等；
\item
  这种谨慎、反复核对是必要的，不是病态。
\end{itemize}

据此可理解为：

\begin{itemize}
\tightlist
\item
  \textbf{强迫性人格特质（obsessive-compulsive traits）}不等于
  \textbf{强迫症（Obsessive-Compulsive Disorder, OCD）}
\end{itemize}

\subsubsection{3.
健康人格的特点}\label{ux5065ux5eb7ux4ebaux683cux7684ux7279ux70b9}

老师概括为：

\begin{itemize}
\tightlist
\item
  相对均衡
\item
  中庸
\item
  中正
\item
  平和
\end{itemize}

\subsubsection{4. 人格问题/人格障碍（personality
disorder）}\label{ux4ebaux683cux95eeux9898ux4ebaux683cux969cux788dpersonality-disorder}

老师指出：

\begin{itemize}
\tightlist
\item
  当某种人格素质比例过高，导致整体人格结构失衡时，
\item
  就可能发展为人格问题甚至人格障碍。
\end{itemize}

\begin{center}\rule{0.5\linewidth}{0.5pt}\end{center}

\subsection{三十七、【举例】超前消费、网贷与人格问题}\label{ux4e09ux5341ux4e03ux4e3eux4f8bux8d85ux524dux6d88ux8d39ux7f51ux8d37ux4e0eux4ebaux683cux95eeux9898}

老师提到：

\begin{itemize}
\tightlist
\item
  过去有``卡奴''，现在则更多表现为网贷滚债；
\item
  这种反复借贷、超前消费、缺乏还款能力评估的行为，
\item
  不只是经济问题，也体现出某种人格和应对方式的问题。
\end{itemize}

\subsubsection{【课堂提醒】}\label{ux8bfeux5802ux63d0ux9192}

老师顺带解释了为什么辅导员要排查网贷：

\begin{itemize}
\tightlist
\item
  不只是怕欠钱；
\item
  更是为了预防由债务、人际、羞耻感等引发的心理危机事件。
\end{itemize}

\begin{center}\rule{0.5\linewidth}{0.5pt}\end{center}

\subsection{三十八、【举例】马加爵案：人格长期压抑与极端爆发}\label{ux4e09ux5341ux516bux4e3eux4f8bux9a6cux52a0ux7235ux6848ux4ebaux683cux957fux671fux538bux6291ux4e0eux6781ux7aefux7206ux53d1}

老师讲到马加爵案，说明人格问题与长期压抑情绪可能带来的严重后果。

\subsubsection{1. 事件概括}\label{ux4e8bux4ef6ux6982ux62ec}

\begin{itemize}
\tightlist
\item
  马加爵为云南大学学生；
\item
  宿舍打牌时被室友嘲讽``作弊''\,``人格有问题''等；
\item
  长期压抑的负性情绪爆发；
\item
  最终实施极端暴力行为，案件轰动全国。
\end{itemize}

\subsubsection{2.
老师借此强调}\label{ux8001ux5e08ux501fux6b64ux5f3aux8c03}

一些看似``小摩擦''的事情，若叠加在长期压抑、失衡的人格结构之上，可能造成极严重后果。

\subsubsection{3.
【学习建议/做人提醒】}\label{ux5b66ux4e60ux5efaux8baeux505aux4ebaux63d0ux9192}

老师由此给出一个非常生活化的提醒：

\begin{itemize}
\tightlist
\item
  将来毕业时，记得拥抱室友，感谢彼此``不杀之恩''。
\end{itemize}

虽然是玩笑口吻，但核心意思是：

\begin{itemize}
\tightlist
\item
  宿舍关系、日常相处、互相体谅非常重要。
\end{itemize}

\subsubsection{4.
【举例】``一饭之恩''救命}\label{ux4e3eux4f8bux4e00ux996dux4e4bux6069ux6551ux547d}

老师提到：

\begin{itemize}
\tightlist
\item
  马加爵唯独放过了一位曾在他困难时给过他饭吃的同学；
\item
  说明一个小小善举，有时真的可能在关键时刻改变结局。
\end{itemize}

\subsubsection{5. 【学习建议】}\label{ux5b66ux4e60ux5efaux8bae-2}

老师最后提醒：

\begin{itemize}
\tightlist
\item
  看到身边同学有困难，能帮就帮；
\item
  ``勿以善小而不为''。
\end{itemize}

\begin{center}\rule{0.5\linewidth}{0.5pt}\end{center}

\subsection{三十九、课末补充：当前教育要求}\label{ux4e09ux5341ux4e5dux8bfeux672bux8865ux5145ux5f53ux524dux6559ux80b2ux8981ux6c42}

录音最后，老师提到：

\begin{itemize}
\item
  \textbf{2026 年 2 月 25 日，教育部举行``落实健康第一''工作部署会}
\item
  强调：

  \begin{itemize}
  \tightlist
  \item
    让学生动起来
  \item
    心理强起来
  \end{itemize}
\end{itemize}

说明心理健康已成为当前教育中的重要要求。

\subsubsection{录音结尾情况}\label{ux5f55ux97f3ux7ed3ux5c3eux60c5ux51b5}

老师最后一句提到：

\begin{itemize}
\tightlist
\item
  ``大概几十秒之后七点四十会在雨课堂发布我们今天的\ldots\ldots{}''
\end{itemize}

录音到这里中断，后半句未录全。
推测这里即将发布\textbf{当堂任务/签到/问卷/课堂活动}之类内容，但原录音未完整保留，不能强行补写。

\begin{center}\rule{0.5\linewidth}{0.5pt}\end{center}

\subsection{本节英文术语索引}\label{ux672cux8282ux82f1ux6587ux672fux8bedux7d22ux5f15-1}

\begin{itemize}
\tightlist
\item
  危机 --- crisis
\item
  心理危机 --- psychological crisis
\item
  危机状态 --- crisis state / state of disequilibrium
\item
  应对能力 --- coping capacity
\item
  精神活性物质 --- psychoactive substance
\item
  物质滥用 --- substance abuse
\item
  酒精 --- alcohol
\item
  匿名戒酒者协会 --- Alcoholics Anonymous（AA）
\item
  生物---心理---社会医学模式 --- biopsychosocial medical model
\item
  松果体 --- pineal gland
\item
  5-羟色胺 --- 5-hydroxytryptamine（5-HT）/ serotonin
\item
  褪黑素 --- melatonin
\item
  季节性情感障碍 --- Seasonal Affective Disorder（SAD）
\item
  碳水化合物 --- carbohydrates
\item
  色氨酸 --- tryptophan
\item
  胰岛素 --- insulin
\item
  血脑屏障 --- blood-brain barrier（BBB）
\item
  创伤后应激障碍 --- post-traumatic stress disorder（PTSD）
\item
  非典型肺炎 --- severe acute respiratory syndrome（SARS）
\item
  新型冠状病毒病 --- COVID-19
\item
  员工帮助计划 --- Employee Assistance Program（EAP）
\item
  人格 --- personality
\item
  应对方式 --- coping style
\item
  气质 --- temperament
\item
  性格 --- character
\item
  能力 --- ability
\item
  人格障碍 --- personality disorder
\item
  强迫症 --- obsessive-compulsive disorder（OCD）
\item
  强迫特质 --- obsessive-compulsive traits
\item
  胆汁质 --- choleric temperament
\item
  多血质 --- sanguine temperament
\item
  黏液质 --- phlegmatic temperament
\item
  抑郁质 --- melancholic temperament
\item
  慢性疾病 --- chronic disease
\item
  精神疾病 --- mental illness
\item
  行为问题 --- behavioral problems
\item
  执业医师考试 --- physician licensing examination
\end{itemize}

如果你愿意，我下一步可以把这份笔记继续整理成``更适合考前背诵的精简版''，或者再给你单独抽一版``本节可能会考的关键词和论述题框架''。

\clearpage

\section{第3讲：心理危机的表现、应激与自杀}\label{ux7b2c3ux8bb2ux5fc3ux7406ux5371ux673aux7684ux8868ux73b0ux5e94ux6fc0ux4e0eux81eaux6740}

\begin{itemize}
\tightlist
\item
  课程：心理危机干预与预防
\item
  老师：彭睿
\item
  日期：2026-03-17
\end{itemize}

\begin{center}\rule{0.5\linewidth}{0.5pt}\end{center}

\subsection{一、课前说明与课堂安排}\label{ux4e00ux8bfeux524dux8bf4ux660eux4e0eux8bfeux5802ux5b89ux6392}

【课程说明】本课为选修课，老师说明不强求所有同学全程高度集中听讲；如果觉得内容有用可以认真听，如果想在下面自习也可以，但\textbf{不要影响课堂纪律}，不要交头接耳、影响他人。

【课堂安排】本节课中间会有：

\begin{itemize}
\tightlist
\item
  \textbf{签到}
\item
  \textbf{随堂测验}
\end{itemize}

两项都通过\textbf{雨课堂}完成，且需在\textbf{当晚 12 点前}完成。

\begin{center}\rule{0.5\linewidth}{0.5pt}\end{center}

\subsection{二、课程导入：什么是``危机（crisis）''与``心理危机（psychological
crisis）''}\label{ux4e8cux8bfeux7a0bux5bfcux5165ux4ec0ux4e48ux662fux5371ux673acrisisux4e0eux5fc3ux7406ux5371ux673apsychological-crisis}

老师先从``危机''的一般概念讲起。

\subsubsection{1.
危机的两个含义}\label{ux5371ux673aux7684ux4e24ux4e2aux542bux4e49}

危机（crisis）有两个层面的含义：

\begin{enumerate}
\def\labelenumi{\arabic{enumi}.}
\tightlist
\item
  \textbf{突发事件本身} 如地震、水灾、空难、疾病、战争等重大应激事件。
\item
  \textbf{人所处的紧张状态}
  即个体在遭遇这些重大应激事件后，进入的一种紧张、失衡、难以把握的心理状态。
\end{enumerate}

\subsubsection{2.
心理危机的定义}\label{ux5fc3ux7406ux5371ux673aux7684ux5b9aux4e49}

心理危机（psychological crisis）是指：
当个体遭遇重大问题或变化时，感到\textbf{难以解决、难以把握}，原有心理平衡被打破，从而出现\textbf{思维或行为的紊乱}。

也就是说，当个体意识到：

\begin{itemize}
\tightlist
\item
  当前困难\textbf{超过了自己的应对能力}
\item
  原有平衡已经无法维持
\end{itemize}

就会出现一系列反应，包括：

\begin{itemize}
\tightlist
\item
  \textbf{生理反应}
\item
  \textbf{情绪反应}
\item
  \textbf{认知反应}
\item
  \textbf{行为反应}
\end{itemize}

常见表现如：

\begin{itemize}
\tightlist
\item
  焦虑（anxiety）
\item
  紧张
\item
  恐惧
\item
  压抑
\item
  痛苦等负性情绪
\end{itemize}

\begin{center}\rule{0.5\linewidth}{0.5pt}\end{center}

\subsection{三、心理危机干预（crisis
intervention）的基本含义}\label{ux4e09ux5fc3ux7406ux5371ux673aux5e72ux9884crisis-interventionux7684ux57faux672cux542bux4e49}

老师强调，心理危机干预是一种\textbf{心理治疗方式}，主要面向处于困境或遭受挫折的人，提供：

\begin{itemize}
\tightlist
\item
  心理关怀
\item
  短程帮助（brief help）
\item
  倾听
\item
  提问
\item
  情绪释放
\item
  干预措施
\end{itemize}

其目标是让当事人能够：

\begin{itemize}
\tightlist
\item
  正确认识自己的危机
\item
  释放被压抑的情绪
\item
  逐步纠正内心``错位''
\item
  恢复心理平衡
\end{itemize}

\subsubsection{1.
判断是否达到``心理危机''的三个标准}\label{ux5224ux65adux662fux5426ux8fbeux5230ux5fc3ux7406ux5371ux673aux7684ux4e09ux4e2aux6807ux51c6}

是否构成心理危机，一般看三个方面：

\begin{enumerate}
\def\labelenumi{\arabic{enumi}.}
\item
  \textbf{有重大影响性的事件}

  \begin{itemize}
  \tightlist
  \item
    如水灾、火灾、地震
  \item
    或重大生活应激事件
  \end{itemize}
\item
  \textbf{出现严重不适感}

  \begin{itemize}
  \tightlist
  \item
    包括生理反应
  \item
    也包括心理反应
  \end{itemize}
\item
  \textbf{原有应对方式已无效}

  \begin{itemize}
  \tightlist
  \item
    个人过去的应对方法已经无法解决问题
  \end{itemize}
\end{enumerate}

满足这些情况时，就需要进行心理危机干预。

\subsubsection{2.
心理危机干预的分类}\label{ux5fc3ux7406ux5371ux673aux5e72ux9884ux7684ux5206ux7c7b}

心理危机干预可分为：

\begin{itemize}
\tightlist
\item
  \textbf{预防性干预（preventive intervention）}
\item
  \textbf{治疗性干预（therapeutic intervention）}
\item
  \textbf{补救性干预（remedial intervention）}
\end{itemize}

【老师观点】目前多数心理危机干预仍停留在\textbf{补救}层面；
真正系统的\textbf{预防性干预}还不够普遍，日常心理健康教育做得也相对较少。

\begin{center}\rule{0.5\linewidth}{0.5pt}\end{center}

\subsection{四、按刺激来源划分的危机类型}\label{ux56dbux6309ux523aux6fc0ux6765ux6e90ux5212ux5206ux7684ux5371ux673aux7c7bux578b}

根据危机刺激来源，可分为三类：

\subsubsection{1. 发展性危机（developmental
crisis）}\label{ux53d1ux5c55ux6027ux5371ux673adevelopmental-crisis}

又称：

\begin{itemize}
\tightlist
\item
  \textbf{内源性危机}
\item
  \textbf{常规性危机}
\end{itemize}

指个体在成长发展过程中，由于阶段性急剧变化而产生的异常反应。

老师强调：人生由一连串发展阶段组成，当一个人从一个阶段进入另一个阶段时，如果原有的适应能力不足，就可能出现发展性危机。

【举例】

\begin{itemize}
\tightlist
\item
  儿童与父母分离后的焦虑
\item
  青少年情感困惑
\item
  中年职场压力、下岗
\item
  父母衰老或死亡
\item
  配偶离去等
\end{itemize}

【要点】 发展性危机具有以下特点：

\begin{itemize}
\tightlist
\item
  常规发生
\item
  可预期
\item
  在生命各阶段都会存在
\end{itemize}

【学习理解】如果个体有足够时间和机会去适应，如：

\begin{itemize}
\tightlist
\item
  获取相关信息
\item
  学习新技能
\item
  承担新角色
\end{itemize}

则发展性危机对个体的冲击会减少。

\subsubsection{2. 境遇性危机（situational
crisis）}\label{ux5883ux9047ux6027ux5371ux673asituational-crisis}

又称：

\begin{itemize}
\tightlist
\item
  \textbf{外源性危机}
\item
  \textbf{环境性危机}
\end{itemize}

指由外部突发、异常事件引起的危机。

【举例】

\begin{itemize}
\tightlist
\item
  地震
\item
  战争
\item
  疾病
\item
  恐怖袭击等
\end{itemize}

\subsubsection{3. 存在性危机（existential
crisis）}\label{ux5b58ux5728ux6027ux5371ux673aexistential-crisis}

与人生重大问题相关，涉及：

\begin{itemize}
\tightlist
\item
  人生目的（purpose of life）
\item
  人生责任（responsibility）
\item
  独立性（independence）
\item
  自由（freedom）
\item
  承诺（commitment）等
\end{itemize}

它表现为围绕这些问题产生的\textbf{内部冲突与焦虑}。

【举例】

\begin{itemize}
\tightlist
\item
  一个人到 40 岁还一事无成，可能产生强烈空虚感和存在性危机
\end{itemize}

【补充】
存在性危机既可能基于现实问题，也可能基于后悔、空虚感等主观体验。

\begin{center}\rule{0.5\linewidth}{0.5pt}\end{center}

\subsection{五、按发生过程划分危机；危机后的四个阶段}\label{ux4e94ux6309ux53d1ux751fux8fc7ux7a0bux5212ux5206ux5371ux673aux5371ux673aux540eux7684ux56dbux4e2aux9636ux6bb5}

\subsubsection{1.
危机按进程可分为}\label{ux5371ux673aux6309ux8fdbux7a0bux53efux5206ux4e3a}

\begin{itemize}
\tightlist
\item
  \textbf{急性危机（acute crisis）}
\item
  \textbf{慢性危机（chronic crisis）}
\item
  \textbf{混合性危机（mixed crisis）}
\end{itemize}

主要依据危机发生的快慢进行区分。

\subsubsection{2.
危机发生后的四个阶段}\label{ux5371ux673aux53d1ux751fux540eux7684ux56dbux4e2aux9636ux6bb5}

老师将危机后的过程分为四期：

\paragraph{（1）冲击期（impact
phase）}\label{ux51b2ux51fbux671fimpact-phase}

危机刚发生不久，个体表现为：

\begin{itemize}
\tightlist
\item
  震惊
\item
  惊慌失措
\end{itemize}

\paragraph{（2）防御期（defensive
phase）}\label{ux9632ux5fa1ux671fdefensive-phase}

个体开始试图恢复心理平衡，努力控制焦虑，恢复受损认知功能，但不知道如何做，因此常出现心理防御机制：

\begin{itemize}
\tightlist
\item
  否认（denial）
\item
  合理化（rationalization）
\end{itemize}

\paragraph{（3）解决期（resolution
phase）}\label{ux89e3ux51b3ux671fresolution-phase}

个体开始：

\begin{itemize}
\tightlist
\item
  采取积极办法
\item
  寻求资源
\item
  减轻焦虑
\item
  恢复社会功能
\end{itemize}

\paragraph{（4）成长期（growth
phase）}\label{ux6210ux957fux671fgrowth-phase}

危机过后，部分人会：

\begin{itemize}
\tightlist
\item
  变得更成熟
\item
  学会处理危机的技巧
\end{itemize}

【老师提醒】并不是每个人都会成长。
有的人危机之后会成长，有的人反而会出现不健康的心理和行为。

\begin{center}\rule{0.5\linewidth}{0.5pt}\end{center}

\subsection{六、本节课聚焦的三个重点问题}\label{ux516dux672cux8282ux8bfeux805aux7126ux7684ux4e09ux4e2aux91cdux70b9ux95eeux9898}

老师说明：并不是所有受灾/受创人群都会发展为明确的精神障碍。
在整个``受灾人群''中，只有一部分进入亚健康状态；在亚健康状态的人群中，也只有一部分进一步出现心理危机或精神障碍。

这些问题之间还会互相重叠，而不是彼此独立。

【复习/引入】老师提到常见相关问题包括：

\begin{itemize}
\tightlist
\item
  急性应激障碍（Acute Stress Disorder, \textbf{ASD}）
\item
  创伤后应激障碍（Post-Traumatic Stress Disorder, \textbf{PTSD}）
\item
  抑郁（depression）
\item
  焦虑（anxiety）
\item
  自杀（suicide）
\item
  酒依赖/物质滥用（alcohol dependence / substance abuse）
\end{itemize}

【本节重点】今天主要讲三个内容：

\begin{enumerate}
\def\labelenumi{\arabic{enumi}.}
\tightlist
\item
  创伤后应激障碍 PTSD
\item
  急性应激障碍 ASD
\item
  自杀（suicide）
\end{enumerate}

【考点理解】老师特别指出：

\begin{itemize}
\tightlist
\item
  \textbf{自杀并不是精神病诊断系统中的单独疾病诊断}
\item
  它会出现在许多精神障碍中
\item
  但由于社会意义重大，所以单独拿出来讲
\end{itemize}

\begin{center}\rule{0.5\linewidth}{0.5pt}\end{center}

\subsection{七、精神科两套常见诊断系统}\label{ux4e03ux7cbeux795eux79d1ux4e24ux5957ux5e38ux89c1ux8bcaux65adux7cfbux7edf}

老师提到精神科常见有两套诊断标准：

\begin{enumerate}
\def\labelenumi{\arabic{enumi}.}
\item
  \textbf{ICD 诊断系统} 国际疾病分类（International Classification of
  Diseases, \textbf{ICD}） 最新为 \textbf{ICD-11}
\item
  \textbf{DSM 诊断系统} 美国精神障碍诊断与统计手册（Diagnostic and
  Statistical Manual of Mental Disorders, \textbf{DSM}） 最新为
  \textbf{DSM-5}
\end{enumerate}

【复习】在 \textbf{DSM-5} 中：

\begin{itemize}
\tightlist
\item
  急性应激障碍（ASD）
\item
  创伤后应激障碍（PTSD）
\end{itemize}

都属于\textbf{创伤及应激相关障碍}（Trauma- and Stressor-Related
Disorders）这一大类。

\begin{center}\rule{0.5\linewidth}{0.5pt}\end{center}

\subsection{八、急性应激障碍（Acute Stress Disorder,
ASD）}\label{ux516bux6025ux6027ux5e94ux6fc0ux969cux788dacute-stress-disorder-asd}

\subsubsection{1. 定义}\label{ux5b9aux4e49}

急性应激障碍又称：

\begin{itemize}
\tightlist
\item
  \textbf{急性心因反应}
\end{itemize}

指个体在突然遭遇强烈应激事件后，\textbf{数分钟到数小时内出现}的急性应激反应，且：

\begin{itemize}
\tightlist
\item
  持续 \textbf{3 天以上}
\item
  不超过 \textbf{1 个月}
\end{itemize}

【考点】 \textbf{``3 天以上、不超过 1 个月''是诊断 ASD 的关键病程标准。}

如果超过 1 个月，就要考虑是否转为 \textbf{PTSD}。

\subsubsection{2.
常见应激事件}\label{ux5e38ux89c1ux5e94ux6fc0ux4e8bux4ef6}

包括：

\begin{itemize}
\tightlist
\item
  重大交通事故
\item
  亲人死亡
\item
  歹徒袭击
\item
  强奸
\item
  虐待
\item
  自然灾害
\item
  战争
\item
  恐怖袭击
\item
  SARS/非典等重大公共卫生事件
\item
  实际接触到死亡或死亡威胁
\end{itemize}

\subsubsection{3. 临床表现}\label{ux4e34ux5e8aux8868ux73b0}

老师强调，精神科疾病主要通过\textbf{临床表现}进行识别和诊断。

\paragraph{（1）意识范围狭窄、意识清晰度下降}\label{ux610fux8bc6ux8303ux56f4ux72edux7a84ux610fux8bc6ux6e05ux6670ux5ea6ux4e0bux964d}

可出现：

\begin{itemize}
\tightlist
\item
  定向困难（disorientation）
\item
  注意力不集中
\item
  对外界刺激理解困难
\item
  言语凌乱破碎
\end{itemize}

【解释】

\begin{itemize}
\tightlist
\item
  \textbf{意识范围狭窄}：只对狭窄范围内的刺激有反应，超出范围可能没有反应
\item
  \textbf{意识清晰度下降}：可表现为嗜睡、意识浑浊、昏睡，甚至昏迷
\item
  \textbf{定向障碍}：分不清时间、地点、人物，严重时甚至不知道自己是谁
\end{itemize}

【举例】刚接触创伤事件时，人常处于``整个人都懵了''的状态。

\paragraph{（2）闯入性再体验}\label{ux95efux5165ux6027ux518dux4f53ux9a8c}

\begin{itemize}
\tightlist
\item
  反复出现创伤性回忆
\item
  在睡梦中反复出现与创伤有关的痛苦体验
\item
  伴随明显精神痛苦或生理反应
\end{itemize}

\paragraph{（3）精神运动性抑制或兴奋}\label{ux7cbeux795eux8fd0ux52a8ux6027ux6291ux5236ux6216ux5174ux594b}

\textbf{抑制时}可表现为：

\begin{itemize}
\tightlist
\item
  目光呆滞
\item
  情感迟钝
\item
  少语少动
\item
  甚至木僵（stupor）
\end{itemize}

【举例】整天端坐、卧床，对外界刺激毫无反应。

\textbf{兴奋时}可表现为：

\begin{itemize}
\tightlist
\item
  激越（agitation）
\item
  喊叫
\item
  情绪爆发
\item
  冲动伤人或毁物行为
\end{itemize}

\paragraph{（4）警觉性增高（hyperarousal）}\label{ux8b66ux89c9ux6027ux589eux9ad8hyperarousal}

包括：

\begin{itemize}
\tightlist
\item
  过分警觉
\item
  反复噩梦
\item
  触景生情式痛苦
\item
  自主神经功能紊乱
\end{itemize}

常见症状：

\begin{itemize}
\tightlist
\item
  心神不宁
\item
  头痛
\item
  出汗
\item
  失眠
\item
  紧张
\item
  坐立不安
\end{itemize}

\paragraph{（5）回避与选择性遗忘}\label{ux56deux907fux4e0eux9009ux62e9ux6027ux9057ux5fd8}

患者会：

\begin{itemize}
\tightlist
\item
  尽量不去想创伤相关的人和事
\item
  回避能引发触景生情的活动和场所
\item
  对创伤事件出现\textbf{选择性失忆}
\end{itemize}

\paragraph{（6）分离症状：人格解体与现实解体}\label{ux5206ux79bbux75c7ux72b6ux4ebaux683cux89e3ux4f53ux4e0eux73b0ux5b9eux89e3ux4f53}

\begin{itemize}
\tightlist
\item
  \textbf{人格解体（depersonalization）}
\item
  \textbf{现实解体（derealization）}
\end{itemize}

本质上都是一种对自身或周围环境的\textbf{不真实感}。

【老师提示】后面视频会帮助理解这两种体验。

\subsubsection{4. 病程演变}\label{ux75c5ux7a0bux6f14ux53d8}

上述症状一般在 \textbf{24～48 小时}后开始减弱。 若症状\textbf{持续超过 1
个月}，则要考虑 \textbf{PTSD}。

\subsubsection{5.
急性应激障碍的亚型}\label{ux6025ux6027ux5e94ux6fc0ux969cux788dux7684ux4e9aux578b}

有一个亚型叫：

\begin{itemize}
\tightlist
\item
  \textbf{急性应激性精神病}
\end{itemize}

特点：

\begin{itemize}
\tightlist
\item
  由强烈心理创伤引起
\item
  以\textbf{幻觉（hallucination）}、\textbf{妄想（delusion）}和严重情感障碍为主要表现
\item
  内容与应激事件密切相关，且一般可以被人理解
\item
  虽然病情看起来严重，但病程较短，通常 \textbf{1 个月内可恢复}
\end{itemize}

\subsubsection{6. DSM-5
诊断标准（课堂版梳理）}\label{dsm-5-ux8bcaux65adux6807ux51c6ux8bfeux5802ux7248ux68b3ux7406}

老师按 \textbf{A-B-C-D-E-F} 标准讲解。

\paragraph{A
标准：暴露于创伤事件}\label{a-ux6807ux51c6ux66b4ux9732ux4e8eux521bux4f24ux4e8bux4ef6}

包括：

\begin{itemize}
\tightlist
\item
  直接经历创伤
\item
  亲眼目睹他人遭遇创伤
\item
  得知亲密家人/朋友遭遇创伤
\item
  反复暴露于创伤事件的令人痛苦细节中
\end{itemize}

【举例】

\begin{itemize}
\tightlist
\item
  急救人员反复收集残骸
\item
  警察反复接触虐待儿童案件细节
\end{itemize}

\paragraph{B 标准：五大类症状中总计至少 9
条}\label{b-ux6807ux51c6ux4e94ux5927ux7c7bux75c7ux72b6ux4e2dux603bux8ba1ux81f3ux5c11-9-ux6761}

五大类为：

\begin{enumerate}
\def\labelenumi{\arabic{enumi}.}
\tightlist
\item
  侵入性症状（intrusion）
\item
  负性心境（negative mood）
\item
  分离症状（dissociation）
\item
  回避（avoidance）
\item
  唤起（arousal）
\end{enumerate}

【要点】在相关症状中，\textbf{至少满足 9 条}。

\subparagraph{侵入性症状}\label{ux4fb5ux5165ux6027ux75c7ux72b6}

\begin{itemize}
\tightlist
\item
  反复、非自愿、侵入性的痛苦回忆
\item
  反复做相关噩梦
\item
  感觉再次亲临事件现场，即\textbf{闪回（flashback）}
\item
  接触相关线索后出现强烈心理/生理反应
\end{itemize}

\subparagraph{负性心境}\label{ux8d1fux6027ux5fc3ux5883}

\begin{itemize}
\tightlist
\item
  无法体验正性情绪
\end{itemize}

\subparagraph{分离症状}\label{ux5206ux79bbux75c7ux72b6}

\begin{itemize}
\tightlist
\item
  人格解体
\item
  现实解体
\item
  无法回忆创伤事件的重要方面
\end{itemize}

\subparagraph{回避症状}\label{ux56deux907fux75c7ux72b6}

\begin{itemize}
\tightlist
\item
  回避相关记忆、思想、感觉
\item
  回避相关外部线索
\end{itemize}

\subparagraph{唤起症状}\label{ux5524ux8d77ux75c7ux72b6}

\begin{itemize}
\tightlist
\item
  睡眠紊乱
\item
  易激惹
\item
  愤怒爆发
\item
  过度警觉
\item
  注意力不集中
\item
  过度惊跳反应
\end{itemize}

\paragraph{C
标准：病程标准}\label{c-ux6807ux51c6ux75c5ux7a0bux6807ux51c6}

\begin{itemize}
\tightlist
\item
  暴露后持续 \textbf{3 天到 1 个月}
\end{itemize}

\paragraph{D
标准：社会功能损害}\label{d-ux6807ux51c6ux793eux4f1aux529fux80fdux635fux5bb3}

\begin{itemize}
\tightlist
\item
  引起明显临床痛苦
\item
  或导致明显社会功能受损
\end{itemize}

\paragraph{E
标准：排除标准}\label{e-ux6807ux51c6ux6392ux9664ux6807ux51c6}

\begin{itemize}
\tightlist
\item
  排除酒精、药物或其他精神疾病所致
\end{itemize}

【考点】 ABCDE 满足后，方可诊断 ASD。

\subsubsection{7. 预防与治疗}\label{ux9884ux9632ux4e0eux6cbbux7597}

\paragraph{（1）预防}\label{ux9884ux9632}

\begin{itemize}
\tightlist
\item
  加强公众心理健康教育
\item
  对创伤事件开展危机干预
\end{itemize}

\paragraph{（2）支持疗法（supportive
therapy）}\label{ux652fux6301ux7597ux6cd5supportive-therapy}

急性应激障碍治疗首先强调支持疗法。

【学习建议】面对危机事件时：

\begin{itemize}
\tightlist
\item
  要学会自助
\item
  依靠社会支持网络
\item
  促进自我修复
\end{itemize}

【干预重点】

\begin{itemize}
\tightlist
\item
  减轻患者的\textbf{自罪自责}
\item
  帮助其认识：在极端事件中，很难做出``最完美选择''
\item
  让其意识到自己已经尽力把损害降到较低水平
\end{itemize}

此外还应：

\begin{itemize}
\tightlist
\item
  设置温馨、安全环境
\item
  帮助其面对和认识最近的经历与感受
\item
  在客观危险消失后，允许适度宣泄情绪
\item
  保证饮水、进食、睡眠
\item
  维持水电解质平衡
\item
  防止身体衰竭
\end{itemize}

【临床关联】若创伤事件导致残疾，还需帮助患者适应残疾后的生活。

\paragraph{（3）心理治疗}\label{ux5fc3ux7406ux6cbbux7597}

\begin{itemize}
\tightlist
\item
  首先是\textbf{认知行为治疗}（Cognitive Behavioral Therapy,
  \textbf{CBT}）
\end{itemize}

\paragraph{（4）药物治疗}\label{ux836fux7269ux6cbbux7597}

目的主要是减少创伤相关恐惧与激越。

可考虑：

\begin{itemize}
\item
  \textbf{β 受体阻滞剂（beta-blockers）}
\item
  若有明显失眠、激越、躁动、自残行为，可用：

  \begin{itemize}
  \tightlist
  \item
    抗精神病药物（antipsychotics）
  \item
    苯二氮卓类（benzodiazepines）
  \end{itemize}
\item
  若伴焦虑、抑郁，可给予相应抗焦虑/抗抑郁药物
\end{itemize}

【用药原则】

\begin{itemize}
\tightlist
\item
  以\textbf{小剂量}为主
\item
  \textbf{疗程不宜过长}
\end{itemize}

【对比】这一点与后面讲的 PTSD 不同。

\begin{center}\rule{0.5\linewidth}{0.5pt}\end{center}

\subsection{九、创伤后应激障碍（Post-Traumatic Stress Disorder,
PTSD）}\label{ux4e5dux521bux4f24ux540eux5e94ux6fc0ux969cux788dpost-traumatic-stress-disorder-ptsd}

\subsubsection{1. 定义}\label{ux5b9aux4e49-1}

PTSD
指个体遭遇异常强烈创伤应激事件后，\textbf{较迟出现}的一类应激相关障碍。

【与 ASD 的区别】

\begin{itemize}
\tightlist
\item
  ASD：事件后很快出现
\item
  PTSD：\textbf{较迟出现}
\item
  PTSD 的病程要求通常\textbf{超过 1 个月}
\end{itemize}

\subsubsection{2. 核心症状}\label{ux6838ux5fc3ux75c7ux72b6}

PTSD 的核心表现包括：

\begin{enumerate}
\def\labelenumi{\arabic{enumi}.}
\tightlist
\item
  \textbf{反复闯入意识或梦境的创伤性体验}
\item
  \textbf{显著而持续的焦虑状态}
\item
  \textbf{对创伤性事件的持续回避}
\end{enumerate}

\subsubsection{3.
流行病学与危险因素}\label{ux6d41ux884cux75c5ux5b66ux4e0eux5371ux9669ux56e0ux7d20}

老师提到：

\begin{itemize}
\tightlist
\item
  PTSD 的患病率约 \textbf{5\%～12\%}
\item
  常与\textbf{酒依赖}共病
\end{itemize}

风险因素从高到低包括：

\begin{itemize}
\tightlist
\item
  \textbf{强奸（rape）}
\item
  强烈暴力行为
\item
  性虐待
\item
  严重事故或重伤
\item
  被刺伤、被枪击
\item
  亲密朋友/亲属突然去世
\item
  童年期受到威胁或患病
\item
  目睹杀人现场
\item
  严重事故
\item
  自然灾害等
\end{itemize}

\subsubsection{4. 结合视频理解
PTSD}\label{ux7ed3ux5408ux89c6ux9891ux7406ux89e3-ptsd}

老师播放了美国退伍军人的英文访谈，用来帮助理解 PTSD 的真实体验。

【临床关联】视频中可见的线索包括：

\begin{itemize}
\tightlist
\item
  长期参战后出现 chronic PTSD
\item
  曾住精神病院
\item
  有严重攻击冲动，甚至准备出去杀人
\item
  反复闪现暴力计划
\item
  需要多种药物控制焦虑、睡眠和激越
\item
  存在极重的内疚、自责、创伤回忆及社会疏离
\end{itemize}

【药物术语补充】 老师借视频顺带讲了几个药名：

\begin{itemize}
\tightlist
\item
  \textbf{Clozapine 氯氮平}：抗精神病药，有较强镇静作用
\item
  \textbf{Abilify
  阿立哌唑（aripiprazole）}：抗精神病药，可改善古怪念头、攻击暴力行为等
\item
  \textbf{Xanax
  阿普唑仑（alprazolam）}：苯二氮卓类抗焦虑药，兼具镇静催眠作用
\end{itemize}

\subsubsection{5. PTSD
的临床表现}\label{ptsd-ux7684ux4e34ux5e8aux8868ux73b0}

\paragraph{（1）创伤性再体验（re-experiencing）}\label{ux521bux4f24ux6027ux518dux4f53ux9a8cre-experiencing}

包括：

\begin{itemize}
\tightlist
\item
  思维、记忆或梦中反复不自主出现创伤情境
\item
  遇到相关线索时触景生情、极度痛苦
\item
  闪回（flashback）
\item
  做与创伤相关的噩梦
\item
  从梦中尖叫、惊醒
\item
  醒后恐惧情绪仍持续
\item
  有时还会出现冲动行为
\item
  甚至出现惊恐发作（panic attack）
\end{itemize}

【解释】惊恐发作常表现为：

\begin{itemize}
\tightlist
\item
  突如其来的急性焦虑
\item
  濒死感
\item
  失控感
\item
  ``快疯了''的感觉
\item
  一般持续十几分钟到半小时后缓解
\end{itemize}

\paragraph{（2）回避与麻木（avoidance and
numbing）}\label{ux56deux907fux4e0eux9ebbux6728avoidance-and-numbing}

回避可以是：

\begin{itemize}
\tightlist
\item
  \textbf{有意识回避}
\item
  \textbf{无意识回避}
\end{itemize}

\textbf{有意识回避}：

\begin{itemize}
\tightlist
\item
  极力不去想相关人和事
\item
  不愿谈起创伤事件
\end{itemize}

【临床关联】媒体采访、法医取证等要求患者回忆创伤事件时，可能给患者带来巨大痛苦。

\textbf{无意识回避}：

\begin{itemize}
\tightlist
\item
  对创伤事件出现选择性遗忘
\item
  与创伤无关的记忆保持完好
\item
  也可能表现为拼命工作，用工作压住创伤回忆
\end{itemize}

【举例】视频中的美国大兵就是不停工作来回避。

\textbf{情感麻木}还可表现为：

\begin{itemize}
\tightlist
\item
  对任何事失去兴趣
\item
  与外界疏离
\item
  缺乏信任感
\item
  难以建立亲密关系
\item
  不愿进行情感交流
\item
  得过且过
\item
  对任何事无动于衷
\item
  甚至出现``生不如死''、自杀观念和行为
\end{itemize}

【复习】老师特别指出：这些表现和\textbf{抑郁症}很相似。

\paragraph{（3）分离症状：人格解体/现实解体}\label{ux5206ux79bbux75c7ux72b6ux4ebaux683cux89e3ux4f53ux73b0ux5b9eux89e3ux4f53}

老师又播放了一段英文视频，专门讲：

\begin{itemize}
\tightlist
\item
  \textbf{人格解体（depersonalization）}
\item
  \textbf{现实解体（derealization）}
\end{itemize}

视频中所描述的体验包括：

\begin{itemize}
\tightlist
\item
  感觉自己不像一个真实的人
\item
  像个``幽灵''
\item
  像在身体外面看着自己
\item
  感觉一切动作不是自己在做
\item
  像与现实脱离
\item
  通过 grounding exercises（落地化训练/现实定向训练）才能慢慢缓回来
\end{itemize}

\paragraph{（4）警觉性增高}\label{ux8b66ux89c9ux6027ux589eux9ad8}

包括：

\begin{itemize}
\tightlist
\item
  入睡困难
\item
  睡眠浅
\item
  容易惊醒
\item
  梦中反复出现创伤场景
\item
  惊醒后情绪激动
\item
  躁动不安
\item
  易激惹
\item
  难以集中注意力
\end{itemize}

\paragraph{（5）伴随问题}\label{ux4f34ux968fux95eeux9898}

老师强调 PTSD 常与以下问题密切相关：

\begin{itemize}
\tightlist
\item
  酒精/药物滥用
\item
  攻击行为
\item
  伤人毁物
\item
  自杀行为
\item
  性格改变（如孤僻、不再信任他人）
\end{itemize}

\subsubsection{6. PTSD
的诊断标准（课堂版）}\label{ptsd-ux7684ux8bcaux65adux6807ux51c6ux8bfeux5802ux7248}

老师将其与 ASD 对照讲解。

\paragraph{A 标准}\label{a-ux6807ux51c6}

与 ASD 基本相同：必须有创伤暴露。

\paragraph{B 标准：侵入性症状至少 1
条}\label{b-ux6807ux51c6ux4fb5ux5165ux6027ux75c7ux72b6ux81f3ux5c11-1-ux6761}

包括：

\begin{itemize}
\tightlist
\item
  反复、非自愿、侵入性的痛苦回忆
\item
  反复做相关噩梦
\item
  闪回
\item
  接触相关线索时出现明显心理痛苦或生理反应
\end{itemize}

\paragraph{C 标准：回避}\label{c-ux6807ux51c6ux56deux907f}

至少满足以下之一：

\begin{itemize}
\tightlist
\item
  回避相关记忆、思想、感觉
\item
  回避相关线索
\end{itemize}

\paragraph{D 标准：认知和心境负性改变，至少 2
条}\label{d-ux6807ux51c6ux8ba4ux77e5ux548cux5fc3ux5883ux8d1fux6027ux6539ux53d8ux81f3ux5c11-2-ux6761}

包括：

\begin{itemize}
\tightlist
\item
  无法回忆创伤的重要部分
\item
  持续负性信念或预期 例如：``我很坏''``没人可信''``世界很危险''
\item
  对创伤原因/结果产生扭曲认知
\item
  持续自罪自责
\item
  恐惧、愤怒、内疚等负性情绪
\item
  兴趣减少
\item
  与他人疏离
\item
  无法体验正性情绪
\end{itemize}

\paragraph{E 标准：警觉性和反应性改变，至少 2
条}\label{e-ux6807ux51c6ux8b66ux89c9ux6027ux548cux53cdux5e94ux6027ux6539ux53d8ux81f3ux5c11-2-ux6761}

包括：

\begin{itemize}
\tightlist
\item
  易激惹
\item
  愤怒爆发
\item
  不顾后果的自毁行为
\item
  过度警觉
\item
  过分惊跳反应
\item
  注意力不集中
\item
  睡眠改变
\end{itemize}

\paragraph{F 标准：病程}\label{f-ux6807ux51c6ux75c5ux7a0b}

\begin{itemize}
\tightlist
\item
  持续时间\textbf{超过 1 个月}
\end{itemize}

\paragraph{G 标准}\label{g-ux6807ux51c6}

\begin{itemize}
\tightlist
\item
  导致明显社会功能损害
\end{itemize}

\paragraph{H 标准}\label{h-ux6807ux51c6}

\begin{itemize}
\tightlist
\item
  排除酒药或躯体疾病所致精神障碍
\end{itemize}

\subsubsection{7. PTSD 的治疗}\label{ptsd-ux7684ux6cbbux7597}

\paragraph{（1）药物治疗}\label{ux836fux7269ux6cbbux7597-1}

如果伴有抑郁表现，可使用：

\begin{itemize}
\item
  \textbf{SSRI}：选择性 5-羟色胺再摄取抑制剂（Selective Serotonin
  Reuptake Inhibitors）

  \begin{itemize}
  \tightlist
  \item
    舍曲林（sertraline）
  \item
    氟西汀（fluoxetine）
  \item
    帕罗西汀（paroxetine）
  \end{itemize}
\item
  \textbf{SNRI}：5-羟色胺-去甲肾上腺素再摄取抑制剂（Serotonin-Norepinephrine
  Reuptake Inhibitors）

  \begin{itemize}
  \tightlist
  \item
    文拉法辛（venlafaxine）
  \end{itemize}
\end{itemize}

【考点】PTSD 药物起效通常较慢：

\begin{itemize}
\tightlist
\item
  一般 \textbf{4～6 周}开始改善
\item
  有时需 \textbf{8 周或更长}才达到充分治疗效果
\end{itemize}

【用药原则】

\begin{itemize}
\tightlist
\item
  从小剂量开始
\item
  逐渐加至治疗剂量
\end{itemize}

其他可合并药物包括：

\begin{itemize}
\tightlist
\item
  哌唑嗪（prazosin）------对噩梦较有效
\item
  曲唑酮（trazodone）
\item
  米氮平（mirtazapine）
\item
  奎硫平（quetiapine）等
\end{itemize}

若表现突出，还可考虑：

\begin{itemize}
\tightlist
\item
  \textbf{抗肾上腺素能药物}：用于警觉高、多动、分离症状明显者
\item
  \textbf{心境稳定剂（mood
  stabilizers）}：用于攻击性、冲动性、行为不稳明显者
\item
  \textbf{新型抗惊厥药}：辅助改善睡眠
\item
  \textbf{非典型抗精神病药}：用于明显恐惧、多疑、过度警觉等
\end{itemize}

【易错/对比】 PTSD
\textbf{一般不常规使用苯二氮卓类药物}，因为研究提示其对长期预后帮助不大。
这点与 ASD 的处理不同。

\paragraph{（2）心理治疗}\label{ux5fc3ux7406ux6cbbux7597-1}

首选：

\begin{itemize}
\tightlist
\item
  \textbf{认知行为治疗（CBT）}
\end{itemize}

核心方法：

\begin{itemize}
\tightlist
\item
  \textbf{暴露疗法（exposure therapy）}
\end{itemize}

【举例】
如果患者因车祸后害怕坐车，可在情绪缓解后，逐步再次进入乘车环境，让其重新认识到：

\begin{itemize}
\tightlist
\item
  乘车总体是相对安全的
\item
  车祸是小概率事件
\end{itemize}

从而逐渐恢复信心。

\paragraph{（3）焦虑处理训练}\label{ux7126ux8651ux5904ux7406ux8badux7ec3}

包括：

\begin{itemize}
\tightlist
\item
  放松训练
\item
  肌肉放松
\item
  腹式呼吸
\item
  正性思维训练
\item
  自信训练
\item
  学会表达感受、意见和愿望
\end{itemize}

\paragraph{（4）其他治疗}\label{ux5176ux4ed6ux6cbbux7597}

\begin{itemize}
\tightlist
\item
  \textbf{生物反馈治疗（biofeedback）}
\item
  \textbf{无抽搐电休克治疗（modified electroconvulsive therapy, MECT）}
\item
  \textbf{重复经颅磁刺激（repetitive transcranial magnetic stimulation,
  rTMS）}
\end{itemize}

【老师解释】

\begin{itemize}
\tightlist
\item
  生物反馈：通过监测脑电等指标，把情绪状态可视化，帮助患者学会调节
\item
  MECT：是在传统 ECT 基础上改良，加入麻醉剂、肌松剂等，减少痛苦和副作用
\item
  rTMS：利用外加磁场影响脑内神经元放电节律，帮助恢复较正常的脑功能节律
\end{itemize}

【易错】老师特别指出： ``电磁辐射对人完全没有影响''这种说法是错误的。 从
rTMS 的治疗原理就能看出，电磁刺激是可以影响神经元放电的。

\begin{center}\rule{0.5\linewidth}{0.5pt}\end{center}

\subsection{十、第一节课结尾：引入``自杀''话题与教学项目招募}\label{ux5341ux7b2cux4e00ux8282ux8bfeux7ed3ux5c3eux5f15ux5165ux81eaux6740ux8bddux9898ux4e0eux6559ux5b66ux9879ux76eeux62dbux52df}

老师说，接下来将进入``自杀（suicide）''这一重要主题。
在正式讲之前，原本准备播放一段关于抑郁症患者自杀自述的 AI
视频，但因第一节课快结束，放到第二节课再看。

【课堂信息】老师顺带介绍自己正在做一个课题：

\begin{itemize}
\tightlist
\item
  用 \textbf{AI 短视频}来呈现精神科中较抽象的病例和症状
\item
  目标是辅助精神病学教学
\item
  当前主要是老师自己在做，非常缺人手
\item
  如果有同学对相关工作感兴趣，可以联系老师
\end{itemize}

\begin{center}\rule{0.5\linewidth}{0.5pt}\end{center}

\subsection{十一、第二节课开始：AI
视频中的自杀未遂与抑郁症表现}\label{ux5341ux4e00ux7b2cux4e8cux8282ux8bfeux5f00ux59cbai-ux89c6ux9891ux4e2dux7684ux81eaux6740ux672aux9042ux4e0eux6291ux90c1ux75c7ux8868ux73b0}

第二节课一开始，老师播放了一段由 AI
制作的、整合真实案例的教学视频，内容围绕：

\begin{itemize}
\tightlist
\item
  自杀未遂后的身体创伤
\item
  抑郁症患者的主观体验
\item
  抑郁症的部分核心症状
\end{itemize}

\subsubsection{1.
视频中展示的自杀未遂体验}\label{ux89c6ux9891ux4e2dux5c55ux793aux7684ux81eaux6740ux672aux9042ux4f53ux9a8c}

【临床关联】跳楼后可出现严重多发伤：

\begin{itemize}
\tightlist
\item
  脚踝骨折
\item
  腿骨折
\item
  肋骨骨折
\item
  血气胸（hemopneumothorax）
\item
  肾脏/骨盆出血
\item
  腰椎粉碎等
\end{itemize}

视频还描述了：

\begin{itemize}
\tightlist
\item
  当时意识模糊、耳鸣、说不出话
\item
  明知自己快死
\item
  喘不上气
\item
  只想睡，但强迫自己不能闭眼
\item
  抢救后反复梦见自己一遍遍跳楼
\item
  醒后持续``生不如死''
\item
  身体疼痛极重
\item
  即使活下来了，仍感到对不起亲人，甚至``还是想死''
\end{itemize}

\subsubsection{2.
视频中呈现的抑郁症体验}\label{ux89c6ux9891ux4e2dux5448ux73b0ux7684ux6291ux90c1ux75c7ux4f53ux9a8c}

【临床关联】视频展现的典型抑郁表现包括：

\paragraph{（1）兴趣丧失 /
快感缺失}\label{ux5174ux8da3ux4e27ux5931-ux5febux611fux7f3aux5931}

\begin{itemize}
\tightlist
\item
  对几乎所有事情失去兴趣
\item
  以前喜欢的旅行、聚餐、工作成就都引不起波澜
\end{itemize}

老师借视频强调：
\textbf{抑郁的反面不一定是``快乐''，而可能是``生命力''。}

\paragraph{（2）意志活动减退}\label{ux610fux5fd7ux6d3bux52a8ux51cfux9000}

原本很简单的小事，也变成巨大负担：

\begin{itemize}
\tightlist
\item
  回消息觉得麻烦
\item
  吃饭都觉得流程太复杂
\item
  倒一杯水都难如登天
\item
  穿衣、洗头都做不到
\item
  可以一个月不洗澡
\item
  长时间躺床不起
\end{itemize}

\paragraph{（3）注意力与决策功能下降}\label{ux6ce8ux610fux529bux4e0eux51b3ux7b56ux529fux80fdux4e0bux964d}

\begin{itemize}
\tightlist
\item
  看文件一下午看不进去
\item
  看了就忘
\item
  穿哪件衣服、点什么外卖都能纠结一两个小时
\end{itemize}

\paragraph{（4）精神运动性迟滞/僵住}\label{ux7cbeux795eux8fd0ux52a8ux6027ux8fdfux6edeux50f5ux4f4f}

【举例】视频中描述：

\begin{itemize}
\tightlist
\item
  明知该打电话求助
\item
  却伸不出手、拿不起电话、拨不了号
\item
  就这样盯着电话看了四个小时
\end{itemize}

\paragraph{（5）睡眠紊乱}\label{ux7761ux7720ux7d0aux4e71}

\begin{itemize}
\tightlist
\item
  先是长期失眠
\item
  后来变成早醒
\item
  凌晨三四点突然醒来
\item
  醒后满脑子负面念头
\item
  再也睡不着
\end{itemize}

【临床关联】老师提到，很多时候精神症状出现前，\textbf{身体症状已经先发出信号}。

\subsubsection{3.
课堂说明：视频并未囊括抑郁全部症状}\label{ux8bfeux5802ux8bf4ux660eux89c6ux9891ux5e76ux672aux56caux62ecux6291ux90c1ux5168ux90e8ux75c7ux72b6}

老师说明：

\begin{itemize}
\tightlist
\item
  目前视频只做出了抑郁症症状的一半
\item
  另一半还在制作中
\end{itemize}

\begin{center}\rule{0.5\linewidth}{0.5pt}\end{center}

\subsection{十二、老师再次介绍 AI
教学视频项目}\label{ux5341ux4e8cux8001ux5e08ux518dux6b21ux4ecbux7ecd-ai-ux6559ux5b66ux89c6ux9891ux9879ux76ee}

【学习建议/课程外拓展】老师再次``打广告''：

如果同学对以下内容感兴趣，可以联系老师：

\begin{itemize}
\tightlist
\item
  精神疾病临床表现
\item
  AI 视频本地化制作
\item
  开源模型部署
\item
  医学教学与 AI 结合
\end{itemize}

老师补充说明：

\begin{itemize}
\tightlist
\item
  目前没有课题经费
\item
  没有现成``免费模型/算力''福利
\item
  现在主要用开源模型在本地跑
\item
  也在联系学校计算机教研室，看看能否借助天津科技大学高性能计算平台，为电脑性能不足的同学提供硬件支持
\end{itemize}

\begin{center}\rule{0.5\linewidth}{0.5pt}\end{center}

\subsection{十三、自杀（suicide）话题正式开始：先纠正``觉得好笑''的态度}\label{ux5341ux4e09ux81eaux6740suicideux8bddux9898ux6b63ux5f0fux5f00ux59cbux5148ux7ea0ux6b63ux89c9ux5f97ux597dux7b11ux7684ux6001ux5ea6}

老师提到，放视频时有同学在笑。老师明确指出：

【课堂提醒】不要把自杀当成笑话。 视频虽然是 AI
合成的教学视频，但素材来自\textbf{三个真实案例}。

老师强调：

\begin{itemize}
\tightlist
\item
  自杀并不是离学生很远的事
\item
  每天都可能有大学生正在自杀
\item
  甚至是\textbf{已经自杀死亡}，不是单指自杀未遂
\end{itemize}

\begin{center}\rule{0.5\linewidth}{0.5pt}\end{center}

\subsection{十四、自杀的流行病学：总体情况}\label{ux5341ux56dbux81eaux6740ux7684ux6d41ux884cux75c5ux5b66ux603bux4f53ux60c5ux51b5}

\subsubsection{1. 全球情况}\label{ux5168ux7403ux60c5ux51b5}

老师提到 2019 年全球年龄标化自杀率地图：

\begin{itemize}
\tightlist
\item
  颜色越深，自杀率越高
\item
  俄罗斯、部分南非国家、美国、加拿大、澳大利亚等相对较高
\item
  中国处于中等水平
\end{itemize}

\subsubsection{2.
中国城乡差异}\label{ux4e2dux56fdux57ceux4e61ux5deeux5f02}

1998---2020 年中国自杀城乡分布显示：

\begin{itemize}
\tightlist
\item
  \textbf{农村高于城市}
\end{itemize}

\subsubsection{3. 性别差异}\label{ux6027ux522bux5deeux5f02-1}

【考点】

\begin{itemize}
\tightlist
\item
  \textbf{自杀死亡（completed suicide）男女比约 3:1}
\item
  \textbf{自杀未遂（suicide attempt）男女比约 1:3}
\end{itemize}

老师用很口语化的话解释：

\begin{itemize}
\tightlist
\item
  女性``要死要活''更多，但不一定成功
\item
  男性尝试少一些，但一旦实施，成功率更高
\end{itemize}

\subsubsection{4. 家族聚集性}\label{ux5bb6ux65cfux805aux96c6ux6027}

自杀具有一定\textbf{家族聚集倾向}，提示：

\begin{itemize}
\tightlist
\item
  遗传因素
\item
  家庭环境因素
\end{itemize}

都很重要。

\subsubsection{5.
全球青年死因排序}\label{ux5168ux7403ux9752ux5e74ux6b7bux56e0ux6392ux5e8f}

在 \textbf{15～29 岁}年龄段，全球前四位死因包括：

\begin{enumerate}
\def\labelenumi{\arabic{enumi}.}
\tightlist
\item
  道路交通事故
\item
  肺结核
\item
  人际暴力
\item
  \textbf{自杀}
\end{enumerate}

【考点】 自杀在青年人中是非常重要的死亡原因。

\begin{center}\rule{0.5\linewidth}{0.5pt}\end{center}

\subsection{十五、大学生自杀：老师重点展开}\label{ux5341ux4e94ux5927ux5b66ux751fux81eaux6740ux8001ux5e08ux91cdux70b9ux5c55ux5f00}

老师引用关于中国大学生自杀的流行病学研究。

\subsubsection{1.
大学生自杀率并不一定高于同龄总体人群}\label{ux5927ux5b66ux751fux81eaux6740ux7387ux5e76ux4e0dux4e00ux5b9aux9ad8ux4e8eux540cux9f84ux603bux4f53ux4ebaux7fa4}

根据以 2008---2010 年中国高校自杀数据为基础的研究：

\begin{itemize}
\tightlist
\item
  我国大学生年均自杀率约为 \textbf{10 万分之 1.24}
\item
  低于青少年总体自杀率 \textbf{10 万分之 3.01}
\item
  也低于我国 20～24 岁总体人群自杀率 （2009 年该年龄段为 \textbf{10
  万分之 3.79}）
\end{itemize}

与美国相比：

\begin{itemize}
\tightlist
\item
  美国大学生自杀率约 \textbf{10 万分之 7.5}
\item
  中国大学生不到其五分之一
\end{itemize}

\subsubsection{2.
但仍绝不能掉以轻心}\label{ux4f46ux4ecdux7eddux4e0dux80fdux6389ux4ee5ux8f7bux5fc3}

【考点/老师强调】 即便按 \textbf{10 万分之 1.24}
计算，因为我国高校在校大学生总数超过 \textbf{3500 万}，所以：

\begin{itemize}
\tightlist
\item
  每年自杀大学生绝对数量仍会\textbf{超过 400 人}
\item
  平均每天都有\textbf{1 名以上}大学生自杀死亡
\end{itemize}

因此：

\begin{itemize}
\tightlist
\item
  不能放松预防和干预
\item
  应持续加强大学生心理危机干预工作
\end{itemize}

\subsubsection{3.
大学生自杀的一些特点}\label{ux5927ux5b66ux751fux81eaux6740ux7684ux4e00ux4e9bux7279ux70b9}

教育部资助项目研究显示：

\begin{itemize}
\tightlist
\item
  女生自杀率高于男生
\item
  家庭经济贫困者比例较高
\item
  研究生自杀现象在多省市、重点高校均有发生
\item
  自杀原因多元，但相对集中
\end{itemize}

主要相关因素包括：

\begin{itemize}
\tightlist
\item
  学习竞争
\item
  就业竞争
\item
  失恋（仍是重要原因之一）
\end{itemize}

\subsubsection{4.
常见方式与``网约自杀''}\label{ux5e38ux89c1ux65b9ux5f0fux4e0eux7f51ux7ea6ux81eaux6740}

大学生自杀方式趋于多元，但：

\begin{itemize}
\tightlist
\item
  \textbf{校内跳楼}仍是最主要方式
\end{itemize}

自 2014 年以来还出现：

\begin{itemize}
\tightlist
\item
  \textbf{网约自杀}
\end{itemize}

即有自杀意念者通过网络社交平台相约一起赴死。 老师指出：

\begin{itemize}
\tightlist
\item
  这些人彼此可能根本不认识
\item
  因为怕一个人死、对死亡又恐惧，所以``找人作伴''
\item
  社交软件会强化自杀意念，起到催化作用
\end{itemize}

【老师观点】

\begin{itemize}
\tightlist
\item
  关闭相关 QQ 群、微信群等是必要的
\item
  但关闭群并不能解决这些人的根本问题
\item
  关键仍是帮助他们解决自身难以处理的人生问题
\end{itemize}

\subsubsection{5. 纪念日}\label{ux7eaaux5ff5ux65e5}

\begin{itemize}
\tightlist
\item
  \textbf{世界预防自杀日：9 月 10 日}
\item
  与教师节是同一天
\end{itemize}

\begin{center}\rule{0.5\linewidth}{0.5pt}\end{center}

\subsection{十六、自杀相关概念辨析}\label{ux5341ux516dux81eaux6740ux76f8ux5173ux6982ux5ff5ux8fa8ux6790}

老师专门讲了几个容易混淆的概念。

\subsubsection{1. 自杀意念（suicidal
ideation）}\label{ux81eaux6740ux610fux5ff5suicidal-ideation}

\begin{itemize}
\tightlist
\item
  有``想死''的念头
\item
  但尚未转化为行为
\end{itemize}

\subsubsection{2. 自杀计划（suicide
plan）}\label{ux81eaux6740ux8ba1ux5212suicide-plan}

比意念更进一步，已经想好：

\begin{itemize}
\tightlist
\item
  什么时间
\item
  什么地点
\item
  用什么方式
\end{itemize}

\subsubsection{3.
自伤行为（self-harm）}\label{ux81eaux4f24ux884cux4e3aself-harm}

如：

\begin{itemize}
\tightlist
\item
  割腕留下疤痕
\end{itemize}

老师提醒：自伤行为有时容易被忽视，但也属于非常重要的危险信号。

\subsubsection{4. 自杀未遂（suicide
attempt）}\label{ux81eaux6740ux672aux9042suicide-attempt}

已经实施自杀，但未成功：

\begin{itemize}
\tightlist
\item
  可能自己中止
\item
  也可能被别人救下
\end{itemize}

\subsubsection{5. 自杀企图}\label{ux81eaux6740ux4f01ux56fe}

老师特别澄清：

\begin{itemize}
\tightlist
\item
  \textbf{``自杀企图''其实是``自杀未遂''的另一种说法}
\item
  并不是单纯只有``想法''
\end{itemize}

\subsubsection{6. 自杀死亡（completed
suicide）}\label{ux81eaux6740ux6b7bux4ea1completed-suicide}

最终因自杀而死亡。

\begin{center}\rule{0.5\linewidth}{0.5pt}\end{center}

\subsection{十七、两个特殊概念：扩大性自杀与间接自杀}\label{ux5341ux4e03ux4e24ux4e2aux7279ux6b8aux6982ux5ff5ux6269ux5927ux6027ux81eaux6740ux4e0eux95f4ux63a5ux81eaux6740}

\subsubsection{1. 扩大性自杀}\label{ux6269ux5927ux6027ux81eaux6740}

又称：

\begin{itemize}
\tightlist
\item
  \textbf{怜悯性杀亲}
\end{itemize}

指个体不仅自己想死，还会认为：

\begin{itemize}
\tightlist
\item
  ``我死后孩子/家人怎么办？''
\item
  ``不如先把他们也带走''
\end{itemize}

于是：

\begin{itemize}
\tightlist
\item
  先杀家人
\item
  再自杀
\end{itemize}

\subsubsection{2. 间接自杀}\label{ux95f4ux63a5ux81eaux6740}

又称：

\begin{itemize}
\tightlist
\item
  \textbf{曲线自杀}
\end{itemize}

指个体自己想死，但又害怕直接自杀，于是：

\begin{itemize}
\tightlist
\item
  通过杀人等行为
\item
  希望获得法律严惩
\item
  借此实现``死亡目的''
\end{itemize}

其特点是：

\begin{itemize}
\tightlist
\item
  被害人往往和施害者没有重大利益纠葛
\item
  甚至可能是陌生人
\end{itemize}

\begin{center}\rule{0.5\linewidth}{0.5pt}\end{center}

\subsection{十八、两个案例：扩大性自杀与间接自杀}\label{ux5341ux516bux4e24ux4e2aux6848ux4f8bux6269ux5927ux6027ux81eaux6740ux4e0eux95f4ux63a5ux81eaux6740}

\subsubsection{1.
扩大性自杀案例}\label{ux6269ux5927ux6027ux81eaux6740ux6848ux4f8b}

老师放了关于``重庆母子三人家中死亡''的报道。

【临床关联】该案特点：

\begin{itemize}
\item
  母亲患抑郁症多年
\item
  在卧室烧炭自杀
\item
  两个孩子也一氧化碳中毒死亡
\item
  留有遗书，表达：

  \begin{itemize}
  \tightlist
  \item
    自杀与长期抑郁有关
  \item
    老公很好
  \item
    对不起家人
  \item
    希望不要解剖、不要调查、不要上新闻
  \end{itemize}
\end{itemize}

【老师借案例强调】

\begin{itemize}
\tightlist
\item
  这类悲剧原本有机会预防
\item
  对长期抑郁患者的陪伴与监护非常重要
\end{itemize}

\subsubsection{2.
间接自杀案例}\label{ux95f4ux63a5ux81eaux6740ux6848ux4f8b}

老师播放了``自己想死但没勇气，于是杀人求死刑''的报道。

【案例要点】

\begin{itemize}
\tightlist
\item
  行为人因感情受挫、自卑、心理崩溃
\item
  杀害与自己无仇怨的儿童
\item
  一开始表示就是想借法律极刑来实现死亡
\item
  但到审判后期又请求从轻处罚
\end{itemize}

【老师分析】 这恰恰反映出自杀者常处于一种：

\begin{itemize}
\tightlist
\item
  犹豫
\item
  摇摆
\item
  矛盾
\end{itemize}

的心理状态。

【危机干预启示】

\begin{itemize}
\tightlist
\item
  自杀者并非总是``坚定求死''
\item
  外界刺激、某些报道、某种认同，都可能让他突然更靠近自杀行为
\item
  这正是危机干预的重要切入点
\end{itemize}

【举例】老师还举到：

\begin{itemize}
\tightlist
\item
  济南抑郁症患者杀全家后自杀
\item
  浙江男子为情所困杀死女友后自杀
\item
  深圳男子割喉杀邻居求死刑
\item
  上海男子故意杀人求死刑
\end{itemize}

这些都属于扩大性自杀或间接自杀的相关案例。

\begin{center}\rule{0.5\linewidth}{0.5pt}\end{center}

\subsection{十九、自杀方式的分布特点}\label{ux5341ux4e5dux81eaux6740ux65b9ux5f0fux7684ux5206ux5e03ux7279ux70b9}

老师讲到不同人群的自杀方式有差异。

\subsubsection{1. 年龄差异}\label{ux5e74ux9f84ux5deeux5f02}

\begin{itemize}
\tightlist
\item
  服农药、跳楼：随年龄增长而减少
\item
  上吊、开煤气：随年龄增长而增加
\end{itemize}

\subsubsection{2. 性别差异}\label{ux6027ux522bux5deeux5f02-2}

\begin{itemize}
\item
  男性更多采用\textbf{激烈方式}

  \begin{itemize}
  \tightlist
  \item
    上吊
  \item
    跳楼
  \end{itemize}
\item
  女性更多采用相对``温和''的方式

  \begin{itemize}
  \tightlist
  \item
    服药/服农药
  \item
    跳河
  \end{itemize}
\end{itemize}

\subsubsection{3. 城乡差异}\label{ux57ceux4e61ux5deeux5f02}

\begin{itemize}
\item
  城市居民：

  \begin{itemize}
  \tightlist
  \item
    开煤气
  \item
    跳楼
  \end{itemize}
\item
  农村居民：

  \begin{itemize}
  \tightlist
  \item
    服农药
  \item
    跳河
  \end{itemize}
\end{itemize}

【考点/干预启示】 自杀方式与``工具获取便利性''密切相关。
因此，\textbf{减少自杀工具的可得性}，也是自杀危机干预的重要切入点。

\begin{center}\rule{0.5\linewidth}{0.5pt}\end{center}

\subsection{二十、自杀的成因：生物---心理---社会多因素模型}\label{ux4e8cux5341ux81eaux6740ux7684ux6210ux56e0ux751fux7269ux5fc3ux7406ux793eux4f1aux591aux56e0ux7d20ux6a21ux578b}

\subsubsection{1.
医学遗传学方面}\label{ux533bux5b66ux9057ux4f20ux5b66ux65b9ux9762}

有自杀或自杀未遂家族史者，其亲属中：

\begin{itemize}
\tightlist
\item
  有自杀行为
\item
  有自杀观念
\end{itemize}

的比例明显更高，提示存在家族聚集性。

\subsubsection{2. 应激-易感模型（stress-diathesis
model）}\label{ux5e94ux6fc0-ux6613ux611fux6a21ux578bstress-diathesis-model}

老师提到，应激易感模型认为：

\begin{itemize}
\tightlist
\item
  遗传因素
\item
  后天环境因素
\end{itemize}

共同导致个体形成某种易感素质。

相关因素包括：

\begin{itemize}
\tightlist
\item
  早年创伤
\item
  慢性疾病
\item
  长期物质滥用
\item
  内外环境共同作用
\end{itemize}

\subsubsection{3.
神经生化方面}\label{ux795eux7ecfux751fux5316ux65b9ux9762}

在应激环境下：

\begin{itemize}
\tightlist
\item
  \textbf{蓝斑（locus coeruleus）}神经活动增强
\item
  \textbf{去甲肾上腺素（norepinephrine, NE）}释放和代谢增加
\end{itemize}

提示 NE 在应激处理里作用重要。

老师还提到：

\begin{itemize}
\tightlist
\item
  \textbf{酪氨酸羟化酶（tyrosine hydroxylase）}是 NE 合成酶
\item
  其基因与适应障碍（adjustment disorder）存在相关性
\end{itemize}

此外：

\begin{itemize}
\tightlist
\item
  脑脊液中\textbf{5-羟色胺（serotonin, 5-HT）}及其代谢产物
  \textbf{5-羟吲哚乙酸（5-HIAA）}水平低下，与自杀显著相关
\item
  可能因中枢 5-HT 降低，导致冲动、攻击性增加
\item
  血浆胆固醇降低，也可能与自杀意念相关
\end{itemize}

\begin{center}\rule{0.5\linewidth}{0.5pt}\end{center}

\subsection{二十一、社会心理因素：哪些人更容易走向自杀}\label{ux4e8cux5341ux4e00ux793eux4f1aux5fc3ux7406ux56e0ux7d20ux54eaux4e9bux4ebaux66f4ux5bb9ux6613ux8d70ux5411ux81eaux6740}

\subsubsection{1.
社会支持不足}\label{ux793eux4f1aux652fux6301ux4e0dux8db3}

高危因素包括：

\begin{itemize}
\tightlist
\item
  社会隔离
\item
  孤独感强
\item
  社会支持差
\item
  缺乏可依赖的亲密关系
\end{itemize}

\subsubsection{2.
负性生活事件}\label{ux8d1fux6027ux751fux6d3bux4e8bux4ef6}

自杀前往往存在各种负性事件。

\textbf{青年人}常见：

\begin{itemize}
\tightlist
\item
  人际冲突
\item
  人际关系丧失
\end{itemize}

也可有：

\begin{itemize}
\tightlist
\item
  法律问题
\item
  经济问题
\end{itemize}

\textbf{老年人}常见：

\begin{itemize}
\tightlist
\item
  健康恶化
\end{itemize}

\subsubsection{3.
儿童青少年早期经历}\label{ux513fux7ae5ux9752ux5c11ux5e74ux65e9ux671fux7ecfux5386}

以下因素会影响人格发展，增加成年后自杀风险：

\begin{itemize}
\tightlist
\item
  父母离异
\item
  父母去世
\item
  虐待
\item
  忽视
\item
  严厉体罚
\item
  情感虐待
\item
  校园暴力/恐吓
\end{itemize}

\subsubsection{4.
家庭精神病理背景}\label{ux5bb6ux5eadux7cbeux795eux75c5ux7406ux80ccux666f}

自杀者父母中，以下问题比例更高：

\begin{itemize}
\tightlist
\item
  情绪障碍
\item
  物质滥用
\item
  攻击行为
\end{itemize}

\begin{center}\rule{0.5\linewidth}{0.5pt}\end{center}

\subsection{二十二、自杀者的常见心理特征}\label{ux4e8cux5341ux4e8cux81eaux6740ux8005ux7684ux5e38ux89c1ux5fc3ux7406ux7279ux5f81}

\subsubsection{1. 认知特点}\label{ux8ba4ux77e5ux7279ux70b9}

\begin{itemize}
\tightlist
\item
  认知不良
\item
  非此即彼
\item
  以偏概全
\item
  容易走极端
\item
  容易钻牛角尖
\item
  宿命论倾向
\item
  认为痛苦不可避免、问题无法解决、无法承受
\item
  习惯从阴暗面看待人、事、己、社会
\item
  容易自卑、自尊过强、怀有偏见和敌意
\item
  缺乏洞察、分析和处理问题能力
\end{itemize}

\subsubsection{2. 情感特点}\label{ux60c5ux611fux7279ux70b9}

\begin{itemize}
\tightlist
\item
  长期慢性痛苦
\item
  焦虑
\item
  抑郁
\item
  愤怒
\item
  厌倦
\item
  内疚
\item
  缺乏精神支柱
\item
  强烈无望感
\end{itemize}

\subsubsection{3.
意志与行为特点}\label{ux610fux5fd7ux4e0eux884cux4e3aux7279ux70b9}

\begin{itemize}
\tightlist
\item
  冲动性
\item
  盲目性
\item
  不计后果
\item
  人际交往少
\item
  回避社交
\item
  难以获得支持
\item
  适应能力差
\item
  对新环境不易适应
\item
  具有一定攻击性
\end{itemize}

\begin{center}\rule{0.5\linewidth}{0.5pt}\end{center}

\subsection{二十三、自杀的社会性因素}\label{ux4e8cux5341ux4e09ux81eaux6740ux7684ux793eux4f1aux6027ux56e0ux7d20}

\subsubsection{1. 年龄}\label{ux5e74ux9f84}

\begin{itemize}
\item
  自杀率总体随年龄增长而升高
\item
  进入老年后更明显
\item
  10 岁以下儿童少见，但有\textbf{低龄化趋势}
\item
  多数国家有两个高峰：

  \begin{itemize}
  \tightlist
  \item
    \textbf{15～35 岁}
  \item
    \textbf{65 岁以上}
  \end{itemize}
\end{itemize}

【考点】大学生群体正处于其中一个高峰年龄段。

\subsubsection{2. 婚姻与家庭}\label{ux5a5aux59fbux4e0eux5bb6ux5ead}

\begin{itemize}
\item
  独身、离婚、丧偶者自杀率高于婚姻稳定者
\item
  已婚者中：

  \begin{itemize}
  \tightlist
  \item
    \textbf{无子女者}自杀率高于有子女者
  \end{itemize}
\end{itemize}

\subsubsection{3. 社会阶层}\label{ux793eux4f1aux9636ux5c42}

老师提到美国研究显示：

\begin{itemize}
\item
  蓝领工人自杀率较低
\item
  两端人群更高：

  \begin{itemize}
  \tightlist
  \item
    高社会地位高压力职业：医生、律师、作家等
  \item
    社会底层：贫困、无固定职业者
  \end{itemize}
\end{itemize}

老师的解释是：

\begin{itemize}
\tightlist
\item
  中间层生活较稳定，心理冲击相对没那么极端
\item
  两端人群要么经济压力过大，要么社会地位高、精神压力极大
\end{itemize}

\subsubsection{4. 宗教与文化}\label{ux5b97ux6559ux4e0eux6587ux5316}

【举例】

\begin{itemize}
\tightlist
\item
  基督教文化中，自杀被视为罪
\item
  某些文化中，自杀可能被赋予特殊意义
\end{itemize}

老师用日本切腹文化举例，说明文化环境也会影响自杀发生。

\begin{center}\rule{0.5\linewidth}{0.5pt}\end{center}

\subsection{二十四、与自杀密切相关的精神障碍与躯体疾病}\label{ux4e8cux5341ux56dbux4e0eux81eaux6740ux5bc6ux5207ux76f8ux5173ux7684ux7cbeux795eux969cux788dux4e0eux8eafux4f53ux75beux75c5}

\subsubsection{1. 精神障碍}\label{ux7cbeux795eux969cux788d}

\paragraph{（1）抑郁症（depression）}\label{ux6291ux90c1ux75c7depression-1}

由于：

\begin{itemize}
\tightlist
\item
  无助
\item
  无望
\item
  无用
\end{itemize}

而出现自杀观念。

\paragraph{（2）双相障碍（bipolar
disorder）}\label{ux53ccux76f8ux969cux788dbipolar-disorder}

因既可有躁狂发作，也可有抑郁发作，因此也与自杀密切相关。

\paragraph{（3）PTSD}\label{ptsd}

前面已讲，是重要相关疾病。

\paragraph{（4）精神分裂症（schizophrenia）}\label{ux7cbeux795eux5206ux88c2ux75c7schizophrenia-1}

自杀可来自两个方面：

\begin{itemize}
\item
  \textbf{命令性幻听（command hallucination）}

  \begin{itemize}
  \tightlist
  \item
    比如幻听命令其跳楼
  \end{itemize}
\item
  \textbf{精神分裂症后抑郁}

  \begin{itemize}
  \tightlist
  \item
    患者在病情部分恢复后意识到自己患病
  \item
    因病耻感、长期治疗、对预后失去信心等而抑郁并自杀
  \end{itemize}
\end{itemize}

\paragraph{（5）焦虑障碍（anxiety
disorders）}\label{ux7126ux8651ux969cux788danxiety-disorders}

严重焦虑、烦躁难忍时，也可能出现自杀行为。

\paragraph{（6）老年期痴呆}\label{ux8001ux5e74ux671fux75f4ux5446}

尤其老年男性，因从``家庭支柱''变为``功能下降、需要照顾的人''，容易产生：

\begin{itemize}
\tightlist
\item
  自卑
\item
  拖累家人的感觉
\item
  自杀观念
\end{itemize}

\subsubsection{2. 躯体疾病}\label{ux8eafux4f53ux75beux75c5}

在自杀死亡者中，有相当比例同时存在躯体疾病。

高风险疾病包括：

\begin{itemize}
\tightlist
\item
  脑损伤
\item
  癫痫
\item
  帕金森病（Parkinson disease）
\item
  癌症（cancer）
\item
  糖尿病（diabetes mellitus）等
\end{itemize}

【原因】

\begin{itemize}
\tightlist
\item
  功能受限
\item
  社交受影响
\item
  慢性疼痛
\item
  毁容
\item
  对治疗失去信心
\item
  经济负担重
\item
  觉得拖累家人
\end{itemize}

\begin{center}\rule{0.5\linewidth}{0.5pt}\end{center}

\subsection{二十五、关于自杀的常见误解，老师逐条纠正}\label{ux4e8cux5341ux4e94ux5173ux4e8eux81eaux6740ux7684ux5e38ux89c1ux8befux89e3ux8001ux5e08ux9010ux6761ux7ea0ux6b63}

\subsubsection{1.
``自杀不可预防''}\label{ux81eaux6740ux4e0dux53efux9884ux9632}

错。 大多数自杀是\textbf{可以预防}的。

\subsubsection{2.
``自杀没有预兆''}\label{ux81eaux6740ux6ca1ux6709ux9884ux5146}

错。 多数人在自杀前都有各种预兆和表达。

\subsubsection{3.
``自杀的人就是真的一心求死''}\label{ux81eaux6740ux7684ux4ebaux5c31ux662fux771fux7684ux4e00ux5fc3ux6c42ux6b7b}

不完全对。 很多人其实处于\textbf{矛盾、犹豫、摇摆}状态。

\subsubsection{4.
``谈自杀的人不会真的去死''}\label{ux8c08ux81eaux6740ux7684ux4ebaux4e0dux4f1aux771fux7684ux53bbux6b7b}

错。 很多人会在自杀前直接或间接表达过其意愿。

\subsubsection{5.
``有自杀想法的人就是精神病''}\label{ux6709ux81eaux6740ux60f3ux6cd5ux7684ux4ebaux5c31ux662fux7cbeux795eux75c5}

这是错误贴标签。 会使其感到侮辱和歧视，反而加重风险。

\subsubsection{6.
``不能和有自杀念头的人谈自杀''}\label{ux4e0dux80fdux548cux6709ux81eaux6740ux5ff5ux5934ux7684ux4ebaux8c08ux81eaux6740}

错。 恰恰应该适当沟通，以便：

\begin{itemize}
\tightlist
\item
  识别风险
\item
  评估企图
\item
  让对方感受到被理解、支持、关爱
\end{itemize}

\subsubsection{7.
``自杀危机缓解了就没事了''}\label{ux81eaux6740ux5371ux673aux7f13ux89e3ux4e86ux5c31ux6ca1ux4e8bux4e86}

错。 自杀危机干预需要\textbf{持续进行}，不能只看眼前好转。

\begin{center}\rule{0.5\linewidth}{0.5pt}\end{center}

\subsection{二十六、自杀的先兆表现：老师反复强调要会识别}\label{ux4e8cux5341ux516dux81eaux6740ux7684ux5148ux5146ux8868ux73b0ux8001ux5e08ux53cdux590dux5f3aux8c03ux8981ux4f1aux8bc6ux522b}

\subsubsection{1. 高危线索}\label{ux9ad8ux5371ux7ebfux7d22}

\begin{itemize}
\tightlist
\item
  最近曾自杀过
\item
  有明确自杀计划
\item
  家庭中刚有人自杀
\item
  近期重大失落
\item
  对生活失去兴趣
\item
  感到无助、绝望
\item
  突然饮酒/吸毒
\item
  生活习惯突然改变
\item
  把珍贵、纪念性物品送人
\item
  情绪明显不稳定
\item
  经常谈到死亡
\end{itemize}

\subsubsection{2. 言语线索}\label{ux8a00ux8bedux7ebfux7d22}

\textbf{直接表达}

\begin{itemize}
\tightlist
\item
  ``我不想活了''
\item
  ``我想自杀''
\end{itemize}

\textbf{间接表达}

\begin{itemize}
\tightlist
\item
  ``没人能帮我''
\item
  ``没有我大家都会更好''
\item
  ``我受不了了''
\item
  ``生活没意义''
\item
  ``再也忍不了这种痛苦了''
\end{itemize}

还包括：

\begin{itemize}
\tightlist
\item
  谈论自杀
\item
  拿自杀开玩笑
\item
  谈自己准备怎么死
\item
  与他人告别
\item
  主动获取自杀方式或工具等
\end{itemize}

\subsubsection{3. 行为线索}\label{ux884cux4e3aux7ebfux7d22}

\begin{itemize}
\tightlist
\item
  严重抑郁
\item
  突然行为改变
\item
  白日梦增多
\item
  幻觉出现等
\end{itemize}

【老师提醒】
新闻里有人跳楼/跳桥但最后没跳，不应仅仅理解为``闹事''\,``讨债''``获利手段''。
在危机干预中，仍要把这类表现按\textbf{真实自杀风险}来对待。

\begin{center}\rule{0.5\linewidth}{0.5pt}\end{center}

\subsection{二十七、几个特别需要提高警惕的人群与时段}\label{ux4e8cux5341ux4e03ux51e0ux4e2aux7279ux522bux9700ux8981ux63d0ux9ad8ux8b66ux60d5ux7684ux4ebaux7fa4ux4e0eux65f6ux6bb5}

老师再次强调，自杀高危人群包括：

\begin{itemize}
\tightlist
\item
  抑郁症患者
\item
  精神分裂症患者
\item
  酒依赖患者
\end{itemize}

【易错】尤其要警惕一种情况：

如果这类患者的情绪\textbf{突然``好转''}，不一定真的是病好了，
有时反而可能是：

\begin{itemize}
\tightlist
\item
  已经下定决心准备自杀
\item
  情感上出现了``解脱感''
\end{itemize}

老师还提到：

\begin{itemize}
\tightlist
\item
  抑郁症自杀不一定发生在最重的时候
\item
  在疾病缓解期也可能风险很高
\item
  某些抗抑郁药说明书中甚至会提示：早期可能\textbf{不降低反而增加}部分患者自杀风险
\end{itemize}

【原因理解】可能是：

\begin{itemize}
\tightlist
\item
  患者情绪尚未明显改善
\item
  但行动能力先恢复了一部分
\item
  于是更有能力把自杀想法转化为行为
\end{itemize}

\begin{center}\rule{0.5\linewidth}{0.5pt}\end{center}

\subsection{二十八、本节内容的结尾情况}\label{ux4e8cux5341ux516bux672cux8282ux5185ux5bb9ux7684ux7ed3ux5c3eux60c5ux51b5}

本次转写在最后一段中断，老师最后正在讲：

\begin{itemize}
\tightlist
\item
  抑郁症患者在疾病高峰期与缓解期的自杀风险差异
\item
  抗抑郁药与自杀风险之间的复杂关系
\end{itemize}

因此这部分内容在原始录音中应还有后续，但当前总稿到此结束。

\begin{center}\rule{0.5\linewidth}{0.5pt}\end{center}

\subsection{本节英文术语索引}\label{ux672cux8282ux82f1ux6587ux672fux8bedux7d22ux5f15-2}

\begin{itemize}
\tightlist
\item
  危机 --- crisis
\item
  心理危机 --- psychological crisis
\item
  心理危机干预 --- crisis intervention
\item
  发展性危机 --- developmental crisis
\item
  境遇性危机 --- situational crisis
\item
  存在性危机 --- existential crisis
\item
  冲击期 --- impact phase
\item
  防御期 --- defensive phase
\item
  解决期 --- resolution phase
\item
  成长期 --- growth phase
\item
  国际疾病分类 --- International Classification of Diseases（ICD）
\item
  ICD 第 11 版 --- ICD-11
\item
  精神障碍诊断与统计手册 --- Diagnostic and Statistical Manual of Mental
  Disorders（DSM）
\item
  DSM 第 5 版 --- DSM-5
\item
  创伤及应激相关障碍 --- Trauma- and Stressor-Related Disorders
\item
  急性应激障碍 --- Acute Stress Disorder（ASD）
\item
  创伤后应激障碍 --- Post-Traumatic Stress Disorder（PTSD）
\item
  焦虑 --- anxiety
\item
  抑郁 --- depression
\item
  自杀 --- suicide
\item
  自杀意念 --- suicidal ideation
\item
  自杀计划 --- suicide plan
\item
  自杀未遂 / 自杀企图 --- suicide attempt
\item
  自伤行为 --- self-harm
\item
  幻觉 --- hallucination
\item
  妄想 --- delusion
\item
  木僵 --- stupor
\item
  定向障碍 --- disorientation
\item
  闪回 --- flashback
\item
  人格解体 --- depersonalization
\item
  现实解体 --- derealization
\item
  警觉性增高 --- hyperarousal
\item
  支持疗法 --- supportive therapy
\item
  认知行为治疗 --- Cognitive Behavioral Therapy（CBT）
\item
  暴露疗法 --- exposure therapy
\item
  β受体阻滞剂 --- beta-blocker
\item
  苯二氮卓类 --- benzodiazepines
\item
  选择性 5-羟色胺再摄取抑制剂 --- Selective Serotonin Reuptake
  Inhibitors（SSRI）
\item
  5-羟色胺-去甲肾上腺素再摄取抑制剂 --- Serotonin-Norepinephrine
  Reuptake Inhibitors（SNRI）
\item
  舍曲林 --- sertraline
\item
  氟西汀 --- fluoxetine
\item
  帕罗西汀 --- paroxetine
\item
  文拉法辛 --- venlafaxine
\item
  哌唑嗪 --- prazosin
\item
  曲唑酮 --- trazodone
\item
  米氮平 --- mirtazapine
\item
  奎硫平 --- quetiapine
\item
  氯氮平 --- clozapine
\item
  阿立哌唑 --- aripiprazole（Abilify）
\item
  阿普唑仑 --- alprazolam（Xanax）
\item
  生物反馈治疗 --- biofeedback
\item
  无抽搐电休克治疗 --- modified electroconvulsive therapy（MECT）
\item
  重复经颅磁刺激 --- repetitive transcranial magnetic
  stimulation（rTMS）
\item
  惊恐发作 --- panic attack
\item
  双相障碍 --- bipolar disorder
\item
  精神分裂症 --- schizophrenia
\item
  命令性幻听 --- command hallucination
\item
  心境稳定剂 --- mood stabilizer
\item
  酒依赖 --- alcohol dependence
\item
  物质滥用 --- substance abuse
\item
  应激-易感模型 --- stress-diathesis model
\item
  适应障碍 --- adjustment disorder
\item
  蓝斑 --- locus coeruleus
\item
  去甲肾上腺素 --- norepinephrine（NE）
\item
  酪氨酸羟化酶 --- tyrosine hydroxylase
\item
  5-羟色胺 --- serotonin（5-HT）
\item
  5-羟吲哚乙酸 --- 5-hydroxyindoleacetic acid（5-HIAA）
\item
  帕金森病 --- Parkinson disease
\item
  血气胸 --- hemopneumothorax
\end{itemize}

如果你愿意，我下一步可以把这份笔记再压缩成一版``考前速记版''。

\clearpage

\section{第4讲：预防、依恋与危机干预}\label{ux7b2c4ux8bb2ux9884ux9632ux4f9dux604bux4e0eux5371ux673aux5e72ux9884}

\subsection{一、课程导入：什么是心理危机，为什么要讲``预防''}\label{ux4e00ux8bfeux7a0bux5bfcux5165ux4ec0ux4e48ux662fux5fc3ux7406ux5371ux673aux4e3aux4ec0ux4e48ux8981ux8bb2ux9884ux9632}

老师先说明：本节课重点不是单纯重复``什么是心理危机''，而是从\textbf{预防（prevention）}的角度来理解心理危机。

\begin{itemize}
\item
  所谓\textbf{心理危机（psychological
  crisis）}，本质上是一种\textbf{心理失衡状态}：

  \begin{itemize}
  \tightlist
  \item
    正常状态下，人能较好处理学习、人际关系和日常任务。
  \item
    危机状态下，原有的内在平衡和生活秩序被打破。
  \end{itemize}
\end{itemize}

【举例】老师用大学生日常情境说明危机如何被触发：

\begin{itemize}
\item
  考试挂科、多门不及格，学业亮红灯。
\item
  面临分流、四六级不过。
\item
  就业压力大，反复面试却不断被刷。
\item
  对多数同学来说，这些可能只是``人生的一道坎''，经过短期调整后可以重新恢复平衡。
\item
  但对少部分同学而言，这些事件会被体验为\textbf{重大挫折}，进而引发：

  \begin{itemize}
  \tightlist
  \item
    对自我能力的怀疑；
  \item
    对未来发展的悲观；
  \item
    对自我价值的否定；
  \item
    甚至出现抑郁倾向，或已达到\textbf{抑郁症（depression）}诊断标准；
  \item
    严重时可能无法上课、无法与他人建立关系。
  \end{itemize}
\end{itemize}

老师强调：

\begin{itemize}
\item
  如果一个人能把危机事件处理好，重新恢复平衡，这一过程也可能带来成长：

  \begin{itemize}
  \tightlist
  \item
    重新思考``我是谁''；
  \item
    重新思考``我适合做什么''；
  \item
    重新发现``我在哪些方面有价值''。
  \end{itemize}
\item
  但如果：

  \begin{itemize}
  \tightlist
  \item
    缺乏有效应对方式；
  \item
    缺乏朋友、家庭等支持系统；
  \item
    周围人不能理解、支持和关怀；
  \item
    仅靠自己又无法解决眼前困难；
  \end{itemize}

  就更容易发展为真正的心理危机。
\end{itemize}

\begin{center}\rule{0.5\linewidth}{0.5pt}\end{center}

\subsection{二、心理危机的影响因素：先看哪些人风险更高}\label{ux4e8cux5fc3ux7406ux5371ux673aux7684ux5f71ux54cdux56e0ux7d20ux5148ux770bux54eaux4e9bux4ebaux98ceux9669ux66f4ux9ad8}

老师提出：要做预防，首先要知道\textbf{哪些人更容易陷入心理危机}。

\subsubsection{1. 生物学因素（biological
factors）}\label{ux751fux7269ux5b66ux56e0ux7d20biological-factors}

老师指出，心理危机并不完全是后天造成的，\textbf{先天差异}也很重要。

\begin{itemize}
\item
  有调查显示：\textbf{30\%～50\% 的自杀风险与遗传因素有关。}
\item
  人与人之间本来就有差异，包括：

  \begin{itemize}
  \tightlist
  \item
    \textbf{神经冲动水平}不同；
  \item
    \textbf{敏感性（sensitivity）}不同；
  \item
    有些人属于\textbf{高敏感人群}。
  \end{itemize}
\end{itemize}

【举例】同样被人批评，有的人``脸皮厚一点''就过去了；但高敏感个体会：

\begin{itemize}
\tightlist
\item
  当场感觉崩溃；
\item
  事后反复回想；
\item
  一直纠结``是不是我哪里不好''。
\end{itemize}

老师进一步指出：

\begin{itemize}
\tightlist
\item
  有的人天生对刺激更敏感；
\item
  有的人天生``神经比较大条''，对刺激不太敏感；
\item
  有的人先天更容易\textbf{冲动（impulsivity）}；
\item
  有的人则更稳重，遇事能先让自己平复。
\end{itemize}

【举例】同样是被老师批评、或与室友争吵：

\begin{itemize}
\tightlist
\item
  有的人会先稳住自己；
\item
  有的人则容易在高唤起状态下做出\textbf{不计后果}的行为。
\end{itemize}

结论： 这些差异并不只是后天教养造成，也与\textbf{遗传和生物学基础}有关。

\begin{center}\rule{0.5\linewidth}{0.5pt}\end{center}

\subsection{三、心理因素：哪些人格与认知模式容易把人推向危机}\label{ux4e09ux5fc3ux7406ux56e0ux7d20ux54eaux4e9bux4ebaux683cux4e0eux8ba4ux77e5ux6a21ux5f0fux5bb9ux6613ux628aux4ebaux63a8ux5411ux5371ux673a}

老师接着从心理层面分析心理危机的高危因素。

\subsubsection{1.
易损人格特质：完美主义（perfectionism）}\label{ux6613ux635fux4ebaux683cux7279ux8d28ux5b8cux7f8eux4e3bux4e49perfectionism}

老师指出，\textbf{完美主义人格}是重要的风险因素。

\begin{itemize}
\tightlist
\item
  一般来说，越是高成就群体，越容易出现较高比例的完美主义特质。
\item
  这也解释了为什么一些``看起来很优秀''的人，心理危机发生率反而不低。
\end{itemize}

老师强调，完美主义有两个维度：

\paragraph{（1）积极维度：高标准、高期待}\label{ux79efux6781ux7ef4ux5ea6ux9ad8ux6807ux51c6ux9ad8ux671fux5f85}

\begin{itemize}
\tightlist
\item
  希望自己做得很好；
\item
  对自己有要求；
\item
  这一面能够促进成就。
\end{itemize}

\paragraph{（2）消极维度：高批评（self-criticism）}\label{ux6d88ux6781ux7ef4ux5ea6ux9ad8ux6279ux8bc4self-criticism}

\begin{itemize}
\item
  一旦做不好，就强烈批评自己：

  \begin{itemize}
  \tightlist
  \item
    ``我怎么可以做不好？''
  \item
    ``我怎么这么差劲？''
  \end{itemize}
\item
  如果这一维度高，就更容易出现：

  \begin{itemize}
  \tightlist
  \item
    心理压力；
  \item
    焦虑（anxiety）；
  \item
    抑郁（depression）；
  \item
    自杀倾向（suicidal tendency）风险升高。
  \end{itemize}
\end{itemize}

【举例】同样是挂科或分数不理想：

\begin{itemize}
\item
  一般人可能会想``这次没发挥好''\,``我没准备好''；
\item
  高批评型完美主义者会想：

  \begin{itemize}
  \tightlist
  \item
    ``我怎么能允许自己考这么差？''
  \item
    ``别人会怎么看我？''
  \item
    ``我是不是根本不行？''
  \end{itemize}
\end{itemize}

老师总结：
完美主义，尤其是其中\textbf{高自我批评}的部分，与抑郁、焦虑、自杀倾向都呈\textbf{高相关}。

\begin{center}\rule{0.5\linewidth}{0.5pt}\end{center}

\subsubsection{2.
冲动性（impulsivity）与情绪调节失调}\label{ux51b2ux52a8ux6027impulsivityux4e0eux60c5ux7eeaux8c03ux8282ux5931ux8c03}

老师说明，危机与\textbf{冲动性}也高度相关。

冲动性有两层来源：

\paragraph{（1）先天神经冲动性差异}\label{ux5148ux5929ux795eux7ecfux51b2ux52a8ux6027ux5deeux5f02}

\begin{itemize}
\tightlist
\item
  有的人本来就更容易冲动。
\end{itemize}

\paragraph{（2）情绪调节功能失调}\label{ux60c5ux7eeaux8c03ux8282ux529fux80fdux5931ux8c03}

\begin{itemize}
\item
  正常成年人在面对刺激时，通常有一定\textbf{情绪调节（emotion
  regulation）}能力：

  \begin{itemize}
  \tightlist
  \item
    能换个角度思考；
  \item
    能自我安抚；
  \item
    不至于陷入持续自责。
  \end{itemize}
\item
  但冲动特质明显的人，往往缺乏这种调节能力。
\end{itemize}

老师提到一个重要概念：

\begin{itemize}
\tightlist
\item
  \textbf{反刍性思维（rumination）}：反复陷在负性思维循环里，持续自责、自我否定。
\end{itemize}

这类个体常见特点：

\begin{itemize}
\tightlist
\item
  思维容易钻牛角尖；
\item
  看问题僵化；
\item
  觉得``事情就是这样，没法改变''；
\item
  很难从辩证、发展的角度理解问题。
\end{itemize}

【老师概括危机中常见状态】

很多陷入危机的同学会说：

\begin{itemize}
\tightlist
\item
  ``没有希望了。''
\item
  ``未来没有希望了。''
\item
  ``这件事没办法解决了。''
\end{itemize}

老师认为，很多危机并非因为``真的毫无出路''，而是因为：

\begin{itemize}
\tightlist
\item
  思维僵化；
\item
  只看到唯一的狭窄出口；
\item
  看不到其他可能性。
\end{itemize}

\begin{center}\rule{0.5\linewidth}{0.5pt}\end{center}

\subsubsection{3. 失调的认知图式（cognitive
schema）}\label{ux5931ux8c03ux7684ux8ba4ux77e5ux56feux5f0fcognitive-schema}

老师接着引入``\textbf{认知图式（schema）}''概念。

所谓认知图式，是指一个人：

\begin{itemize}
\tightlist
\item
  如何看待自己；
\item
  如何看待他人；
\item
  如何看待世界。
\end{itemize}

这些图式来自成长过程中的长期经验积累。

\paragraph{（1）关于自我的图式}\label{ux5173ux4e8eux81eaux6211ux7684ux56feux5f0f}

老师说明：

\begin{itemize}
\item
  早期我们往往是通过别人的反馈来认识自己。
\item
  如果经常被表扬，就更容易觉得：

  \begin{itemize}
  \tightlist
  \item
    自己是好的；
  \item
    自己是可爱的。
  \end{itemize}
\item
  如果长期被批评、被否定，就可能形成：

  \begin{itemize}
  \tightlist
  \item
    ``我很差劲''；
  \item
    ``我连小事都做不好''；
  \item
    ``我不值得被肯定''。
  \end{itemize}
\end{itemize}

\paragraph{（2）关于他人的图式}\label{ux5173ux4e8eux4ed6ux4ebaux7684ux56feux5f0f}

\begin{itemize}
\tightlist
\item
  别人是否可靠？
\item
  我能不能信任别人？
\item
  别人靠近我是安全的还是危险的？
\end{itemize}

\paragraph{（3）关于世界的图式}\label{ux5173ux4e8eux4e16ux754cux7684ux56feux5f0f}

\begin{itemize}
\tightlist
\item
  外界是安全的还是危险的？
\item
  他人是可依靠的还是会伤害我的？
\end{itemize}

老师强调，如果一个人形成了\textbf{负性认知图式}，那么在接受外界信息时，就更容易：

\begin{itemize}
\tightlist
\item
  选择性关注负面信息；
\item
  用原有模板解释世界；
\item
  反复验证自己的负性假设。
\end{itemize}

【举例】

如果一个人内心已有``我不可爱''\,``我不值得被爱''的图式，那么：

\begin{itemize}
\item
  别人聚会没叫自己，就会自动解读为：

  \begin{itemize}
  \tightlist
  \item
    ``他们是不是不喜欢我？''
  \item
    ``我是不是不值得被关注？''
  \end{itemize}
\item
  这会进一步强化其负性自我认识，形成\textbf{恶性循环}。
\end{itemize}

老师总结：

\begin{itemize}
\tightlist
\item
  容易抑郁的人，更容易用抑郁的方式解读世界；
\item
  容易焦虑的人，更容易用焦虑的方式解读世界；
\item
  问题不一定在外界，而可能在于\textbf{加工外界信息的方式存在偏差}。
\end{itemize}

\begin{center}\rule{0.5\linewidth}{0.5pt}\end{center}

\subsubsection{4. 问题解决不足（poor problem
solving）}\label{ux95eeux9898ux89e3ux51b3ux4e0dux8db3poor-problem-solving}

老师强调，很多危机个体还有一个共同点：\textbf{问题解决策略不足}。

\begin{itemize}
\tightlist
\item
  他们看不到外界资源；
\item
  看不到帮助；
\item
  觉得别人不会理解自己；
\item
  也不相信别人能帮到自己。
\end{itemize}

【举例：考试焦虑】

老师提到，有些大学生在考试前焦虑到：

\begin{itemize}
\tightlist
\item
  根本无法去考试；
\item
  觉得``天塌了''；
\item
  因为从小到大，自己一直依靠考试来获得认可；
\item
  到大学发现自己不再有明显优势时，就开始严重怀疑自我价值。
\end{itemize}

但老师举例说，其实这类问题也可以有解决策略，例如：

\begin{itemize}
\tightlist
\item
  申请\textbf{缓考}；
\item
  暂时回家休息；
\item
  调整心态后再回来考试。
\end{itemize}

老师指出：

\begin{itemize}
\tightlist
\item
  虽然这种做法带有一定回避色彩，但它依然是一种\textbf{问题解决策略}；
\item
  许多同学陷入危机，不是因为绝对没有办法，而是因为\textbf{看不到办法}。
\end{itemize}

老师的核心观点是：

\begin{itemize}
\tightlist
\item
  在任何糟糕的状态下，都仍然存在某种解决问题的方法；
\item
  危机中最可怕的，不是困难本身，而是\textbf{看不到出路}。
\end{itemize}

\begin{center}\rule{0.5\linewidth}{0.5pt}\end{center}

\subsubsection{5.
对自身问题解决能力缺乏信心}\label{ux5bf9ux81eaux8eabux95eeux9898ux89e3ux51b3ux80fdux529bux7f3aux4e4fux4fe1ux5fc3}

老师进一步指出，有些人即使已经看到策略，仍然会怀疑自己：

\begin{itemize}
\tightlist
\item
  ``这样行吗？''
\item
  ``我能做到吗？''
\item
  ``这个办法靠谱吗？''
\end{itemize}

【举例】初中、高中拒学的孩子会认为：

\begin{itemize}
\tightlist
\item
  不上学，别的路也走不通；
\item
  上学这条路又被自己堵死了；
\item
  其他事自己也做不成；
\item
  好像每一条路都没有出口。
\end{itemize}

老师在这里强调一个重要观点：

\begin{itemize}
\tightlist
\item
  每个人都有个体差异；
\item
  不是每个人都擅长学习；
\item
  但每个人都一定有自己擅长的领域；
\item
  只是在危机中，人往往\textbf{看不到自己的价值}和可能性。
\end{itemize}

\begin{center}\rule{0.5\linewidth}{0.5pt}\end{center}

\subsection{四、应激事件：危机的``扳机点''（trigger）}\label{ux56dbux5e94ux6fc0ux4e8bux4ef6ux5371ux673aux7684ux6273ux673aux70b9trigger}

老师指出：

\begin{itemize}
\tightlist
\item
  生物学因素、心理因素是\textbf{易感因素}；
\item
  \textbf{应激事件（stressful
  event）}则像一个``\textbf{扳机点（trigger）}''，会把这些潜在风险激活。
\end{itemize}

【举例：失恋】

\begin{itemize}
\item
  对大多数人来说，失恋虽痛苦，但经过两三个月、半年，通常可以逐渐恢复；
\item
  但有些人会因为失恋出现：

  \begin{itemize}
  \tightlist
  \item
    跑去对方学校或单位大吵大闹；
  \item
    寻死觅活；
  \item
    做出极端行为。
  \end{itemize}
\end{itemize}

老师用比喻说明：

\begin{itemize}
\tightlist
\item
  同样一块石头扔进水里，每个人激起的涟漪不一样。
\item
  不同反应，和一个人的人格特质、心理状态、神经冲动性、依恋状况等都有关。
\end{itemize}

【举例：四级没过】

老师提到一个真实类型案例：

\begin{itemize}
\item
  某同学因为\textbf{英语四级没过}，最终做出极端行为。
\item
  问题不只是``四级没过''本身，而是这个同学把：

  \begin{itemize}
  \tightlist
  \item
    自我价值；
  \item
    能力感；
  \item
    未来前途；
  \end{itemize}

  全都和四级绑定在一起。
\end{itemize}

因此，某个具体事件就成了危机的启动按钮。

\begin{center}\rule{0.5\linewidth}{0.5pt}\end{center}

\subsection{五、应激事件的常见类型}\label{ux4e94ux5e94ux6fc0ux4e8bux4ef6ux7684ux5e38ux89c1ux7c7bux578b}

老师继续列举多种可能激发危机的应激事件。

\subsubsection{1. 突发事件}\label{ux7a81ux53d1ux4e8bux4ef6}

包括：

\begin{itemize}
\tightlist
\item
  家庭重大变故；
\item
  失恋；
\item
  自然灾害；
\item
  社会公共意外事件；
\item
  火灾、地震等；
\item
  疫情等公共卫生事件。
\end{itemize}

【举例】老师提到：

\begin{itemize}
\tightlist
\item
  地震、火灾、社会公共事件；
\item
  疫情最初阶段，特别是身处疫区或一线的医护人员，也会面临较大的心理危机。
\end{itemize}

\subsubsection{2.
严重身体疾病}\label{ux4e25ux91cdux8eabux4f53ux75beux75c5}

老师强调：

\begin{itemize}
\tightlist
\item
  身体健康与心理健康密切相关；
\item
  如果一个人身患重病，尤其是绝症，其心理状态通常会受到重大影响。
\end{itemize}

并指出，这类人往往会经历一个过程：

\begin{itemize}
\tightlist
\item
  最开始难以置信、难以接受；
\item
  之后逐渐接受；
\item
  再慢慢修复。
\end{itemize}

\subsubsection{3.
环境适应不良}\label{ux73afux5883ux9002ux5e94ux4e0dux826f}

老师特别强调了``\textbf{适应}''本身就是危机因素。

【复习/延伸】老师提到一项经验性观察：

\begin{itemize}
\item
  高中、初中毕业班学生的心理健康状况，反而可能好于刚入学的新生。
\item
  即：

  \begin{itemize}
  \tightlist
  \item
    高一、初一等刚进入新环境的学生，
  \item
    有时比高三、初三更需要关注。
  \end{itemize}
\end{itemize}

原因在于：\textbf{适应压力}。

\begin{center}\rule{0.5\linewidth}{0.5pt}\end{center}

\subsection{六、大学适应问题：宿舍、学习、人际、自律都可能成为危机来源}\label{ux516dux5927ux5b66ux9002ux5e94ux95eeux9898ux5bbfux820dux5b66ux4e60ux4ebaux9645ux81eaux5f8bux90fdux53efux80fdux6210ux4e3aux5371ux673aux6765ux6e90}

老师把重点落在大学生常见危机情境上。

\subsubsection{1.
第一次住校、过集体生活的适应}\label{ux7b2cux4e00ux6b21ux4f4fux6821ux8fc7ux96c6ux4f53ux751fux6d3bux7684ux9002ux5e94}

老师提醒大家不要小看``住校''这件事。

\begin{itemize}
\tightlist
\item
  有人喜欢热闹，适应得很好；
\item
  有人却觉得非常折磨。
\end{itemize}

主要冲突包括：

\begin{itemize}
\tightlist
\item
  熄灯时间与作息不同；
\item
  室友夜间活跃、讲话、听音乐、打游戏；
\item
  自己想睡觉却被影响；
\item
  宿舍文化和个人习惯冲突。
\end{itemize}

\subsubsection{2.
私人边界与宿舍摩擦}\label{ux79c1ux4ebaux8fb9ux754cux4e0eux5bbfux820dux6469ux64e6}

【举例】老师提到常见矛盾：

\begin{itemize}
\item
  室友未经同意使用自己的肥皂、盆、耳机、充电器等。
\item
  有的人会觉得：

  \begin{itemize}
  \tightlist
  \item
    ``都这么熟了，用一下还要打招呼吗？''
  \end{itemize}
\item
  但另一些同学真正介意的不是``东西本身''，而是：

  \begin{itemize}
  \tightlist
  \item
    自己的私人空间是否被尊重；
  \item
    自己作为个体是否被尊重。
  \end{itemize}
\end{itemize}

老师指出：

\begin{itemize}
\tightlist
\item
  看似``鸡毛蒜皮''的小事，
\item
  往往正是宿舍关系破裂、冲突升级的根源。
\end{itemize}

\subsubsection{3.
大学生活的自由与空虚}\label{ux5927ux5b66ux751fux6d3bux7684ux81eaux7531ux4e0eux7a7aux865a}

老师指出，从高中进入大学后，许多外部管束突然减少：

\begin{itemize}
\tightlist
\item
  没有人天天催你学习；
\item
  老师上完课就不再紧盯；
\item
  家长也不在身边时时管着。
\end{itemize}

短期看会觉得很自由，但长期可能出现：

\begin{itemize}
\tightlist
\item
  空虚；
\item
  缺乏规划；
\item
  自律不足。
\end{itemize}

【举例】老师说，有同学大学四年一直打游戏，直到找工作时突然觉醒，发现自己：

\begin{itemize}
\tightlist
\item
  好像一无所获；
\item
  没学到什么；
\item
  成长很慢；
\end{itemize}

于是开始出现抑郁和焦虑。

\subsubsection{4.
学业压力与学习困难}\label{ux5b66ux4e1aux538bux529bux4e0eux5b66ux4e60ux56f0ux96be}

老师再次提到：

\begin{itemize}
\tightlist
\item
  多门不及格；
\item
  学业压力大；
\item
  学习困难；
\item
  学习不再是自己优势时的自我怀疑。
\end{itemize}

\subsubsection{5.
家庭贫困与经济负担}\label{ux5bb6ux5eadux8d2bux56f0ux4e0eux7ecfux6d4eux8d1fux62c5}

\begin{itemize}
\tightlist
\item
  担心家庭经济压力；
\item
  想通过兼职减轻父母负担；
\item
  又难以平衡学业与打工之间的关系。
\end{itemize}

\subsubsection{6.
人际关系失调、就业困难、考研失败}\label{ux4ebaux9645ux5173ux7cfbux5931ux8c03ux5c31ux4e1aux56f0ux96beux8003ux7814ux5931ux8d25}

包括：

\begin{itemize}
\tightlist
\item
  投简历石沉大海；
\item
  毕业求职困难；
\item
  考研失败；
\item
  对未来极度焦虑。
\end{itemize}

\subsubsection{7.
同伴危机的``传染性''影响}\label{ux540cux4f34ux5371ux673aux7684ux4f20ux67d3ux6027ux5f71ux54cd}

老师指出，身边同学出现危机，也会影响自己。

【举例】

\begin{itemize}
\tightlist
\item
  有同学长期接触一个抑郁状态很重的朋友；
\item
  对方天天发大量负面、抑郁图片；
\item
  时间久了，自己的情绪也被明显拖低。
\end{itemize}

老师借此说明： \textbf{同龄人对同龄人的影响是非常大的。}

\begin{center}\rule{0.5\linewidth}{0.5pt}\end{center}

\subsection{七、创伤（trauma）：心理危机背后的深层因素}\label{ux4e03ux521bux4f24traumaux5fc3ux7406ux5371ux673aux80ccux540eux7684ux6df1ux5c42ux56e0ux7d20}

老师继续分析心理危机的另一个重要来源：\textbf{创伤（trauma）}。

\subsubsection{1.
创伤的基本含义}\label{ux521bux4f24ux7684ux57faux672cux542bux4e49}

老师给出的核心理解是：

\begin{itemize}
\item
  创伤不是单纯``遇到困难''；
\item
  而是个体处于\textbf{威胁性环境}中，且其防御和应对能力不足，导致一种：

  \begin{itemize}
  \tightlist
  \item
    无助感；
  \item
    无法防备感；
  \item
    内在理解被撼动的经历。
  \end{itemize}
\end{itemize}

创伤会长期影响一个人对以下内容的理解：

\begin{itemize}
\tightlist
\item
  对自己的理解；
\item
  对他人的理解；
\item
  对世界的理解。
\end{itemize}

\subsubsection{2.
创伤会动摇人的基本信念}\label{ux521bux4f24ux4f1aux52a8ux6447ux4ebaux7684ux57faux672cux4fe1ux5ff5}

老师指出，人在正常状态下通常有一些基本信念，例如：

\begin{itemize}
\tightlist
\item
  生活是可预测的；
\item
  我对生活有一定掌控感；
\item
  世界基本是公平的；
\item
  我是有价值的；
\item
  危险离我很远。
\end{itemize}

而在创伤状态下，这些信念会被动摇：

\begin{itemize}
\tightlist
\item
  生活变得不可预测；
\item
  世界不再公平；
\item
  自我价值感下降；
\item
  危险感显著上升。
\end{itemize}

【举例】

\begin{itemize}
\tightlist
\item
  发生地震的人，会觉得世界突然失控；
\item
  得绝症的人，会觉得``为什么这种事会落到我头上''；
\item
  经历战争的人，不会再觉得``危险离我很远''。
\end{itemize}

\begin{center}\rule{0.5\linewidth}{0.5pt}\end{center}

\subsection{八、依恋（attachment）与依恋创伤：危机的关键背景}\label{ux516bux4f9dux604battachmentux4e0eux4f9dux604bux521bux4f24ux5371ux673aux7684ux5173ux952eux80ccux666f}

在第一段末尾，老师引出``\textbf{依恋（attachment）}''概念，第二段接着重点展开。

\subsubsection{1.
为什么要讲依恋创伤}\label{ux4e3aux4ec0ux4e48ux8981ux8bb2ux4f9dux604bux521bux4f24}

老师说明：

\begin{itemize}
\tightlist
\item
  很多创伤、很多心理危机，和\textbf{依恋创伤（attachment
  trauma）}密切相关；
\item
  在各种创伤中，依恋创伤占有很大比重。
\end{itemize}

\subsubsection{2. 依恋的定义}\label{ux4f9dux604bux7684ux5b9aux4e49}

依恋是指：

\begin{itemize}
\tightlist
\item
  人在出生后与\textbf{重要抚养者（primary
  caregiver）}建立起来的一种情感连接。
\end{itemize}

老师强调，依恋对整个人的发展都很重要，因为它是：

\begin{itemize}
\tightlist
\item
  安全感（sense of security）的基础；
\item
  基本信任（basic trust）的基础。
\end{itemize}

\subsubsection{3.
零到一岁的核心心理任务：建立信任感}\label{ux96f6ux5230ux4e00ux5c81ux7684ux6838ux5fc3ux5fc3ux7406ux4efbux52a1ux5efaux7acbux4fe1ux4efbux611f}

老师指出，在心理发展中：

\begin{itemize}
\tightlist
\item
  \textbf{0～1 岁}的重要任务就是建立\textbf{信任感（trust）}。
\end{itemize}

原因是：

\begin{itemize}
\item
  婴儿无法独立生存；
\item
  需要喂养、保暖、安抚；
\item
  如果其需要能够被及时、稳定地回应，就会逐渐形成：

  \begin{itemize}
  \tightlist
  \item
    世界是可信的；
  \item
    世界是安全的；
  \item
    我是可爱的；
  \item
    我是有能力的。
  \end{itemize}
\end{itemize}

也就是说，早期被如何对待，会逐渐沉淀为：

\begin{itemize}
\tightlist
\item
  我怎么看自己；
\item
  我怎么看别人；
\item
  我怎么看这个世界。
\end{itemize}

\begin{center}\rule{0.5\linewidth}{0.5pt}\end{center}

\subsection{九、依恋实验：小猴子更需要``情感抚慰''而不只是喂养}\label{ux4e5dux4f9dux604bux5b9eux9a8cux5c0fux7334ux5b50ux66f4ux9700ux8981ux60c5ux611fux629aux6170ux800cux4e0dux53eaux662fux5582ux517b}

【举例】老师讲到经典心理学实验：

\begin{itemize}
\item
  把刚出生的小猴子放在两个``模拟母猴''前面：

  \begin{itemize}
  \tightlist
  \item
    一个是\textbf{有奶瓶的铁丝母猴}；
  \item
    一个是\textbf{包着柔软毛巾的母猴}。
  \end{itemize}
\item
  结果是：

  \begin{itemize}
  \tightlist
  \item
    小猴子只有在饿得不行时，才去有奶瓶的铁丝母猴那里吃奶；
  \item
    其余大部分时间，都待在柔软的``毛巾母猴''身边。
  \end{itemize}
\end{itemize}

老师借此说明：

\begin{itemize}
\tightlist
\item
  婴儿需要的并不仅仅是食物，
\item
  还需要\textbf{情感上的依附、安全与安抚}。
\end{itemize}

\begin{center}\rule{0.5\linewidth}{0.5pt}\end{center}

\subsection{十、依恋的形成：生命早期``跟谁长大''非常重要}\label{ux5341ux4f9dux604bux7684ux5f62ux6210ux751fux547dux65e9ux671fux8ddfux8c01ux957fux5927ux975eux5e38ux91cdux8981}

老师指出：

\begin{itemize}
\tightlist
\item
  在生命最早期，尤其\textbf{0～3 岁}，
\item
  跟谁一起长大，对后续人格和关系模式影响很大。
\end{itemize}

可能是：

\begin{itemize}
\tightlist
\item
  母亲；
\item
  父亲；
\item
  姥姥姥爷；
\item
  爷爷奶奶；
\item
  或其他主要照顾者。
\end{itemize}

通常：

\begin{itemize}
\tightlist
\item
  早期长期陪伴、照顾自己的人，
\item
  往往就是日后自己最亲近、最有情感连接的人。
\end{itemize}

老师把这概括为： 这是人生命早期形成的、与重要他人的情感连接。

\begin{center}\rule{0.5\linewidth}{0.5pt}\end{center}

\subsection{十一、依恋如何影响成人后的恋爱与人际关系}\label{ux5341ux4e00ux4f9dux604bux5982ux4f55ux5f71ux54cdux6210ux4ebaux540eux7684ux604bux7231ux4e0eux4ebaux9645ux5173ux7cfb}

老师指出：

\begin{itemize}
\tightlist
\item
  婴儿期的依恋经历，会成为儿童期乃至\textbf{成人期关系模式的模板}；
\item
  包括恋爱关系，也是一种依恋关系。
\end{itemize}

【举例】老师说，很多人会发现：

\begin{itemize}
\tightlist
\item
  自己谈恋爱时，总在找``同一类人''；
\item
  即使换了对象，某些人格特质、心理特征还是很相似。
\end{itemize}

这说明：

\begin{itemize}
\tightlist
\item
  人在关系中有自己的\textbf{内在模板（internal working model）}；
\item
  会反复被某类关系吸引。
\end{itemize}

【举例】

有些人在亲密关系中会觉得：

\begin{itemize}
\tightlist
\item
  ``太近了我就想跑。''
\item
  ``关系一靠近，我就觉得不安全。''
\item
  ``我会觉得被吞没、被控制。''
\end{itemize}

老师认为，这类表现常与早期依恋经验有关。

\begin{center}\rule{0.5\linewidth}{0.5pt}\end{center}

\subsection{十二、四种依恋类型（attachment
styles）}\label{ux5341ux4e8cux56dbux79cdux4f9dux604bux7c7bux578battachment-styles}

老师接着系统讲了依恋类型。

\subsubsection{1. 安全型依恋（secure
attachment）}\label{ux5b89ux5168ux578bux4f9dux604bsecure-attachment}

在儿童期的表现：

\begin{itemize}
\tightlist
\item
  母亲离开时会难过；
\item
  重新相聚后很快恢复开心。
\end{itemize}

在成人后的表现：

\begin{itemize}
\tightlist
\item
  在浪漫关系中较自由、自主；
\item
  进入婚姻或家庭关系后，距离感较灵活；
\item
  既能享受亲密，也能享受独处；
\item
  人际关系总体更适应性。
\end{itemize}

老师指出：

\begin{itemize}
\tightlist
\item
  \textbf{安全型依恋在人群中比例最高}。
\end{itemize}

\begin{center}\rule{0.5\linewidth}{0.5pt}\end{center}

\subsubsection{2. 回避型依恋（avoidant
attachment）}\label{ux56deux907fux578bux4f9dux604bavoidant-attachment}

在儿童期的表现：

\begin{itemize}
\tightlist
\item
  照顾者离开或回来，表面上都``无所谓''；
\item
  不是内心真的不在乎，而是表现出退缩、忽视和拒绝依恋需求。
\end{itemize}

在成人后的表现：

\begin{itemize}
\tightlist
\item
  关系中容易疏远、退缩；
\item
  情感回避；
\item
  不愿谈内心感受；
\item
  较回避身体和情感上的接近。
\end{itemize}

\begin{center}\rule{0.5\linewidth}{0.5pt}\end{center}

\subsubsection{3. 焦虑型 / 抵抗型依恋（anxious / resistant
attachment）}\label{ux7126ux8651ux578b-ux62b5ux6297ux578bux4f9dux604banxious-resistant-attachment}

在儿童期的表现：

\begin{itemize}
\tightlist
\item
  对照顾者离开非常敏感、警觉；
\item
  强烈黏着照顾者；
\item
  特别害怕分离。
\end{itemize}

在成人后的表现：

\begin{itemize}
\tightlist
\item
  亲密关系中容易过度接近；
\item
  纠缠；
\item
  侵入性强；
\item
  边界不清；
\item
  非常害怕被抛弃。
\end{itemize}

\begin{center}\rule{0.5\linewidth}{0.5pt}\end{center}

\subsubsection{4. 紊乱型依恋（disorganized
attachment）}\label{ux7d0aux4e71ux578bux4f9dux604bdisorganized-attachment}

特点是：

\begin{itemize}
\tightlist
\item
  同时带有回避型和焦虑型的特征；
\item
  人际边界更混乱；
\item
  关系模式更不稳定。
\end{itemize}

\begin{center}\rule{0.5\linewidth}{0.5pt}\end{center}

\subsection{十三、依恋创伤如何在危机中被激活}\label{ux5341ux4e09ux4f9dux604bux521bux4f24ux5982ux4f55ux5728ux5371ux673aux4e2dux88abux6fc0ux6d3b}

老师强调：

\begin{itemize}
\tightlist
\item
  创伤，尤其是依恋创伤，常常会在后来的关系事件中被重新激活。
\end{itemize}

【举例：失恋】

有些人失恋后反应极端，不只是因为``分手''本身，而是因为：

\begin{itemize}
\tightlist
\item
  分手激活了早年被拒绝、被抛弃的体验；
\item
  再次唤起了那种原始的崩溃感。
\end{itemize}

老师指出：

\begin{itemize}
\tightlist
\item
  很多危机看起来只是``危机事件''，
\item
  其实背后是\textbf{旧创伤被点燃}了。
\end{itemize}

并进一步强调：

\begin{itemize}
\tightlist
\item
  依恋不仅影响正常心理发展，
\item
  也会调节一个人的\textbf{生理唤醒水平}和对刺激的反应；
\item
  同样的应激事件，之所以有人反应更大，往往与早年依恋创伤被触发有关。
\end{itemize}

\begin{center}\rule{0.5\linewidth}{0.5pt}\end{center}

\subsection{十四、创伤会让人更渴求关系，但又可能不信任关系}\label{ux5341ux56dbux521bux4f24ux4f1aux8ba9ux4ebaux66f4ux6e34ux6c42ux5173ux7cfbux4f46ux53c8ux53efux80fdux4e0dux4fe1ux4efbux5173ux7cfb}

老师在这里提出一个关键判断：

\begin{itemize}
\tightlist
\item
  \textbf{任何创伤都会激活人的依恋系统。}
\end{itemize}

因为当一个人陷入创伤时，他会感到：

\begin{itemize}
\tightlist
\item
  生活失衡；
\item
  内心失控；
\item
  安全感被破坏。
\end{itemize}

而在这种情况下，人最本能的需要就是：

\begin{itemize}
\tightlist
\item
  抓住一个关系；
\item
  寻找支持。
\end{itemize}

但问题在于：

\begin{itemize}
\item
  如果此人拥有充足社会支持，

  \begin{itemize}
  \tightlist
  \item
    大概率不至于陷入严重危机；
  \end{itemize}
\item
  如果其过去依恋经历不安全，

  \begin{itemize}
  \tightlist
  \item
    会觉得关系靠不住；
  \item
    会觉得关系具有侵入性；
  \item
    最终还是陷入危机和孤立。
  \end{itemize}
\end{itemize}

老师总结一句很重的话：

\begin{itemize}
\tightlist
\item
  \textbf{基于人际关系的创伤最难以忍受，关系越亲密，伤害越重。}
\end{itemize}

【举例】

\begin{itemize}
\tightlist
\item
  有些人一分手就寻死觅活；
\item
  有些人会伤害自己，或者通过服药等激进行为表达痛苦；
\item
  本质上是依恋系统被强烈破坏了。
\end{itemize}

\begin{center}\rule{0.5\linewidth}{0.5pt}\end{center}

\subsection{十五、强迫性重复（repetition
compulsion）：为什么总是找``同一类人''}\label{ux5341ux4e94ux5f3aux8febux6027ux91cdux590drepetition-compulsionux4e3aux4ec0ux4e48ux603bux662fux627eux540cux4e00ux7c7bux4eba}

老师提到一个重要心理现象：\textbf{强迫性重复（repetition compulsion）}。

表现为：

\begin{itemize}
\tightlist
\item
  明明上一段关系很痛苦；
\item
  但下一段仍会被类似的人吸引；
\item
  关系模式不断重演。
\end{itemize}

老师解释：

\begin{itemize}
\tightlist
\item
  熟悉的关系模式会带来一种``可掌控感''；
\item
  哪怕这种熟悉未必健康，它仍会让人感觉安全；
\item
  甚至对于一些人而言，\textbf{糟糕的关系也胜过没有关系}。
\end{itemize}

【举例】

有些同学明明觉得当前恋爱关系并不满意，却还是不愿离开，因为：

\begin{itemize}
\tightlist
\item
  大家都在谈恋爱；
\item
  自己害怕孤单；
\item
  哪怕关系不好，也比完全没有关系更容易承受。
\end{itemize}

\begin{center}\rule{0.5\linewidth}{0.5pt}\end{center}

\subsection{十六、依恋创伤的后果：基本信任被破坏，情绪与认同都会受影响}\label{ux5341ux516dux4f9dux604bux521bux4f24ux7684ux540eux679cux57faux672cux4fe1ux4efbux88abux7834ux574fux60c5ux7eeaux4e0eux8ba4ux540cux90fdux4f1aux53d7ux5f71ux54cd}

老师总结依恋创伤的影响时提到：

\begin{itemize}
\tightlist
\item
  会破坏\textbf{基本信任感（basic trust）}；
\item
  包括对世界、对重要他人的信任。
\end{itemize}

当个体受到刺激时，创伤还可能带来：

\begin{itemize}
\tightlist
\item
  情绪紊乱；
\item
  自我认同下降；
\item
  知觉障碍；
\item
  记忆损害等。
\end{itemize}

\begin{center}\rule{0.5\linewidth}{0.5pt}\end{center}

\subsection{十七、创伤不一定只发生在童年，成年后的重大创伤同样深重}\label{ux5341ux4e03ux521bux4f24ux4e0dux4e00ux5b9aux53eaux53d1ux751fux5728ux7ae5ux5e74ux6210ux5e74ux540eux7684ux91cdux5927ux521bux4f24ux540cux6837ux6df1ux91cd}

老师特别澄清一个常见误解：

\begin{itemize}
\tightlist
\item
  不是说只有童年创伤才重要；
\item
  成年后的重大创伤，同样可能深刻改变一个人。
\end{itemize}

【举例：战争与 PTSD】

老师举例说：

\begin{itemize}
\item
  很多经历过战争的老兵，即使战争结束后已经回归平静生活，
\item
  仍会长期受创伤影响，
\item
  出现\textbf{创伤后应激障碍（Post-Traumatic Stress Disorder,
  PTSD）}表现，如：

  \begin{itemize}
  \tightlist
  \item
    做噩梦；
  \item
    闪回（flashback）；
  \item
    好像再次回到战场。
  \end{itemize}
\end{itemize}

老师借此说明：

\begin{itemize}
\tightlist
\item
  创伤可以贯穿整个人的生命周期；
\item
  只是生命早期因应对能力弱，更容易形成深层影响。
\end{itemize}

\begin{center}\rule{0.5\linewidth}{0.5pt}\end{center}

\subsection{十八、创伤如何在后来事件中被重新唤醒}\label{ux5341ux516bux521bux4f24ux5982ux4f55ux5728ux540eux6765ux4e8bux4ef6ux4e2dux88abux91cdux65b0ux5524ux9192}

【举例】老师举了两个非常典型的例子：

\subsubsection{1.
童年羞耻创伤被研究生阶段激活}\label{ux7ae5ux5e74ux7f9eux803bux521bux4f24ux88abux7814ux7a76ux751fux9636ux6bb5ux6fc0ux6d3b}

\begin{itemize}
\tightlist
\item
  一个孩子8岁时，在比赛中罚丢了关键球；
\item
  教练在队友和父母面前大吼大叫；
\item
  这个羞耻体验未被处理好，埋在心里。
\item
  到了读研后，导师在课题组公开说其研究``是一团垃圾''，
\item
  就可能再次激活当年的羞耻创伤，
\item
  让人回到当年被当众羞辱的情境中。
\end{itemize}

\subsubsection{2.
失恋激活早年被抛弃经历}\label{ux5931ux604bux6fc0ux6d3bux65e9ux5e74ux88abux629bux5f03ux7ecfux5386}

\begin{itemize}
\item
  表面上只是``男朋友/女朋友提出分手''；
\item
  但对有些人而言，真正被唤起的是：

  \begin{itemize}
  \tightlist
  \item
    早年被抛弃；
  \item
    被送人；
  \item
    被拒绝；
  \end{itemize}

  那种深层体验。
\end{itemize}

老师想强调的是：

\begin{itemize}
\tightlist
\item
  很多危机表面上是``现在发生的事''，
\item
  实际上背后连着``很早以前没处理完的事''。
\end{itemize}

\begin{center}\rule{0.5\linewidth}{0.5pt}\end{center}

\subsection{十九、如何预防心理危机：从可改变的部分入手}\label{ux5341ux4e5dux5982ux4f55ux9884ux9632ux5fc3ux7406ux5371ux673aux4eceux53efux6539ux53d8ux7684ux90e8ux5206ux5165ux624b}

老师进入本节核心的``预防''部分，指出：

\begin{itemize}
\tightlist
\item
  生物学因素改变不了；
\item
  但很多心理与行为层面的因素，是可以逐渐调整的。
\end{itemize}

\begin{center}\rule{0.5\linewidth}{0.5pt}\end{center}

\subsection{二十、【学习建议】调整完美主义：允许自己做不到``样样都好''}\label{ux4e8cux5341ux5b66ux4e60ux5efaux8baeux8c03ux6574ux5b8cux7f8eux4e3bux4e49ux5141ux8bb8ux81eaux5df1ux505aux4e0dux5230ux6837ux6837ux90fdux597d}

老师先谈完美主义的调整。

\subsubsection{1.
人格会逐渐固化，但大学阶段仍很有可塑性}\label{ux4ebaux683cux4f1aux9010ux6e10ux56faux5316ux4f46ux5927ux5b66ux9636ux6bb5ux4ecdux5f88ux6709ux53efux5851ux6027}

老师指出：

\begin{itemize}
\tightlist
\item
  人格到成年后会逐渐固化；
\item
  老年人的人格更难改变；
\item
  但大学生还处在很有可塑性的阶段。
\end{itemize}

\subsubsection{2.
不能要求自己所有维度都完美}\label{ux4e0dux80fdux8981ux6c42ux81eaux5df1ux6240ux6709ux7ef4ux5ea6ux90fdux5b8cux7f8e}

【举例】老师列举大学生常见``全都想做好''的状态：

\begin{itemize}
\tightlist
\item
  学习要顶尖；
\item
  社团要活跃；
\item
  还想兼职赚钱；
\item
  还想当班干部；
\item
  还想做老师助手；
\item
  还想学才艺。
\end{itemize}

老师提醒：

\begin{itemize}
\tightlist
\item
  每个人都只有 24 小时；
\item
  时间和精力并不会因为你要求高就自动增加；
\item
  如果什么都想做到很好，只会把自己累垮。
\end{itemize}

\subsubsection{3.
允许自己在``不重要的方面''不那么优秀}\label{ux5141ux8bb8ux81eaux5df1ux5728ux4e0dux91cdux8981ux7684ux65b9ux9762ux4e0dux90a3ux4e48ux4f18ux79c0}

老师的建议是：

\begin{itemize}
\tightlist
\item
  不必在每个方面都做到最好；
\item
  在一些不重要的方面，可以允许自己犯错、允许自己没那么好；
\item
  适度的不完美，反而更容易被人接纳。
\end{itemize}

【学习建议】老师原话核心意思可整理为：

\begin{itemize}
\tightlist
\item
  允许自己犯错；
\item
  承认自己有不足；
\item
  不必处处苛求；
\item
  一个人适当不完美，往往更可爱，也更轻松。
\end{itemize}

\begin{center}\rule{0.5\linewidth}{0.5pt}\end{center}

\subsection{二十一、【学习建议】减少冲动性：先给自己一个``缓冲地带''}\label{ux4e8cux5341ux4e00ux5b66ux4e60ux5efaux8baeux51cfux5c11ux51b2ux52a8ux6027ux5148ux7ed9ux81eaux5df1ux4e00ux4e2aux7f13ux51b2ux5730ux5e26}

对于冲动性较高的人，老师提出非常具体的建议：

\begin{itemize}
\tightlist
\item
  在情绪高唤起的时候，不要立刻行动；
\item
  给自己几秒钟缓冲；
\item
  深呼吸；
\item
  先冷静下来；
\item
  再考虑怎么解决问题。
\end{itemize}

【举例】

如果自己特别容易：

\begin{itemize}
\tightlist
\item
  和别人冲突；
\item
  一激动就想动手；
\item
  一生气就想把对方狠狠干一顿；
\end{itemize}

那么更要训练自己：

\begin{itemize}
\tightlist
\item
  暂停；
\item
  冷静；
\item
  延迟反应。
\end{itemize}

老师指出：

\begin{itemize}
\tightlist
\item
  两个人吵架，往往是越吵越凶；
\item
  如果先分开，各自找空间冷静，很多架其实是吵不起来的。
\end{itemize}

\begin{center}\rule{0.5\linewidth}{0.5pt}\end{center}

\subsection{二十二、【学习建议】修正失调认知图式：重新看待自己、他人和世界}\label{ux4e8cux5341ux4e8cux5b66ux4e60ux5efaux8baeux4feeux6b63ux5931ux8c03ux8ba4ux77e5ux56feux5f0fux91cdux65b0ux770bux5f85ux81eaux5df1ux4ed6ux4ebaux548cux4e16ux754c}

老师强调：

\begin{itemize}
\tightlist
\item
  图式虽然早年形成，但不是永远不能变。
\end{itemize}

\subsubsection{1.
重新看待自己}\label{ux91cdux65b0ux770bux5f85ux81eaux5df1}

如果你总觉得自己很差劲，就要反思：

\begin{itemize}
\tightlist
\item
  我真的一无是处吗？
\item
  我身上真的没有任何闪光点吗？
\end{itemize}

老师指出：

\begin{itemize}
\tightlist
\item
  一个人能够一路活到现在、走到现在，
\item
  一定有一些地方是做得好的；
\item
  只是很多人看不到自己身上的``资源''和``宝藏''。
\end{itemize}

\subsubsection{2.
重新看待别人}\label{ux91cdux65b0ux770bux5f85ux522bux4eba}

如果一个人曾被霸凌、被欺骗，就可能形成：

\begin{itemize}
\tightlist
\item
  人不可靠；
\item
  人让人害怕；
\item
  我要远离人群；
\item
  我要用盔甲把自己包起来。
\end{itemize}

老师承认：

\begin{itemize}
\tightlist
\item
  这种模式曾经可能是保护自己的方式；
\end{itemize}

但老师也提醒：

\begin{itemize}
\tightlist
\item
  当你未来想建立亲密关系时，这种高度防御就会变成障碍；
\item
  别人走不进你，你也无法真正走向别人。
\end{itemize}

【学习建议】老师给出的方向是：

\begin{itemize}
\tightlist
\item
  选择相对安全的人际关系；
\item
  在安全的基础上，一点点``松绑''；
\item
  逐渐放下一部分防御；
\item
  学会在关系中感受安全，而非一直封闭自己。
\end{itemize}

\begin{center}\rule{0.5\linewidth}{0.5pt}\end{center}

\subsection{二十三、【学习建议】提高问题解决能力：不要只靠自己死扛}\label{ux4e8cux5341ux4e09ux5b66ux4e60ux5efaux8baeux63d0ux9ad8ux95eeux9898ux89e3ux51b3ux80fdux529bux4e0dux8981ux53eaux9760ux81eaux5df1ux6b7bux625b}

老师强调，很多危机其实和``不会求助''有关。

\subsubsection{1.
要看到其他人的视角}\label{ux8981ux770bux5230ux5176ux4ed6ux4ebaux7684ux89c6ux89d2}

【举例】老师讲了``多级火箭''灵感故事：

\begin{itemize}
\tightlist
\item
  发明者苦思冥想，想不出怎样靠足够燃料把卫星送上轨道；
\item
  后来受到``吹笛子合奏''的启发，想到多级装置逐步脱落的办法。
\end{itemize}

老师用这个故事想说明：

\begin{itemize}
\tightlist
\item
  很多问题靠一个人是想不到突破口的；
\item
  别人未必是你这个专业的同行；
\item
  但他可能正好能看到你看不到的角度。
\end{itemize}

老师借用古诗意思说明：

\begin{itemize}
\tightlist
\item
  身在局中，反而容易看不清；
\item
  换个视角，可能就能看到新的出路。
\end{itemize}

【学习建议】

\begin{itemize}
\tightlist
\item
  遇到困难要学会求助；
\item
  不要只困在自己的视角里；
\item
  可以向同学、朋友、老师、学长学姐请教；
\item
  通过``多视角''寻找解决策略。
\end{itemize}

\begin{center}\rule{0.5\linewidth}{0.5pt}\end{center}

\subsection{二十四、【学习建议】理解创伤对自己的影响，并尝试疗愈}\label{ux4e8cux5341ux56dbux5b66ux4e60ux5efaux8baeux7406ux89e3ux521bux4f24ux5bf9ux81eaux5df1ux7684ux5f71ux54cdux5e76ux5c1dux8bd5ux7597ux6108}

老师指出：

\begin{itemize}
\item
  如果一个人曾经历过创伤，就要逐渐学会觉察：

  \begin{itemize}
  \tightlist
  \item
    过去的经历是如何影响现在的；
  \item
    自己是否在重复旧模式；
  \item
    是否总在进入类似的关系和困境。
  \end{itemize}
\end{itemize}

【学习建议】

\begin{itemize}
\tightlist
\item
  识别自己是否在重复过去的关系模式；
\item
  一点点从旧创伤里走出来；
\item
  尝试建立更健康的关系；
\item
  允许自己被修复，而不是一直被旧创伤牵着走。
\end{itemize}

\begin{center}\rule{0.5\linewidth}{0.5pt}\end{center}

\subsection{二十五、【学习建议】必要时寻求心理咨询（psychological
counseling）}\label{ux4e8cux5341ux4e94ux5b66ux4e60ux5efaux8baeux5fc5ux8981ux65f6ux5bfbux6c42ux5fc3ux7406ux54a8ux8be2psychological-counseling}

老师明确提到：

\begin{itemize}
\tightlist
\item
  心理咨询是疗愈心理创伤的一种方式。
\end{itemize}

因为：

\begin{itemize}
\item
  很多创伤是人不愿轻易让别人看到的；
\item
  但它们又确实在影响：

  \begin{itemize}
  \tightlist
  \item
    危机发生；
  \item
    亲密关系建立；
  \item
    当下生活功能。
  \end{itemize}
\end{itemize}

老师认为，在一个\textbf{安全的咨询关系}中，可以逐渐做到：

\begin{itemize}
\tightlist
\item
  看见过往经历对自己的影响；
\item
  理解自己当前关系模式的问题；
\item
  学会改善现有关系与应对方式。
\end{itemize}

\begin{center}\rule{0.5\linewidth}{0.5pt}\end{center}

\subsection{二十六、【学习建议】同辈互助：同学、室友、学长学姐都是重要资源}\label{ux4e8cux5341ux516dux5b66ux4e60ux5efaux8baeux540cux8f88ux4e92ux52a9ux540cux5b66ux5ba4ux53cbux5b66ux957fux5b66ux59d0ux90fdux662fux91cdux8981ux8d44ux6e90}

老师转到大学校园现实场景，指出：

\begin{itemize}
\item
  每个班通常都有心理委员、生活委员；
\item
  但更普遍、更容易用到的其实是：

  \begin{itemize}
  \tightlist
  \item
    室友；
  \item
    同伴；
  \item
    学姐学长；
  \item
    学弟学妹；
  \item
    朋友。
  \end{itemize}
\end{itemize}

老师认为：

\begin{itemize}
\tightlist
\item
  对大学生来说，遇到困难时最常见的求助对象就是同龄人；
\item
  因为他们更有类似的生活经历，也更容易理解你的处境。
\end{itemize}

【学习建议】

遇到心理危机时，要学会求助于：

\begin{itemize}
\tightlist
\item
  同伴；
\item
  室友；
\item
  朋友；
\item
  学长学姐；
\item
  学弟学妹等现实支持系统。
\end{itemize}

\begin{center}\rule{0.5\linewidth}{0.5pt}\end{center}

\subsection{二十七、【考点/提醒】作为同伴，也要学会识别危险信号并及时转介}\label{ux4e8cux5341ux4e03ux8003ux70b9ux63d0ux9192ux4f5cux4e3aux540cux4f34ux4e5fux8981ux5b66ux4f1aux8bc6ux522bux5371ux9669ux4fe1ux53f7ux5e76ux53caux65f6ux8f6cux4ecb}

老师特别提醒：

\begin{itemize}
\tightlist
\item
  如果我们自己是被求助对象，
\item
  尤其当发现对方出现\textbf{自杀相关危险征兆}时，
\item
  要有警觉意识。
\end{itemize}

关键原则是：

\begin{itemize}
\tightlist
\item
  识别危险信号；
\item
  及时转介；
\item
  面对紧急情况不能只靠``陪聊''，要升级处理。
\end{itemize}

\begin{center}\rule{0.5\linewidth}{0.5pt}\end{center}

\subsection{二十八、危机干预（crisis
intervention）的基本技巧}\label{ux4e8cux5341ux516bux5371ux673aux5e72ux9884crisis-interventionux7684ux57faux672cux6280ux5de7}

老师随后从``预防''过渡到``危机干预''。

\subsubsection{1. 倾听（listening）}\label{ux503eux542clistening}

【举例】如果一个失恋、情绪极差，甚至有轻生念头的同学来找你：

首先要做的是：\textbf{做一个好的倾听者}。

老师强调：

\begin{itemize}
\tightlist
\item
  不急着评判；
\item
  不抢着讲大道理；
\item
  不在对方没说完时就打断和否定。
\end{itemize}

并指出，倾听不仅适用于危机情境，也适用于日常关系。

\paragraph{倾听包括两部分：}\label{ux503eux542cux5305ux62ecux4e24ux90e8ux5206}

\subparagraph{（1）言语倾听}\label{ux8a00ux8bedux503eux542c}

\begin{itemize}
\tightlist
\item
  用语言鼓励对方继续表达。
\end{itemize}

\subparagraph{（2）非言语倾听}\label{ux975eux8a00ux8bedux503eux542c}

\begin{itemize}
\tightlist
\item
  你的身体在场，心也要在场；
\item
  不是一边玩手机一边敷衍听；
\item
  而是真正面对着对方，认真听他说。
\end{itemize}

\begin{center}\rule{0.5\linewidth}{0.5pt}\end{center}

\subsubsection{2. 共情（empathy）}\label{ux5171ux60c5empathy}

老师强调：\textbf{共情（empathy）}是非常重要的心理学概念。

共情的意思是：

\begin{itemize}
\tightlist
\item
  设身处地理解对方的想法和感受。
\end{itemize}

老师特别区分了：

\begin{itemize}
\tightlist
\item
  \textbf{共情 ≠ 同情（sympathy）}
\end{itemize}

区别在于：

\begin{itemize}
\tightlist
\item
  同情往往带有一点``我高你低、我觉得你很惨''的意味；
\item
  共情则是站到对方的位置上去感受他的处境。
\end{itemize}

\subsubsection{3. 情感反应（emotional
reflection）}\label{ux60c5ux611fux53cdux5e94emotional-reflection}

老师说，在共情基础上，还要做出适当情感回应，例如：

\begin{itemize}
\tightlist
\item
  ``我知道你现在很不容易。''
\item
  ``我知道你现在很难。''
\item
  ``我知道你现在很委屈。''
\item
  ``我知道你现在很愤怒。''
\item
  ``我猜你现在可能会觉得很绝望。''
\end{itemize}

这类回应能让对方感到：

\begin{itemize}
\tightlist
\item
  自己被理解；
\item
  情绪被看见；
\item
  不是孤立无援的。
\end{itemize}

\begin{center}\rule{0.5\linewidth}{0.5pt}\end{center}

\subsubsection{4.
具体化（concretization）}\label{ux5177ux4f53ux5316concretization}

老师提出：要对对方保持\textbf{好奇}，并帮助其把问题说清楚。

例如：

\begin{itemize}
\tightlist
\item
  ``你这两天看起来情绪很低落，是发生什么事情了吗？''
\item
  ``可以跟我讲一讲吗？''
\item
  ``你现在最难受的是什么？''
\end{itemize}

具体化的作用是：

\begin{itemize}
\tightlist
\item
  澄清事件；
\item
  澄清情绪；
\item
  让问题更可被理解和处理。
\end{itemize}

\begin{center}\rule{0.5\linewidth}{0.5pt}\end{center}

\subsubsection{5.
帮助寻找资源（resources）}\label{ux5e2eux52a9ux5bfbux627eux8d44ux6e90resources}

老师强调：

\begin{itemize}
\tightlist
\item
  危机中的人常常只看到绝境；
\item
  旁人可以帮他看到资源。
\end{itemize}

这些资源可能包括：

\begin{itemize}
\tightlist
\item
  同学；
\item
  室友；
\item
  老师；
\item
  家人；
\item
  学长学姐；
\item
  制度性安排；
\item
  现实中的替代方案。
\end{itemize}

\begin{center}\rule{0.5\linewidth}{0.5pt}\end{center}

\subsubsection{6. 拓宽可能性（expand
possibilities）}\label{ux62d3ux5bbdux53efux80fdux6027expand-possibilities}

老师反复强调一个核心思想：

\begin{itemize}
\tightlist
\item
  到达罗马不止一条路；
\item
  当事人看不到的路，别人也许能看到。
\end{itemize}

可以通过：

\begin{itemize}
\tightlist
\item
  多找几个人商量；
\item
  头脑风暴；
\item
  重新评估还有哪些出口。
\end{itemize}

\begin{center}\rule{0.5\linewidth}{0.5pt}\end{center}

\subsubsection{7. 灌注希望（instill
hope）}\label{ux704cux6ce8ux5e0cux671binstill-hope}

老师认为，任何危机干预里都不能缺少``希望''。

核心观点包括：

\begin{itemize}
\tightlist
\item
  不要轻易放弃自己；
\item
  山重水复疑无路，柳暗花明又一村；
\item
  任何时候都有对应时候的解决办法；
\item
  危机不是绝对死局，而是需要``见招拆招''。
\end{itemize}

【学习建议】老师这里延伸出一种对生活的态度：

\begin{itemize}
\tightlist
\item
  不要过度追求一切都可控；
\item
  接受未知；
\item
  接受不确定性；
\item
  用开放心态面对生活的意外；
\item
  学会``见招拆招''，焦虑水平反而会下降。
\end{itemize}

\begin{center}\rule{0.5\linewidth}{0.5pt}\end{center}

\subsection{二十九、什么时候做危机干预：不是越早越好}\label{ux4e8cux5341ux4e5dux4ec0ux4e48ux65f6ux5019ux505aux5371ux673aux5e72ux9884ux4e0dux662fux8d8aux65e9ux8d8aux597d}

老师特别讲了危机干预的时机问题。

【举例：地震、火灾】

老师说：

\begin{itemize}
\item
  危机事件刚发生时，最紧急的不是心理干预；
\item
  而是先满足生存和安全需要，例如：

  \begin{itemize}
  \tightlist
  \item
    食物；
  \item
    水；
  \item
    帐篷；
  \item
    绷带；
  \item
    救援队；
  \item
    搜索与医疗支持。
  \end{itemize}
\end{itemize}

因此：

\begin{itemize}
\tightlist
\item
  心理援助/危机干预\textbf{不是第一秒就上去做}；
\item
  一般在事件发生后 \textbf{24～48 小时} 左右就可以开始介入；
\item
  一周后、一个月后也可以；
\item
  甚至一年后仍可能需要介入。
\end{itemize}

老师解释：

\begin{itemize}
\tightlist
\item
  因为有些人会出现\textbf{PTSD}；
\item
  创伤反应不一定立即出现，
\item
  可能延后很久才显现。
\end{itemize}

\begin{center}\rule{0.5\linewidth}{0.5pt}\end{center}

\subsection{三十、危机/创伤干预的三个阶段}\label{ux4e09ux5341ux5371ux673aux521bux4f24ux5e72ux9884ux7684ux4e09ux4e2aux9636ux6bb5}

老师最后讲到心理创伤和心理危机干预的\textbf{三阶段模型}。

\subsubsection{第一阶段：稳定化（stabilization）}\label{ux7b2cux4e00ux9636ux6bb5ux7a33ux5b9aux5316stabilization}

这是最先要做的。

目标是：

\begin{itemize}
\tightlist
\item
  让一个人先安定下来；
\item
  建立安全感；
\item
  防止情绪完全崩溃。
\end{itemize}

老师提到，心理学里会用一些技巧帮助稳定，例如：

\begin{itemize}
\tightlist
\item
  放松训练；
\item
  稳定化技术；
\item
  让恐惧记忆暂时被``放进保险箱''之类的意象技术。
\end{itemize}

老师特别强调：

\begin{itemize}
\tightlist
\item
  在个体尚未稳定前，不能急着``拉开伤口''处理创伤；
\item
  就像受伤先要\textbf{止血}，而不是马上深挖伤口。
\end{itemize}

\begin{center}\rule{0.5\linewidth}{0.5pt}\end{center}

\subsubsection{第二阶段：创伤暴露与处理（trauma exposure /
processing）}\label{ux7b2cux4e8cux9636ux6bb5ux521bux4f24ux66b4ux9732ux4e0eux5904ux7406trauma-exposure-processing}

当个体已有足够安全感和稳定度后，才可以：

\begin{itemize}
\tightlist
\item
  非常温和地；
\item
  一点一点地；
\item
  慢慢接触和处理创伤内容。
\end{itemize}

老师强调：

\begin{itemize}
\tightlist
\item
  创伤终究需要被处理；
\item
  但处理必须有时机、有节奏，不能粗暴推进。
\end{itemize}

\begin{center}\rule{0.5\linewidth}{0.5pt}\end{center}

\subsubsection{第三阶段：重建意义（meaning
reconstruction）}\label{ux7b2cux4e09ux9636ux6bb5ux91cdux5efaux610fux4e49meaning-reconstruction}

老师指出：

\begin{itemize}
\tightlist
\item
  创伤不只是``忘掉''就结束了；
\item
  还需要被整合、被重新理解。
\end{itemize}

也就是说：

\begin{itemize}
\tightlist
\item
  创伤确实会带来负面影响；
\item
  但在整合过程中，也可能产生新的意义。
\end{itemize}

【举例】老师提到疫情等重大公共事件，虽然是负性的，但也会促使人：

\begin{itemize}
\tightlist
\item
  重新思考生活；
\item
  重新理解自己和世界；
\item
  重新安排人生。
\end{itemize}

本段录音到此处结束，老师关于``重建意义''的展开应当还未完全讲完。

\begin{center}\rule{0.5\linewidth}{0.5pt}\end{center}

\subsection{本节英文术语索引}\label{ux672cux8282ux82f1ux6587ux672fux8bedux7d22ux5f15-3}

\begin{itemize}
\tightlist
\item
  心理危机 --- psychological crisis
\item
  预防 --- prevention
\item
  生物学因素 --- biological factors
\item
  遗传因素 --- genetic factors / hereditary factors
\item
  高敏感人群 --- highly sensitive person / high sensitivity group
\item
  冲动性 --- impulsivity
\item
  情绪调节 --- emotion regulation
\item
  反刍性思维 --- rumination
\item
  认知图式 --- cognitive schema
\item
  完美主义 --- perfectionism
\item
  自我批评 --- self-criticism
\item
  焦虑 --- anxiety
\item
  抑郁 / 抑郁症 --- depression
\item
  自杀倾向 --- suicidal tendency
\item
  应激事件 --- stressful event
\item
  扳机点 --- trigger
\item
  创伤 --- trauma
\item
  依恋 --- attachment
\item
  依恋创伤 --- attachment trauma
\item
  重要抚养者 / 主要照顾者 --- primary caregiver
\item
  安全感 --- sense of security
\item
  基本信任 --- basic trust
\item
  安全型依恋 --- secure attachment
\item
  回避型依恋 --- avoidant attachment
\item
  焦虑型依恋 --- anxious attachment
\item
  抵抗型依恋 --- resistant attachment
\item
  紊乱型依恋 --- disorganized attachment
\item
  亲密关系 --- intimate relationship
\item
  基本信任感 --- basic trust
\item
  创伤后应激障碍 --- Post-Traumatic Stress Disorder（PTSD）
\item
  闪回 --- flashback
\item
  心理咨询 --- psychological counseling
\item
  危机干预 --- crisis intervention
\item
  倾听 --- listening
\item
  共情 --- empathy
\item
  同情 --- sympathy
\item
  情感反应 --- emotional reflection
\item
  具体化 --- concretization
\item
  资源 --- resources
\item
  稳定化 --- stabilization
\item
  创伤暴露 --- trauma exposure
\item
  意义重建 --- meaning reconstruction
\end{itemize}

如果你愿意，我下一步可以把这份笔记再压缩成一版适合期末背诵的``考前提纲版''。

\clearpage

\section{第5讲：法律伦理、系统视角、树理论与自杀识别}\label{ux7b2c5ux8bb2ux6cd5ux5f8bux4f26ux7406ux7cfbux7edfux89c6ux89d2ux6811ux7406ux8bbaux4e0eux81eaux6740ux8bc6ux522b}

\subsection{一、上节课回顾与本节课导入}\label{ux4e00ux4e0aux8282ux8bfeux56deux987eux4e0eux672cux8282ux8bfeux5bfcux5165}

【复习】上节课主要讲了\textbf{心理危机的预防（prevention of
psychological crisis）}，老师回顾了影响心理危机产生的因素，包括：

\begin{itemize}
\tightlist
\item
  \textbf{生物学因素（biological factors）}
\item
  \textbf{心理因素（psychological factors）}
\item
  \textbf{后天创伤（trauma）}
\item
  其他会影响个体心理危机形成的因素
\end{itemize}

本节课转入\textbf{心理危机干预（psychological crisis
intervention）}，重点讨论： 当危机事件已经发生时，如何进行有效干预。

\begin{center}\rule{0.5\linewidth}{0.5pt}\end{center}

\subsection{二、心理危机干预中的法律与伦理问题}\label{ux4e8cux5fc3ux7406ux5371ux673aux5e72ux9884ux4e2dux7684ux6cd5ux5f8bux4e0eux4f26ux7406ux95eeux9898}

老师首先强调：心理危机干预工作不是``想怎么做就怎么做''，而是始终受到：

\begin{itemize}
\tightlist
\item
  \textbf{法律（law）}
\item
  \textbf{伦理（ethics）}
\end{itemize}

的双重约束。

\subsubsection{1.
保密原则不是绝对的}\label{ux4fddux5bc6ux539fux5219ux4e0dux662fux7eddux5bf9ux7684}

老师通过几个典型情境说明：\textbf{心理咨询中的保密原则（confidentiality）并非在任何情况下都必须绝对遵守}。

\paragraph{情境一：来访者明确表达自杀计划}\label{ux60c5ux5883ux4e00ux6765ux8bbfux8005ux660eux786eux8868ux8fbeux81eaux6740ux8ba1ux5212}

如果一位同学告诉咨询老师：

\begin{itemize}
\tightlist
\item
  自己已有明确的\textbf{自杀计划（suicide plan）}
\item
  或经评估属于高\textbf{自杀风险（suicide risk）}
\end{itemize}

那么此时需要\textbf{突破保密原则}。原因是：

\begin{itemize}
\tightlist
\item
  \textbf{生命安全第一（safety first / life safety first）}
\item
  心理工作者首先要保证来访者有机会被保护、被帮助
\end{itemize}

此时通常会告知：

\begin{itemize}
\tightlist
\item
  \textbf{监护人（guardian）}
\item
  \textbf{临时监护人}
\end{itemize}

老师特别说明，在学校环境中，同学们的临时监护人通常是：

\begin{itemize}
\tightlist
\item
  \textbf{辅导员}
\item
  \textbf{学院分管学生工作的副书记}
\end{itemize}

而正式监护人通常是父母。

\paragraph{情境二：来访者表达伤害他人的计划}\label{ux60c5ux5883ux4e8cux6765ux8bbfux8005ux8868ux8fbeux4f24ux5bb3ux4ed6ux4ebaux7684ux8ba1ux5212}

如果来访者明确告诉咨询老师，自己准备以某种方式对他人实施报复、危及对方生命安全，那么咨询老师同样需要\textbf{突破保密原则}，及时通知可能受害的人并采取保护措施。

\paragraph{情境三：涉及违法犯罪}\label{ux60c5ux5883ux4e09ux6d89ux53caux8fddux6cd5ux72afux7f6a}

如果来访者触犯法律，例如属于在逃人员，则咨询工作也必须受到法律条款约束，此时也可能需要突破保密。

\paragraph{情境四：未成年人正在遭受侵害}\label{ux60c5ux5883ux56dbux672aux6210ux5e74ux4ebaux6b63ux5728ux906dux53d7ux4fb5ux5bb3}

如果未成年人正在遭受：

\begin{itemize}
\tightlist
\item
  \textbf{家庭暴力（domestic violence）}
\item
  或其他危及人身安全的侵害
\end{itemize}

同样可能需要突破保密原则。

\begin{center}\rule{0.5\linewidth}{0.5pt}\end{center}

\subsubsection{2.
心理咨询受哪些法律约束}\label{ux5fc3ux7406ux54a8ux8be2ux53d7ux54eaux4e9bux6cd5ux5f8bux7ea6ux675f}

老师提到，心理咨询与危机干预是在一个更大的法律框架下展开的，例如：

\begin{itemize}
\tightlist
\item
  \textbf{宪法}
\item
  \textbf{未成年人保护法}
\item
  \textbf{精神卫生法（Mental Health Law）}
\end{itemize}

这些法律共同规定了咨询工作可以如何开展、能与哪些对象工作、在何种情况下必须采取行动。

\begin{center}\rule{0.5\linewidth}{0.5pt}\end{center}

\subsubsection{3.
心理咨询还受到伦理守则约束}\label{ux5fc3ux7406ux54a8ux8be2ux8fd8ux53d7ux5230ux4f26ux7406ux5b88ux5219ux7ea6ux675f}

除法律外，咨询师还要遵循\textbf{伦理守则（ethical
code）}。老师提到的重点包括：

\paragraph{（1）信息保密}\label{ux4fe1ux606fux4fddux5bc6}

咨询师需要保护当事人的信息，哪怕在：

\begin{itemize}
\tightlist
\item
  个案报告中
\item
  课堂讲述案例时
\end{itemize}

也必须尽量隐去能够暴露身份的信息。

\paragraph{（2）披露最小化}\label{ux62abux9732ux6700ux5c0fux5316}

在必须突破保密原则时，也要做到：

\begin{itemize}
\tightlist
\item
  \textbf{最小披露（minimal disclosure）}
\item
  只披露必要的信息
\item
  尽量保护当事人隐私
\end{itemize}

\paragraph{（3）避免双重关系}\label{ux907fux514dux53ccux91cdux5173ux7cfb}

老师特别解释了\textbf{双重关系（dual relationship）}：

在中国文化中常有``找熟人好办事''的倾向，但心理咨询中通常\textbf{应避免和熟人建立咨询关系}。例如：

\begin{itemize}
\tightlist
\item
  父母认识某位咨询师
\item
  便直接让孩子去找这位认识的咨询师
\end{itemize}

在专业伦理上通常并不合适，因为双重关系会影响咨询的边界与专业性。

\begin{center}\rule{0.5\linewidth}{0.5pt}\end{center}

\subsection{三、危机干预的核心前提：生命安全第一}\label{ux4e09ux5371ux673aux5e72ux9884ux7684ux6838ux5fc3ux524dux63d0ux751fux547dux5b89ux5168ux7b2cux4e00}

老师反复强调，心理危机干预最根本的原则是：

\subsection{\texorpdfstring{\textbf{生命安全第一位}}{生命安全第一位}}\label{ux751fux547dux5b89ux5168ux7b2cux4e00ux4f4d}

这也是为什么在必要时要突破保密原则。

老师将其和医学生熟悉的职业使命作了类比：

\begin{itemize}
\tightlist
\item
  医生强调\textbf{救死扶伤}
\item
  心理工作者同样强调\textbf{生命至上（life first）}
\end{itemize}

因此，在来访者出现危机时，首要任务不是深层分析人格，而是先判断：

\begin{itemize}
\tightlist
\item
  有没有生命危险
\item
  危险等级多高
\item
  是否需要立刻保护起来
\end{itemize}

\begin{center}\rule{0.5\linewidth}{0.5pt}\end{center}

\subsection{四、从系统视角理解心理危机}\label{ux56dbux4eceux7cfbux7edfux89c6ux89d2ux7406ux89e3ux5fc3ux7406ux5371ux673a}

【复习】老师再次回到上节课内容，指出心理危机不能只从个体内部去理解，不能简单归因为``这个人自己有问题''。

\subsubsection{1.
个体因素只是其中一部分}\label{ux4e2aux4f53ux56e0ux7d20ux53eaux662fux5176ux4e2dux4e00ux90e8ux5206}

个体层面的影响因素包括：

\begin{itemize}
\tightlist
\item
  遗传因素（genetic factors）
\item
  气质特点（temperament）
\item
  神经冲动水平
\item
  人格特质（personality traits）
\item
  当前遭遇的情境事件（situational events）
\end{itemize}

\begin{center}\rule{0.5\linewidth}{0.5pt}\end{center}

\subsubsection{2.
更大的社会与文化系统同样重要}\label{ux66f4ux5927ux7684ux793eux4f1aux4e0eux6587ux5316ux7cfbux7edfux540cux6837ux91cdux8981}

老师强调，心理危机还受到更宏观系统的影响。

\paragraph{（1）经济环境}\label{ux7ecfux6d4eux73afux5883}

经济环境的变化会显著影响个体和家庭的危机风险。

【举例】老师以疫情后部分行业受冲击为例，指出：

\begin{itemize}
\tightlist
\item
  一些行业赚钱变难
\item
  某些从业者可能从原本经济条件很好，突然陷入破产或重大财务危机
\item
  对个人和家庭来说，这不仅是经济危机，也是心理危机
\end{itemize}

老师举到房地产行业的例子：
行业在不同时期经历大起大落，对从业者造成巨大心理冲击。

\begin{center}\rule{0.5\linewidth}{0.5pt}\end{center}

\paragraph{（2）文化环境}\label{ux6587ux5316ux73afux5883}

不同文化背景下，人面对灾难后的恢复方式不同。

【举例】老师提到\textbf{斯里兰卡海啸（Sri Lanka tsunami）}后的现象：

\begin{itemize}
\tightlist
\item
  一些欧美心理工作者到当地救援
\item
  但他们惊讶地发现，当地民众在彼此帮助的社会氛围中，恢复得很快
\end{itemize}

【举例】老师又提到\textbf{汶川地震（Wenchuan earthquake）}：

\begin{itemize}
\tightlist
\item
  许多心理救援工作者进入灾区
\item
  但部分当地群众并不那么依赖专业心理援助
\item
  四川人民通过聚在一起打麻将等方式，反而慢慢恢复了秩序感与支持感
\end{itemize}

老师的解释是：

\begin{itemize}
\item
  \textbf{文化会影响人如何理解灾难、如何恢复、如何互相支持}
\item
  像打麻将这种活动，在心理学视角下不仅是娱乐，也有：

  \begin{itemize}
  \tightlist
  \item
    团体合作
  \item
    相互支持
  \item
    重建秩序感
  \end{itemize}
\end{itemize}

\begin{center}\rule{0.5\linewidth}{0.5pt}\end{center}

\paragraph{（3）中国传统文化的影响}\label{ux4e2dux56fdux4f20ux7edfux6587ux5316ux7684ux5f71ux54cd}

老师认为，中国传统文化中有许多与心理韧性有关的智慧，例如：

\begin{itemize}
\tightlist
\item
  \textbf{中庸}
\item
  \textbf{辩证看待问题}
\item
  ``\textbf{塞翁失马，焉知非福}''
\end{itemize}

这些文化资源会影响人们如何理解人生轨迹、如何应对灾难与变化。

\begin{center}\rule{0.5\linewidth}{0.5pt}\end{center}

\subsubsection{3.
家庭与成长环境的重要性}\label{ux5bb6ux5eadux4e0eux6210ux957fux73afux5883ux7684ux91cdux8981ux6027}

老师强调：

\subsection{家庭和成长环境对人的影响极大}\label{ux5bb6ux5eadux548cux6210ux957fux73afux5883ux5bf9ux4ebaux7684ux5f71ux54cdux6781ux5927}

并提出一个核心观点：

\begin{itemize}
\tightlist
\item
  家庭是一个人的``\textbf{根}''
\item
  很多性格、创伤、应对方式，存在\textbf{代际传递（intergenerational
  transmission）}
\end{itemize}

也就是说：

\begin{itemize}
\tightlist
\item
  一代影响一代
\item
  某些创伤与性格模式，会在家族中不断重复
\end{itemize}

\begin{center}\rule{0.5\linewidth}{0.5pt}\end{center}

\subsubsection{4.
成长环境塑造社交能力与人格}\label{ux6210ux957fux73afux5883ux5851ux9020ux793eux4ea4ux80fdux529bux4e0eux4ebaux683c}

老师指出，不同成长环境对人格塑造差异很大。例如：

\begin{itemize}
\tightlist
\item
  多子女家庭 vs 独自成长
\item
  农村/人口密集、邻里互动多的环境 vs 城市中邻里关系淡薄的环境
\item
  从小社交训练多 vs 社交训练匮乏
\end{itemize}

老师借此解释为何现在很多同学会觉得自己``不会社交''：

\begin{itemize}
\tightlist
\item
  并不一定只是个人能力差
\item
  很可能和成长环境中缺乏社交训练有关
\end{itemize}

\begin{center}\rule{0.5\linewidth}{0.5pt}\end{center}

\subsubsection{5.
每一代人所处的大时代不同，心理议题也不同}\label{ux6bcfux4e00ux4ee3ux4ebaux6240ux5904ux7684ux5927ux65f6ux4ee3ux4e0dux540cux5fc3ux7406ux8baeux9898ux4e5fux4e0dux540c}

老师引导大家从代际角度看问题：

\begin{itemize}
\item
  祖辈可能经历过战争、饥荒，因此更节俭
\item
  父母一代经历改革开放，更强调``知识改变命运''
\item
  当代年轻人物质相对不匮乏，但可能更常感到：

  \begin{itemize}
  \tightlist
  \item
    缺少玩伴
  \item
    缺少理解
  \item
    竞争激烈
  \item
    未来不确定
  \end{itemize}
\end{itemize}

因此，每一代人面临的心理冲突并不相同。

\begin{center}\rule{0.5\linewidth}{0.5pt}\end{center}

\subsubsection{6.
人是``系统中的人''}\label{ux4ebaux662fux7cfbux7edfux4e2dux7684ux4eba}

老师总结：

\begin{quote}
一个人出现心理问题或心理危机，绝不是只看他自己。
\end{quote}

要把人放回完整系统中理解，包括：

\begin{itemize}
\tightlist
\item
  大的社会环境
\item
  家庭
\item
  学校
\item
  社区
\item
  亲朋好友
\item
  个体的生理与心理因素
\item
  所经历的事件
\end{itemize}

\begin{center}\rule{0.5\linewidth}{0.5pt}\end{center}

\subsubsection{7.
心理危机干预也必须是系统性的}\label{ux5fc3ux7406ux5371ux673aux5e72ux9884ux4e5fux5fc5ux987bux662fux7cfbux7edfux6027ux7684}

危机干预不能只干预个人，而要同时考虑：

\begin{itemize}
\tightlist
\item
  个体解决问题的能力
\item
  个体看待世界与看待自己的方式
\item
  家庭系统
\item
  学校系统
\item
  社区支持
\item
  辅导员、老师、周围同学等资源
\end{itemize}

老师提到，心理工作中常会使用：

\begin{itemize}
\tightlist
\item
  \textbf{家庭治疗（family therapy）}
\end{itemize}

正是因为干预常常需要进入系统层面，而不是只盯着一个人。

\begin{center}\rule{0.5\linewidth}{0.5pt}\end{center}

\subsection{五、心理危机的``树理论''}\label{ux4e94ux5fc3ux7406ux5371ux673aux7684ux6811ux7406ux8bba}

老师接下来讲到一个非常重要的比喻：\textbf{树理论}。
原转写中``数理论''应为``\textbf{树理论}''。

\subsubsection{1.
根：家庭与依恋关系}\label{ux6839ux5bb6ux5eadux4e0eux4f9dux604bux5173ux7cfb}

老师认为，一个人想成长得好，首先要有发达的``根系''。

这里的``根''在心理学上主要指：

\begin{itemize}
\tightlist
\item
  \textbf{家庭（family）}
\item
  \textbf{依恋关系（attachment / attachment relationship）}
\end{itemize}

意义在于：

\begin{itemize}
\tightlist
\item
  人生一定会有风吹雨打
\item
  如果有稳固的家庭和依恋关系
\item
  遭遇挫折时，知道背后有人托举自己
\item
  知道有避风港和爱的港湾
\end{itemize}

这能带来：

\begin{itemize}
\tightlist
\item
  对自己和他人的基础信任感（basic trust）
\end{itemize}

\begin{center}\rule{0.5\linewidth}{0.5pt}\end{center}

\subsubsection{2.
树干：信仰、价值观、社会支持}\label{ux6811ux5e72ux4fe1ux4ef0ux4ef7ux503cux89c2ux793eux4f1aux652fux6301}

树要长得好，还需要粗壮的树干。

老师把树干对应为：

\begin{itemize}
\tightlist
\item
  \textbf{信仰（belief / faith）}
\item
  \textbf{价值观（values）}
\item
  \textbf{社会支持（social support）}
\end{itemize}

\paragraph{（1）信仰}\label{ux4fe1ux4ef0}

老师提醒大家思考：

\begin{itemize}
\tightlist
\item
  你是否真正相信自己在做的事
\item
  你是否相信学习能给你带来意义和幸福
\end{itemize}

如果没有信仰，就容易动摇、容易被外界影响。

\paragraph{（2）价值观}\label{ux4ef7ux503cux89c2}

不同的价值观会影响一个人面对困难时的应对方式。 例如：

\begin{itemize}
\tightlist
\item
  想让世界变得更好
\item
  想为自己而活
\end{itemize}

这些都是价值观。

\paragraph{（3）社会支持}\label{ux793eux4f1aux652fux6301}

社会支持包括：

\begin{itemize}
\tightlist
\item
  朋友
\item
  家人
\item
  老师
\item
  亲朋好友
\item
  困难时可以求助的人
\end{itemize}

老师引用一句中国话说明社会支持的重要性：

\begin{quote}
平时肯帮人，急时有人帮。
\end{quote}

\begin{center}\rule{0.5\linewidth}{0.5pt}\end{center}

\subsubsection{3.
枝叶与外部环境：社会环境也决定成长状态}\label{ux679dux53f6ux4e0eux5916ux90e8ux73afux5883ux793eux4f1aux73afux5883ux4e5fux51b3ux5b9aux6210ux957fux72b6ux6001}

一棵树能否枝繁叶茂，还取决于它长在什么环境里：

\begin{itemize}
\tightlist
\item
  是在阳光、水分、土壤充足的地方
\item
  还是在悬崖缝隙中、资源稀缺的地方
\end{itemize}

同样，一个人的发展也受到外部社会环境的影响，例如：

\begin{itemize}
\tightlist
\item
  国家政策
\item
  行业发展
\item
  经济背景
\item
  社会机遇
\end{itemize}

【举例】老师再次以改革开放后的社会机会为例，说明：

\begin{itemize}
\tightlist
\item
  有些人的发展并不只是个人努力，也与时代机会密切相关
\end{itemize}

\begin{center}\rule{0.5\linewidth}{0.5pt}\end{center}

\subsubsection{4.
树理论与危机承受力}\label{ux6811ux7406ux8bbaux4e0eux5371ux673aux627fux53d7ux529b}

老师用树理论解释危机的后果：

\paragraph{如果一个人拥有：}\label{ux5982ux679cux4e00ux4e2aux4ebaux62e5ux6709}

\begin{itemize}
\tightlist
\item
  发达的根系
\item
  粗壮的树干
\item
  较好的外部环境
\end{itemize}

那么即使遭遇``狂风暴雨''般的危机，也可能只是：

\begin{itemize}
\tightlist
\item
  折断一些枝条
\item
  掉落一些叶子
\end{itemize}

但整棵树仍能活下来，并继续生长。

\paragraph{如果一个人缺乏：}\label{ux5982ux679cux4e00ux4e2aux4ebaux7f3aux4e4f}

\begin{itemize}
\tightlist
\item
  家庭支持
\item
  稳定依恋
\item
  信仰与意义感
\item
  社会支持
\end{itemize}

那么一场危机就可能动摇他的生存根基。

【举例】老师举到求职失败的情境：
如果一个人没有信仰、没有支持系统、家庭关系也差，那么``找不到工作''不只是现实挫折，还会立刻演变成：

\begin{itemize}
\tightlist
\item
  我活着的意义是什么？
\item
  没有人会帮助我
\item
  我不能信任任何人
\end{itemize}

这就是危机迅速加深的过程。

\begin{center}\rule{0.5\linewidth}{0.5pt}\end{center}

\subsection{六、危机干预的三个阶段}\label{ux516dux5371ux673aux5e72ux9884ux7684ux4e09ux4e2aux9636ux6bb5}

在树理论之后，老师顺势讲到危机干预的三个阶段。

\subsubsection{1. 第一阶段：危机评估（crisis
assessment）}\label{ux7b2cux4e00ux9636ux6bb5ux5371ux673aux8bc4ux4f30crisis-assessment}

首先要评估：

\begin{itemize}
\tightlist
\item
  自杀风险
\item
  攻击他人的风险
\end{itemize}

老师强调，不同风险等级差异很大：

\begin{itemize}
\tightlist
\item
  只有一闪而过的自杀念头
\item
  已经形成明确计划
\item
  已购买工具
\item
  已踩点准备实施
\end{itemize}

这些不是一回事。

\begin{center}\rule{0.5\linewidth}{0.5pt}\end{center}

\subsubsection{2. 第二阶段：行动干预（action
intervention）}\label{ux7b2cux4e8cux9636ux6bb5ux884cux52a8ux5e72ux9884action-intervention}

如果风险已经较高，例如已经有明确自杀计划，就必须立即采取现实层面的保护行动，例如：

\begin{itemize}
\tightlist
\item
  突破保密原则
\item
  通知父母
\item
  让辅导员或其他人员进行24小时看护
\item
  先把人保护起来
\end{itemize}

\begin{center}\rule{0.5\linewidth}{0.5pt}\end{center}

\subsubsection{3.
第三阶段：心理干预与后续心理治疗}\label{ux7b2cux4e09ux9636ux6bb5ux5fc3ux7406ux5e72ux9884ux4e0eux540eux7eedux5fc3ux7406ux6cbbux7597}

\paragraph{（1）心理干预}\label{ux5fc3ux7406ux5e72ux9884}

心理干预的直接目标是：

\begin{itemize}
\tightlist
\item
  \textbf{稳定情绪（stabilization）}
\item
  \textbf{处理当前问题}
\item
  \textbf{帮助问题解决}
\end{itemize}

【举例】老师举了``研究生选导师落空''的例子：

\begin{itemize}
\tightlist
\item
  原以为导师会收自己
\item
  但政策变化，导师能带的人数减少
\item
  导致该同学研究生阶段安排受阻，产生危机
\end{itemize}

此时心理干预要做的，就是帮助当事人讨论还有哪些替代路径和现实解决方案。

\paragraph{（2）后续心理治疗（psychotherapy）}\label{ux540eux7eedux5fc3ux7406ux6cbbux7597psychotherapy}

如果危机已经解除，后续咨询工作的重点就转向：

\begin{itemize}
\tightlist
\item
  帮助来访者成长
\item
  发展应对技巧（coping skills）
\item
  为今后遇到类似情景做准备
\end{itemize}

\begin{center}\rule{0.5\linewidth}{0.5pt}\end{center}

\subsection{七、什么是危机，以及危机常见反应}\label{ux4e03ux4ec0ux4e48ux662fux5371ux673aux4ee5ux53caux5371ux673aux5e38ux89c1ux53cdux5e94}

老师给出危机的定义：

\begin{quote}
危机是人在紧急状态下，原有心理平衡被打破，继而出现情感、认知、行为等功能失调。
\end{quote}

\begin{center}\rule{0.5\linewidth}{0.5pt}\end{center}

\subsubsection{1. 认知反应（cognitive
reactions）}\label{ux8ba4ux77e5ux53cdux5e94cognitive-reactions}

危机中的常见认知表现包括：

\begin{itemize}
\tightlist
\item
  否认（denial）
\item
  注意力不集中
\item
  危机场景反复再现
\item
  失去信心
\item
  内疚自责
\item
  丧失安全感
\end{itemize}

老师特别强调： 危机中的人，最突出的认知特点是：

\subsection{\texorpdfstring{\textbf{认知变窄（narrowed
cognition）}}{认知变窄（narrowed cognition）}}\label{ux8ba4ux77e5ux53d8ux7a84narrowed-cognition}

也就是：

\begin{itemize}
\tightlist
\item
  眼里只有问题和困难
\item
  看不到其他可能性、资源和希望
\end{itemize}

【举例】例如重大疾病或绝症诊断刚出现时，人的第一反应常常是否认：

\begin{itemize}
\tightlist
\item
  是不是搞错了？
\item
  怎么会发生在我身上？
\end{itemize}

\begin{center}\rule{0.5\linewidth}{0.5pt}\end{center}

\subsubsection{2. 情感/情绪反应（emotional
reactions）}\label{ux60c5ux611fux60c5ux7eeaux53cdux5e94emotional-reactions}

常见情绪表现包括：

\begin{itemize}
\tightlist
\item
  混乱
\item
  害怕
\item
  恐惧
\item
  沮丧
\item
  麻木
\item
  怀疑
\item
  悲伤
\item
  绝望
\item
  无助
\item
  羞耻
\item
  易怒
\item
  情绪难以平静
\item
  情绪失调
\end{itemize}

老师以地震、火灾等情景说明，人会明显体验到：

\begin{itemize}
\tightlist
\item
  环境失控
\item
  内在失控
\item
  对危险的强烈反应
\end{itemize}

\begin{center}\rule{0.5\linewidth}{0.5pt}\end{center}

\subsubsection{3. 行为反应（behavioral
reactions）}\label{ux884cux4e3aux53cdux5e94behavioral-reactions}

危机中的行为表现可能有：

\begin{itemize}
\tightlist
\item
  攻击
\item
  社交退缩
\item
  逃避
\item
  食欲异常
\item
  哭泣
\item
  饮酒增加
\item
  药物使用增加
\item
  坐立不安
\item
  过度警觉（hypervigilance）
\end{itemize}

老师强调： 危机中的行为特点常表现为：

\begin{itemize}
\tightlist
\item
  \textbf{反常}
\item
  \textbf{易失控}
\end{itemize}

【临床关联】老师举了父亲脑梗后的变化：

\begin{itemize}
\tightlist
\item
  不愿意见人
\item
  容易易怒
\end{itemize}

说明疾病本身也可能构成一个人的心理危机。

老师提醒大家，将来进入医院后面对重病患者时，也应理解：

\begin{itemize}
\tightlist
\item
  这些情绪和行为往往是正常危机反应的一部分
\end{itemize}

\begin{center}\rule{0.5\linewidth}{0.5pt}\end{center}

\subsection{八、危机干预的基本原则}\label{ux516bux5371ux673aux5e72ux9884ux7684ux57faux672cux539fux5219}

老师总结了几条危机干预的原则：

\subsubsection{1. 保障安全}\label{ux4fddux969cux5b89ux5168}

首要目的是保证被干预者的人身安全。

\subsubsection{2. 聚焦问题}\label{ux805aux7126ux95eeux9898}

危机干预主要聚焦于：

\begin{itemize}
\tightlist
\item
  当前的情绪冲突
\item
  当前的问题解决
\item
  当前的情绪调节困难
\end{itemize}

而不是去深挖人格深层问题。

\subsubsection{3. 激活资源（resource
activation）}\label{ux6fc0ux6d3bux8d44ux6e90resource-activation}

危机干预要帮助来访者看见并调动：

\begin{itemize}
\tightlist
\item
  自身资源
\item
  家庭资源
\item
  学校资源
\item
  社交资源
\end{itemize}

老师再次强调：

\begin{quote}
危机既是危险，也是机会。
\end{quote}

``机会''在于： 如果一个人走过危机，他会习得新的能力，例如：

\begin{itemize}
\tightlist
\item
  生存技能
\item
  应对挫折的能力
\item
  自我重新定义的能力
\item
  社交求助能力
\end{itemize}

【举例】老师提到：

\begin{itemize}
\tightlist
\item
  有的同学以前在高中一直是佼佼者
\item
  进入大学后学习优势消失
\item
  自我价值严重受挫
\end{itemize}

这时危机的意义在于： 要学会重新审视自己------除了学习，我还有什么价值？

【举例】老师还提到研究生与导师相处的问题，说明危机也可能倒逼一个人学习：

\begin{itemize}
\tightlist
\item
  主动沟通
\item
  求助
\item
  寻找资源
\item
  发展社交能力
\end{itemize}

\begin{center}\rule{0.5\linewidth}{0.5pt}\end{center}

\subsection{九、后半节：危机干预前的准备与工作状态}\label{ux4e5dux540eux534aux8282ux5371ux673aux5e72ux9884ux524dux7684ux51c6ux5907ux4e0eux5de5ux4f5cux72b6ux6001}

第二个片段开始，老师进一步讲危机干预的流程。

\subsubsection{1.
干预前必须有充分准备}\label{ux5e72ux9884ux524dux5fc5ux987bux6709ux5145ux5206ux51c6ux5907}

老师把危机干预类比为外科手术：

\begin{itemize}
\tightlist
\item
  外科医生上手术台前需要做好身心准备
\item
  危机干预者同样需要进入专业状态
\end{itemize}

要求包括：

\begin{itemize}
\tightlist
\item
  心理上准备充分
\item
  生理状态到位
\item
  能把注意力真正放在当事人身上
\item
  而不是被自己当天的私人事务占据
\end{itemize}

\begin{center}\rule{0.5\linewidth}{0.5pt}\end{center}

\subsubsection{2.
危机干预必须依靠团队}\label{ux5371ux673aux5e72ux9884ux5fc5ux987bux4f9dux9760ux56e2ux961f}

老师再次强调：

\subsection{心理危机干预不是一个人能完成的}\label{ux5fc3ux7406ux5371ux673aux5e72ux9884ux4e0dux662fux4e00ux4e2aux4ebaux80fdux5b8cux6210ux7684}

需要团队合作，例如在学校场景中，需要：

\begin{itemize}
\tightlist
\item
  心理健康老师
\item
  辅导员
\item
  任课老师
\item
  班主任
\item
  父母
\item
  其他支持者
\end{itemize}

\begin{center}\rule{0.5\linewidth}{0.5pt}\end{center}

\subsubsection{3.
干预者要保持松弛而稳定的状态}\label{ux5e72ux9884ux8005ux8981ux4fddux6301ux677eux5f1bux800cux7a33ux5b9aux7684ux72b6ux6001}

面对一个说``我没有活下去勇气、想结束生命''的人，干预者自己也会紧张、害怕、担心。
但老师强调：

\begin{itemize}
\tightlist
\item
  不能先被吓倒
\item
  要能陪在这个人身边
\item
  当事人能感受到你的状态
\end{itemize}

\begin{center}\rule{0.5\linewidth}{0.5pt}\end{center}

\subsubsection{4.
干预中的态度要求}\label{ux5e72ux9884ux4e2dux7684ux6001ux5ea6ux8981ux6c42}

老师强调危机干预者应具备：

\begin{itemize}
\tightlist
\item
  真诚（genuineness）
\item
  关心（care）
\item
  热爱/善意
\item
  温和而坚定（gentle but firm）
\item
  专业高效
\item
  依法依规
\end{itemize}

尤其重要的是：

\begin{quote}
你关心的是``这个人好不好''，而不是只关心他的行为是否符合社会规范。
\end{quote}

\begin{center}\rule{0.5\linewidth}{0.5pt}\end{center}

\subsection{十、人的基本需要与危机的关系}\label{ux5341ux4ebaux7684ux57faux672cux9700ux8981ux4e0eux5371ux673aux7684ux5173ux7cfb}

老师接着从人的需要谈危机形成机制，提到\textbf{马斯洛需要层次理论（Maslow's
hierarchy of needs）}。

人有许多基本需要，包括：

\begin{itemize}
\tightlist
\item
  生理需要（physiological needs）
\item
  安全需要（safety needs）
\item
  尊重需要（esteem needs）
\item
  自我实现需要（self-actualization）
\end{itemize}

当某些重要需要长期得不到满足时，人就容易产生心理痛苦。

\begin{center}\rule{0.5\linewidth}{0.5pt}\end{center}

\subsubsection{1. 关系的需要（need for
relationship）}\label{ux5173ux7cfbux7684ux9700ux8981need-for-relationship}

老师认为：

\begin{itemize}
\tightlist
\item
  人需要建立关系
\item
  如果在环境中无法建立关系，就容易出问题
\end{itemize}

【举例】大学宿舍、班级环境里能否建立关系非常重要。 如果一个人长期：

\begin{itemize}
\tightlist
\item
  独来独往
\item
  感到孤立无援
\item
  感受不到支持、认可、帮助
\end{itemize}

就更容易陷入无助。

\begin{center}\rule{0.5\linewidth}{0.5pt}\end{center}

\subsubsection{2. 自尊的需要（need for
self-esteem）}\label{ux81eaux5c0aux7684ux9700ux8981need-for-self-esteem}

每个人都希望被尊重。 老师举到学校资助工作的例子：

【举例】如果对经济困难学生的支持能够通过大数据识别、无声无息地给予，而不必让学生反复公开申请、暴露贫困身份，这本身就是一种尊重。

老师还讲到社交和教育中的具体技巧：

\begin{itemize}
\tightlist
\item
  \textbf{表扬尽量在人前}
\item
  \textbf{批评尽量一对一、私下进行}
\end{itemize}

将来无论是：

\begin{itemize}
\tightlist
\item
  为人父母
\item
  做老师
\item
  与人交往
\end{itemize}

都应注意这一点，因为每个人都有``面子''和尊严需求。

\begin{center}\rule{0.5\linewidth}{0.5pt}\end{center}

\subsubsection{3. 安全感的需要（need for
security）}\label{ux5b89ux5168ux611fux7684ux9700ux8981need-for-security}

安全感既来自外部，也来自内部。

外部包括：

\begin{itemize}
\tightlist
\item
  社会是否稳定
\item
  是否有战争
\item
  环境是否可预期
\end{itemize}

内部则包括：

\begin{itemize}
\tightlist
\item
  我能否通过努力达到想要的生活状态
\item
  我是否对未来有掌控感
\end{itemize}

老师指出，很多同学之所以特别努力，某种程度上也是在争取自己的安全感。

\begin{center}\rule{0.5\linewidth}{0.5pt}\end{center}

\subsubsection{4. 亲密关系的需要（need for
intimacy）}\label{ux4eb2ux5bc6ux5173ux7cfbux7684ux9700ux8981need-for-intimacy}

即便有些人嘴上说不需要恋爱、不需要结婚，但从人的心理深处看，大多数人都仍然渴望：

\begin{itemize}
\tightlist
\item
  亲密感
\item
  被走近
\item
  被理解
\item
  被连接
\end{itemize}

\begin{center}\rule{0.5\linewidth}{0.5pt}\end{center}

\subsubsection{5. 成就感与自我认同（achievement and
self-identity）}\label{ux6210ux5c31ux611fux4e0eux81eaux6211ux8ba4ux540cachievement-and-self-identity}

老师强调：

\begin{itemize}
\tightlist
\item
  每个人都需要在某个方面感到``我能做成点事''
\item
  每个人的自我认同方式不同
\end{itemize}

例如：

\begin{itemize}
\tightlist
\item
  学习好
\item
  口才好
\item
  社交强
\item
  钢琴弹得好
\end{itemize}

因此，自我价值不能只压缩在一个非常狭窄的维度里。

\begin{center}\rule{0.5\linewidth}{0.5pt}\end{center}

\subsubsection{6.
情绪与感受幸福的能力}\label{ux60c5ux7eeaux4e0eux611fux53d7ux5e78ux798fux7684ux80fdux529b}

老师讲到一个重要观点：

\begin{itemize}
\tightlist
\item
  有些人更容易感受到快乐
\item
  有些人即便客观上已很成功，也感受不到意义
\end{itemize}

【举例】老师讲了一个案例：
一个已经``在别人眼里很有成就''的人，却非常不快乐，甚至觉得活着没有意义。后来有人建议他：

\begin{itemize}
\tightlist
\item
  给街头乞讨者、吃不上饭的人送盒饭
\end{itemize}

结果他持续做这件事后，逐渐感到自己的人生重新有了意义，甚至后来还组织更多人去灾区送饭。

老师借此说明：

\begin{quote}
真正持久地让人快乐的，不一定是功成名就，而是做自己觉得有意义、有价值的事。
\end{quote}

\begin{center}\rule{0.5\linewidth}{0.5pt}\end{center}

\subsection{十一、不同需要受挫后的心理表现}\label{ux5341ux4e00ux4e0dux540cux9700ux8981ux53d7ux632bux540eux7684ux5fc3ux7406ux8868ux73b0}

老师接着从不同心理需要受挫的角度，分析危机的表现。

\subsubsection{1. 关系受挫}\label{ux5173ux7cfbux53d7ux632b}

例如被抛弃后，可能遭受强烈心理冲击。
成年后可能会拼命寻找``不被抛弃的关系''。

\subsubsection{2. 自尊受挫}\label{ux81eaux5c0aux53d7ux632b}

如果一个人自尊很脆弱，被老师当众批评，可能立刻陷入强烈羞耻感，甚至转而攻击对方。

\subsubsection{3. 安全感不足}\label{ux5b89ux5168ux611fux4e0dux8db3}

安全感低的人，容易长期焦虑，通过不停做事来获得控制感。
老师提到这类人可能表现为``\textbf{工作狂（workaholic）}''。

\subsubsection{4.
亲密关系不足}\label{ux4eb2ux5bc6ux5173ux7cfbux4e0dux8db3}

容易体验到孤独，严重时会感到这个世界没人支持自己，从而出现绝望。

\subsubsection{5. 成就感受挫}\label{ux6210ux5c31ux611fux53d7ux632b}

权力、掌控、成就受挫时，也可能带来强烈内在挣扎。

\subsubsection{6.
自我认同受挫}\label{ux81eaux6211ux8ba4ux540cux53d7ux632b}

容易陷入迷茫，不知道自己是谁、价值何在。

\begin{center}\rule{0.5\linewidth}{0.5pt}\end{center}

\subsubsection{7.
针对不同受挫类型的干预方向}\label{ux9488ux5bf9ux4e0dux540cux53d7ux632bux7c7bux578bux7684ux5e72ux9884ux65b9ux5411}

老师概括提到：

\begin{itemize}
\tightlist
\item
  对关系受挫者：帮助其在关系中获得安全感、稳定关系
\item
  对自尊受挫者：处理羞耻感（shame）
\item
  对缺乏尊重者：深入共情、给予安全与抱持（holding）
\item
  对焦虑和缺乏安全感者：容纳其焦虑
\item
  对特别需要权力和控制者：赋权（empowerment）、理解、稳定关系
\end{itemize}

老师也提示，这些原则比较概括，大家一时未必都能完全理解。

\begin{center}\rule{0.5\linewidth}{0.5pt}\end{center}

\subsection{十二、危机干预的具体步骤}\label{ux5341ux4e8cux5371ux673aux5e72ux9884ux7684ux5177ux4f53ux6b65ux9aa4}

老师再次把危机干预归纳为三个步骤：

\subsubsection{1.
明确问题和目标}\label{ux660eux786eux95eeux9898ux548cux76eeux6807}

先弄清问题是什么、目标是什么，不能在问题未明时就开始盲目干预。

老师特别提醒：

\begin{itemize}
\tightlist
\item
  学业问题往往只是表层问题
\item
  更深层可能是自我价值、自我人生意义受到挑战
\item
  亲密关系破裂也不只是``分手''本身，而常激活被抛弃感等更深层需求
\end{itemize}

\begin{center}\rule{0.5\linewidth}{0.5pt}\end{center}

\subsubsection{2. 快速而准确地评估（rapid and accurate
assessment）}\label{ux5febux901fux800cux51c6ux786eux5730ux8bc4ux4f30rapid-and-accurate-assessment}

评估贯穿干预全过程，尤其包括：

\begin{itemize}
\tightlist
\item
  自杀风险
\item
  自伤风险
\item
  攻击他人风险
\item
  现实压力
\item
  支持资源
\item
  是否存在临床诊断
\end{itemize}

老师举到风险评估中可能会参考的因素：

\begin{itemize}
\tightlist
\item
  自杀计划高低
\item
  既往是否有自杀史
\item
  次数、频率
\item
  现实压力程度
\item
  资源多少
\item
  是否有抑郁症、精神分裂症、躁狂、双相情感障碍等诊断
\end{itemize}

\begin{center}\rule{0.5\linewidth}{0.5pt}\end{center}

\subsubsection{3.
给予希望与问题解决}\label{ux7ed9ux4e88ux5e0cux671bux4e0eux95eeux9898ux89e3ux51b3}

危机中的人认知狭窄，容易认为：

\begin{itemize}
\tightlist
\item
  没路可走了
\item
  只能这样了
\end{itemize}

因此干预者要做的是：

\begin{itemize}
\tightlist
\item
  帮助其看见被忽略的其他路径
\item
  扩展人生选择视野
\item
  发现周围资源
\item
  一步一步找到问题解决方法
\end{itemize}

老师还提到有时会签订：

\begin{itemize}
\tightlist
\item
  \textbf{不伤害协议（no-harm agreement / safety contract）}
\end{itemize}

即要求当事人承诺：

\begin{itemize}
\tightlist
\item
  不实施自伤/自杀行为
\item
  若出现冲动，先联系咨询师或院系老师
\end{itemize}

\begin{center}\rule{0.5\linewidth}{0.5pt}\end{center}

\subsection{十三、危机干预中的安全措施}\label{ux5341ux4e09ux5371ux673aux5e72ux9884ux4e2dux7684ux5b89ux5168ux63aaux65bd}

老师反复强调，整个干预前后都必须：

\begin{itemize}
\tightlist
\item
  持续评估风险
\item
  足够重视安全问题
\item
  制定行动策略
\end{itemize}

对于有自杀风险的人，要尽可能使其远离：

\begin{itemize}
\tightlist
\item
  药品
\item
  刀具
\item
  其他危险器械
\item
  危险地点
\end{itemize}

必要时：

\begin{itemize}
\tightlist
\item
  实施监护
\item
  在法律允许范围内送医治疗
\end{itemize}

【举例】如果学生已经购买了可能危及生命的药品、器械，就应让他人暂时代为保管，使工具变得``不可获得''。

\begin{center}\rule{0.5\linewidth}{0.5pt}\end{center}

\subsection{十四、建立关系：危机干预有效的前提}\label{ux5341ux56dbux5efaux7acbux5173ux7cfbux5371ux673aux5e72ux9884ux6709ux6548ux7684ux524dux63d0}

老师特别强调，危机干预能否有效，第一步不是技巧，而是：

\subsection{\texorpdfstring{\textbf{关系（therapeutic
relationship）}}{关系（therapeutic relationship）}}\label{ux5173ux7cfbtherapeutic-relationship}

来访者必须感到：

\begin{itemize}
\tightlist
\item
  你愿意听
\item
  你真诚接纳
\item
  你尊重他
\item
  你能理解他的核心困难
\end{itemize}

如果当事人根本不相信你，那后续干预就很难真正有效。

\begin{center}\rule{0.5\linewidth}{0.5pt}\end{center}

\subsection{十五、学校心理健康资源与心理咨询的意义}\label{ux5341ux4e94ux5b66ux6821ux5fc3ux7406ux5065ux5eb7ux8d44ux6e90ux4e0eux5fc3ux7406ux54a8ux8be2ux7684ux610fux4e49}

接近课程后段时，老师开始介绍学校资源。

\subsubsection{1. 课程资源}\label{ux8bfeux7a0bux8d44ux6e90}

学校里获取心理学知识的便捷方式包括：

\begin{itemize}
\tightlist
\item
  《大学生心理健康》
\item
  《心理危机干预》
\item
  《压力与健康》
\item
  其他相关通识课程
\end{itemize}

老师认为，这类课程是最简单快捷的心理学知识入口。

\begin{center}\rule{0.5\linewidth}{0.5pt}\end{center}

\subsubsection{2.
心理社团与活动资源}\label{ux5fc3ux7406ux793eux56e2ux4e0eux6d3bux52a8ux8d44ux6e90}

如果参加心理社团，也有机会更近距离接触：

\begin{itemize}
\tightlist
\item
  心理活动
\item
  学校心理健康资源
\end{itemize}

\begin{center}\rule{0.5\linewidth}{0.5pt}\end{center}

\subsubsection{3.
心理咨询的作用：像一根拐杖}\label{ux5fc3ux7406ux54a8ux8be2ux7684ux4f5cux7528ux50cfux4e00ux6839ux62d0ux6756}

老师用一个非常形象的比喻说明\textbf{心理咨询（psychological
counseling）}：

\begin{quote}
心理咨询就像人生中的一根拐杖。
\end{quote}

意思是：

\begin{itemize}
\tightlist
\item
  每个人在人生路上都可能``崴脚''\,``摔跟头''
\item
  当自己暂时走不动时，需要一个外在支撑
\item
  在支撑下，人一边往前走，一边慢慢恢复
\end{itemize}

当伤好了，拐杖就可以放下。

\begin{center}\rule{0.5\linewidth}{0.5pt}\end{center}

\subsubsection{4.
心理咨询不是``立竿见影''的药}\label{ux5fc3ux7406ux54a8ux8be2ux4e0dux662fux7acbux7affux89c1ux5f71ux7684ux836f}

老师提醒大家不要对心理咨询抱有不切实际的期待：

\begin{itemize}
\tightlist
\item
  不是一次咨询就能解决多年积累的问题
\item
  不是像吃感冒药一样立刻见效
\item
  它是一个缓慢渗透的过程
\end{itemize}

它的作用在于：

\begin{itemize}
\tightlist
\item
  帮助你更全面、深刻地理解自己
\item
  看见既往经历如何塑造了现在的自己
\item
  调整思维方式
\item
  逐渐把新理念、新方法渗透到生活、学习、工作中
\end{itemize}

\begin{center}\rule{0.5\linewidth}{0.5pt}\end{center}

\subsection{十六、自杀识别：课程最后的重要补充}\label{ux5341ux516dux81eaux6740ux8bc6ux522bux8bfeux7a0bux6700ux540eux7684ux91cdux8981ux8865ux5145}

老师说，这部分在危机干预中非常重要，因为：

\begin{itemize}
\tightlist
\item
  自杀相关问题风险等级最高
\item
  是咨询老师经常处理的重点内容
\end{itemize}

\begin{center}\rule{0.5\linewidth}{0.5pt}\end{center}

\subsubsection{1. 自杀的语言征兆（verbal warning
signs）}\label{ux81eaux6740ux7684ux8bedux8a00ux5f81ux5146verbal-warning-signs}

如果一个人经常说：

\begin{itemize}
\tightlist
\item
  我不想活了
\item
  活着特别没劲
\item
  太累了
\item
  活着没有意义
\end{itemize}

这些都不是随便说说，而是重要的\textbf{求助信号（help-seeking signal /
SOS）}，需要高度警觉。

老师特别强调：

\begin{itemize}
\tightlist
\item
  很多自杀成功案例，在事后回看时，前期其实有很多可以被识别、被救援的机会
\item
  只是周围人没有意识到严重性
\end{itemize}

\begin{center}\rule{0.5\linewidth}{0.5pt}\end{center}

\subsubsection{2. 自杀的行为征兆（behavioral warning
signs）}\label{ux81eaux6740ux7684ux884cux4e3aux5f81ux5146behavioral-warning-signs}

老师提到的行为信号包括：

\begin{itemize}
\tightlist
\item
  频繁搜索与自杀相关的信息
\item
  浏览自杀相关网站
\item
  搜索致命方式
\item
  准备自杀工具
\item
  进行踩点
\item
  物品整理与分发
\item
  出现道别、告别行为
\end{itemize}

【举例】如果一个人开始把自己珍视的东西分给别人，并说``以后我也用不着了''，这就是很危险的信号。

\begin{center}\rule{0.5\linewidth}{0.5pt}\end{center}

\subsubsection{3.
自杀过程中的情绪变化}\label{ux81eaux6740ux8fc7ux7a0bux4e2dux7684ux60c5ux7eeaux53d8ux5316}

老师指出，自杀风险并不是``越激动越危险''，真正要警惕的是一个阶段性的变化：

\paragraph{前期：}\label{ux524dux671f}

\begin{itemize}
\tightlist
\item
  易怒
\item
  烦躁
\item
  情绪剧烈波动
\end{itemize}

\paragraph{后期：}\label{ux540eux671f}

\begin{itemize}
\tightlist
\item
  突然进入``冷静期''
\end{itemize}

这个``冷静期''非常危险，因为这往往意味着：

\begin{itemize}
\tightlist
\item
  当事人已经做了决定
\item
  不再纠结
\item
  风险反而升高
\end{itemize}

老师特别提醒大家：

\begin{quote}
不要误以为``突然冷静下来''就是好了。
\end{quote}

\begin{center}\rule{0.5\linewidth}{0.5pt}\end{center}

\subsubsection{4.
社交媒体与失联信号}\label{ux793eux4ea4ux5a92ux4f53ux4e0eux5931ux8054ux4fe1ux53f7}

老师提到，一些外在线索也要引起警觉，例如：

\begin{itemize}
\tightlist
\item
  朋友圈背景或内容突然全部变黑
\item
  发布异常字眼
\item
  明显失联
\end{itemize}

这些都可能是危险信号。

\begin{center}\rule{0.5\linewidth}{0.5pt}\end{center}

\subsection{十七、关于自杀的几个常见误解}\label{ux5341ux4e03ux5173ux4e8eux81eaux6740ux7684ux51e0ux4e2aux5e38ux89c1ux8befux89e3}

\subsubsection{1.
误解一：有过自杀尝试的人，之后风险会降低}\label{ux8befux89e3ux4e00ux6709ux8fc7ux81eaux6740ux5c1dux8bd5ux7684ux4ebaux4e4bux540eux98ceux9669ux4f1aux964dux4f4e}

老师明确指出：\textbf{错。}

既往有自杀史（history of suicide
attempt）的人，之后的自杀风险反而\textbf{更高}。 原因是：

\begin{itemize}
\tightlist
\item
  既往尝试让其获得了某种``自杀能力''
\item
  包括对死亡、疼痛的恐惧降低
\end{itemize}

\begin{center}\rule{0.5\linewidth}{0.5pt}\end{center}

\subsubsection{2.
误解二：和有自杀风险的人谈论自杀，会诱发他去做}\label{ux8befux89e3ux4e8cux548cux6709ux81eaux6740ux98ceux9669ux7684ux4ebaux8c08ux8bbaux81eaux6740ux4f1aux8bf1ux53d1ux4ed6ux53bbux505a}

老师明确指出：\textbf{通常不是这样。}

如果你是以一种：

\begin{itemize}
\tightlist
\item
  真正关心他
\item
  很在意他生命安全
\item
  严肃而真诚
\end{itemize}

的态度去讨论，他反而可能会感到：

\begin{itemize}
\tightlist
\item
  如释重负
\item
  终于可以把无法说出口的压力讲出来
\end{itemize}

老师认为，这样的讨论常常有助于减轻其内在包袱，而不是增加风险。

\begin{center}\rule{0.5\linewidth}{0.5pt}\end{center}

\subsubsection{3.
误解三：说``我不想活了''的人只是说说而已}\label{ux8befux89e3ux4e09ux8bf4ux6211ux4e0dux60f3ux6d3bux4e86ux7684ux4ebaux53eaux662fux8bf4ux8bf4ux800cux5df2}

老师明确反对这种看法。

她认为，这通常是一个：

\begin{itemize}
\tightlist
\item
  \textbf{SOS求救信号}
\item
  说明此人至少已经认真考虑过这件事
\item
  也可能正在观察：周围是否有人愿意帮助自己
\end{itemize}

因此，周围人的态度至关重要。

\begin{center}\rule{0.5\linewidth}{0.5pt}\end{center}

\subsection{十八、普通同学遇到有自杀倾向的人时该怎么做}\label{ux5341ux516bux666eux901aux540cux5b66ux9047ux5230ux6709ux81eaux6740ux503eux5411ux7684ux4ebaux65f6ux8be5ux600eux4e48ux505a}

老师最后专门讲了``非专业人员能做什么''。

\subsubsection{1.
不要承诺保密}\label{ux4e0dux8981ux627fux8bfaux4fddux5bc6}

这是老师反复强调的重点。

如果对方说：

\begin{itemize}
\tightlist
\item
  我很信任你
\item
  但你一定要替我保密
\end{itemize}

这时不能答应。 老师提醒：

\begin{itemize}
\tightlist
\item
  一旦你承诺保密，而对方后来真的实施了严重行为
\item
  你可能会背负长期的内疚与自责
\end{itemize}

更合适的说法是：

\begin{itemize}
\tightlist
\item
  我非常担心你
\item
  我非常在意你
\item
  正因为我不希望失去你，所以我不能答应完全保密
\item
  我需要让更多人参与进来，先保证你的安全
\end{itemize}

\begin{center}\rule{0.5\linewidth}{0.5pt}\end{center}

\subsubsection{2.
立刻调动身边资源}\label{ux7acbux523bux8c03ux52a8ux8eabux8fb9ux8d44ux6e90}

老师强调，不要一个人硬扛。 普通同学做不到：

\begin{itemize}
\tightlist
\item
  24小时盯护
\item
  独自承担全部风险
\end{itemize}

因此要第一时间告诉：

\begin{itemize}
\item
  \textbf{辅导员}
\item
  必要时由老师进一步通知：

  \begin{itemize}
  \tightlist
  \item
    家长
  \item
    学校心理健康教育中心
  \item
    宿舍同学
  \item
    好朋友
  \item
    男女朋友
  \item
    其他支持者
  \end{itemize}
\end{itemize}

核心原则是：

\subsection{先救人，先让他活下来}\label{ux5148ux6551ux4ebaux5148ux8ba9ux4ed6ux6d3bux4e0bux6765}

活下来，才有后续帮助的可能。

\begin{center}\rule{0.5\linewidth}{0.5pt}\end{center}

\subsubsection{3.
后续可能的走向}\label{ux540eux7eedux53efux80fdux7684ux8d70ux5411}

老师指出，危机之后不一定都要长期心理咨询：

\begin{itemize}
\tightlist
\item
  有些同学是``一过性''的，危机过去就好了
\item
  有些同学则需要长期跟进、持续支持、心理帮扶
\end{itemize}

这取决于问题性质、人格基础和支持情况。

\begin{center}\rule{0.5\linewidth}{0.5pt}\end{center}

\subsection{十九、老师对``危机''的总体态度}\label{ux5341ux4e5dux8001ux5e08ux5bf9ux5371ux673aux7684ux603bux4f53ux6001ux5ea6}

课程最后，老师再次总结自己对危机的看法：

\subsubsection{1. 危机是危险}\label{ux5371ux673aux662fux5371ux9669}

因为它是人生的一道坎，可能直接威胁安全。

\subsubsection{2.
危机也是机会}\label{ux5371ux673aux4e5fux662fux673aux4f1a}

如果一个人顺利走过危机，他常常会获得：

\begin{itemize}
\tightlist
\item
  更强的问题解决能力
\item
  对自我的重新定义
\item
  更深的自我觉察（self-awareness）
\item
  更成熟的应对方式
\end{itemize}

【举例】老师再次回到前面``送饭后重新找到生命意义''的案例，说明：

\begin{itemize}
\tightlist
\item
  危机可能迫使一个人重新寻找生命意义
\item
  重新建立和世界、和他人的连接
\end{itemize}

因此，老师鼓励大家：

\begin{itemize}
\tightlist
\item
  不要只用悲观态度看待危机
\item
  要看到其中``学习与成长''的可能性
\end{itemize}

\begin{center}\rule{0.5\linewidth}{0.5pt}\end{center}

\subsection{二十、课堂结尾与随堂任务}\label{ux4e8cux5341ux8bfeux5802ux7ed3ux5c3eux4e0eux968fux5802ux4efbux52a1}

【课堂说明】老师最后提醒：

\begin{itemize}
\tightlist
\item
  本节课内容到此结束
\item
  有一个\textbf{随堂测试}
\item
  截止时间为\textbf{今晚12点}
\item
  需要在\textbf{雨课堂}完成
\item
  题目数量为\textbf{5道}
\end{itemize}

\begin{center}\rule{0.5\linewidth}{0.5pt}\end{center}

\subsection{本节英文术语索引}\label{ux672cux8282ux82f1ux6587ux672fux8bedux7d22ux5f15-4}

\begin{itemize}
\tightlist
\item
  心理危机 ------ psychological crisis
\item
  心理危机预防 ------ prevention of psychological crisis
\item
  心理危机干预 ------ psychological crisis intervention
\item
  法律 ------ law
\item
  伦理 ------ ethics
\item
  保密原则 ------ confidentiality
\item
  保密例外 / 突破保密 ------ breach of confidentiality / exception to
  confidentiality
\item
  最小披露 ------ minimal disclosure
\item
  双重关系 ------ dual relationship
\item
  监护人 ------ guardian
\item
  生命安全第一 ------ safety first / life safety first
\item
  生物学因素 ------ biological factors
\item
  心理因素 ------ psychological factors
\item
  创伤 ------ trauma
\item
  家庭治疗 ------ family therapy
\item
  依恋 / 依恋关系 ------ attachment / attachment relationship
\item
  社会支持 ------ social support
\item
  信仰 ------ belief / faith
\item
  价值观 ------ values
\item
  危机评估 ------ crisis assessment
\item
  自杀风险 ------ suicide risk
\item
  自杀计划 ------ suicide plan
\item
  行动干预 ------ action intervention
\item
  心理治疗 ------ psychotherapy
\item
  认知反应 ------ cognitive reactions
\item
  情绪反应 ------ emotional reactions
\item
  行为反应 ------ behavioral reactions
\item
  否认 ------ denial
\item
  过度警觉 ------ hypervigilance
\item
  资源激活 ------ resource activation
\item
  应对技巧 ------ coping skills
\item
  马斯洛需要层次理论 ------ Maslow's hierarchy of needs
\item
  生理需要 ------ physiological needs
\item
  安全需要 ------ safety needs
\item
  尊重需要 ------ esteem needs
\item
  自我实现需要 ------ self-actualization
\item
  自尊 ------ self-esteem
\item
  亲密关系 ------ intimacy / intimate relationship
\item
  自我认同 ------ self-identity
\item
  羞耻感 ------ shame
\item
  抱持 ------ holding
\item
  赋权 ------ empowerment
\item
  不伤害协议 ------ no-harm agreement / safety contract
\item
  治疗关系 / 咨询关系 ------ therapeutic relationship
\item
  心理咨询 ------ psychological counseling
\item
  求助信号 ------ help-seeking signal / SOS
\item
  自杀尝试史 ------ history of suicide attempt
\item
  自我觉察 ------ self-awareness
\item
  双相情感障碍 ------ bipolar disorder
\item
  精神分裂症 ------ schizophrenia
\item
  躁狂 ------ mania
\item
  家庭暴力 ------ domestic violence
\end{itemize}

你要是愿意，我下一条可以继续把这份笔记再压缩成一版``考试复习提纲版''，专门提炼成更适合临考背诵的版本。

\clearpage

\section{第6讲：学生心理危机三级预防}\label{ux7b2c6ux8bb2ux5b66ux751fux5fc3ux7406ux5371ux673aux4e09ux7ea7ux9884ux9632}

\textbf{课程}：心理危机干预与预防 \textbf{老师}：毛富强
\textbf{日期}：2026-04-07

\subsection{一、课前说明与本节导入}\label{ux4e00ux8bfeux524dux8bf4ux660eux4e0eux672cux8282ux5bfcux5165}

\begin{itemize}
\item
  老师先说明：这次上课通知存在混乱，课表与现场通知不一致，因此\textbf{本次不点名}，但仍按教学任务继续讲课。
\item
  老师提到，单讲抽象理论往往``不容易抓住要领''，因此本节改用\textbf{案例分析}的方式来讲\textbf{学生心理危机（psychological
  crisis in students）}。
\item
  老师说明案例选择的原因：

  \begin{itemize}
  \tightlist
  \item
    身边学校的真实案例通常涉及保密，不能直接讲。
  \item
    因此只能讲已经被网络公开、社会高度关注的案例。
  \end{itemize}
\item
  老师强调：\textbf{不同个体、不同事件背后的心理学原因，往往具有相似性}，所以个案分析有助于理解一般规律。
\end{itemize}

\begin{center}\rule{0.5\linewidth}{0.5pt}\end{center}

\subsection{二、案例引入：胡鑫宇事件的基本经过}\label{ux4e8cux6848ux4f8bux5f15ux5165ux80e1ux946bux5b87ux4e8bux4ef6ux7684ux57faux672cux7ecfux8fc7}

【案例】本节以\textbf{胡鑫宇事件}为例分析学生心理危机。

\begin{itemize}
\item
  胡鑫宇，江西省上饶市铅山县致远中学高一男生。
\item
  老师用``阳光大男孩''来形容其外在形象，指出这类学生放在日常校园环境中``并不违和''，因此\textbf{心理危机并不总是从外表上能看出来}。
\item
  事件时间线：

  \begin{itemize}
  \tightlist
  \item
    \textbf{2022年9月1日}入学。
  \item
    入学一个多月后，\textbf{10月14日}在校内失踪。
  \end{itemize}
\item
  学校背景：

  \begin{itemize}
  \tightlist
  \item
    为寄宿制学校。
  \item
    校内有教学楼、围墙、小山包、人工湖等区域。
  \item
    监控并非全覆盖，只覆盖部分楼道、教室、食堂等区域。
  \end{itemize}
\item
  失踪前最后轨迹：

  \begin{itemize}
  \tightlist
  \item
    监控显示其在\textbf{17:49}仍出现在宿舍门口。
  \item
    \textbf{18:00}应上晚自习，但没有到教室。
  \end{itemize}
\item
  班主任发现其未到课后，及时向学校报告；学校组织校内寻找，但未找到。
\item
  当晚约\textbf{23:00}，班主任联系家长，告知孩子失踪。
\end{itemize}

\begin{center}\rule{0.5\linewidth}{0.5pt}\end{center}

\subsection{三、家属介入、学校搜索与网络舆情发酵}\label{ux4e09ux5bb6ux5c5eux4ecbux5165ux5b66ux6821ux641cux7d22ux4e0eux7f51ux7edcux8206ux60c5ux53d1ux9175}

\begin{itemize}
\item
  胡鑫宇母亲次日早晨赶到学校，在宿舍中发现线索：

  \begin{itemize}
  \tightlist
  \item
    家里此前为其购买过\textbf{录音笔（voice
    recorder）}，原本用于学英语；
  \item
    这支录音笔后来被发现不在原处。
  \end{itemize}
\item
  学校扩大搜索范围：

  \begin{itemize}
  \tightlist
  \item
    地下室、阁楼、水池、人工湖、化粪池等区域都进行了寻找。
  \end{itemize}
\item
  之后报警，警方介入，通过走访与调取监控展开调查，但初期仍无结果。
\item
  家属随后开始在网络上持续发声，悬赏寻人金额从\textbf{10万}不断增加，最终到\textbf{50万}。
\item
  随着时间推移，网络上出现大量谣言，包括但不限于：

  \begin{itemize}
  \tightlist
  \item
    被藏于车辆中带离学校；
  \item
    被人抽血、贩卖器官；
  \item
    被教师杀害并处理尸体；
  \item
    校内曾发生多起类似失踪事件等。
  \end{itemize}
\item
  老师强调：\textbf{事后证明上述传言大多为谣言}，其中有人借机蹭热度、博眼球，造成了恶劣社会影响。
\end{itemize}

【提示】这里老师实际上借案例提醒：

\begin{itemize}
\tightlist
\item
  心理危机事件不仅有个体层面的伤害，
\item
  还会伴随\textbf{舆论失真、社会恐慌、次生伤害}等问题。
\end{itemize}

\begin{center}\rule{0.5\linewidth}{0.5pt}\end{center}

\subsection{四、案件结果：警方结论与录音证据}\label{ux56dbux6848ux4ef6ux7ed3ux679cux8b66ux65b9ux7ed3ux8bbaux4e0eux5f55ux97f3ux8bc1ux636e}

\begin{itemize}
\item
  失踪两个多月后，\textbf{2023年1月28日}，有人报警称在树林中发现一具尸体。
\item
  经\textbf{DNA（deoxyribonucleic acid）比对}确认，死者即为胡鑫宇。
\item
  现场还发现了重要物证：\textbf{录音笔}。
\item
  该录音笔被送往专业机构进行\textbf{声音鉴定（voice identification /
  forensic audio analysis）}与音频恢复。
\item
  公安部门最终公布结论：

  \begin{itemize}
  \tightlist
  \item
    胡鑫宇系\textbf{自缢身亡（death by suicide by hanging）}；
  \item
    现场为\textbf{第一现场（primary scene）}；
  \item
    无搏斗、拖拽痕迹；
  \item
    录音内容未被剪辑、未被篡改。
  \end{itemize}
\item
  录音时间覆盖\textbf{2022年10月14日17:40至23:08}之间的多个时段。
\item
  其中有两段录音较清晰地表达了\textbf{自杀意愿（suicidal intent）}。
\item
  老师据此归纳：胡鑫宇在失踪前已做出明确的自杀准备，并在临死前经历了较长时间的犹豫与挣扎。
\end{itemize}

\begin{center}\rule{0.5\linewidth}{0.5pt}\end{center}

\subsection{五、案例的心理学分析：胡鑫宇身上的危险特征}\label{ux4e94ux6848ux4f8bux7684ux5fc3ux7406ux5b66ux5206ux6790ux80e1ux946bux5b87ux8eabux4e0aux7684ux5371ux9669ux7279ux5f81}

\subsubsection{1.
性格特征：典型内向、孤僻、压抑}\label{ux6027ux683cux7279ux5f81ux5178ux578bux5185ux5411ux5b64ux50fbux538bux6291}

【考点】老师明确把``内向性格''视为该案例的重要风险因素。

\begin{itemize}
\item
  胡鑫宇被描述为\textbf{典型内向型人格（introverted personality）}：

  \begin{itemize}
  \tightlist
  \item
    有事不愿交流；
  \item
    不愿倾诉；
  \item
    不会主动求助；
  \item
    更倾向于``自己憋着''。
  \end{itemize}
\item
  老师认为，这类学生表面看``不显山不露水''，但\textbf{风险在于外界不容易及时发现其内心危机}。
\item
  其情绪宣泄方式较为隐蔽，只能在作业本边角写一些文字表达情绪。
\end{itemize}

\subsubsection{2.
对新环境的不适应}\label{ux5bf9ux65b0ux73afux5883ux7684ux4e0dux9002ux5e94}

【复习】老师借此回到心理危机发生机制：\textbf{环境变化 +
个体脆弱特征}容易触发危机。

\begin{itemize}
\item
  胡鑫宇在作业本上写过类似内容，反映：

  \begin{itemize}
  \tightlist
  \item
    对新环境难以适应；
  \item
    对自己内向的性格感到烦恼；
  \item
    只能通过写字来缓解情绪。
  \end{itemize}
\item
  老师由此判断：

  \begin{itemize}
  \tightlist
  \item
    他在进入高中后，对\textbf{学习环境}、\textbf{人际关系}等均不适应；
  \item
    性格特点加重了这种不适应。
  \end{itemize}
\end{itemize}

\subsubsection{3.
睡眠与节律问题}\label{ux7761ux7720ux4e0eux8282ux5f8bux95eeux9898}

【考点】老师把\textbf{睡眠障碍（sleep
disturbance）}列为判断其抑郁状态的重要证据之一。

\begin{itemize}
\item
  胡鑫宇写到：

  \begin{itemize}
  \tightlist
  \item
    早晨很早醒来；
  \item
    醒来后又躺着，不敢起床，怕影响别人。
  \end{itemize}
\item
  老师由此分析：

  \begin{itemize}
  \tightlist
  \item
    已出现\textbf{早醒（early morning awakening）}；
  \item
    也存在\textbf{入睡困难（difficulty falling asleep / insomnia）}；
  \item
    形成\textbf{节律紊乱（circadian rhythm disturbance）}。
  \end{itemize}
\item
  睡眠问题会进一步影响：

  \begin{itemize}
  \tightlist
  \item
    白天清醒度；
  \item
    上课状态；
  \item
    学习效率。
  \end{itemize}
\end{itemize}

\subsubsection{4.
对学习成就的高期待与现实落差}\label{ux5bf9ux5b66ux4e60ux6210ux5c31ux7684ux9ad8ux671fux5f85ux4e0eux73b0ux5b9eux843dux5dee}

\begin{itemize}
\tightlist
\item
  胡鑫宇原本中考成绩较好，并获得学校\textbf{1500元奖学金}入学。
\item
  家庭和本人都对其高中阶段抱有较高期待，希望通过努力考上重点大学。
\item
  但入学后因各种不适应因素，学习状态明显下降。
\item
  老师认为，这种``\textbf{高期待---低适应---高挫败}''的落差，是其心理危机加剧的重要背景。
\end{itemize}

\begin{center}\rule{0.5\linewidth}{0.5pt}\end{center}

\subsection{六、老师对其``重度抑郁''判断的六条依据}\label{ux516dux8001ux5e08ux5bf9ux5176ux91cdux5ea6ux6291ux90c1ux5224ux65adux7684ux516dux6761ux4f9dux636e}

【考点】老师明确提出：胡鑫宇在自杀前已表现出\textbf{重度抑郁（major
depression / severe depressive state）}的特征。

\subsubsection{（1）严重睡眠障碍}\label{ux4e25ux91cdux7761ux7720ux969cux788d}

\begin{itemize}
\tightlist
\item
  入睡困难；
\item
  早醒；
\item
  昼夜节律被破坏。
\end{itemize}

\subsubsection{（2）学习功能受损}\label{ux5b66ux4e60ux529fux80fdux53d7ux635f}

\begin{itemize}
\tightlist
\item
  \textbf{注意力不集中（inattention）}；
\item
  \textbf{记忆力下降（memory decline）}；
\item
  \textbf{思维迟钝（slowed thinking）}；
\item
  理解力下降；
\item
  白天上课``听不进去''。
\end{itemize}

\subsubsection{（3）明显社交障碍}\label{ux660eux663eux793eux4ea4ux969cux788d}

\begin{itemize}
\tightlist
\item
  自卑；
\item
  在意他人评价；
\item
  害怕被嫌弃；
\item
  采取\textbf{回避（avoidance）}策略；
\item
  越回避越孤独，越孤独越回避，形成恶性循环。
\end{itemize}

\subsubsection{（4）负性认知与自我评价下降}\label{ux8d1fux6027ux8ba4ux77e5ux4e0eux81eaux6211ux8bc4ux4ef7ux4e0bux964d}

\begin{itemize}
\tightlist
\item
  出现明显的\textbf{自责（self-blame）}；
\item
  \textbf{自我否定（self-denial）}；
\item
  觉得自己性格不行、学习能力不行、未来``没戏了''。
\item
  老师强调：负性认知越强，情绪越低落；情绪越低落，又会进一步强化负性认知。
\end{itemize}

\subsubsection{（5）心身症状}\label{ux5fc3ux8eabux75c7ux72b6}

\begin{itemize}
\tightlist
\item
  出现\textbf{厌食（loss of appetite / anorexia）}；
\item
  \textbf{呕吐（vomiting）}；
\item
  老师将其视为抑郁相关的\textbf{心身症状（psychosomatic
  symptoms）}表现。
\end{itemize}

\subsubsection{（6）持续情绪低落与轻生念头}\label{ux6301ux7eedux60c5ux7eeaux4f4eux843dux4e0eux8f7bux751fux5ff5ux5934}

\begin{itemize}
\tightlist
\item
  经常在无人时独自哭泣；
\item
  已出现``活着没意思''等想法；
\item
  说明已存在明确的\textbf{轻生观念（suicidal ideation）}。
\end{itemize}

【结论】老师认为：上述六条足以支持``胡鑫宇在自杀前已陷入严重抑郁状态''的判断。

\begin{center}\rule{0.5\linewidth}{0.5pt}\end{center}

\subsection{七、谁本可以帮助他：家长、老师、同学与本人}\label{ux4e03ux8c01ux672cux53efux4ee5ux5e2eux52a9ux4ed6ux5bb6ux957fux8001ux5e08ux540cux5b66ux4e0eux672cux4eba}

老师接下来围绕``胡鑫宇能不能不死''展开讨论，重点不是追责本身，而是从中提炼\textbf{危机干预（crisis
intervention）}与\textbf{预防（prevention）}的启示。

\begin{center}\rule{0.5\linewidth}{0.5pt}\end{center}

\subsection{八、家长层面的责任与失误}\label{ux516bux5bb6ux957fux5c42ux9762ux7684ux8d23ux4efbux4e0eux5931ux8bef}

【考点】老师认为：\textbf{家长是未成年学生心理危机中的第一道防线}。

\begin{itemize}
\item
  胡鑫宇曾多次向家中求助，尤其多次联系母亲。
\item
  \textbf{9月27日}，一天内与母亲通话三次，核心诉求是：

  \begin{itemize}
  \tightlist
  \item
    想回家；
  \item
    学不进去；
  \item
    想休息一段时间。
  \end{itemize}
\item
  \textbf{10月5日}，又连续与母亲通话三次，总时长近40分钟，仍在表达：

  \begin{itemize}
  \tightlist
  \item
    压力太大；
  \item
    希望回家休息；
  \item
    甚至提出休学、退学。
  \end{itemize}
\item
  老师指出：

  \begin{itemize}
  \tightlist
  \item
    母亲并未把这些求助理解为心理危机信号；
  \item
    而是将其理解为``娇气''\,``意志不坚定''。
  \end{itemize}
\item
  老师进一步分析，母亲拒绝孩子回家的原因包括：

  \begin{enumerate}
  \def\labelenumi{\arabic{enumi}.}
  \tightlist
  \item
    \textbf{经济因素}：家庭经济紧张，休学意味着额外成本；
  \item
    \textbf{期待压力}：全家把``改变命运''的希望寄托在孩子身上；
  \item
    \textbf{学校光环}：认为进入``好学校''就意味着必须坚持到底，不能退出。
  \end{enumerate}
\item
  老师的核心判断是：

  \begin{itemize}
  \tightlist
  \item
    家长缺乏心理危机识别知识；
  \item
    在孩子明确求助时未能提供支持；
  \item
    若当时能带孩子就医、休息、接受治疗，结局可能不同。
  \end{itemize}
\end{itemize}

【学习建议】老师借此提醒：

\begin{itemize}
\tightlist
\item
  当青少年连续、多次、情绪明显地表达``撑不住了''\,``想回家''``学不进去''时，不能只按``娇气''处理；
\item
  应把它当作\textbf{危机求助信号（help-seeking signal）}看待。
\end{itemize}

\begin{center}\rule{0.5\linewidth}{0.5pt}\end{center}

\subsection{九、老师层面的责任与失误}\label{ux4e5dux8001ux5e08ux5c42ux9762ux7684ux8d23ux4efbux4e0eux5931ux8bef}

【考点】寄宿制学校中的老师，某种程度上承担了\textbf{代行监护（in loco
parentis）}的职责。

\begin{itemize}
\item
  老师指出，学校教师的关注点往往被成绩、升学率、刷题压力占据。
\item
  对于胡鑫宇，班主任虽曾找他谈话，但重点放在：

  \begin{itemize}
  \tightlist
  \item
    成绩下滑；
  \item
    排名下降；
  \item
    学习状态不对。
  \end{itemize}
\item
  而没有真正深入关心：

  \begin{itemize}
  \tightlist
  \item
    他为什么难受；
  \item
    他是否存在情绪问题；
  \item
    他是否需要帮助。
  \end{itemize}
\item
  有教师甚至曾劝其``转学''，这可能进一步加重其挫败感与排斥感。
\item
  老师特别提到：

  \begin{itemize}
  \tightlist
  \item
    胡鑫宇在\textbf{10月10日至13日}之间，曾\textbf{7次}到宿舍楼三楼阳台``踩点''；
  \item
    但这些异常行为没有引起足够警觉。
  \end{itemize}
\item
  老师据此认为，这是学校在危机识别方面的明显缺失。
\end{itemize}

【易错】只关注成绩，不关注状态，是学校心理危机管理中的典型误区。

\begin{center}\rule{0.5\linewidth}{0.5pt}\end{center}

\subsection{十、同学层面的责任与失误}\label{ux5341ux540cux5b66ux5c42ux9762ux7684ux8d23ux4efbux4e0eux5931ux8bef}

\begin{itemize}
\item
  胡鑫宇在生前最后几天也曾向同学表达痛苦。
\item
  但同学们大多处于``自顾不暇''的学习状态，没有足够耐心或能力去承接这些信息。
\item
  他生前有很多异常表现，例如：

  \begin{itemize}
  \tightlist
  \item
    最后一天在教室坐了一整天；
  \item
    几乎不抬头；
  \item
    书也没有打开。
  \end{itemize}
\item
  但这些异常并未引起同伴、老师的有效反应。
\item
  老师强调：

  \begin{itemize}
  \tightlist
  \item
    只要其中某位同学、老师及时上报、反馈、干预，都可能改变结局。
  \item
    因此同伴系统也是学生危机预防中的重要环节。
  \end{itemize}
\end{itemize}

\begin{center}\rule{0.5\linewidth}{0.5pt}\end{center}

\subsection{十一、本人层面的因素：退避型应对与缺乏主动求助}\label{ux5341ux4e00ux672cux4ebaux5c42ux9762ux7684ux56e0ux7d20ux9000ux907fux578bux5e94ux5bf9ux4e0eux7f3aux4e4fux4e3bux52a8ux6c42ux52a9}

【考点】老师认为，个体自身的\textbf{应对方式（coping
style）}也是危机发展的关键因素。

\begin{itemize}
\item
  胡鑫宇在面对压力时，采用的是\textbf{退避（avoidant coping）}方式：

  \begin{itemize}
  \tightlist
  \item
    回避问题；
  \item
    假装没事；
  \item
    不愿麻烦别人；
  \item
    缺乏主动求医和主动求助。
  \end{itemize}
\item
  老师指出：

  \begin{itemize}
  \tightlist
  \item
    他其实有手机；
  \item
    理论上仍存在对外求助的机会；
  \item
    但他最终没有把求助真正推进到行动层面。
  \end{itemize}
\item
  老师借此说明：

  \begin{itemize}
  \tightlist
  \item
    个体不是完全被动的；
  \item
    每个人都需要学习在危机中主动寻求帮助。
  \end{itemize}
\end{itemize}

【学习建议】不要把``不给别人添麻烦''当成美德。

\begin{itemize}
\tightlist
\item
  在心理危机中，\textbf{主动求助（help-seeking）}本身就是一种能力。
\end{itemize}

\begin{center}\rule{0.5\linewidth}{0.5pt}\end{center}

\subsection{十二、由案例引出的生命意义教育}\label{ux5341ux4e8cux7531ux6848ux4f8bux5f15ux51faux7684ux751fux547dux610fux4e49ux6559ux80b2}

【学习建议】老师由此转向``生命意义（meaning of life）''教育。

\begin{itemize}
\item
  老师认为，我们从小较少被系统地教育``人为什么活着''。
\item
  常见的单一逻辑是：

  \begin{itemize}
  \tightlist
  \item
    学习为了成绩，
  \item
    成绩为了大学，
  \item
    大学为了工作，
  \item
    但最终``活着为了什么''反而不清楚。
  \end{itemize}
\item
  老师提到社会上对人生意义有不同理解，例如：

  \begin{itemize}
  \tightlist
  \item
    感受生活；
  \item
    创造价值；
  \item
    为社会做贡献；
  \item
    活得精彩。
  \end{itemize}
\item
  他强调：\textbf{人生意义没有标准答案}，它应当来自个体自己的体会。
\end{itemize}

\begin{center}\rule{0.5\linewidth}{0.5pt}\end{center}

\subsection{十三、卡伦·霍尼与``被爱的渴求''}\label{ux5341ux4e09ux5361ux4f26ux970dux5c3cux4e0eux88abux7231ux7684ux6e34ux6c42}

【考点】老师在这里引入\textbf{卡伦·霍尼（Karen Horney）}的观点。

\begin{itemize}
\item
  卡伦·霍尼是\textbf{新精神分析学派（Neo-Freudian / neo-psychoanalytic
  school）}的重要代表人物。
\item
  老师引用她的思想来解释：

  \begin{itemize}
  \tightlist
  \item
    人的心理健康不仅需要物质营养，还需要\textbf{精神营养（emotional
    nourishment）}；
  \item
    其中最关键的是\textbf{真挚、温暖、持续的母爱（maternal love）}。
  \end{itemize}
\item
  老师特别强调``\textbf{被爱的渴求（need to be loved）}''：

  \begin{itemize}
  \tightlist
  \item
    一个人若从小缺乏稳定、无条件的爱，
  \item
    可能在成长中形成不安全感、自卑感、人际困难，
  \item
    更容易在受挫时把失败理解为``没人爱我''\,``我被抛弃了''。
  \end{itemize}
\item
  老师把这一理论应用到胡鑫宇案例中：

  \begin{itemize}
  \tightlist
  \item
    胡鑫宇在困境中向家长求助却被拒绝，
  \item
    可能进一步强化了``只有我学习好，家人才爱我''的条件性信念，
  \item
    最终发展成绝望感。
  \end{itemize}
\item
  老师据此判断：

  \begin{itemize}
  \tightlist
  \item
    胡鑫宇的社交障碍不只是单纯内向，
  \item
    也可能与成长中情感支持不足有关。
  \end{itemize}
\end{itemize}

【学习建议】缺乏被爱体验的人，更需要通过后天学习来培养：

\begin{itemize}
\tightlist
\item
  主动爱人的能力；
\item
  建立关系的能力；
\item
  感受支持、接受支持的能力。
\end{itemize}

\begin{center}\rule{0.5\linewidth}{0.5pt}\end{center}

\subsection{十四、第二节：学生心理危机的三级预防}\label{ux5341ux56dbux7b2cux4e8cux8282ux5b66ux751fux5fc3ux7406ux5371ux673aux7684ux4e09ux7ea7ux9884ux9632}

【考点】本节后半部分正式转入\textbf{三级预防（three levels of
prevention）}。

\subsubsection{1. 一级预防（primary
prevention）}\label{ux4e00ux7ea7ux9884ux9632primary-prevention}

\begin{itemize}
\tightlist
\item
  含义：\textbf{消除病因或减少致病因素}；
\item
  核心：在危机发生前进行预防，防患于未然；
\item
  目标：尽量不让心理危机发生。
\end{itemize}

\subsubsection{2. 二级预防（secondary
prevention）}\label{ux4e8cux7ea7ux9884ux9632secondary-prevention}

\begin{itemize}
\item
  核心是``三早''：

  \begin{itemize}
  \tightlist
  \item
    \textbf{早发现（early detection）}
  \item
    \textbf{早诊断（early diagnosis）}
  \item
    \textbf{早治疗（early treatment）}
  \end{itemize}
\item
  含义：危机已经出现苗头或已经发生，此时要及时识别、及时干预，争取快速缓解、改善预后（prognosis）。
\end{itemize}

\subsubsection{3. 三级预防（tertiary
prevention）}\label{ux4e09ux7ea7ux9884ux9632tertiary-prevention}

\begin{itemize}
\item
  核心：\textbf{康复（rehabilitation）}与\textbf{防复发（relapse
  prevention）}；
\item
  目标：

  \begin{itemize}
  \tightlist
  \item
    减少再次发生；
  \item
    减少心理残疾或长期功能损害；
  \item
    帮助个体尽快恢复社会功能。
  \end{itemize}
\end{itemize}

【可能考简答】学生心理危机三级预防的定义与区别。

\begin{center}\rule{0.5\linewidth}{0.5pt}\end{center}

\subsection{十五、老师提到的社会层面防自杀共识}\label{ux5341ux4e94ux8001ux5e08ux63d0ux5230ux7684ux793eux4f1aux5c42ux9762ux9632ux81eaux6740ux5171ux8bc6}

老师提到，世界卫生组织认为并非所有自杀都能完全避免，但可通过公共策略降低风险。

\subsubsection{1.
减少获取自杀手段的机会}\label{ux51cfux5c11ux83b7ux53d6ux81eaux6740ux624bux6bb5ux7684ux673aux4f1a}

\begin{itemize}
\tightlist
\item
  例如限制高致死性农药的使用；
\item
  强调管制枪支、危险器具的重要性。
\item
  老师借中美差异举例，说明\textbf{高致死性工具的可得性（access to lethal
  means）}会显著影响非正常死亡风险。
\end{itemize}

\subsubsection{2.
积极治疗精神障碍患者}\label{ux79efux6781ux6cbbux7597ux7cbeux795eux969cux788dux60a3ux8005}

\begin{itemize}
\tightlist
\item
  \textbf{精神分裂症（schizophrenia）}、\textbf{抑郁症（depression）}等患者是高风险人群；
\item
  因此要做到及时诊断、及时治疗；
\item
  老师提到我国对\textbf{六类重性精神疾病（severe mental
  disorders）}进行登记管理。
\end{itemize}

\subsubsection{3.
对自杀未遂者进行随访}\label{ux5bf9ux81eaux6740ux672aux9042ux8005ux8fdbux884cux968fux8bbf}

\begin{itemize}
\tightlist
\item
  \textbf{自杀未遂（suicide attempt）}者后续再自杀风险较高；
\item
  因此不能把``救回来''当作问题终点，而应持续随访与干预。
\end{itemize}

\subsubsection{4.
负责任的媒体报道}\label{ux8d1fux8d23ux4efbux7684ux5a92ux4f53ux62a5ux9053}

【考点】老师强调：\textbf{自杀具有一定``传染性''或模仿效应（contagion /
imitation effect）}。

\begin{itemize}
\tightlist
\item
  媒体过度渲染具体事件、地点、方式，可能引发模仿；
\item
  但完全不报道也会带来隐患，因为公众可能缺乏警觉和重视。
\item
  这体现出危机传播中的平衡问题。
\end{itemize}

\subsubsection{5.
培训基层卫生保健人员}\label{ux57f9ux8badux57faux5c42ux536bux751fux4fddux5065ux4ebaux5458}

\begin{itemize}
\tightlist
\item
  需要建立心理危机识别和干预队伍；
\item
  基层人员是学生心理问题早期识别的重要力量。
\end{itemize}

\begin{center}\rule{0.5\linewidth}{0.5pt}\end{center}

\subsection{十六、学生心理危机重点人群：七类高风险对象}\label{ux5341ux516dux5b66ux751fux5fc3ux7406ux5371ux673aux91cdux70b9ux4ebaux7fa4ux4e03ux7c7bux9ad8ux98ceux9669ux5bf9ux8c61}

【考点】老师明确列出学生中的\textbf{七类重点人群（high-risk groups）}。

\subsubsection{1. 有家族史者（family
history）}\label{ux6709ux5bb6ux65cfux53f2ux8005family-history}

\begin{itemize}
\tightlist
\item
  直系亲属中有精神疾病者，个体心理脆弱性可能更高。
\end{itemize}

\subsubsection{2. 有既往史者（past history / previous psychiatric
history）}\label{ux6709ux65e2ux5f80ux53f2ux8005past-history-previous-psychiatric-history}

\begin{itemize}
\tightlist
\item
  过去曾发生过精神疾病或严重心理问题者，再次失衡风险更高。
\end{itemize}

\subsubsection{3. 单亲家庭学生（single-parent
family）}\label{ux5355ux4eb2ux5bb6ux5eadux5b66ux751fsingle-parent-family}

\begin{itemize}
\item
  可能面临：

  \begin{itemize}
  \tightlist
  \item
    经济压力；
  \item
    家庭角色功能不足；
  \item
    教养支持不完整。
  \end{itemize}
\end{itemize}

\subsubsection{4. 经济困难学生（financial
hardship）}\label{ux7ecfux6d4eux56f0ux96beux5b66ux751ffinancial-hardship}

\begin{itemize}
\tightlist
\item
  容易自卑、压力大，物质支持不足会放大心理困境。
\end{itemize}

\subsubsection{5. 存在人格缺陷者（personality
vulnerabilities）}\label{ux5b58ux5728ux4ebaux683cux7f3aux9677ux8005personality-vulnerabilities}

\begin{itemize}
\item
  老师列举：

  \begin{itemize}
  \tightlist
  \item
    内向；
  \item
    敏感；
  \item
    脆弱；
  \item
    完美主义（perfectionism）；
  \item
    高自尊；
  \item
    孤僻；
  \item
    执着、死心眼等。
  \end{itemize}
\end{itemize}

\subsubsection{6. 遭遇应激事件者（stressful life
events）}\label{ux906dux9047ux5e94ux6fc0ux4e8bux4ef6ux8005stressful-life-events}

\begin{itemize}
\item
  如：

  \begin{itemize}
  \tightlist
  \item
    重大考试失败；
  \item
    失恋；
  \item
    财产损失；
  \item
    慢性疾病；
  \item
    其他重大挫折与丧失。
  \end{itemize}
\end{itemize}

\subsubsection{7.
已出现精神症状者}\label{ux5df2ux51faux73b0ux7cbeux795eux75c7ux72b6ux8005}

\begin{itemize}
\item
  如：

  \begin{itemize}
  \tightlist
  \item
    焦虑（anxiety）；
  \item
    抑郁（depression）；
  \item
    幻觉（hallucination）；
  \item
    妄想（delusion）等。
  \end{itemize}
\end{itemize}

【可能考简答】学生心理危机重点人群有哪些。

\begin{center}\rule{0.5\linewidth}{0.5pt}\end{center}

\subsection{十七、心理危机的高发时间、年级与场所}\label{ux5341ux4e03ux5fc3ux7406ux5371ux673aux7684ux9ad8ux53d1ux65f6ux95f4ux5e74ux7ea7ux4e0eux573aux6240}

\subsubsection{1. 高发时间}\label{ux9ad8ux53d1ux65f6ux95f4}

\begin{itemize}
\tightlist
\item
  \textbf{春夏之交（4---5月）}
\item
  \textbf{秋冬之交}
\item
  老师认为换季时部分敏感个体更容易出问题。
\end{itemize}

\subsubsection{2. 高发年级}\label{ux9ad8ux53d1ux5e74ux7ea7}

\begin{itemize}
\item
  新生入学阶段问题较多，因为适应困难；
\item
  但总体上，\textbf{临近毕业年级问题更严重}。
\item
  原因在于：

  \begin{itemize}
  \tightlist
  \item
    毕业压力；
  \item
    毕不了业；
  \item
    考研失败；
  \item
    就业困难；
  \item
    异地恋与未来选择冲突等多种矛盾交织。
  \end{itemize}
\end{itemize}

\subsubsection{3. 高发场所}\label{ux9ad8ux53d1ux573aux6240}

\begin{itemize}
\item
  老师引用调查资料指出，大学生校内自杀中：

  \begin{itemize}
  \tightlist
  \item
    以\textbf{跳楼}最常见；
  \item
    主要发生在\textbf{宿舍楼、教学楼}等场所；
  \item
    许多事件发生在校内高层建筑。
  \end{itemize}
\item
  老师强调：

  \begin{itemize}
  \tightlist
  \item
    六楼以上高层建筑应加强管理；
  \item
    楼道、窗户等处应做好安全防护。
  \end{itemize}
\end{itemize}

【考点】危机防控不仅要关注``人''，也要关注``地点''和``环境''。

\begin{center}\rule{0.5\linewidth}{0.5pt}\end{center}

\subsection{十八、学校层面的一级预防：课程与活动}\label{ux5341ux516bux5b66ux6821ux5c42ux9762ux7684ux4e00ux7ea7ux9884ux9632ux8bfeux7a0bux4e0eux6d3bux52a8}

\subsubsection{1.
心理健康教育课程的重要性}\label{ux5fc3ux7406ux5065ux5eb7ux6559ux80b2ux8bfeux7a0bux7684ux91cdux8981ux6027}

【学习建议】老师强调：

\begin{itemize}
\tightlist
\item
  大学课堂教育仍是心理健康知识传播的\textbf{主渠道（main channel）}；
\item
  应用学生喜闻乐见的方式进行宣讲；
\item
  要树立``\textbf{健康第一，学业第二}''的观念。
\end{itemize}

老师传达的核心价值包括：

\begin{itemize}
\tightlist
\item
  热爱生命；
\item
  热爱生活；
\item
  了解自己；
\item
  接纳困难；
\item
  明白``办法总比困难多''。
\end{itemize}

\subsubsection{2.
参与心理健康教育活动}\label{ux53c2ux4e0eux5fc3ux7406ux5065ux5eb7ux6559ux80b2ux6d3bux52a8}

\begin{itemize}
\tightlist
\item
  老师认为，单靠理论课印象往往不深；
\item
  还需要通过体验式活动提升认知与自我察觉能力。
\end{itemize}

\begin{center}\rule{0.5\linewidth}{0.5pt}\end{center}

\subsection{十九、学校特色活动之一：内观认知疗法训练班}\label{ux5341ux4e5dux5b66ux6821ux7279ux8272ux6d3bux52a8ux4e4bux4e00ux5185ux89c2ux8ba4ux77e5ux7597ux6cd5ux8badux7ec3ux73ed}

【考点】老师重点介绍了学校的\textbf{内观认知疗法（Naikan-based cognitive
training / introspective cognitive therapy）}培训班。

\subsubsection{1. 基本情况}\label{ux57faux672cux60c5ux51b5-3}

\begin{itemize}
\tightlist
\item
  已举办\textbf{60多期}；
\item
  截至去年，完成学员\textbf{2600多人}；
\item
  今年学校``525心理健康月''仍将继续推广。
\end{itemize}

\subsubsection{2.
理论来源与形式}\label{ux7406ux8bbaux6765ux6e90ux4e0eux5f62ux5f0f}

\begin{itemize}
\tightlist
\item
  老师认为其源于\textbf{佛教禅修（Buddhist meditation / Zen-inspired
  introspection）}；
\item
  主要形式是\textbf{选择性回忆（selective recollection）}；
\item
  采用体验式训练；
\item
  一般持续\textbf{一周}。
\end{itemize}

\subsubsection{3. 组织方式}\label{ux7ec4ux7ec7ux65b9ux5f0f}

\begin{itemize}
\tightlist
\item
  通过简易屏风等形式营造体验环境；
\item
  每天集中数小时进行体验；
\item
  让学生通过回顾自身经历，察觉自身问题并促进成长。
\end{itemize}

\subsubsection{4. 效果统计}\label{ux6548ux679cux7edfux8ba1}

老师给出了对300多名学生的统计结果：

\begin{itemize}
\tightlist
\item
  三个问题中解决\textbf{两个}的占 \textbf{63.4\%}；
\item
  解决\textbf{一个}的占 \textbf{19.1\%}；
\item
  \textbf{一个都没解决}的占 \textbf{17.5\%}；
\item
  总有效率约 \textbf{82.5\%}。
\end{itemize}

【考点】这里可能涉及对活动效果的理解与记忆。

\subsubsection{5.
老师举的适用问题例子}\label{ux8001ux5e08ux4e3eux7684ux9002ux7528ux95eeux9898ux4f8bux5b50}

\begin{itemize}
\tightlist
\item
  自卑；
\item
  疑病倾向 / 过度担心自己患病；
\item
  人际关系问题；
\item
  情绪导致的身体不适。
\end{itemize}

\subsubsection{6.
老师总结其核心作用}\label{ux8001ux5e08ux603bux7ed3ux5176ux6838ux5fc3ux4f5cux7528}

\begin{itemize}
\tightlist
\item
  \textbf{学习感恩（gratitude）}
\item
  \textbf{学会体谅（empathy / understanding others）}
\item
  改变认知方式，学会珍惜已有、理解他人、减少怨恨与冲突。
\end{itemize}

【举例】老师用``追求异性成功或失败''举例，说明：

\begin{itemize}
\tightlist
\item
  得到后容易不珍惜；
\item
  得不到时容易怨恨；
\item
  通过内观训练，应学会珍惜拥有、尊重他人选择、减少执念与伤害性反应。
\end{itemize}

\begin{center}\rule{0.5\linewidth}{0.5pt}\end{center}

\subsection{二十、学校特色活动之二：中国森田疗法研修班}\label{ux4e8cux5341ux5b66ux6821ux7279ux8272ux6d3bux52a8ux4e4bux4e8cux4e2dux56fdux68eeux7530ux7597ux6cd5ux7814ux4feeux73ed}

【考点】老师重点介绍了\textbf{森田疗法（Morita therapy）}。

\subsubsection{1.
创立者与地位}\label{ux521bux7acbux8005ux4e0eux5730ux4f4d}

\begin{itemize}
\tightlist
\item
  创立者：日本学者\textbf{森田正马（Shoma Morita）}
\item
  创立时间：\textbf{1919年}
\item
  老师称其为东方心理疗法的重要代表，可与精神分析并列讨论。
\end{itemize}

\subsubsection{2. 理论来源}\label{ux7406ux8bbaux6765ux6e90}

\begin{itemize}
\item
  老师认为其思想本质上与\textbf{中国道家哲学（Taoist
  philosophy）}密切相关。
\item
  核心理念可概括为：

  \begin{itemize}
  \tightlist
  \item
    \textbf{接纳客观（accept reality）}
  \item
    \textbf{为所当为（do what should be done）}
  \item
    \textbf{顺其自然（let feelings be, act with purpose）}
  \end{itemize}
\end{itemize}

\subsubsection{3. 适应范围}\label{ux9002ux5e94ux8303ux56f4}

\begin{itemize}
\tightlist
\item
  神经症（neurosis）
\item
  酒依赖（alcohol dependence）
\item
  失眠（insomnia）
\item
  适应障碍（adjustment disorder）
\item
  应激障碍（stress-related disorder）
\item
  抑郁症等
\end{itemize}

\subsubsection{4.
老师对森田疗法原理的解释}\label{ux8001ux5e08ux5bf9ux68eeux7530ux7597ux6cd5ux539fux7406ux7684ux89e3ux91ca}

\begin{itemize}
\tightlist
\item
  人出现心理问题，本质上往往是\textbf{主观与客观不一致}；
\item
  个体总想用主观去强行改变客观，结果失败；
\item
  森田疗法的作用是让主观逐渐适应客观，从而减轻痛苦。
\end{itemize}

\subsubsection{5.
老师的学术活动与案例}\label{ux8001ux5e08ux7684ux5b66ux672fux6d3bux52a8ux4e0eux6848ux4f8b}

\begin{itemize}
\item
  老师提到自己曾赴日本参加相关会议；
\item
  介绍了与\textbf{糖尿病视网膜病变（diabetic
  retinopathy）}患者合作开展森田疗法干预的研究：

  \begin{itemize}
  \tightlist
  \item
    该类患者常伴有抑郁、焦虑；
  \item
    情绪差会降低治疗依从性（adherence/compliance）；
  \item
    接受森田疗法后，患者依从性及病情控制优于对照组。
  \end{itemize}
\end{itemize}

\subsubsection{6.
学校近期活动安排}\label{ux5b66ux6821ux8fd1ux671fux6d3bux52a8ux5b89ux6392}

\begin{itemize}
\item
  老师计划在\textbf{525心理健康月}期间开办\textbf{免费森田疗法体验班}；
\item
  形式包括：

  \begin{itemize}
  \tightlist
  \item
    线下讲座；
  \item
    线上40天体验；
  \item
    每天推送短课；
  \item
    研究生担任指导老师；
  \item
    每天约体验1小时；
  \item
    结业后发证书。
  \end{itemize}
\end{itemize}

\subsubsection{7.
老师概括的森田疗法最终目标}\label{ux8001ux5e08ux6982ux62ecux7684ux68eeux7530ux7597ux6cd5ux6700ux7ec8ux76eeux6807}

\begin{itemize}
\tightlist
\item
  \textbf{自我管理（self-management）}
\item
  规律生活；
\item
  吃好、睡好；
\item
  接纳情绪；
\item
  顺其自然；
\item
  为所当为；
\item
  改善人际关系；
\item
  获得社会支持；
\item
  及时识别症状并主动求助。
\end{itemize}

\begin{center}\rule{0.5\linewidth}{0.5pt}\end{center}

\subsection{二十一、三级预防的落脚点：接纳困难，及时求助}\label{ux4e8cux5341ux4e00ux4e09ux7ea7ux9884ux9632ux7684ux843dux811aux70b9ux63a5ux7eb3ux56f0ux96beux53caux65f6ux6c42ux52a9}

【学习建议】结尾部分，老师把三级预防落到现实生活中的态度和行动上。

\begin{itemize}
\item
  每个人都不是``神仙''，都会遇到暂时过不去的困难。
\item
  关键不在于``永远没问题''，而在于：

  \begin{itemize}
  \tightlist
  \item
    能否接纳自己的困难；
  \item
    能否及时求助；
  \item
    能否在问题早期进行处理。
  \end{itemize}
\item
  老师明确反对：

  \begin{itemize}
  \tightlist
  \item
    好面子；
  \item
    打肿脸充胖子；
  \item
    讳疾忌医（hide illness and avoid treatment）。
  \end{itemize}
\item
  他强调：

  \begin{itemize}
  \tightlist
  \item
    有些问题不是``硬扛''就能过去；
  \item
    及时求助并不丢人，反而是成熟和负责的表现。
  \end{itemize}
\end{itemize}

\begin{center}\rule{0.5\linewidth}{0.5pt}\end{center}

\subsection{本节英文术语索引}\label{ux672cux8282ux82f1ux6587ux672fux8bedux7d22ux5f15-5}

\begin{itemize}
\tightlist
\item
  心理危机 --- psychological crisis
\item
  危机干预 --- crisis intervention
\item
  预防 --- prevention
\item
  一级预防 --- primary prevention
\item
  二级预防 --- secondary prevention
\item
  三级预防 --- tertiary prevention
\item
  早发现、早诊断、早治疗 --- early detection, early diagnosis, early
  treatment
\item
  康复 --- rehabilitation
\item
  预后 --- prognosis
\item
  防复发 --- relapse prevention
\item
  内向型人格 --- introverted personality
\item
  睡眠障碍 --- sleep disturbance
\item
  早醒 --- early morning awakening
\item
  入睡困难 / 失眠 --- difficulty falling asleep / insomnia
\item
  节律紊乱 --- circadian rhythm disturbance
\item
  重度抑郁 --- major depression / severe depressive state
\item
  社交障碍 --- social impairment / social dysfunction
\item
  回避 --- avoidance
\item
  自责 --- self-blame
\item
  自我否定 --- self-denial
\item
  心身症状 --- psychosomatic symptoms
\item
  厌食 / 食欲下降 --- loss of appetite / anorexia
\item
  情绪低落 --- depressed mood
\item
  轻生观念 / 自杀意念 --- suicidal ideation
\item
  自杀意愿 --- suicidal intent
\item
  自杀未遂 --- suicide attempt
\item
  家族史 --- family history
\item
  既往史 --- past history / previous psychiatric history
\item
  应激事件 --- stressful life events
\item
  焦虑 --- anxiety
\item
  抑郁 --- depression
\item
  幻觉 --- hallucination
\item
  妄想 --- delusion
\item
  负性认知 --- negative cognition
\item
  主动求助 --- help-seeking
\item
  生命意义 --- meaning of life
\item
  被爱的渴求 --- need to be loved
\item
  母爱 --- maternal love
\item
  新精神分析学派 --- neo-psychoanalytic school
\item
  卡伦·霍尼 --- Karen Horney
\item
  声音鉴定 --- forensic audio analysis / voice identification
\item
  DNA比对 --- DNA matching
\item
  录音笔 --- voice recorder
\item
  内观认知疗法 --- Naikan-based cognitive training / introspective
  cognitive therapy
\item
  选择性回忆 --- selective recollection
\item
  感恩 --- gratitude
\item
  体谅 / 共情理解 --- empathy / understanding others
\item
  森田疗法 --- Morita therapy
\item
  森田正马 --- Shoma Morita
\item
  神经症 --- neurosis
\item
  酒依赖 --- alcohol dependence
\item
  适应障碍 --- adjustment disorder
\item
  应激障碍 --- stress-related disorder
\item
  自我管理 --- self-management
\item
  社会支持 --- social support
\item
  糖尿病视网膜病变 --- diabetic retinopathy
\item
  依从性 --- adherence / compliance
\end{itemize}

这节课如果你还要的话，我也可以继续给你补一个\textbf{``极简背诵版''}，专门提炼出老师最可能考、最值得期末前速看的内容。

\clearpage

\section{第7讲：心理问题的客观存在、主观加工与应对思路}\label{ux7b2c7ux8bb2ux5fc3ux7406ux95eeux9898ux7684ux5ba2ux89c2ux5b58ux5728ux4e3bux89c2ux52a0ux5de5ux4e0eux5e94ux5bf9ux601dux8def}

\begin{itemize}
\tightlist
\item
  课程：心理危机干预与预防
\item
  老师：李建伟
\item
  日期：2026-04-14
\end{itemize}

\begin{center}\rule{0.5\linewidth}{0.5pt}\end{center}

\subsection{一、课前说明：下次考试与本节课程安排}\label{ux4e00ux8bfeux524dux8bf4ux660eux4e0bux6b21ux8003ux8bd5ux4e0eux672cux8282ux8bfeux7a0bux5b89ux6392}

【考务】老师一开始先说明下次课是\textbf{考试}，要求如下：

\begin{itemize}
\tightlist
\item
  下次课一定要\textbf{准时到达}。
\item
  考试形式是\textbf{开卷考试（open-book exam）}。
\item
  但\textbf{不允许带电子设备}，只能带\textbf{纸质资料}。
\item
  参加考试时带\textbf{学生证}即可。
\item
  没来的同学要互相通知，下周记得来参加考试。
\item
  老师强调：只要认真作答，\textbf{一般不会故意为难大家}，但\textbf{不能缺考}。
\end{itemize}

【课后任务】
每次课后会有几道题，这次课后也会布置\textbf{两道非常简单的题}，提醒大家不要忘记完成。

【课程安排】
本来这节课计划是本课程最后一次课的\textbf{电影赏析}，但因为播放器出了问题，影片无法播放，所以改为继续讲课。老师顺带推荐了一个比较符合心理学口味的电影系列：

【举例】《头脑特工队》（\emph{Inside Out}）

\begin{itemize}
\tightlist
\item
  老师认为，这类作品体现了西方思维的一种特点：喜欢把人的情绪拆分、分析其构成与原因。
\item
  这种思维方式有其优势：能够把问题描述得比较清楚。
\item
  但它也有局限：\textbf{把问题说清楚了，不等于就有办法解决问题}。
\item
  老师借此引到医学中的类似现象：病因、机制可能讲得很清楚，但真正可用的治疗办法并不一定很多。
\end{itemize}

\begin{center}\rule{0.5\linewidth}{0.5pt}\end{center}

\subsection{二、导入：心理问题``看不见、摸不着''，但客观存在}\label{ux4e8cux5bfcux5165ux5fc3ux7406ux95eeux9898ux770bux4e0dux89c1ux6478ux4e0dux7740ux4f46ux5ba2ux89c2ux5b58ux5728}

老师指出，心理学要帮助大家换一个角度理解心理问题的产生与存在方式。

\subsubsection{1.
每个人都在散发不同的``心理状态''}\label{ux6bcfux4e2aux4ebaux90fdux5728ux6563ux53d1ux4e0dux540cux7684ux5fc3ux7406ux72b6ux6001}

老师让大家体会：虽然大家都坐在教室里，但每个人其实都带着不同的信息、状态或``气场''：

\begin{itemize}
\tightlist
\item
  有些人焦虑；
\item
  有些人没有动力；
\item
  有些人静不下来，必须不断干点别的事。
\end{itemize}

【举例】老师提到一个来访者：

\begin{itemize}
\tightlist
\item
  一会儿看电视，一会儿刷手机；
\item
  很难专注于做一件事情；
\item
  自己怀疑是不是 \textbf{ADHD（Attention-Deficit/Hyperactivity
  Disorder，注意缺陷多动障碍）}。
\end{itemize}

老师的判断是：

\begin{itemize}
\tightlist
\item
  这人高中以前学习还不错，因此\textbf{未必真的是 ADHD}；
\item
  更重要的问题在于：\textbf{无法和自己独处}，一旦独处就受不了，像要``发疯''一样；
\item
  所以他不是单纯``坐不住''，而是\textbf{难以承受与自己相处时的内在状态}。
\end{itemize}

老师借这个例子强调：心理困扰不是一句``你坐那儿看书不就行了吗''就能解决的，因为当事人\textbf{不是不想，而是做不到}。

\subsubsection{2.
心理问题的第一个核心认识}\label{ux5fc3ux7406ux95eeux9898ux7684ux7b2cux4e00ux4e2aux6838ux5fc3ux8ba4ux8bc6}

【重点】心理问题有一个非常重要的特点：

\begin{itemize}
\tightlist
\item
  \textbf{看不见}
\item
  \textbf{摸不着}
\item
  但它在每个人内在都是一种\textbf{客观存在（objective existence）}
\end{itemize}

老师进一步解释：

\begin{itemize}
\tightlist
\item
  对于看得见、摸得着的东西，人容易找到调控对象、找到``靶点''。
\item
  但如果一个对象既看不见又摸不着，你还要去调控它，就会变得特别困难。
\item
  这正是心理学面临的特点：它明明存在，却不易直接把握。
\end{itemize}

\subsubsection{3.
每个人都不一样}\label{ux6bcfux4e2aux4ebaux90fdux4e0dux4e00ux6837}

老师再次强调个体差异：

\begin{itemize}
\tightlist
\item
  有人喜欢孤独；
\item
  有人无法忍受孤独；
\item
  有人渴望被赞许；
\item
  有人一接触他人就紧张。
\end{itemize}

这些都说明：心理状态确实存在，而且\textbf{个体之间差异极大}。

\begin{center}\rule{0.5\linewidth}{0.5pt}\end{center}

\subsection{三、为什么这个时代心理问题更多}\label{ux4e09ux4e3aux4ec0ux4e48ux8fd9ux4e2aux65f6ux4ee3ux5fc3ux7406ux95eeux9898ux66f4ux591a}

老师提出一个问题：为什么现在这个年代心理问题特别多？并追问大家：

\begin{itemize}
\tightlist
\item
  你们现在还高兴吗？
\item
  还快乐吗？
\end{itemize}

老师的判断是：现在这个时代，敢明确说自己很高兴、很快乐的人，其实很少。

\subsubsection{1.
外部信息不断侵入精神世界}\label{ux5916ux90e8ux4fe1ux606fux4e0dux65adux4fb5ux5165ux7cbeux795eux4e16ux754c}

老师认为，现在每个人基本都有：

\begin{itemize}
\tightlist
\item
  手机
\item
  电脑
\end{itemize}

而我们的精神世界不断被外部信息占据。这些信息往往带有：

\begin{itemize}
\tightlist
\item
  情绪化
\item
  不安稳
\item
  焦灼
\item
  刺激性
\end{itemize}

【重点】任何信息都会在大脑中留下痕迹。
长期如此，人的整体精神状态就会变得：

\begin{itemize}
\tightlist
\item
  不平稳
\item
  更容易情绪化
\item
  更容易焦虑和不安
\end{itemize}

\begin{center}\rule{0.5\linewidth}{0.5pt}\end{center}

\subsection{四、案例一：不会做题就崩溃------内在批评声音的发现}\label{ux56dbux6848ux4f8bux4e00ux4e0dux4f1aux505aux9898ux5c31ux5d29ux6e83ux5185ux5728ux6279ux8bc4ux58f0ux97f3ux7684ux53d1ux73b0}

【举例】老师讲了一个高中生的案例。

\subsubsection{1. 表面表现}\label{ux8868ux9762ux8868ux73b0}

这个学生的特点是：

\begin{itemize}
\tightlist
\item
  一道题不会做，就会\textbf{崩溃}
\item
  所谓崩溃，就是必须\textbf{大哭一场}
\end{itemize}

家人会劝他：

\begin{itemize}
\tightlist
\item
  不会做题很正常；
\item
  哭也解决不了问题。
\end{itemize}

但老师强调：这恰恰说明一个事实------\textbf{心理问题是客观存在的}，不是靠讲道理就能立刻改变。

\subsubsection{2.
崩溃后的反应链}\label{ux5d29ux6e83ux540eux7684ux53cdux5e94ux94fe}

这个学生的典型反应是：

\begin{enumerate}
\def\labelenumi{\arabic{enumi}.}
\tightlist
\item
  题不会做；
\item
  崩溃大哭；
\item
  哭完以后头晕；
\item
  倒头睡觉，睡很久；
\item
  睡完以后才稍微缓过来。
\end{enumerate}

\subsubsection{3.
心理咨询中的分析路径}\label{ux5fc3ux7406ux54a8ux8be2ux4e2dux7684ux5206ux6790ux8defux5f84}

老师借此说明心理咨询的思路：

\begin{itemize}
\tightlist
\item
  不是只看外在行为；
\item
  而是顺着``当时你是什么感觉''这个线索往下找。
\end{itemize}

老师指出，心理学某种意义上就是帮助人\textbf{认识自己}。

于是继续往下追问：

\begin{itemize}
\tightlist
\item
  当题不会做的时候，你内心到底发生了什么？
\end{itemize}

\subsubsection{4.
找到``内在批评者''}\label{ux627eux5230ux5185ux5728ux6279ux8bc4ux8005}

在咨询中，这个孩子逐渐意识到：

\begin{itemize}
\tightlist
\item
  当题不会做的时候，
\item
  他总感觉旁边有一个\textbf{强烈的声音在批评他}。
\end{itemize}

这个声音像什么？

\begin{itemize}
\tightlist
\item
  像一个一直盯着自己的``第三只眼''；
\item
  像一个不断审视、指责自己的存在。
\end{itemize}

老师借此和同学们建立共鸣：

【提示】如果你自律性特别强，常常感觉总有一个声音在盯着自己、批评自己、要求自己，那说明你活得并不轻松，你的内在是被掌控住的。

\subsubsection{5.
这个声音具体在说什么}\label{ux8fd9ux4e2aux58f0ux97f3ux5177ux4f53ux5728ux8bf4ux4ec0ux4e48}

在进一步还原时，这个学生感受到那个声音像在说：

\begin{itemize}
\tightlist
\item
  ``你这个题又不会做。''
\item
  ``你这个题都不会做。''
\end{itemize}

问题在于：

\begin{itemize}
\tightlist
\item
  这个声音一来，他\textbf{没法反驳}；
\item
  也\textbf{没法对抗}；
\item
  所以会一下子崩溃。
\end{itemize}

\subsubsection{6. 双椅技术（empty-chair / two-chair
technique）}\label{ux53ccux6905ux6280ux672fempty-chair-two-chair-technique}

老师在这里提到心理咨询中的一个技术：

【术语】双椅技术（two-chair technique / empty-chair technique）

大意是：

\begin{itemize}
\tightlist
\item
  把内在那个批评自己的部分``具象化''出来；
\item
  让来访者尝试和那个声音、那个``人''对话；
\item
  例如对它说：``你别总这样批评我。''
\end{itemize}

但在实际操作中，这个学生一旦被要求去``面对''那个内在批评者时，立刻表现出强烈恐惧：

\begin{itemize}
\tightlist
\item
  整个人一下子转过去；
\item
  明明没有外人；
\item
  但他就是\textbf{不敢看``它''}。
\end{itemize}

老师借此再次强调：

【重点】心理问题是真实不虚的。
虽然你看不见、摸不着它，但你可能\textbf{根本控制不了它}。

\begin{center}\rule{0.5\linewidth}{0.5pt}\end{center}

\subsection{五、精神世界像``心田''：长什么，决定你活成什么样}\label{ux4e94ux7cbeux795eux4e16ux754cux50cfux5fc3ux7530ux957fux4ec0ux4e48ux51b3ux5b9aux4f60ux6d3bux6210ux4ec0ux4e48ux6837}

老师随后引入一个比喻：

\begin{itemize}
\tightlist
\item
  人的内在精神世界，像一块\textbf{田地（心田）}。
\end{itemize}

\subsubsection{1.
心田里如果是石头和杂草，就长不出好庄稼}\label{ux5fc3ux7530ux91ccux5982ux679cux662fux77f3ux5934ux548cux6742ux8349ux5c31ux957fux4e0dux51faux597dux5e84ux7a3c}

如果一块田里有很多：

\begin{itemize}
\tightlist
\item
  石头
\item
  杂草
\end{itemize}

那么它自然长不出好的庄稼。

老师将其类比到人的内在：

\begin{itemize}
\item
  一个人很痛苦时，心里长出来的常常是：

  \begin{itemize}
  \tightlist
  \item
    痛苦
  \item
    抱怨
  \item
    不满
  \end{itemize}
\item
  一个人若内在能长出：

  \begin{itemize}
  \tightlist
  \item
    感恩
  \item
    快乐
  \item
    高兴
  \end{itemize}
\end{itemize}

那他活得就会轻松自在。

\subsubsection{2.
抱怨型人格的理解}\label{ux62b1ux6028ux578bux4ebaux683cux7684ux7406ux89e3}

老师指出：

\begin{itemize}
\tightlist
\item
  如果一个人总在抱怨别人哪里不对、哪里不好，
\item
  可以高度怀疑：他的``心田''里长不出丰盛的庄稼，只能长出``杂草''。
\end{itemize}

\subsubsection{\texorpdfstring{3.
``我必须\ldots\ldots{}''带来的僵化}{3. ``我必须\ldots\ldots''带来的僵化}}\label{ux6211ux5fc5ux987bux5e26ux6765ux7684ux50f5ux5316}

老师还提到一种常见的内在状态：

\begin{itemize}
\tightlist
\item
  ``我必须要怎么样''
\item
  ``我一定要怎么样''
\end{itemize}

【易错理解】这类念头表面看像上进、自律，实际上会让人的精神世界变得：

\begin{itemize}
\tightlist
\item
  固定
\item
  僵硬
\item
  不灵活
\end{itemize}

而这恰恰会带来更大的痛苦。

\subsubsection{4.
一个颠覆性的观点：痛苦未必主要来自外在环境}\label{ux4e00ux4e2aux98a0ux8986ux6027ux7684ux89c2ux70b9ux75dbux82e6ux672aux5fc5ux4e3bux8981ux6765ux81eaux5916ux5728ux73afux5883}

老师特别强调一个他认为很重要的观点：

\begin{itemize}
\tightlist
\item
  很多人以为自己的痛苦都是外在环境造成的；
\item
  但实际上，\textbf{即使没有这个事情，也会有别的事情冒出来困扰你}。
\end{itemize}

也就是说：

【重点】现在的痛苦，不一定主要取决于当前的外在环境；更核心的是，你的内在本身就有让你痛苦的``种子''或模式。

\begin{center}\rule{0.5\linewidth}{0.5pt}\end{center}

\subsection{六、案例二：强迫症的迁移------困扰会不断换对象}\label{ux516dux6848ux4f8bux4e8cux5f3aux8febux75c7ux7684ux8fc1ux79fbux56f0ux6270ux4f1aux4e0dux65adux6362ux5bf9ux8c61}

【举例】老师讲了一个强迫症（Obsessive-Compulsive Disorder,
OCD）患者的经历，用来说明痛苦背后是同一种内在焦虑在不断``发芽''。

\subsubsection{1.
小时候：担心脚卷入自行车后轮}\label{ux5c0fux65f6ux5019ux62c5ux5fc3ux811aux5377ux5165ux81eaux884cux8f66ux540eux8f6e}

\begin{itemize}
\item
  父亲骑自行车送他上学；
\item
  他坐在后面，总想：

  \begin{itemize}
  \tightlist
  \item
    如果自己的脚塞进后车轮、车辐条，会怎样？
  \end{itemize}
\item
  因此非常痛苦，还埋怨父亲不要骑车送他。
\end{itemize}

后来：

\begin{itemize}
\tightlist
\item
  住校了；
\item
  不骑车了；
\item
  这个问题暂时没了。
\end{itemize}

\subsubsection{2.
上高中：开始被物理``受力分析''困住}\label{ux4e0aux9ad8ux4e2dux5f00ux59cbux88abux7269ux7406ux53d7ux529bux5206ux6790ux56f0ux4f4f}

高中学物理后，他又开始出现新困扰：

\begin{itemize}
\tightlist
\item
  走路时不停分析脚受了哪些力；
\item
  这个力是怎么产生的；
\item
  脑子停不下来。
\end{itemize}

他当时以为：

\begin{itemize}
\tightlist
\item
  只要不学物理了，
\item
  自己就自由了。
\end{itemize}

\subsubsection{3.
上大学：又被马克思主义原理困住}\label{ux4e0aux5927ux5b66ux53c8ux88abux9a6cux514bux601dux4e3bux4e49ux539fux7406ux56f0ux4f4f}

后来不学物理了，又进入大学。结果又出现新的强迫性困扰：

\begin{itemize}
\tightlist
\item
  为什么第一章讲这个？
\item
  为什么第二章讲那个？
\item
  为什么这个内容要放在最后一章？
\end{itemize}

他会反复琢磨：

\begin{itemize}
\tightlist
\item
  \textbf{必须}把这个逻辑彻底搞明白；
\item
  搞不明白，别的内容都看不进去。
\end{itemize}

而且偏偏考研还会考这部分，于是他和这个问题较劲了四五年。

\subsubsection{4.
工作后：困扰再次换壳}\label{ux5de5ux4f5cux540eux56f0ux6270ux518dux6b21ux6362ux58f3}

后来工作了，不考这些理论了，但他当了老师以后，新的困扰又冒出来：

\begin{itemize}
\tightlist
\item
  物理力学该怎么给学生讲？
\item
  这个力学关系到底怎么来？
\end{itemize}

【结论】老师借这个案例说明：

\begin{itemize}
\tightlist
\item
  表面看，人的困扰对象一直在变；
\item
  但核心不在这些具体对象本身；
\item
  核心在于你内在有一个\textbf{焦虑、痛苦的种子}，它会不断寻找新的载体发芽。
\end{itemize}

所以：

【备考/理解重点】不要以为``等十年后离开这个环境我就好了''。
老师的意思是：如果根本模式不变，换了环境也可能继续痛苦。

\begin{center}\rule{0.5\linewidth}{0.5pt}\end{center}

\subsection{七、情绪幸福感相对稳定：心理状态具有持续性}\label{ux4e03ux60c5ux7eeaux5e78ux798fux611fux76f8ux5bf9ux7a33ux5b9aux5fc3ux7406ux72b6ux6001ux5177ux6709ux6301ux7eedux6027}

老师提到一个调查研究，大意是连续测量人的情绪幸福指数，结果发现：

\begin{itemize}
\tightlist
\item
  假设满分 10 分；
\item
  有的人长期大约在 3 分；
\item
  有的人长期大约在 10 分；
\end{itemize}

即便过了很多年，这种总体倾向也相对稳定。

老师借此说明：

\begin{itemize}
\tightlist
\item
  如果你现在总觉得活得痛苦，
\item
  十年以后可能仍然容易有类似感觉；
\item
  如果你现在比较愉快、关系也好，
\item
  十年以后也可能保持这种倾向。
\end{itemize}

他借此再次说明：

【重点】心理问题有其\textbf{客观实在性}，不是别人说一句``想开点''就能马上改变。

\begin{center}\rule{0.5\linewidth}{0.5pt}\end{center}

\subsection{八、老师对心理问题的总结：主观性、必然性、主导性}\label{ux516bux8001ux5e08ux5bf9ux5fc3ux7406ux95eeux9898ux7684ux603bux7ed3ux4e3bux89c2ux6027ux5fc5ux7136ux6027ux4e3bux5bfcux6027}

老师说，这三个特点是他自己的总结。

\subsubsection{（一）主观性（subjectivity）}\label{ux4e00ux4e3bux89c2ux6027subjectivity}

\begin{itemize}
\tightlist
\item
  每个人的痛苦程度和痛苦方式都不一样；
\item
  没有两个人的痛苦是完全相同的。
\end{itemize}

【举例】社交恐惧 / 社交焦虑障碍（Social Anxiety Disorder, SAD）

\paragraph{1. 女生常见表现}\label{ux5973ux751fux5e38ux89c1ux8868ux73b0}

\begin{itemize}
\tightlist
\item
  一说话就脸红；
\item
  脸红后觉得很丢人；
\item
  越觉得丢人越紧张；
\item
  脑子一片空白；
\item
  更说不出话；
\item
  最终形成闭环。
\end{itemize}

\paragraph{2. 男生常见表现}\label{ux7537ux751fux5e38ux89c1ux8868ux73b0}

老师举了一个门诊案例：

\begin{itemize}
\item
  这个男生不敢出去开会；
\item
  起初老师以为可能是\textbf{广场恐惧症（agoraphobia）}，因为有些人会不敢坐火车、不敢去公共场所；
\item
  但对方说不是，他敢坐火车，问题在于：

  \begin{itemize}
  \tightlist
  \item
    开会休息时大家会去厕所；
  \item
    他不敢去公共厕所。
  \end{itemize}
\end{itemize}

老师指出：

\begin{itemize}
\tightlist
\item
  男性的社交焦虑障碍，有时会表现为\textbf{不敢去公共厕所}；
\item
  这可以是其典型表现之一。
\end{itemize}

\paragraph{3.
另一位社交焦虑来访者：上班像``上坟''}\label{ux53e6ux4e00ux4f4dux793eux4ea4ux7126ux8651ux6765ux8bbfux8005ux4e0aux73edux50cfux4e0aux575f}

这个来访者每天上班非常痛苦，原因不是工作本身，而是\textbf{怎么和人相处}让他极度焦虑：

\begin{itemize}
\item
  坐地铁时担心遇到同事：

  \begin{itemize}
  \tightlist
  \item
    要不要打招呼？
  \item
    遇到领导怎么打？
  \item
    遇到普通同事怎么打？
  \end{itemize}
\item
  拿着豆浆和领导打招呼合不合适？
\item
  干脆不坐电梯，去爬楼梯；
\item
  到工位还要经过几排桌子，又开始纠结：

  \begin{itemize}
  \tightlist
  \item
    要不要和同事打招呼？
  \item
    如果对方没看我，我要不要主动打？
  \end{itemize}
\end{itemize}

白天好不容易熬过去了，下班后还不轻松，还要不断\textbf{复盘（rumination）}：

\begin{itemize}
\tightlist
\item
  我早上说的话妥不妥当？
\item
  别人会不会觉得我蠢？
\end{itemize}

结果每天筋疲力尽。

老师借此说明：

【重点】对有些人来说，和人交往根本不是``小事''，而是一个巨大的痛苦来源。这就是心理问题的主观性。

\begin{center}\rule{0.5\linewidth}{0.5pt}\end{center}

\subsubsection{（二）必然性}\label{ux4e8cux5fc5ux7136ux6027}

老师接着举了\textbf{疑病症 /
疾病焦虑障碍}这一类例子（按语境应为疑病相关表现）。

【举例】一个高中女生担心自己得了狂犬病：

\begin{itemize}
\item
  她常去喂小区里的流浪猫；
\item
  有一天发现原来 5 只猫，现在只剩 4 只；
\item
  她立刻联想到：

  \begin{itemize}
  \tightlist
  \item
    少的那只是不是得狂犬病死了？
  \item
    那会不会传染给我？
  \end{itemize}
\end{itemize}

结果：

\begin{itemize}
\tightlist
\item
  一夜没睡；
\item
  恐惧袭来以后，怎么劝都不行。
\end{itemize}

父母带她去查：

\begin{itemize}
\tightlist
\item
  调监控；
\item
  去防疫站咨询；
\item
  医生说没有抓伤、没有外伤，问题不大；
\item
  她仍然不放心。
\end{itemize}

然后她使出焦虑者的``必杀技''：

\begin{itemize}
\tightlist
\item
  ``\textbf{万一呢？}''
\end{itemize}

因为医生也不可能绝对担保，她就继续陷在恐惧里。

即使后来去做抗原抗体检测、结果没事，她还是会在网上继续查资料。结果又看到一些信息，说某些阶段可能检测也未必能完全排除，于是焦虑仍然存在。

她每次见老师都很绝望，甚至说：

\begin{itemize}
\tightlist
\item
  狂犬病潜伏期大概多久；
\item
  过几天可能你就见不着我了。
\end{itemize}

\paragraph{老师借此说明``必然性''的意思：}\label{ux8001ux5e08ux501fux6b64ux8bf4ux660eux5fc5ux7136ux6027ux7684ux610fux601d}

\begin{itemize}
\tightlist
\item
  家长会觉得：如果她不去喂流浪猫，就不会有这个问题；
\item
  但老师的判断是：\textbf{即使不是这件事，她也会被别的事情困住}；
\item
  因为她身体里、内在里本来就存在一个等待爆发的恐惧模式。
\end{itemize}

也就是说：

【重点】心理问题不是单纯偶然被某件事``制造''出来的，而是内在模式遇到机会后\textbf{必然会寻找出口}。

\paragraph{另一个例子：收入问题并不是焦虑真正原因}\label{ux53e6ux4e00ux4e2aux4f8bux5b50ux6536ux5165ux95eeux9898ux5e76ux4e0dux662fux7126ux8651ux771fux6b63ux539fux56e0}

另一个来访者曾说：

\begin{itemize}
\tightlist
\item
  我的问题你解决不了；
\item
  因为我现在焦虑，是因为工作不好、没有收入；
\item
  只要收入好了，我就会好多了。
\end{itemize}

但后来他真的收入变好了，焦虑反而更大。于是他才意识到：

\begin{itemize}
\tightlist
\item
  自己过去以为的``原因''，并不是真正原因；
\item
  心理问题更像是身体内部早就存在的东西。
\end{itemize}

\begin{center}\rule{0.5\linewidth}{0.5pt}\end{center}

\subsubsection{（三）主导性}\label{ux4e09ux4e3bux5bfcux6027}

老师用强迫症继续解释``主导性''。

【举例】一个大货车司机：

\paragraph{1. 诱发事件}\label{ux8bf1ux53d1ux4e8bux4ef6}

\begin{itemize}
\item
  开大货车拐弯时，把旁边一辆电动车挂倒了；
\item
  从那以后，他开始反复担心：

  \begin{itemize}
  \tightlist
  \item
    我会不会又挂到别人？
  \end{itemize}
\end{itemize}

\paragraph{2. 问题扩展}\label{ux95eeux9898ux6269ux5c55}

\begin{itemize}
\tightlist
\item
  开车时别人往前看，他总忍不住往后看；
\item
  后来家人和别人觉得这样不行，干脆别开大货车了。
\end{itemize}

\paragraph{3. 症状迁移}\label{ux75c7ux72b6ux8fc1ux79fb}

\begin{itemize}
\tightlist
\item
  不开大货车后，改骑电动车；
\item
  骑电动车到了人多的地方，又担心会不会碰到别人；
\item
  于是电动车也不敢骑了。
\end{itemize}

\paragraph{4.
徒步时继续检查}\label{ux5f92ux6b65ux65f6ux7ee7ux7eedux68c0ux67e5}

\begin{itemize}
\item
  家里劝他要``克服困难''，去人多的地方；
\item
  他每次咬牙去，但总是不断检查：

  \begin{itemize}
  \tightlist
  \item
    我是不是撞到别人了？
  \item
    我是不是碰到别人了？
  \end{itemize}
\end{itemize}

老师强调：

【重点】真正让他越来越严重的，不只是最初那个念头，而是后面不断发生的：

\begin{itemize}
\tightlist
\item
  反复检查
\item
  不断对抗
\item
  不断卷入
\end{itemize}

这就是心理问题的``主导性''：

\begin{itemize}
\tightlist
\item
  问题一旦发生，
\item
  你越和它缠斗、越对抗它，
\item
  它往往越厉害。
\end{itemize}

\begin{center}\rule{0.5\linewidth}{0.5pt}\end{center}

\subsection{九、痛苦与``增容''：老师对苦难的理解}\label{ux4e5dux75dbux82e6ux4e0eux589eux5bb9ux8001ux5e08ux5bf9ux82e6ux96beux7684ux7406ux89e3}

接下来老师从更宏观、更价值性的角度谈``痛苦''。

\subsubsection{1.
不要只盯着``我有多痛苦''}\label{ux4e0dux8981ux53eaux76efux7740ux6211ux6709ux591aux75dbux82e6}

【举例】有来访者问老师：

\begin{itemize}
\tightlist
\item
  你见过比我遭受更严重痛苦的人吗？
\end{itemize}

老师后来用\textbf{郑和}做例子来回应：

\begin{itemize}
\tightlist
\item
  郑和原本是官宦家庭出身；
\item
  后来家破人亡，又遭到极大的人生创伤；
\item
  从处境上看，其痛苦远超一般人。
\end{itemize}

老师用这个例子想说明：

\begin{itemize}
\tightlist
\item
  不要一直沉浸在``我最痛苦''里；
\item
  更重要的是：\textbf{你如何把痛苦转化成力量，走出来}。
\end{itemize}

\subsubsection{2.
历史上能成事的人，往往都经历过磨难}\label{ux5386ux53f2ux4e0aux80fdux6210ux4e8bux7684ux4ebaux5f80ux5f80ux90fdux7ecfux5386ux8fc7ux78e8ux96be}

老师的观点是：

\begin{itemize}
\tightlist
\item
  历史上凡是对社会有积极贡献的人，
\item
  往往都经历过许多痛苦和磨难。
\end{itemize}

因为：

\begin{itemize}
\tightlist
\item
  人只有经历事，能力才会被锻炼出来；
\item
  如果一直顺利，容量就不会被扩大。
\end{itemize}

\subsubsection{3.
心理健康与``包容力''}\label{ux5fc3ux7406ux5065ux5eb7ux4e0eux5305ux5bb9ux529b}

老师提出：

【重点】心理健康，某种意义上就是你\textbf{能包容一切}的时候最健康。

\begin{itemize}
\tightlist
\item
  当你包容不了某些东西时，就会非常痛苦。
\item
  所以磨难可以理解为一种``扩容''\,``增容''。
\end{itemize}

老师的说法是：

\begin{itemize}
\tightlist
\item
  过得去，容量变大，层次升高；
\item
  过不去，就容易被打落下去。
\end{itemize}

\subsubsection{4.
换个角度看``难相处的人''}\label{ux6362ux4e2aux89d2ux5ea6ux770bux96beux76f8ux5904ux7684ux4eba}

当来访者总说：

\begin{itemize}
\tightlist
\item
  ``我有没有道理？''
\item
  ``明明就是对方不对。''
\end{itemize}

老师往往会回答：

\begin{itemize}
\tightlist
\item
  你说得有道理；
\item
  但换个角度，那个人其实是来帮你\textbf{增容}的。
\end{itemize}

也就是说：

\begin{itemize}
\tightlist
\item
  你连这样一个人都包容不下，
\item
  那将来更复杂的人和事你怎么面对？
\end{itemize}

老师最后把这个观点总结成一句话：

【观点】人生不是按``道理''活的，而是按``心量''活的。

\begin{center}\rule{0.5\linewidth}{0.5pt}\end{center}

\subsection{十、续讲：对心理问题的宏观理解与对自媒体的提醒}\label{ux5341ux7eedux8bb2ux5bf9ux5fc3ux7406ux95eeux9898ux7684ux5b8fux89c2ux7406ux89e3ux4e0eux5bf9ux81eaux5a92ux4f53ux7684ux63d0ux9192}

下半节课开始，老师希望大家形成对心理问题的宏观印象。

\subsubsection{1. 少看自媒体}\label{ux5c11ux770bux81eaux5a92ux4f53}

【学习建议】老师明确提醒：

\begin{itemize}
\tightlist
\item
  尽量少看自媒体。
\end{itemize}

理由包括：

\begin{itemize}
\tightlist
\item
  做自媒体的人为了生存，需要不断输出；
\item
  会非常消耗精气神；
\item
  平台内容里有大量情绪宣泄，不一定真正有助于解决问题。
\end{itemize}

老师观察到：

\begin{itemize}
\tightlist
\item
  如果网上发积极、阳光的内容，未必有多少人点赞；
\item
  但如果发抱怨、不公平、受委屈的内容，点赞会非常多。
\end{itemize}

这说明：

\begin{itemize}
\item
  自媒体平台上很多内容迎合的是\textbf{情绪宣泄}；
\item
  而对于负面情绪，老师主张：

  \begin{itemize}
  \tightlist
  \item
    \textbf{别和它纠缠}
  \item
    越纠缠，越出不来
  \end{itemize}
\end{itemize}

【举例】老师曾有一个学生在网上反复抱怨父母以前怎么对待自己，下面点赞很多。老师对他说：

\begin{itemize}
\tightlist
\item
  别再发这些东西了；
\item
  发完之后，你的问题并不会真正解决。
\end{itemize}

\begin{center}\rule{0.5\linewidth}{0.5pt}\end{center}

\subsection{十一、对``原生家庭决定论''的批评}\label{ux5341ux4e00ux5bf9ux539fux751fux5bb6ux5eadux51b3ux5b9aux8bbaux7684ux6279ux8bc4}

\subsubsection{1.
网络上常见说法：所有心理问题都来自原生家庭}\label{ux7f51ux7edcux4e0aux5e38ux89c1ux8bf4ux6cd5ux6240ux6709ux5fc3ux7406ux95eeux9898ux90fdux6765ux81eaux539fux751fux5bb6ux5ead}

老师特别点名这个流行观点：

【易错】``所有心理问题都是原生家庭（family of origin）的问题。''

老师对此持明显反对态度。

\subsubsection{2. 老师的理由}\label{ux8001ux5e08ux7684ux7406ux7531}

他举了一个历史社会层面的观察：

\begin{itemize}
\tightlist
\item
  50 年前，中国很多家庭有五六个、七八个孩子；
\item
  如果真的是原生家庭决定一切，
\item
  那同一个家庭里所有孩子都应普遍出问题；
\item
  但现实并非如此。
\end{itemize}

因此老师认为：

\begin{itemize}
\tightlist
\item
  原生家庭可能有影响；
\item
  但\textbf{不能简单建立直接必然因果关系}。
\end{itemize}

\subsubsection{3.
为什么这样说仍然不够}\label{ux4e3aux4ec0ux4e48ux8fd9ux6837ux8bf4ux4ecdux7136ux4e0dux591f}

老师说，仅仅这么讲没有说服力，因为当事人自己会觉得：

\begin{itemize}
\tightlist
\item
  ``你看我父母就是这样对我的。''
\item
  ``我的痛苦就是他们导致的。''
\end{itemize}

于是老师继续用案例来说明：
有时真正的问题，是\textbf{个体自己的内在加工方式}。

\begin{center}\rule{0.5\linewidth}{0.5pt}\end{center}

\subsection{十二、案例三：同样一件事，不同的人加工完全不同}\label{ux5341ux4e8cux6848ux4f8bux4e09ux540cux6837ux4e00ux4ef6ux4e8bux4e0dux540cux7684ux4ebaux52a0ux5de5ux5b8cux5168ux4e0dux540c}

\subsubsection{1.
关门声被体验成``暴力''}\label{ux5173ux95e8ux58f0ux88abux4f53ux9a8cux6210ux66b4ux529b}

有位来访者从老师诊室离开后，老师像平常一样关门。下次对方来时却问：

\begin{itemize}
\tightlist
\item
  ``你怎么这么暴力？''
\end{itemize}

老师很惊讶，对方解释说：

\begin{itemize}
\tightlist
\item
  你关门声音那么大。
\end{itemize}

老师借此说明：

\begin{itemize}
\tightlist
\item
  对他来说，普通关门声都能被体验成极大的冲击；
\item
  那么他平时如何体验父母的一声咳嗽、一次语气变化，就可想而知。
\end{itemize}

【重点】问题不只是外界发生了什么，更在于他的内在以什么方式加工这一切。

\subsubsection{2.
同样的母亲行为，不同叙述方向完全相反}\label{ux540cux6837ux7684ux6bcdux4eb2ux884cux4e3aux4e0dux540cux53d9ux8ff0ux65b9ux5411ux5b8cux5168ux76f8ux53cd}

另一个例子：

\paragraph{例子
A：负向加工}\label{ux4f8bux5b50-aux8d1fux5411ux52a0ux5de5}

有来访者回忆小时候和母亲去旅游：

\begin{itemize}
\tightlist
\item
  到旅馆后，床上的被子不够；
\item
  母亲想去别的床上拿被子，结果和别人发生争执；
\item
  这个来访者由此把母亲加工成``有问题''\,``令人厌恶''的形象。
\end{itemize}

\paragraph{例子
B：正向加工}\label{ux4f8bux5b50-bux6b63ux5411ux52a0ux5de5}

同样地，还有人说：

\begin{itemize}
\tightlist
\item
  大学快毕业时，火车上有人挤了自己一下；
\item
  父亲五十多岁了，还要为自己冲上去和别人打架；
\item
  于是得出结论：父亲对我多好。
\end{itemize}

老师借这两个例子说明：

【重点】同样是``父母做了一件事''，不同的人内在加工出来的意义可以完全不同。
而且最关键的是：\textbf{每个人都觉得自己的加工是对的。}

\begin{center}\rule{0.5\linewidth}{0.5pt}\end{center}

\subsection{十三、再次总结：心理问题与原生家庭``有关''，但不能锁死成因果}\label{ux5341ux4e09ux518dux6b21ux603bux7ed3ux5fc3ux7406ux95eeux9898ux4e0eux539fux751fux5bb6ux5eadux6709ux5173ux4f46ux4e0dux80fdux9501ux6b7bux6210ux56e0ux679c}

老师在这里再次强调一遍：

\begin{itemize}
\item
  心理问题和原生家庭\textbf{可以有关系}；
\item
  但绝不能简单归结为：

  \begin{itemize}
  \tightlist
  \item
    ``我的问题全是父母造成的。''
  \end{itemize}
\end{itemize}

因为一旦你这样归结：

\begin{itemize}
\tightlist
\item
  你会把自己的人生解释权完全交给过去；
\item
  也就更难走出来。
\end{itemize}

老师认为，这类说法之所以``害人不浅''，是因为：

\begin{itemize}
\tightlist
\item
  它\textbf{符合人的痛苦模式}；
\item
  但\textbf{不符合让人真正从痛苦中解脱出来的方法}。
\end{itemize}

\begin{center}\rule{0.5\linewidth}{0.5pt}\end{center}

\subsection{十四、人与世界如何互动，世界就如何回应你}\label{ux5341ux56dbux4ebaux4e0eux4e16ux754cux5982ux4f55ux4e92ux52a8ux4e16ux754cux5c31ux5982ux4f55ux56deux5e94ux4f60}

老师把心理问题进一步提升到``精神世界''的层面来理解。

\subsubsection{1.
你的方式，会塑造你得到的回应}\label{ux4f60ux7684ux65b9ux5f0fux4f1aux5851ux9020ux4f60ux5f97ux5230ux7684ux56deux5e94}

老师提出一个命题：

\begin{itemize}
\tightlist
\item
  你用什么样的方式和世界互动，
\item
  世界就会以什么样的方式回应你。
\end{itemize}

\subsubsection{2.
如果你总觉得自己忙、累，你就不容易快乐}\label{ux5982ux679cux4f60ux603bux89c9ux5f97ux81eaux5df1ux5fd9ux7d2fux4f60ux5c31ux4e0dux5bb9ux6613ux5febux4e50}

老师问大家：

\begin{itemize}
\tightlist
\item
  你们是不是总觉得自己很忙、很累？
\end{itemize}

然后直接指出：

\begin{itemize}
\tightlist
\item
  如果你一直以``忙、累''的方式和世界接触，
\item
  那么这种模式本身就会让你更难快乐。
\end{itemize}

\begin{center}\rule{0.5\linewidth}{0.5pt}\end{center}

\subsection{十五、案例四：小熊一家------``心想事成''的负向验证模式}\label{ux5341ux4e94ux6848ux4f8bux56dbux5c0fux718aux4e00ux5bb6ux5fc3ux60f3ux4e8bux6210ux7684ux8d1fux5411ux9a8cux8bc1ux6a21ux5f0f}

老师讲了一个小故事来说明``负向加工''如何把现实一步步导向自己最害怕的结果。

\subsubsection{1. 故事背景}\label{ux6545ux4e8bux80ccux666f}

\begin{itemize}
\item
  小熊以前没有和父母一起生活；
\item
  后来到了上学年纪，才被接到父母身边；
\item
  因为缺乏长期相处，小熊总有一种感觉：

  \begin{itemize}
  \tightlist
  \item
    \textbf{父母不爱我。}
  \end{itemize}
\end{itemize}

父母也感到：

\begin{itemize}
\tightlist
\item
  和孩子之间总有隔阂；
\item
  相处起来不轻松、不自在。
\end{itemize}

\subsubsection{2.
一家人去野餐}\label{ux4e00ux5bb6ux4ebaux53bbux91ceux9910}

熊爸爸为了改善关系，提议周末一家人去野餐。 一开始大家都很开心。

\subsubsection{3.
关键事件：让小熊回家拿起子}\label{ux5173ux952eux4e8bux4ef6ux8ba9ux5c0fux718aux56deux5bb6ux62ffux8d77ux5b50}

到中午吃饭时，罐头打不开，起子忘带了。 熊爸爸就让小熊跑回家拿起子。

这一刻，小熊立刻冒出一个念头：

\begin{itemize}
\tightlist
\item
  父母是不是故意把我支开？
\item
  我走了，他们是不是要背着我吃别的好吃的？
\end{itemize}

于是他说：

\begin{itemize}
\tightlist
\item
  我不去；
\item
  我去了你们把别的东西吃了怎么办？
\end{itemize}

父母向他保证：

\begin{itemize}
\tightlist
\item
  你不回来，我们不吃。
\end{itemize}

\subsubsection{4. 结局}\label{ux7ed3ux5c40-1}

按理说来回最多一个小时，但：

\begin{itemize}
\tightlist
\item
  一个小时过去了，小熊没回来；
\item
  两个小时过去了，小熊还没回来；
\item
  天快黑了，父母也饿了；
\item
  最后父母一人吃了一口面包，准备回家。
\end{itemize}

这时候，小熊从远处的树后哭着出来，质问：

\begin{itemize}
\tightlist
\item
  ``你们说过什么？你们说过什么？''
\end{itemize}

\subsubsection{5.
老师要说明的核心}\label{ux8001ux5e08ux8981ux8bf4ux660eux7684ux6838ux5fc3}

这个故事不是在说谁``客观上''绝对错，而是在说明：

【重点】一个人如果内心先有了``他们不爱我''的预设，
那么一遇到事情，就会往那个方向加工；
一旦这么加工，就容易把事情一步步推成自己最害怕的样子。

老师把这种现象概括为：

【概念】``心想事成''
也就是：\textbf{你心里害怕什么，往往就更容易发生什么。}

\begin{center}\rule{0.5\linewidth}{0.5pt}\end{center}

\subsection{十六、恋爱中的``验证''会制造痛苦闭环}\label{ux5341ux516dux604bux7231ux4e2dux7684ux9a8cux8bc1ux4f1aux5236ux9020ux75dbux82e6ux95edux73af}

【举例】一个学生谈恋爱，刚开始很好，后来心里的``心魔''出来了：

\begin{itemize}
\tightlist
\item
  总怀疑：对方到底爱不爱我？
\end{itemize}

于是开始验证。

\subsubsection{1.
验证方式：发消息看是否秒回}\label{ux9a8cux8bc1ux65b9ux5f0fux53d1ux6d88ux606fux770bux662fux5426ux79d2ux56de}

\begin{itemize}
\tightlist
\item
  如果秒回，就暂时安心；
\item
  如果对方没秒回，尤其在晚上睡觉时没有回，
\item
  她就开始焦虑，彻夜不安。
\end{itemize}

\subsubsection{2. 反复验证}\label{ux53cdux590dux9a8cux8bc1}

即便对方解释、安慰，过一段时间她又会想到：

\begin{itemize}
\tightlist
\item
  ``万一他变心了呢？''
\end{itemize}

于是继续验证、继续试探。

\subsubsection{3. 结果}\label{ux7ed3ux679c-1}

最终对方被折腾得受不了，说了一句：

\begin{itemize}
\tightlist
\item
  ``你这是无理取闹。''
\end{itemize}

她立刻崩溃，进一步确认：

\begin{itemize}
\tightlist
\item
  你看，你都对我发火了；
\item
  这说明你就是不爱我。
\end{itemize}

老师借此说明：

【重点】如果你总去``验证''别人是不是爱你、在乎你、认可你，
你很可能已经进入了痛苦的自动循环。

\subsubsection{4.
老师给出的建议}\label{ux8001ux5e08ux7ed9ux51faux7684ux5efaux8bae}

【学习建议】如果有这种担心，不要一味验证，而要练习：

\begin{itemize}
\tightlist
\item
  让自己平静下来；
\item
  和那个痛苦\textbf{共处（coexist with distress）}。
\end{itemize}

老师认为：

\begin{itemize}
\tightlist
\item
  和痛苦共处，是人成长的必由之路；
\item
  如果一直追求一种``绝对安全''，那反而是最不安全的事。
\end{itemize}

并进一步总结：

【观点】人生很多痛苦，都是``太认真''造成的。 ``既来之，则安之。''

\begin{center}\rule{0.5\linewidth}{0.5pt}\end{center}

\subsection{十七、为什么人天然更容易往负面加工：老师所谓的``魔''}\label{ux5341ux4e03ux4e3aux4ec0ux4e48ux4ebaux5929ux7136ux66f4ux5bb9ux6613ux5f80ux8d1fux9762ux52a0ux5de5ux8001ux5e08ux6240ux8c13ux7684ux9b54}

老师接着追问：

\begin{itemize}
\tightlist
\item
  为什么小熊总往负面想？
\item
  为什么大家一闲下来，很难自动加工出积极的内容？
\end{itemize}

老师的观点是：

\begin{itemize}
\tightlist
\item
  人的头脑天然更容易往``危险''\,``不安全''``会出事''的方向加工；
\item
  很难自发稳定地产生积极、安稳的内容。
\end{itemize}

他把这种倾向称为：

【概念】``魔''

也就是：

\begin{itemize}
\item
  头脑会自动加工：

  \begin{itemize}
  \tightlist
  \item
    我不这么做会不会危险？
  \item
    我不这么做会不会出问题？
  \item
    我不这样努力将来会不会糟糕？
  \end{itemize}
\end{itemize}

老师的意思是：这种倾向不是个别人有，而是人普遍存在的。

\begin{center}\rule{0.5\linewidth}{0.5pt}\end{center}

\subsection{十八、借《西游记》谈``磨炼''与精神内核}\label{ux5341ux516bux501fux897fux6e38ux8bb0ux8c08ux78e8ux70bcux4e0eux7cbeux795eux5185ux6838}

老师引用《西游记》来说明：

\begin{itemize}
\tightlist
\item
  人生中的种种困难、折磨，本身也可以视作一种``炼''。
\end{itemize}

他提到几个意象：

\begin{itemize}
\tightlist
\item
  孙悟空为什么要在八卦炉里炼？
\item
  为什么会经历那么多磨难？
\item
  为什么经历关键冲突之后，团队反而更稳固？
\end{itemize}

老师把大家的现实生活比作也在``八卦炉''里被炼：

\begin{itemize}
\tightlist
\item
  就像烤串一样，翻过来烤、调过去烤；
\item
  有人能熬过去，有人会``烤糊''。
\end{itemize}

他的核心意思是：

\begin{itemize}
\tightlist
\item
  若能在磨炼中练出``火眼金睛''，
\item
  就能识破生活中很多看似重要、其实并不那么重要的东西；
\item
  真正重要的是培养自己内在的精神品质。
\end{itemize}

\subsubsection{1.
老师强调``精神内核''}\label{ux8001ux5e08ux5f3aux8c03ux7cbeux795eux5185ux6838}

他认为：

\begin{itemize}
\tightlist
\item
  一个社会、一个国家、一个家庭，真正重要的是人有没有精神内核；
\item
  遇到困难时，是选择抱怨，还是选择承担、建设。
\end{itemize}

\subsubsection{2.
家庭为什么会过不好}\label{ux5bb6ux5eadux4e3aux4ec0ux4e48ux4f1aux8fc7ux4e0dux597d}

老师认为很多家庭的问题都在于：

\begin{itemize}
\item
  每个人都理直气壮地抱怨别人不对；
\item
  但很少有人想着：

  \begin{itemize}
  \tightlist
  \item
    我能为这个家付出什么？
  \end{itemize}
\end{itemize}

他的结论是：

【重点】一个家庭若人人都想着付出、建设它，就不容易过不好；
如果人人都盯着别人不对，这个家庭就会一直鸡飞狗跳。

\begin{center}\rule{0.5\linewidth}{0.5pt}\end{center}

\subsection{十九、修身齐家：少责备，多承担}\label{ux5341ux4e5dux4feeux8eabux9f50ux5bb6ux5c11ux8d23ux5907ux591aux627fux62c5}

老师把这一层上升到中国传统文化的表达：

\begin{itemize}
\tightlist
\item
  修身
\item
  齐家
\item
  治国
\item
  平天下
\end{itemize}

\subsubsection{1.
``齐家''不是要求别人，而是先承担自己那部分责任}\label{ux9f50ux5bb6ux4e0dux662fux8981ux6c42ux522bux4ebaux800cux662fux5148ux627fux62c5ux81eaux5df1ux90a3ux90e8ux5206ux8d23ux4efb}

老师认为：

\begin{itemize}
\tightlist
\item
  在家庭中应先承担责任，而不是先责怪别人；
\item
  在集体中应先想着如何让集体更好，而不是总挑别人毛病。
\end{itemize}

\subsubsection{2.
``仁义礼智信''不是僵化教条}\label{ux4ec1ux4e49ux793cux667aux4fe1ux4e0dux662fux50f5ux5316ux6559ux6761}

老师认为，传统文化中很多内容并不是死板教条，而是帮助人：

\begin{itemize}
\tightlist
\item
  找到自己的位置；
\item
  承担自己的角色；
\item
  让社会按合理规则运转。
\end{itemize}

他以``君君臣臣，父父子子''为例解释说：

\begin{itemize}
\item
  不是在强调盲目服从；
\item
  而是在强调：

  \begin{itemize}
  \tightlist
  \item
    处在什么位置，就承担什么位置的责任。
  \end{itemize}
\end{itemize}

\subsubsection{3.
``礼崩乐坏''的理解}\label{ux793cux5d29ux4e50ux574fux7684ux7406ux89e3}

老师的解释是：

\begin{itemize}
\tightlist
\item
  当人心完全被负面、抱怨、挑剔、私欲所控制时，
\item
  社会就不能按合理规则运转，
\item
  就会出现秩序紊乱。
\end{itemize}

\begin{center}\rule{0.5\linewidth}{0.5pt}\end{center}

\subsection{二十、识人之术：远离持续抱怨外界的人}\label{ux4e8cux5341ux8bc6ux4ebaux4e4bux672fux8fdcux79bbux6301ux7eedux62b1ux6028ux5916ux754cux7684ux4eba}

【学习建议】老师提醒大家：

\begin{itemize}
\item
  对于那种整天抱怨：

  \begin{itemize}
  \tightlist
  \item
    这个人不好
  \item
    那个事不对
  \item
    外界全错
  \end{itemize}
\item
  的人，要尽量远离。
\end{itemize}

老师的理由是：

\begin{itemize}
\tightlist
\item
  一个人心里装着什么，就更容易看到外界什么。
\end{itemize}

\subsubsection{1.
苏东坡与佛印的故事}\label{ux82cfux4e1cux5761ux4e0eux4f5bux5370ux7684ux6545ux4e8b}

老师讲了两个与苏东坡、佛印有关的故事。

\paragraph{故事一：八风吹不动，一屁打过江}\label{ux6545ux4e8bux4e00ux516bux98ceux5439ux4e0dux52a8ux4e00ux5c41ux6253ux8fc7ux6c5f}

苏东坡自觉修养很好，觉得自己``八风吹不动''，写给佛印看。
佛印只回了一个``屁''字。 苏东坡顿时生气，跑去理论。

佛印借此指出：

\begin{itemize}
\tightlist
\item
  你不是说自己不动心吗？
\item
  怎么一个``屁''字就把你打过江来了？
\end{itemize}

老师用这个故事说明：

\begin{itemize}
\tightlist
\item
  表面上的``看开''不算真功夫；
\item
  真正的心量，是能不能经得起刺激。
\end{itemize}

\paragraph{故事二：心里有什么，就看别人像什么}\label{ux6545ux4e8bux4e8cux5fc3ux91ccux6709ux4ec0ux4e48ux5c31ux770bux522bux4ebaux50cfux4ec0ux4e48}

苏东坡说：

\begin{itemize}
\tightlist
\item
  我看自己像佛；
\item
  看别人像狗屎。
\end{itemize}

佛印说：

\begin{itemize}
\tightlist
\item
  我看你像佛；
\item
  我看自己像狗屎。
\end{itemize}

后来苏东坡妹妹提醒他：

\begin{itemize}
\tightlist
\item
  你心里有什么，才会看到别人像什么；
\item
  佛印心里有佛，所以看你像佛；
\item
  你心里装着``狗屎''，才会看别人像狗屎。
\end{itemize}

老师借此总结：

【重点】如果你总觉得这个世界到处都不对、别人都不行，往往也说明你自己的内在加工层次出了问题。

\begin{center}\rule{0.5\linewidth}{0.5pt}\end{center}

\subsection{二十一、关于目标、努力与放松}\label{ux4e8cux5341ux4e00ux5173ux4e8eux76eeux6807ux52aaux529bux4e0eux653eux677e}

老师最后又谈到``有心栽花花不开，无心插柳柳成荫''。

\subsubsection{1.
太执着于目标，反而容易痛苦}\label{ux592aux6267ux7740ux4e8eux76eeux6807ux53cdux800cux5bb9ux6613ux75dbux82e6}

老师提醒大家：

\begin{itemize}
\tightlist
\item
  如果你特别执着于某个目标，
\item
  到了``我必须得到它''的程度，
\item
  那往往会让自己很痛苦。
\end{itemize}

因此他建议：

【学习建议】对自己宽容一点，目标放一放，别总和自己过不去。

\begin{center}\rule{0.5\linewidth}{0.5pt}\end{center}

\subsection{二十二、临下课前的两道题提示}\label{ux4e8cux5341ux4e8cux4e34ux4e0bux8bfeux524dux7684ux4e24ux9053ux9898ux63d0ux793a}

【考点/课后题提示】老师提醒大家，课后有两道选择题，大约 10
分，而且都很简单，实际上课上已经讲到答案方向了：

\subsubsection{1. 第一题}\label{ux7b2cux4e00ux9898}

\begin{itemize}
\tightlist
\item
  当事情涉及危机时，
\item
  \textbf{生命最重要}。
\end{itemize}

\subsubsection{2. 第二题}\label{ux7b2cux4e8cux9898}

\begin{itemize}
\item
  当困难超出自己能力范围时，
\item
  \textbf{不要自己硬扛}；
\item
  可以：

  \begin{itemize}
  \tightlist
  \item
    听课学习
  \item
    做心理咨询
  \item
    和朋友倾诉
  \end{itemize}
\end{itemize}

\begin{center}\rule{0.5\linewidth}{0.5pt}\end{center}

\subsection{二十三、人生建议：千里马常有，而伯乐不常有}\label{ux4e8cux5341ux4e09ux4ebaux751fux5efaux8baeux5343ux91ccux9a6cux5e38ux6709ux800cux4f2fux4e50ux4e0dux5e38ux6709}

临下课前，老师又讲了关于人生发展的一些思考。

\subsubsection{1.
大家都是``千里马''}\label{ux5927ux5bb6ux90fdux662fux5343ux91ccux9a6c}

老师认为：

\begin{itemize}
\tightlist
\item
  在座各位都可以是``千里马''；
\item
  但问题在于，未必总能遇到真正识别和使用人才的环境，也就是``伯乐''。
\end{itemize}

\subsubsection{2.
人生中三个``好''很重要}\label{ux4ebaux751fux4e2dux4e09ux4e2aux597dux5f88ux91cdux8981}

老师认为，一个人若有下面三个条件，人生会顺很多：

\begin{enumerate}
\def\labelenumi{\arabic{enumi}.}
\tightlist
\item
  \textbf{好的父母}：给予基础支持；
\item
  \textbf{好的老师}：帮助你找到适合自己的方向；
\item
  \textbf{好的团队/好的老板}：给你发挥能力的舞台。
\end{enumerate}

\subsubsection{3. 因材施教（teaching students in accordance with their
aptitude）}\label{ux56e0ux6750ux65bdux6559teaching-students-in-accordance-with-their-aptitude}

老师谈到``好的老师''时，提到了孔子的一个重要观念：

【术语】因材施教

老师的理解是：

\begin{itemize}
\item
  真正好的老师，不是让所有人都走同一条路；
\item
  而是帮助你发现：

  \begin{itemize}
  \tightlist
  \item
    你的天命在哪里
  \item
    你的天赋在哪里
  \item
    你真正适合做什么
  \end{itemize}
\end{itemize}

\subsubsection{4.
良心拷问：你真的喜欢自己学的专业吗？}\label{ux826fux5fc3ux62f7ux95eeux4f60ux771fux7684ux559cux6b22ux81eaux5df1ux5b66ux7684ux4e13ux4e1aux5417}

老师最后问大家：

\begin{itemize}
\tightlist
\item
  你真的喜欢现在学的专业吗？
\item
  如果实在不喜欢，就不要只是机械地逼迫自己。
\end{itemize}

他的意思是：

\begin{itemize}
\tightlist
\item
  若天赋、兴趣根本不在这里，
\item
  一味苦撑，可能是在浪费生命。
\end{itemize}

\subsubsection{5.
如何判断什么是你真正适合的事}\label{ux5982ux4f55ux5224ux65adux4ec0ux4e48ux662fux4f60ux771fux6b63ux9002ux5408ux7684ux4e8b}

老师给了一个标准：

【判断标准】 如果有一件事情：

\begin{itemize}
\tightlist
\item
  没有外在报酬；
\item
  没有人肯定你；
\item
  暂时也没有结果；
\end{itemize}

但你做起来依然：

\begin{itemize}
\tightlist
\item
  开心
\item
  快乐
\item
  愿意一直做
\end{itemize}

那它很可能就是你真正擅长、真正适合的方向。

老师鼓励大家：

\begin{itemize}
\tightlist
\item
  不要放弃这种事情；
\item
  继续坚持下去，
\item
  你会更轻松地走向自己的目标和成绩。
\end{itemize}

\begin{center}\rule{0.5\linewidth}{0.5pt}\end{center}

\subsection{本节英文术语索引}\label{ux672cux8282ux82f1ux6587ux672fux8bedux7d22ux5f15-6}

\begin{itemize}
\tightlist
\item
  心理危机干预 -- psychological crisis intervention
\item
  开卷考试 -- open-book exam
\item
  注意缺陷多动障碍 -- Attention-Deficit/Hyperactivity Disorder（ADHD）
\item
  心理咨询 -- psychological counseling / psychotherapy
\item
  双椅技术 -- two-chair technique / empty-chair technique
\item
  精神世界 -- mental world / inner psychic world
\item
  社交恐惧症 / 社交焦虑障碍 -- Social Anxiety Disorder（SAD）
\item
  广场恐惧症 -- agoraphobia
\item
  强迫症 -- Obsessive-Compulsive Disorder（OCD）
\item
  疑病症 / 疾病焦虑障碍 -- illness anxiety disorder / hypochondriasis
\item
  主观性 -- subjectivity
\item
  客观存在 -- objective existence
\item
  因材施教 -- teaching students in accordance with their aptitude
\end{itemize}

如果你愿意，我下一条可以继续把这份笔记再压缩成一版更适合考前速看的``提纲版''。

\end{document}
